\documentclass[11pt,a4paper]{article}
\usepackage[T1]{fontenc}
\usepackage[utf8]{inputenc}
\usepackage{lmodern}
\ifdefined\pdfmapfile
  \pdfmapfile{+lm.map}
  \pdfmapfile{+cm.map}
  \pdfmapfile{+symbols.map}
  \pdfmapfile{+euler.map}
  \pdfmapfile{+rsfs.map}
\fi
\usepackage[margin=27mm]{geometry}
\usepackage{amsmath,amssymb,amsthm,mathtools}
\usepackage{mathrsfs}
\usepackage{array,longtable,booktabs,enumitem}
\usepackage{xcolor}
\usepackage{microtype}
\usepackage{hyperref}
\usepackage{xurl}
\hypersetup{colorlinks=true,linkcolor=blue!45!black,urlcolor=blue!45!black,
 pdftitle={ALP CONNES: consequences, constructions, and the cubic gate},
 pdfauthor={Research collaboration, 22 November 2026}}
\newtheorem{theorem}{Theorem}[section]
\newtheorem{proposition}[theorem]{Proposition}
\newtheorem{lemma}[theorem]{Lemma}
\newtheorem{corollary}[theorem]{Corollary}
\theoremstyle{definition}
\newtheorem{definition}[theorem]{Definition}
\theoremstyle{remark}
\newtheorem{remark}[theorem]{Remark}
\newtheorem{example}[theorem]{Example}
\newcommand{\R}{\mathbb R}
\newcommand{\C}{\mathbb C}
\newcommand{\Q}{\mathbb Q}
\newcommand{\Z}{\mathbb Z}
\newcommand{\N}{\mathbb N}
\newcommand{\A}{\mathbb A}
\DeclareMathOperator{\Tr}{Tr}
\DeclareMathOperator{\Spec}{Spec}
\DeclareMathOperator{\Hom}{Hom}
\DeclareMathOperator{\Gal}{Gal}
\DeclareMathOperator{\supp}{supp}
\DeclareMathOperator{\rank}{rank}
\DeclareMathOperator{\End}{End}
\setlength{\emergencystretch}{3em}
\setlist{nosep}
\newcommand{\proofstep}[1]{\par\smallskip\noindent\textbf{#1}\quad}

\title{Alp\"oge's theorem through Connes's programme:\\
proved consequences, explicit constructions, and the cubic gate}
\author{ALP\_CONNES research record}
\date{22 November 2026\quad--\quad working paper}

\begin{document}
\maketitle
\begin{abstract}
We take the supplied theorem A of Alp\"oge (20 November 2026) as an
input: RH for the Riemann zeta function and for all primitive
Dirichlet $L$-functions. We distinguish its consequences from
constructions already unconditional in the sources. The proved
results include rational-character Weil positivity, the exact
scope of the cohomological and Sobolev spectral realizations,
actual finite-prolate defect convergence on compact Mellin-frequency
intervals, closed finite-window operators, quantitative small Weil
eigenvalues, numerical Hodge data, and explicit Artin--Tate and
ray-class models. Follow-up~1 adds the BCM and CM systems,
neutrality, the local-factor moment problem, and the 2026
geometric and trace papers. The finite-order character at
$p=269$ has an explicit maximal integral ray-class host and
an unconditional strictly positive Weil window through
$\lambda=4$. Its positive-weight strict $3$-adic Selmer
group is finite cyclic, with one canonical residual generator;
an integral CM Hom realizes its character. An exact
Poitou--Tate/K-theory order formula identifies the remaining
size invariant, and ramified-inertia descent removes the
integral character-comparison error. Certified logarithmic
class calculations, exact genus theory and conjugation symmetry
prove that its split cyclotomic Iwasawa module is finite.
Consequently these fixed-field three-primary lengths lie
between four and five and are periodic in the weight modulo three.
The only possible residue patterns are $(4,4,4)$, $(5,4,4)$ and
$(4,5,4)$; a specified first-layer conjugation component decides which.
The other coefficient
primes and the growing host comparison remain open.
Unconditional theta quasimodes and strong
resolvent limits extend to the degree-two cubic Hecke sector.
The arithmetic sparse-window inequality H04.4
for that cubic character has not been proved. Nor is a
gap in a proposed inference described as a theorem of
logical independence. Every numbered finding has a full
proof; the check register specifies which readings and
computations have actually been performed.
\end{abstract}

\tableofcontents
\clearpage

\section{Input, normalization, and the state of the questions}
\label{sec:input}

\paragraph{Theorem A (Alp\"oge, supplied theorem, 20 November 2026).}
The Riemann zeta function satisfies RH. Dirichlet $L$-functions
have no nontrivial zero with real part greater than $1/2$;
all nontrivial zeros of every primitive Dirichlet function
have real part $1/2$. The original prompt and its follow-up
are included as \path{TASK_PROMPT.md} and
\path{TASK_FOLLOWUP_1.md}. This paper accepts A exactly in
that role. It does not claim an independent Lean compilation.

The attached map \path{OURS_ALPOGE_LEAN_MAP_20261120.md}
locates the public statements at \path{Main.lean}, lines
4, 6, 8, 10 and 12, and \path{Statements.lean}, lines
3--17; the kernel-level interface is in
\path{Supplements/Kernels/MainTheorems.lean}, lines 7--16.
The map is a textual audit of the supplied archive, not a
build certificate. The generic analytic theorem actually
recovered and read here is \texttt{Deduction.landau\_window}
in \path{Development/AutomorphicArithmetic/DedekindLFunctions/LandauTheorems/landauWindow.lean}.
Its arithmetic premise remains separate from that
generic analytic implication. The four-page
\path{ALPOGE_remark.pdf} was read by text extraction;
no visual scan inspection or proof verification of A is
claimed by that operation.

\paragraph{Prior-run inputs.}
The starting point is
\path{RETURN_ALP_MACHINE_checkpoint_04_20261121T2304Z},
especially AM-H01--AM-H04 and AM-H06--AM-H08, and
\path{RETURN_MAS_ARTIN_checkpoint_10_20261121T2301Z},
especially MA-FT-001--MA-FT-002 and the actual character
Gram and trace interfaces they use. We carry their proved
statements with their hypotheses and distinguish a full
reading in this run from a carried background assertion.
The registers give the check type for each such use.
No numbered assertion from those runs is relabelled here
as a new theorem.

\paragraph{What the status words mean.}
``Proved'' below means the displayed conclusion follows
from the expressly stated inputs; some such inputs include A.
``Unconditional'' means A and cubic RH are not needed.
A checked counterexample refutes only the precise
formulation specified. An open comparison, a missing
construction, or an unavailable archival source is a
gate, not a positive or negative theorem about every
possible approach. Source statements known before this
run are cited as such. The file \path{QUESTIONS.tsv}
tracks the individual questions; \path{CHECK_EVENTS.tsv}
records checks rather than projecting a stronger status
from another statement's verification.

\begin{center}
\begin{tabular}{p{.15\textwidth}p{.38\textwidth}p{.37\textwidth}}
\toprule
Task & Proved or concretely constructed & Remaining comparison or gate\\
\midrule
Trace and windows & All rational unitary sectors, Weil positivity,
closed compact-resolvent windows, precise spectral multiplicities &
High-frequency control for the original C99 cutoff comparison;
local zero-measure convergence is proved; unitarization of nonsimple jet blocks\\
Spectral triples & Fixed-window Galerkin convergence, quantitative
low-energy bands and unconditional theta quasimodes & Selection and comparison of
the specific lowest eigenfunction and the proposed prolate trial\\
Sites and motives & Analytic Hodge inequality, numerical square,
actual Artin--Tate examples and sheaf constructions & A geometric
surface square, its polarized cohomology and surface Riemann--Roch\\
BCM and neutrality & Exact positive/twisted trace distinction, cubic
arithmetic observable, finite unit compensation and norm counts &
No KMS or local-purity statement supplies the cubic arithmetic ramp\\
Cubic Kummer & Minimal conductor host, strict finite test windows,
explicitly bounded, three-periodic fixed-field three-primary strict groups and integral CM
coefficients & H04.4; the other coefficient primes and quantitative
transport to growing hosts\\
\bottomrule
\end{tabular}
\end{center}

\part{Proved findings and source corrections}

\section{All rational characters, the trace pairing, and finite windows}
\label{sec:trace-windows}

This section uses the following source abbreviations. C99 is Connes,
\emph{Trace formula in noncommutative geometry and the zeros of the
Riemann zeta function}, the supplied \texttt{main.tex}
(arXiv:math/9811068). CCM07 is Connes--Consani--Marcolli,
\emph{The Weil proof and the geometry of the adeles class space},
the supplied \texttt{Yuri70v2.tex} (arXiv:math/0703392).
CC20 is Connes--Consani, \emph{Weil positivity and Trace formula:
the archimedean place}, the supplied \texttt{weil-compo.tex}
(arXiv:2006.13771). The statement numbers below have been checked
against the theorem-counter declarations and the complete sequence
of numbered environments in these supplied sources. Source labels
are also supplied, to make the locators independent of typesetting.

\subsection{The character quantifier and the precise positivity theorem}

Write \(C_{\mathbb Q}=\mathbb A_{\mathbb Q}^{\times}/\mathbb Q^\times\),
and give its compact norm-one subgroup Haar measure one. Let
\(\mathcal D(C_{\mathbb Q})\) denote finite sums of locally constant
functions on the compact factor times smooth functions of compact
support on \(\mathbb R_+^\times\). Convolution and adjoint use
multiplicative Haar measure:
\[
 (f*g)(x)=\int f(y)g(y^{-1}x)\,d^\times y,\qquad
 f^*(x)=\overline{f(x^{-1})}.
\]
Put
\[
 \widehat f(\chi,z)=
 \int_{\widehat{\mathbb Z}^{\times}}\int_0^\infty
 f(u,r)\overline{\chi(u)}r^z\,\frac{dr}{r}\,du .
\]
Here \(\chi\) labels a primitive Dirichlet character with
Euler coefficients \(\chi(p)\). Its associated Hecke character
has compact-unit part \(\overline\chi\) in the rational
representative convention used here. It is not counted again
each time it is induced to a larger modulus.

\begin{proposition}[AC-T01: the exact rational character range]
\label{AC-T01}
There is a topological group isomorphism
\[
 C_{\mathbb Q}\simeq\widehat{\mathbb Z}^{\times}\times
 \mathbb R_+^\times .
\]
Every continuous \emph{unitary} character of \(C_{\mathbb Q}\) is
\(\chi(u)r^{it}\), for a finite-order Dirichlet character \(\chi\)
and a real number \(t\). Consequently A implies RH for every
unitarily normalized Gr\"ossencharacter of \(\mathbb Q\).
A nonunitary norm twist changes the center of the critical strip;
it must first be removed before real part \(1/2\) is asserted.
\end{proposition}

\begin{proof}
For an idele \(a=(a_\infty,(a_p))\), choose the positive rational
number \(b=\prod_p p^{-v_p(a_p)}\). Then each \(ba_p\) is a unit.
Multiplication by the rational number \(-1\), if necessary, makes
the real component positive. This gives a representative
\((r,u)\) with \(r>0\) and \(u\in\widehat{\mathbb Z}^{\times}\).
A rational number preserving all finite valuations is \(\pm1\);
preserving positivity at infinity then makes it \(1\). The
representative is unique. The construction and its inverse are
continuous, and multiplication of representatives proves the
group assertion. The idele norm is \(r\).

A continuous character of the compact profinite group
\(\widehat{\mathbb Z}^{\times}\) has open kernel. To see this
directly, take an arc about \(1\) in the circle which contains no
nontrivial subgroup, and an open subgroup whose image lies in that
arc. Its image is trivial. Thus the character factors through
\((\mathbb Z/q\mathbb Z)^\times\) for some \(q\).
A continuous unitary character of \(\mathbb R_+^\times\), after
the logarithm, is a character \(x\mapsto e^{itx}\) of
\((\mathbb R,+)\). These observations give the classification.
The corresponding \(L\)-function is \(L(s+it,\overline\chi)\)
in these compact-unit coordinates.
Translation by \(it\) preserves real parts of zeros, so A gives
the assertion. An imprimitive Euler presentation is not a new
Hecke character; the assertion concerns the primitive completed
\(L\)-function, not extra zeros of omitted Euler factors.
\end{proof}

For each primitive \(\chi\), write its nontrivial zeros as
\(\rho=\frac12+i\gamma\), and their multiplicities as \(m_{\chi,\gamma}\).
Under A, all \(\gamma\) are real. Define the centered Weil
functional by its spectral side,
\[
 \mathcal W(f)=
 \sum_{\chi}\sum_\gamma m_{\chi,\gamma}
       \widehat f(\chi,i\gamma).                         \tag{T.1}
\]
For a test in \(\mathcal D\) the outer sum is finite and the inner
sums are absolutely convergent. This is also the centered
geometric distribution given by C99, Appendix II, Theorem 6.
The same identity on the strong Schwartz space used in CCM07 is
CCM07, Theorem 4.16 (label \texttt{varthetagTrH1}) together with
Theorem 6.1 (label \texttt{TraceformulaThm}), with the factor
\(|\cdot|^{-1/2}\) inserted.

\begin{theorem}[AC-T02: the whole rational Weil form and the CCM trace pairing]
\label{AC-T02}
Assuming A,
\[
 \mathcal W(f*f^*)=
 \sum_{\chi,\gamma}m_{\chi,\gamma}
          |\widehat f(\chi,i\gamma)|^2\geq0.             \tag{T.2}
\]
For the relative cyclic cohomology object \(H^1\) of CCM07,
Definition 4.10 (label \texttt{motH1AKCK}), this is precisely
\[
 \operatorname{Tr}\!
 \left(\underline\vartheta_m
       ((\Delta^{-1/2}f)*(\Delta^{-1/2}f^*))\mid H^1\right)
       \geq0,\qquad \Delta^z f(x)=|x|^z f(x).             \tag{T.3}
\]
Equivalently, the trace pairing satisfies
\(\langle h,h^\sharp\rangle\geq0\), where
\(h^\sharp(x)=|x|^{-1}\overline{h(x^{-1})}\).
These assertions hold on CCM07's strong Schwartz domain as well
as on \(\mathcal D\). In the converse direction, positivity of
(T.2) on all compactly supported tests implies A.
\end{theorem}

\begin{proof}
Multiplicative Fourier transformation takes convolution to
multiplication. On a unitary character \((u,r)\mapsto
\overline{\chi(u)}r^{i\gamma}\), it takes \(f^*\) to the conjugate of the
transform of \(f\). Under A every point occurring in (T.1) is
unitary. Substitution gives (T.2) term by term. Absolute
convergence justifies the substitution; equivalently it follows
from the trace-class assertion in the cited source. The same
argument applies to the strong Schwartz space, whose defining
decay makes the spectral trace converge.

The source normalization is
\(\underline\vartheta_m(x)=|x|^{1/2}W(x)\);
see CCM07 (label \texttt{Wvarthetag}) in the proof of
Theorem 4.16. Hence
\(W(f)=\underline\vartheta_m(\Delta^{-1/2}f)\).
Multiplication by \(|x|^{-1/2}\) is an automorphism of convolution,
so this equality converts (T.2) to (T.3).
Setting \(h=\Delta^{-1/2}f\) gives
\(\Delta^{-1/2}f^*=h^\sharp\), proving the second form.
The converse is Weil's positivity criterion in precisely the
form of CCM07, Proposition 6.2
(label \texttt{RHpositTr}) and Corollary 6.3
(label \texttt{RHpositTrsharp}), followed by AC-T01.
Thus no additional unproved positivity assertion is used here.
\end{proof}

The cohomological/nuclear trace identity of CCM07, Theorem 6.1
was already an unconditional theorem; A proves its positive
Hermitian specialization, not its existence. Similarly, the
radical statement and the ability to modify representatives
in CCM07, Lemma 4.17 and Proposition 6.4
(labels \texttt{vanishV}, \texttt{pairTrvanishV}) are
unconditional. A permits quotienting the now positive form
by its full nullspace and completing it. It does not assert
that the previously identified subspace \(\mathcal V\) is
the entire nullspace.

\subsection{What the original spectral realization becomes}

\begin{proposition}[AC-T03: exhaustion, multiplicities, and Jordan blocks]
\label{AC-T03}
For \(k=\mathbb Q\), every character sector in C99,
Theorem III.1, now has as its spectral set \emph{all} the
nontrivial zeros of the corresponding Dirichlet \(L\)-function,
after subtracting \(1/2\). For its Sobolev parameter
\(\delta>1\), a zero of multiplicity \(m\) occurs with the
multiplicity
\[
 m_\delta=\min(m,N_\delta),\qquad
 N_\delta=\max\{n\in\mathbb Z_{\geq1}:n<(1+\delta)/2\}.    \tag{T.4}
\]
Every prescribed finite multiset of zeros is realized with
its full multiplicities by choosing \(\delta\) sufficiently
large. A makes no change to formula (T.4). The stated source
results supply no uniform multiplicity bound independent
of height, and do not establish that this particular
generator becomes skew-adjoint.
\end{proposition}

\begin{proof}
Theorem III.1 states that the spectrum consists exactly of
the critical zeros, with (T.4). A says that every nontrivial
zero is critical, so the spectral-set conclusion follows.
For a finite multiset, let \(M\) be its largest multiplicity
and choose \(\delta>2M-1\). Then \(N_\delta\geq M\).
The supplied theorem contains no assertion that its
multiplicity formula changes when RH holds. Moreover,
the paragraph immediately after C99, III.(30), explicitly
states that a multiple zero produces a nontrivial Jordan
form when \(\delta\) is sufficiently large. A skew-adjoint
operator is normal and has no nontrivial finite-dimensional
Jordan block: if \((D-i\gamma)v=w\ne0\) and
\((D-i\gamma)w=0\), then skew-adjointness gives
\(\|w\|^2=\langle (D-i\gamma)v,w\rangle
=-\langle v,(D-i\gamma)w\rangle=0\), a contradiction.
A is a zero-location assertion, not a simplicity assertion.
Thus the stated skew-adjointness inference would require
further information, not provided by A or by Theorem III.1.
\end{proof}

\begin{theorem}[AC-T04: an explicit positive realization, with its exact scope]
\label{AC-T04}
A permits the following constructions.

First, the pre-Hilbert space
\[
 \mathcal D(C_{\mathbb Q})/\mathcal N,\qquad
 \mathcal N=\{f:\widehat f(\chi,i\gamma)=0
                    \text{ for every }(\chi,\gamma)\},
 \quad
 (f,g)_{\mathcal W}=\mathcal W(f*g^*)
                                                               \tag{T.5}
\]
has a unitary action of \(C_{\mathbb Q}\) by translation.
Its completion is the positive trace-pairing realization.

Second, a multiplicity-preserving trace model is
\[
 \mathcal H_{\rm mult}
  =\bigoplus_{\chi,\gamma}\mathbb C^{m_{\chi,\gamma}},
 \qquad
 H(v_{\chi,\gamma})=(\gamma v_{\chi,\gamma}),\quad
 \operatorname{Dom}H
  =\left\{v:\sum_{\chi,\gamma}
             \gamma^2\|v_{\chi,\gamma}\|^2<\infty\right\}.        \tag{T.6}
\]
The diagonal action
\[
 U(u,r)|_{\mathbb C^{m_{\chi,\gamma}}}
       =\overline{\chi(u)}r^{i\gamma}\operatorname{id}
\]
is unitary and \(H\) is self-adjoint. For finite-character
smooth compactly supported \(f\), the operator \(U(f)\)
is trace class and
\[
 \operatorname{Tr}U(f)=\mathcal W(f).                            \tag{T.7}
\]
The construction (T.6) uses the zeros as input. It is not
an identification with the original adelic cutoff spaces,
nor does (T.5) automatically retain a multiplicity
\(m_{\chi,\gamma}\)-dimensional eigenspace.
\end{theorem}

\begin{proof}
Equation (T.2), polarization, and Cauchy--Schwarz for a
nonnegative sesquilinear form show that quotienting by
its nullspace gives an inner product. The sum-of-squares
formula identifies that nullspace exactly as in (T.5).
Translation multiplies each transform by a unit-modulus
number, so it preserves the inner product and extends
unitarily to the completion. Strong continuity follows
first for a fixed \(f\) by dominated convergence in the
absolutely convergent square sum and then by density.

For (T.6), the direct sum is a Hilbert space. Multiplication
by a real sequence on its maximal square-summability
domain is self-adjoint: for example the diagonal inverses
\((H\pm i)^{-1}\) have coefficients \((\gamma\pm i)^{-1}\),
are bounded, and have ranges in the stated domain.
The action is unitary; dominated convergence in the
square-summable coordinates proves strong continuity.
Integrated against a test \(f\), it is diagonal with
eigenvalue \(\widehat f(\chi,i\gamma)\), repeated
\(m_{\chi,\gamma}\) times. There are only finitely many
character sectors, and rapid vertical decay together
with the usual \(O_\chi(T\log(T+2))\) zero count makes
\(\sum m_{\chi,\gamma}|\widehat f(\chi,i\gamma)|\) finite.
Equivalently, this is the absolute spectral convergence
in C99, Corollary III.2 and CCM07, Theorem 4.16.
This proves trace class and (T.7).

Finally, (T.5) samples a test at a distinct spectral point
only once, with weight \(m_{\chi,\gamma}\). A weight in
a scalar spectral measure does not produce \(m\)
independent vectors at that point. This explains why
(T.5) and the multiplicity-preserving model (T.6) have
different purposes when a zero is multiple.
\end{proof}

\subsection{A fully explicit theorem on every finite window}

Fix a primitive Dirichlet character \(\chi\) of conductor \(q\),
let \(\epsilon\in\{0,1\}\) satisfy \(\chi(-1)=(-1)^\epsilon\),
and write \(\psi=\Gamma'/\Gamma\).
For \(\lambda>1\) set \(L=\log\lambda\), \(I=[-L,L]\).
For \(v\in C_c^\infty((-L,L))\) put
\[
 g(e^t)=v(t),\qquad
 f=g*g^*,\qquad
 F_+v(\tau)=\int_I v(t)e^{i\tau t}\,dt,\qquad
 \ell_\pm(v)=\int_I v(t)e^{\pm t/2}\,dt .
\]
We use the \(L^2\) inner product linear in the second variable.
Let \(T_a\) be compressed translation on \(L^2(I)\):
\[
 (T_a v)(t)=\widetilde v(t-a)\quad(t\in I),
\]
where \(\widetilde v\) is extended by zero outside \(I\).
Thus \(T_a^*=T_{-a}\), and \(T_a=0\) for \(|a|\geq2L\)
(up to the immaterial endpoints). Define the bounded
operator given by a finite sum of translations
\[
 P_{\chi,\lambda}=
 \sum_{\substack{p\ {\rm prime},\,k\geq1\\p^k\leq\lambda^2}}
 \frac{\log p}{p^{k/2}}
 \left(\chi(p)^k T_{-k\log p}
       +\overline{\chi(p)^k}T_{k\log p}\right).                 \tag{T.8}
\]
Each translation on \(L^2(I)\) need not have finite rank.
As usual, \(\chi(p)=0\) at ramified primes.
Let \(\delta_\chi=1\) for the trivial primitive character
(\(q=1\)) and \(0\) otherwise, and set
\[
 a_\chi(\tau)=
 \log(q/\pi)+
 \operatorname{Re}\psi\!\left(\frac{1/2+\epsilon+i\tau}{2}\right).
                                                               \tag{T.9}
\]

\begin{theorem}[AC-T05: the exact semilocal window identity]
\label{AC-T05}
Assuming A, for every \(\lambda>1\) and every such \(v\),
\[
\begin{split}
 q_{\chi,\lambda}(v)
 &:=
 \int_{\mathbb R}a_\chi(\tau)|F_+v(\tau)|^2\,
                    \frac{d\tau}{2\pi}
 -\langle v,P_{\chi,\lambda}v\rangle\\
 &\hspace{6mm}
 +2\delta_\chi\operatorname{Re}
            \bigl(\ell_+(v)\overline{\ell_-(v)}\bigr)\\
 &=\sum_{\rho=1/2+i\gamma}
               m_{\chi,\gamma}|F_+v(\gamma)|^2\ \geq0 .
\end{split}\tag{T.10}
\]
Only the prime powers \(p^k\leq\lambda^2\) occur.
For \(\chi=1\), imposing either \(\ell_+(v)=0\) or
\(\ell_-(v)=0\) removes the two polar terms. For a nontrivial
primitive character there are no polar terms to remove.
Positivity of all these forms, for all \(\lambda\) and all
primitive \(\chi\), is equivalent to A.
\end{theorem}

\begin{proof}
The completed Dirichlet function has archimedean factor
\[
 (q/\pi)^{(s+\epsilon)/2}
       \Gamma((s+\epsilon)/2).
\]
Twice the real part of its logarithmic derivative on
\(s=1/2+i\tau\) is (T.9). Applying the explicit formula
C99, Appendix II, Theorem 6 in the \(\chi\)-sector
therefore gives, for an arbitrary smooth compactly
supported \(f\) in the centered normalization,
\[
\begin{split}
 \sum_\rho m_\rho\int_0^\infty
             f(x)x^{\rho-1/2}\,\frac{dx}{x}
={}&
 \delta_\chi\left(\int f(x)x^{1/2}\frac{dx}{x}
             +\int f(x)x^{-1/2}\frac{dx}{x}\right)\\
 &+\int_{\mathbb R}a_\chi(\tau)
           \left(\int f(x)x^{i\tau}\frac{dx}{x}\right)
                   \frac{d\tau}{2\pi}\\
 &-\sum_{p,k\geq1}\frac{\log p}{p^{k/2}}
           \bigl(\chi(p)^k f(p^k)
             +\overline{\chi(p)^k}f(p^{-k})\bigr).
\end{split}\tag{T.11}
\]
For clarity, the conductor term comes from
\(2\,d\log(q^{s/2})/ds=\log q\); the real gamma term comes
from pairing the two functional-equation contours; and
the prime terms follow from
\(-L'/L(s,\chi)=\sum_{p,k}(\log p)\chi(p)^k p^{-ks}\)
on \(\operatorname{Re}s>1\).
These are the local components in the cited general
explicit formula. The \(\chi=1\) version, with the same
signs, is also printed in CC20, Appendix B, equations
\texttt{bombieriexplicit}--\texttt{burnolexplicit1}.
The gamma factor has no zeros; its trivial-zero
cancellations are already incorporated in (T.11).

Now \(f(e^a)=\int v(t)\overline{v(t-a)}dt\), so
\(\operatorname{supp}f\subset[\lambda^{-2},\lambda^2]\),
and
\[
 f(e^a)=\langle v,T_{-a}v\rangle,\qquad
 \int f(x)x^{i\tau}\frac{dx}{x}=|F_+v(\tau)|^2 .
\]
Consequently the prime sum is exactly (T.8); terms outside
the stated range vanish. The two polar transforms sum
to \(2\operatorname{Re}(\ell_+\overline{\ell_-})\).
Under A, \(\rho-1/2=i\gamma\) with \(\gamma\) real.
Substitution in (T.11) proves the identity and inequality.

For the converse, every compact logarithmic support is
contained in some interval \([-L,L]\). The hypotheses
therefore give the centered Weil positivity in each
character sector. The sectors are orthogonal under
convolution and finite sums give all tests in
\(\mathcal D(C_{\mathbb Q})\). AC-T02 and Weil's
criterion imply A. The endpoint constraints in the
trivial sector can also be imposed without losing the
criterion: CC20, Proposition 46
(label \texttt{mainprop}), proves exactly this
finite-codimension version.
\end{proof}

\paragraph{Prior result and the extension made here.}
For the trivial character, the closed Weil form and
compact-resolvent operator are already established in
Connes--Consani--Moscovici, \emph{Zeta spectral triples}
(arXiv:2511.22755), \S3, Proposition 3.3
(label \texttt{Hilbert}), Proposition 3.5
(label \texttt{compres}), and Theorem 3.6
(label \texttt{thmsmallest}). The following proof
records the required form for every primitive
Dirichlet character, including odd and complex
characters, and the consequence of A.

\begin{theorem}[AC-T06: closed Dirichlet window operators]
\label{AC-T06}
The forms (T.10) have a canonical closed realization on
\(L^2(I)\). Under A their associated self-adjoint
operators \(A_{\chi,\lambda}\) are nonnegative and have
compact resolvent. In particular their spectra consist
of nonnegative eigenvalues of finite multiplicity tending
to \(+\infty\). Their kernels are precisely
\[
 \ker A_{\chi,\lambda}
 =\{v\in\mathcal D_{\chi,\lambda}:
                 F_+v(\gamma)=0
          \text{ for all zeros }1/2+i\gamma
          \text{ of }L(s,\chi)\},                             \tag{T.12}
\]
where the form domain \(\mathcal D_{\chi,\lambda}\) is
specified in the proof. Every finite matrix obtained by
polarizing (T.10) on test vectors is positive semidefinite.
The same statements hold for the form restricted to
either closed polar-annihilating hyperplane
\(\ker\ell_+\) or \(\ker\ell_-\).
\end{theorem}

\begin{proof}
The digamma function has no poles on the vertical lines
in (T.9). Stirling's formula gives
\[
 a_\chi(\tau)=\log|\tau|+\log(q/(2\pi))+O_\chi(|\tau|^{-2})
       \quad(|\tau|\longrightarrow\infty).
\]
It is continuous, bounded below, and tends to infinity.
Choose \(c_\chi\) such that \(b_\chi=a_\chi+c_\chi\geq1\).
This weight is comparable above and below to
\(1+\log(2+|\tau|)\). Define
\[
 \mathcal D_{\chi,\lambda}
 =\left\{v\in L^2(I):
  \int b_\chi(\tau)|F_+v(\tau)|^2\,\frac{d\tau}{2\pi}
                      <\infty\right\}.                       \tag{T.13}
\]
The positive weighted Fourier form is closed: extension
by zero identifies \(L^2(I)\) with a closed subspace of
\(L^2(\mathbb R)\), and multiplication by
\(b_\chi^{1/2}\) is a closed operator in Fourier space.

Smooth tests in the open interval are a form core.
Indeed first replace \(v(t)\) by
\(v_r(t)=r^{-1/2}v(t/r)\) with \(r<1\) tending to \(1\).
Dilation is strongly continuous in the weighted Fourier
norm. This follows by approximation by compactly supported
Fourier functions and the uniform bound on
\(b_\chi(\tau/r)/b_\chi(\tau)\) for \(1/2\leq r\leq1\).
The support of \(v_r\) lies in the interior of \(I\).
Convolving \(v_r\) with a smooth compactly supported
approximate identity of sufficiently small radius
preserves support in \(I\). In Fourier space it multiplies
by a uniformly bounded function tending pointwise to
\(1\); dominated convergence gives form-norm convergence.
This proves the core assertion.

The operator (T.8) is bounded and self-adjoint, with
\(\|P_{\chi,\lambda}\|\leq
2\sum_{p^k\leq\lambda^2}(\log p)p^{-k/2}\).
The polar term is a bounded Hermitian form because
\(\ell_\pm\) are bounded functionals on \(L^2(I)\).
Hence adding these bounded terms to the closed weighted
Fourier form gives a closed semibounded form on (T.13).
On the core it is (T.10), so A makes it nonnegative
on its whole domain. The representation theorem for
closed nonnegative forms defines
\(A_{\chi,\lambda}\) uniquely.

Here is the compactness argument. A form-norm bounded
set has uniformly small Fourier tails:
\[
 \int_{|\tau|>R}|F_+v(\tau)|^2\frac{d\tau}{2\pi}
 \leq
 \frac{C}{\inf_{|\tau|>R}b_\chi(\tau)}
 \longrightarrow0.
\]
The Fourier transform from \(L^2(I)\) to \(L^2([-R,R])\)
has square-integrable kernel \(e^{i\tau t}\), and is
therefore compact. Low-frequency compactness and the
uniform tail bound show that the embedding of the form
domain into \(L^2(I)\) is compact. The form resolvent
maps bounded sets of \(L^2(I)\) to form-norm bounded
sets, proving compact resolvent.

To extend the sum-of-squares formula, approximate
\(v\in\mathcal D_{\chi,\lambda}\) by core vectors \(v_n\).
The sequences
\((\sqrt{m_{\chi,\gamma}}F_+v_n(\gamma))_\gamma\)
are Cauchy in \(\ell^2\), because their squared distance
is \(q_{\chi,\lambda}(v_n-v_m)\). Each coordinate
evaluation is continuous on \(L^2(I)\), by
Cauchy--Schwarz. The limit is therefore the corresponding
sequence for \(v\), and (T.10) extends to the form domain.
A vector belongs to the kernel of a nonnegative operator
if and only if its form value is zero. This gives (T.12).

Restriction to either hyperplane is again a closed form.
Core density in that hyperplane follows by taking smooth
approximants and subtracting their small moment error
times one fixed smooth test whose corresponding moment
is nonzero. The compactness argument persists under
restriction. Finally a polarized Gram matrix is
positive semidefinite because its quadratic form is
\(q_{\chi,\lambda}(\sum_j c_jv_j)\geq0\).
\end{proof}

These constructed window operators have the Weil form
as their quadratic form. Their spectral conclusion, strengthened to strict positivity
in AC-T13 below, does not identify their eigenvalues with
the zero ordinates. No assertion
of a uniform spectral gap as \(\lambda\to\infty\), nor
a convergence rate for any different spectral-triple
approximation, is included in AC-T06.

\subsection{What the semilocal compressions satisfy}

Let \(\vartheta(r)\xi(x)=r^{-1/2}\xi(r^{-1}x)\) on
\(L^2(\mathbb R)_{\rm ev}\), and let \(\mathbf S\) be
the Sonin projection onto the even functions which,
together with their additive Fourier transforms,
vanish on \([-1,1]\). CC20, Theorem 27
(label \texttt{devil}), is the unconditional identity
\[
 \operatorname{Tr}(\vartheta(f)\mathbf S)
      =W_\infty(f)+E(f),\qquad
 E(f)=\int_0^\infty f(r^{-1})\varepsilon(r)\frac{dr}{r}.          \tag{T.14}
\]
Here \(W_\infty=-W_{\mathbb R}\) in that paper's unitary
normalization, and its explicit convergent prolate
expression is
\[
 \varepsilon(r)=
 \sum_n\frac{\lambda(n)}{\sqrt{1-\lambda(n)^2}}
       \langle \xi_n,\vartheta(r^{-1})\zeta_n\rangle
 \quad(r\geq1),\qquad \varepsilon(r^{-1})=\varepsilon(r).
                                                               \tag{T.15}
\]
All vectors and eigenvalues in this formula are defined
in CC20, \S4, Proposition 25 (label \texttt{devil0});
the convergence treatment is Appendix F, including
Lemma 49 (label \texttt{lemesti}).

\begin{theorem}[AC-T07: corrected Sonin and semilocal trace inequalities]
\label{AC-T07}
Let \(g\in C_c^\infty([\lambda^{-1},\lambda])\),
\(f=g*g^*\), and assume either
\(\int g(r)r^{1/2}d^\times r=0\) or
\(\int g(r)r^{-1/2}d^\times r=0\).
Assuming A,
\[
\begin{split}
 0\leq \sum_{\rho=1/2+i\gamma}m_\rho
                  |\widehat g_+(\gamma)|^2
 ={}&\operatorname{Tr}
          (\vartheta(g)\mathbf S\vartheta(g)^*)-E(f)\\
 &-\sum_{p^k\leq\lambda^2}
      \frac{\log p}{p^{k/2}}
             \bigl(f(p^k)+f(p^{-k})\bigr),
                                                               \\
 \widehat g_+(\tau)&=\int g(r)r^{i\tau}d^\times r .
\end{split}\tag{T.16}
\]
Thus the Sonin trace dominates the explicit sum of
the prolate remainder and these finitely many prime
contributions, for every window.

More generally, let \(S\) contain infinity, the
ramified primes of a character sector, and the primes
occurring in its test. Put \(v(t)=g(e^t)\) and
\(f_\chi(u,r)=\chi(u)f(r)\), viewed on \(C_S\)
after quotienting by the outside-\(S\) compact units.
In the unitary normalization of
C99, Theorem VII.4, the finite part of the semilocal
cutoff trace in that sector is
\[
 \operatorname{FP}_{\Lambda\to\infty}
 \operatorname{Tr}(R_\Lambda\vartheta_S(f_\chi))
     =B_\chi(f)-q_{\chi,\lambda}(v),                           \tag{T.17}
\]
where the subtracted divergent term is
\(2f(1)\log'\Lambda\), and \(B_\chi(f)\) is the polar
term in (T.10). It is therefore at most \(B_\chi(f)\);
on the polar-annihilating subspace it is nonpositive.
This concerns the finite part, not the sign of each
finite-\(\Lambda\) cutoff trace.
\end{theorem}

\begin{proof}
For the trivial character, (T.11), with the polar
terms removed, gives
\(q=W_\infty(f)-\sum_p W_p(f)\), where
\[
 W_p(f)=(\log p)\sum_{k\geq1}p^{-k/2}
                  (f(p^k)+f(p^{-k})).
\]
Insert (T.14). For \(f=g*g^*\), the trace in (T.14)
equals \(\operatorname{Tr}
(\vartheta(g)\mathbf S\vartheta(g)^*)\) by the
trace-class and cyclicity assertions proved there.
This proves the equality in (T.16); A supplies
nonnegativity through (T.10).

For (T.17), write the original nonunitary scaling
action as \(U(a)=|a|^{1/2}\vartheta_S(a)\).
In C99, Theorem VII.4, substitute the test
\(h(a)=|a|^{-1/2}f_\chi(a)\), after averaging the
irrelevant compact-unit coordinates. Then
\(U(h)=\vartheta_S(f_\chi)\) and \(h(1)=f(1)\).
The finite part is exactly the sum of local
geometric terms. Since the chosen \(S\) contains
all nonzero terms for this test, the explicit
formula (T.11) identifies this local sum with
\(B_\chi(f)-q_{\chi,\lambda}(v)\). This proves
(T.17) and its inequalities. The theorem uses
\(R_\Lambda=\widehat P_\Lambda P_\Lambda\), which
need not be an orthogonal projection, so the
argument makes no pointwise-in-\(\Lambda\)
positivity claim.
\end{proof}

CC20's original small-window theorem is stronger than
the bare inequality \(W_\infty(g*g^*)\geq0\) in a
different direction. Theorem 1 assumes support in
\([2^{-1/2},2^{1/2}]\) and two transform vanishings.
Theorem 44 (label \texttt{mainthmfine}) proves
\[
 W_\infty(g*g^*)\geq
 \operatorname{Tr}(\vartheta(g)\mathbf S\vartheta(g)^*)
                   -c|\widehat g(0)|^2,\qquad 13<c<17,
                                                               \tag{T.18}
\]
under its specified one endpoint vanishing. These
were already source theorems. Their quantitative
constant depends on the prolate estimates and
numerical ingredients in CC20; this run does not
certify those numerical computations anew.
A proves (T.16) on all windows, but does not make
the explicit \(E\) and prime terms zero. The quantitative bound in (T.18) still uses the
cited prolate estimates. CC20, Remark 45
(label \texttt{remlowc}), gives the source
obstruction to removing the positive constant
under those hypotheses.

\begin{proposition}[AC-T08: compression positivity does not give a restricted flow]
\label{AC-T08}
For any unitary representation \(U\) and orthogonal
projection \(P\), the compression
\(\Phi(f)=PU(f)P|_{P\mathcal H}\) satisfies
\[
 \Phi(g*g^*)-\Phi(g)\Phi(g)^*
   =PU(g)(1-P)U(g)^*P\geq0.                                  \tag{T.19}
\]
The kernel \((x,y)\mapsto PU(x^{-1}y)P\) is positive
definite. These facts are unconditional. The
compressed family \(PU(x)P\) need not be a
representation, so positivity alone does not
produce the spectral theorem for a restricted
scaling flow.
\end{proposition}

\begin{proof}
Use \(U(g*g^*)=U(g)U(g)^*\), insert \(1=P+(1-P)\),
and subtract. For the kernel, vectors
\(\xi_j\in P\mathcal H\) give
\[
 \sum_{i,j}\langle\xi_i,PU(x_i^{-1}x_j)P\xi_j\rangle
  =\left\|\sum_j U(x_j)\xi_j\right\|^2\geq0.
\]
For a counterexample to the representation assertion,
take \(U(t)=\operatorname{diag}(e^{it},e^{-it})\)
and project onto \((1,1)/\sqrt2\). The compression
is the scalar \(\cos t\), and
\(\cos(t+s)\ne\cos t\cos s\) in general.
This verifies the failure even in a finite-dimensional
setting with no domain issues. For a compression of
an unbounded generator there are, in addition,
domain and self-adjointness questions; (T.19)
does not answer them.
\end{proof}

\begin{proposition}[AC-T09: a compactly supported counterexample to global archimedean positivity]
\label{AC-T09}
The claim that the archimedean term \(W_\infty\)
alone is nonnegative on every window, even after
both polar moments vanish, is false. This failure
is compatible with A.
\end{proposition}

\begin{proof}
In logarithmic coordinates the archimedean form is
the Fourier multiplier
\[
 a_1(\tau)=\operatorname{Re}\psi(1/4+i\tau/2)-\log\pi.
\]
The reflection and duplication formulas for
\(\Gamma\) give
\[
 a_1(0)=-\gamma_{\rm E}-\frac\pi2-3\log2-\log\pi<0.
\]
For example this follows from
\(\psi(1/2)=-\gamma_{\rm E}-2\log2\),
\(\psi(3/4)-\psi(1/4)=\pi\), and
\(\psi(1/4)+\psi(3/4)=2\psi(1/2)-2\log2\).

Choose a nonzero real \(\phi\in C_c^\infty((-1,1))\),
put \(\phi_T(t)=T^{-1/2}\phi(t/T)\), and
\[
 v_T=(\partial_t^2-1/4)\phi_T,\qquad g_T(e^t)=v_T(t).
\]
Integration by parts proves
\(\ell_+(v_T)=\ell_-(v_T)=0\).
The support of \(g_T\) lies in \([e^{-T},e^T]\).
Its archimedean form is
\[
 \int_{\mathbb R} a_1(\eta/T)
    ((\eta/T)^2+1/4)^2
       |F_+\phi(\eta)|^2\frac{d\eta}{2\pi}.
\]
As \(T\to\infty\) this tends to
\(a_1(0)\|\phi\|_2^2/16<0\).
Dominated convergence is valid since \(a_1(\tau)\)
has logarithmic growth and \(F_+\phi\) decays
faster than any power. Thus a sufficiently large
finite \(T\) gives a smooth compactly supported
counterexample. The full form in (T.10) contains
prime terms as well. Consequently there is no
contradiction with A.
\end{proof}

\subsection{The 1999 global cutoff: source scope and a conductor-level reduction}

Theorem VIII.5 of C99 explicitly assumes positive
characteristic. Its implication from the cutoff
trace to positivity uses the exact inclusion
\(E(B_{\Lambda,0})\subset S_\Lambda\), C99 VIII.(23),
and its reverse implication uses Lemma VIII.3.
At the end of \S VIII, equations (30)--(33), C99
discusses the number-field case, restricts to
\(k=\mathbb Q\), uses even real prolate functions,
and states that this gives the analogues of
Lemma 1, Theorem 5, and Lemma 3 for the Riemann
zeta function. Thus, \emph{if that stated
number-field analogue is taken as a cited
theorem}, A gives its zeta cutoff formula:
\[
 \operatorname{Tr}(Q_\Lambda U(h))
 =2h(1)\log'\Lambda+
   \sum_v\int_{\mathbb Q_v^\times}'
        \frac{h(u^{-1})}{|1-u|}\,d^\times u+o(1).               \tag{T.20}
\]
This is a consequence of the asserted analogue,
with the prolate cutoff meant there. It is not a
derivation of a full all-character number-field
cutoff theorem from the positive-characteristic
Theorem VIII.5. The introduction's all-character
formulation and the nuclear formula of CCM07
do not supply the missing cutoff estimate.

\begin{proposition}[AC-T10: why the literal function-field cutoff cannot be copied]
\label{AC-T10}
A nonzero function on \(\mathbb R\) cannot have
both itself and its additive Fourier transform
compactly supported. Hence the number-field
extension must use an approximate/prolate
replacement for the literal simultaneous-support
cutoff in Theorem VIII.5.
\end{proposition}

\begin{proof}
If \(f\in L^2(\mathbb R)\) has compact support,
then it is in \(L^1\), and
\[
 \widehat f(z)=\int f(x)e^{-2\pi i xz}\,dx
\]
extends to an entire function by differentiation
under the compactly supported integral.
If its restriction to the real axis also has
compact support, continuity makes it vanish
on a real open interval. The identity theorem
makes it identically zero, and injectivity
of Fourier transformation gives \(f=0\).
This obstruction is explicitly noted in
C99, \S VIII just before equation (31).
\end{proof}

\begin{proposition}[AC-T11: finite conductor, Gauss phase, and parity]
\label{AC-T11}
Let \(\chi\) be a primitive nontrivial Dirichlet
character of conductor \(q\), and extend it
by zero to nonunits modulo \(q\). On the finite
adeles put
\[
 \varphi_\chi(x)=
   \sum_{a\bmod q}\chi(a)\,
                     1_{a+q\widehat{\mathbb Z}}(x).
                                                               \tag{T.21}
\]
Choose the standard finite additive character
so that for \(y=b/q+\widehat{\mathbb Z}\) its
value on the residue \(a\) is \(e^{2\pi iab/q}\).
With \(\operatorname{vol}(\widehat{\mathbb Z})=1\),
\[
 \widehat\varphi_\chi(y)=
 \begin{cases}
 q^{-1}\tau(\chi)\overline{\chi(b)},&
           y=b/q+\widehat{\mathbb Z},\\
 0,&y\notin q^{-1}\widehat{\mathbb Z},
 \end{cases}
 \qquad
 \tau(\chi)=\sum_{a\bmod q}\chi(a)e^{2\pi ia/q}.                 \tag{T.22}
\]
In particular it vanishes unless \(b\) is a unit
modulo \(q\), and \(|\tau(\chi)|=\sqrt q\).
For an archimedean vector of parity
\(\chi(-1)=(-1)^\epsilon\), simultaneous
\emph{nominal} adelic norm cutoffs \(\Lambda\)
are therefore satisfied by the uniform archimedean rectangle
\[
 |x|\leq\Lambda,\qquad
 |\xi|\leq\Lambda/q.                                         \tag{T.23}
\]
Unitary dilation transforms this to the symmetric
rectangle of side parameter \(\Lambda/\sqrt q\).
This rectangle is a sufficient cutoff for the chosen finite tensor
factor.  It is not an equivalence with an arbitrary adelic norm
cutoff: on fibers where the finite norm is smaller, the latter
permits a larger real coordinate.  The comparison with the full
global cutoff remains part of T-CUT below.
Even and odd prolate sectors handle
\(\epsilon=0,1\), respectively. The scalar
Gauss phase disappears from the geometry of the
two orthogonal projections.
\end{proposition}

\begin{proof}
The Fourier transform of \(1_{a+q\widehat{\mathbb Z}}\)
is \(q^{-1}\) times the indicated additive
character on \(q^{-1}\widehat{\mathbb Z}\),
and zero outside that annihilator. This gives
the finite sum in (T.22).
For a unit \(b\), changing variables \(a\mapsto
b^{-1}a\) turns the sum into
\(\overline{\chi(b)}\tau(\chi)\).
For nonunit \(b\), choose a prime divisor \(p\)
of \((b,q)\). Primitivity implies that the
character is nontrivial on the kernel of the
reduction from the units modulo \(q\) to the
units modulo \(q/p\). Choose \(u\) in that
kernel with \(\chi(u)\ne1\). The additive
character in the sum is unchanged by
\(a\mapsto ua\), since \(p\mid b\);
the multiplicative character is multiplied by
\(\chi(u)\). Thus the sum is zero.
Finite Fourier Plancherel now gives
\[
 q\sum_{a\bmod q}|\chi(a)|^2
   =\sum_{b\bmod q}
        |\tau(\chi)\overline{\chi(b)}|^2,
\]
so \(q\varphi(q)=|\tau(\chi)|^2\varphi(q)\).

The support of (T.21) has finite norm at most
\(1\), while the support of (T.22) has finite
norm at most \(q\); at every prime dividing
\(q\) the latter valuation is exactly the
negative conductor exponent. Combining with
the real factor gives (T.23). If
\(D_r h(x)=r^{-1/2}h(x/r)\), then \(D_r\)
multiplies the time support by \(r\) and
divides the Fourier support by \(r\). Take
\(r=\sqrt q\) and side
\(\Lambda/\sqrt q\).
Fourier transformation preserves parity.
Multiplication of a Fourier transform by a
unit scalar has no effect on its support
projection. These facts prove the last
assertions.
\end{proof}

There is an explicit finite-dimensional version
of the remaining analytic problem. For a finite
span of compactly supported prolate vectors
\(\psi_1,\ldots,\psi_N\) of the required parity
in (T.23), define
\[
 e_j(x)=x^{1/2}\sum_{n\geq1}\chi(n)\psi_j(nx)
       \quad(x>0),\quad
 G_{ij}=\langle e_i,e_j\rangle_{L^2(dx/x)},\quad
 M(h)_{ij}=\langle e_i,V(h)e_j\rangle .                         \tag{T.24}
\]
For fixed nontrivial \(\chi\), its periodic partial
sums are bounded. Summation by parts therefore
bounds the sum defining \(e_j(x)/\sqrt x\)
uniformly as \(x\downarrow0\), for the bounded-
variation prolate cutoff vectors. It vanishes
for \(x>\Lambda\). Thus \(e_j\in L^2(dx/x)\).
Their Mellin transforms in a right half-plane
are \(L(s,\chi)\) times the Mellin transforms
of the \(\psi_j\); uniqueness proves linear
independence. Hence \(G\) is positive definite.
The orthogonal projection \(Q_{\chi,\Lambda}\)
onto their span is concretely
\(E G^{-1}E^*\), where \(Ec=\sum_jc_je_j\), and
\[
 \operatorname{Tr}(Q_{\chi,\Lambda}V(h))
           =\operatorname{tr}(G^{-1}M(h)).                    \tag{T.25}
\]
The expected prolate threshold is the one in
C99, VIII.(32), rescaled as in (T.23):
indices of the required parity up to
\(4\Lambda^2/q\). For the trivial character
one must also account for the two endpoint
conditions and their polar correction.

\paragraph{Exact unproved gate T-CUT.}
For the prolate image spaces just described
(and a cutoff convention whose threshold errors
are controlled), prove
\[
 2h(1)\log\Lambda
 -\operatorname{tr}(G_{\chi,\Lambda}^{-1}M_{\chi,\Lambda}(h))
 \longrightarrow
 \sum_\gamma m_{\chi,\gamma}\widehat h_+(\gamma)               \tag{T.26}
\]
for every smooth compactly supported \(h\),
with the polar correction in the trivial
sector. Before A, the proposed limit is the
sum of the harmonic measures in C99,
Lemma VIII.3; A collapses these measures to
Dirac masses if this limit theorem has been
established. The finite Fourier calculation
AC-T11 settles the conductor and parity
part. It does not estimate the inverse Gram
matrices or the tails of the \(E_\chi\)-images.

Two attempted shortcuts fail at identifiable
steps. First, exact simultaneous support
fails by AC-T10. Second, a prolate vector
has a nonzero Fourier tail, so its \(E_\chi\)
image need not vanish below \(\Lambda^{-1}\).
The exact inclusion \(Q'_{\Lambda,0}\leq
S_\Lambda\) in C99, VIII.(23) is therefore
not available without an error estimate.
Orthogonality of prolate vectors for the
real \(L^2\) norm does not give orthogonality
for the \(E_\chi\) norm: the matrix \(G\)
is essential. Expanding its double arithmetic
sum and interchanging integrals term by
term requires convergence not supplied by
bounded character partial sums. Thus
(T.26) is not proved in this work. The
claim that the supplied positive-characteristic
theorem alone settles every number-field
cutoff has been corrected to this precise gate.

\subsection{Geometric language and the statements that remain}

For a finite-order \(\chi\), the arithmetic
object is a rank-one finite-monodromy
\(\ell\)-adic local system on
\(\operatorname{Spec}\mathbb Z[1/q\ell]\),
or its Artin motive with coefficients
containing the character values. Its
unramified Frobenius eigenvalues already
have complex absolute value \(1\).
A supplies an analytic statement about the
completed \(L\)-function; it does not newly
prove that elementary weight-zero fact.
Equations (T.2) and (T.10) are positive
pairings of test functions attached to
these \(L\)-functions.

CCM07, Definition 4.10 constructs a cokernel
in a category of cyclic modules and
equation \texttt{TorH1mot} applies a derived
functor to it. This is the precise
categorical location of the source's
cohomological terminology. The supplied
construction does not identify it with
finite-dimensional \(\ell\)-adic
cohomology of a proper scheme, or with
a specified pure motive. AC-T04 adds
a positive Hilbert completion and an
explicit spectral trace model. Neither
construction supplies an algebraic
square, a \(K\)-theory comparison map,
or algebraic correspondences realizing
the pairing.

For the intersection notation of CCM07,
Proposition 7.2 (label \texttt{PosTrRHint}),
A gives the stated numerical inequality
in the real test-function convention.
For complex tests, the endpoint term
written there as \(2\widehat f(0)
\widehat f(1)\) must be replaced by
\[
 2\operatorname{Re}
      \bigl(\widehat f(0)\overline{\widehat f(1)}\bigr).
\]
This follows directly by transforming
\(f*f^\sharp\), as in (T.10).
It is a Hermitian correction, not a
new geometric intersection construction.

The exact remaining questions in this
part are: the full all-character
number-field cutoff limit T-CUT; an
identification of its prolate image
spaces with the positive trace-pairing
completion; and any proposed semilocal
Sonin domination stronger than (T.16),
which requires a separate estimate for
the explicit remainder. The local
trace formula C99, Theorem V.3, the
semilocal trace formula Theorem VII.4,
the CCM07 nuclear trace formula, and
CC20's already proved small-window
theorems are not newly proved by A.

\subsection{The inherited slow-band trace theorem and a sharper window tool}

\begin{theorem}[AC-T12: A closes the MAS\_ARTIN slow-band arithmetic gate over \(\mathbb Q\)]
\label{AC-T12}
Fix \(0<\alpha<1/2\) and \(s\geq0\). Use the
actual source operators of MAS\_ARTIN checkpoint 10,
\texttt{modules/full\_test\_slow\_band.tex},
MA-FT-001 and MA-FT-002, in
\[
 H_{\alpha,s}
 =L^2\!\left(\mathbb R,
         e^{2\alpha|x|}(1+x^2)^{s/2}\,dx\right).
\]
In the rational trivial sector and every primitive
nonprincipal Dirichlet sector, A gives, for every
\(g\in C_c^\infty(\mathbb R)\) with \(g(0)=0\),
\[
 L_{\alpha,s}(g)=\sum_\rho m_\rho
        \int_{\mathbb R}g(t)e^{i\gamma_\rho t}\,dt,
 \qquad \Delta_{\alpha,s}(g)=0.                               \tag{T.27}
\]
In particular the paired actual local and global
annular-cutoff limits of MA-FT.(7) agree whenever
\(X,Y\to\infty\) with \(Y/X\to0\), and with the
fixed common annular profile specified there.
The source-thin version additionally retains
its condition \(\delta X\to0\).
This is true at every positive strip width,
with full multiplicities and for every \(s\geq0\).
\end{theorem}

\begin{proof}
The source defines \(d=1/2-\alpha\), \(u=1/2+\alpha\)
and retained multiplicities \(r_\rho^+\) and
\(r_\rho^{\rm st}\), with possible deficiencies
\(\delta_\rho^\pm\), in MA-FT.(2).
Under A every zero has real part \(1/2\), strictly
between \(d\) and \(u\). Hence
\(r_\rho^+=r_\rho^{\rm st}=m_\rho\) and
\(\delta_\rho^-=\delta_\rho^+=0\); no zero lies
on a strip boundary, so the boundary jet
parameter \(b_s\) has no effect.
Substitute these values in MA-FT.(5)--(7).
All deficiency sums disappear, giving (T.27)
and the stated actual operator-limit
identities. The support, profile, and
path hypotheses remain precisely those
of MA-FT-001, rather than a new order
of limits.
\end{proof}

The proof of MA-FT-001/002 was read in full
for this use. Its ambient trace-class,
near-norm-one, and local-quotient results
are retained as the explicitly named
background dependencies of that proof;
they have not all been re-proved in
this section. The statement that A
annihilates its two deficiency divisors
is checked directly above.

This supplies another actual trace realization.
It leaves a concrete comparison problem before
one obtains C99's prolate cutoff: the source
uses positive \(\alpha\), one-sided
Paley--Wiener theta-image closures and
a slow-band joint limit; C99 uses the
unweighted simultaneous time/frequency
prolate geometry. A comparison must
control the trace after the Hilbert
space identification, including passage
\(\alpha\downarrow0\), and must restore
the identity contact term. Indeed
\(g(0)=0\) excludes every nonzero test
of the form \(g=v*v^*\), since
\((v*v^*)(0)=\|v\|_2^2\).
Equality on this test class alone cannot
determine a distribution's coefficient
of \(\delta_0\). Two distributions
differing by \(c\delta_0\) agree on
every such \(g\). This is the precise
missing normalization, not a claim
that the source theorem is wrong.

\begin{corollary}[AC-T13: strictly positive spectrum in every fixed Dirichlet window]
\label{AC-T13}
Assuming A, \(\ker A_{\chi,\lambda}=0\) for
every primitive Dirichlet character and every
finite \(\lambda>1\). Consequently
\(\min\operatorname{Spec}A_{\chi,\lambda}>0\).
This is a statement for each fixed window;
it supplies no uniform lower bound as the
window grows.
\end{corollary}

\begin{proof}
Under GRH, Carneiro--Chandee--Milinovich,
\emph{A note on the zeros of zeta and
\(L\)-functions}, arXiv:1503.00955v1,
Theorem 5, (4.4), and Example 8, give
\(S(t,\chi)=O_\chi(\log|t|/\log\log|t|)\).
The jump at a zero bounds its multiplicity
by the same order. The
zero-count identity (4.1) in that source gives
\(\asymp_\chi T\log T\) zeros with multiplicity; hence there
are \(\gg_\chi T\log\log T\) distinct
ordinates with \(|\gamma|\leq T\).

If \(v\) belonged to (T.12), its entire
transform \(F_+v(z)\) would vanish at all
these ordinates. Cauchy--Schwarz gives
\(|F_+v(z)|\leq\sqrt{2L}\|v\|_2e^{L|\operatorname{Im}z|}\).
For a nonzero such entire function,
Jensen's formula about any point at
which it does not vanish bounds its
number of zeros in \(|z|\leq T\) by
\(O_v(T)\): apply Jensen at radius
\(2T+O(1)\), and use the displayed
growth bound. This contradicts the
distinct-ordinate count. Thus \(v=0\).
Compact resolvent and nonnegativity
in AC-T06 then imply that the least
eigenvalue is strictly positive.
\end{proof}

For the trivial character this is the
same strictness argument as AC-S03
in the spectral-triple part of this
work; AC-T13 supplies its primitive
Dirichlet extension.

\begin{proposition}[AC-T14: the exact norm of a compressed prime translation]
\label{AC-T14}
On \(L^2([-L,L])\), for \(a>0\) and \(z\in\mathbb C\),
\[
 \|zT_a+\overline zT_{-a}\|=
 \begin{cases}
  2|z|\cos\!\left(\dfrac{\pi}
             {\lceil 2L/a\rceil+1}\right),&0<a<2L,\\
  0,&a\geq2L.
 \end{cases}                                                 \tag{T.28}
\]
In particular if \(L\leq a<2L\), the norm
is \(|z|\), rather than the triangle
bound \(2|z|\). For any finite sum
of prime translations, summing the
individual quantities (T.28) gives
a valid operator-norm bound.
\end{proposition}

\begin{proof}
Write almost every point of the interval
uniquely as
\(t=-L+r+ja\), with \(0\leq r<a\) and
\(j=0,\ldots,N(r)-1\), omitting values
outside the interval. This is a unitary
direct-integral decomposition into
finite coordinate chains: integration
in \(t\) becomes integration in \(r\)
followed by summation in \(j\).
On each chain \(T_a\) is the one-step
shift with zero boundary condition.
If \(0<a<2L\), the essential maximum
chain length is
\(N_{\max}=\lceil2L/a\rceil\).
When \(2L/a\) is an integer, the
additional endpoint occurs only
for a measure-zero set of \(r\);
there is no extra coordinate in
the essential maximum.

A diagonal unitary on a chain removes
the phase of \(z\). The remaining
matrix is \(|z|\) times the adjacency
matrix of a path with \(N(r)\)
vertices. Its eigenvectors have
coordinates
\(\sin(jk\pi/(N(r)+1))\),
\(j=1,\ldots,N(r)\), and eigenvalues
\(2|z|\cos(k\pi/(N(r)+1))\),
\(k=1,\ldots,N(r)\), as direct
substitution in the two-term
recurrence verifies. Its norm is
therefore
\(2|z|\cos(\pi/(N(r)+1))\).
Taking the essential supremum in
the direct integral gives (T.28).
For \(a\geq2L\) both compressed
translations vanish almost
everywhere. The last assertion
is the operator-norm triangle
inequality.
\end{proof}

\section{Spectral triples, zeta cycles, and prolate operators}
\label{sec:ACspectral}

\subsection{Sources and scope}
In this section \(A\) is the input theorem of the question.  Its zeta
component is used; the Dirichlet component is not needed here.  The source
abbreviations are as follows.  \(\mathrm{S21}\) is Connes--Consani,
\emph{Spectral triples and zeta-cycles}, arXiv:2106.01715, attached
\texttt{Spectraltriples.tex}.  \(\mathrm{S22}\) is their
\emph{Hochschild homology, trace map and zeta-cycles}, arXiv:2207.10419,
attached \texttt{Hochschild-zeta.tex}.  \(\mathrm{S23}\) is
Connes--Consani--Moscovici, \emph{Zeta zeros and prolate wave operators},
arXiv:2310.18423, attached \texttt{mainc2m24fine.tex}.
\(\mathrm{S25}\) is their \emph{Zeta spectral triples},
arXiv:2511.22755, attached \texttt{mc2arXiv.tex}.  Source labels below are
literal LaTeX labels, so they remain exact locators if the numbering changes
between versions.  The source numbering used for S21, S22 and S25 is that
of the supplied text.

The following distinctions are essential.  A spectrum contained in the
real line does not by itself give a positive inner product making a
specified operator self-adjoint.  A positive quadratic form does not by
itself identify its eigenvectors.  A counting-function asymptotic does
not identify individual spectral points.  The results below give exact
positive conclusions and identify the additional hypotheses at each of
these three transitions.

\subsection{What becomes complete in the zeta-cycle and sheaf realizations}

\begin{proposition}[AC-S01: all zeros are covered]
\label{AC-S01}
Let
\[
 C_L=\mathbb R_+^*/\exp(L\mathbb Z),\qquad
 \mathcal H_L=\overline{\Sigma_{e^L}\mathcal E(\mathcal S_0^{\rm ev})}^{\,\perp}
 \subset L^2(C_L),\quad L>0.
\]
Under \(A\), the spectrum of the self-adjoint generator of scaling on
\(\mathcal H_L\), with the sign convention in which a character is
\(u^{it}\), is exactly
\[
 \{\gamma\in\mathbb R:\zeta(1/2+i\gamma)=0,
                       \ \gamma L\in2\pi\mathbb Z\}.
\]
Each character in a single circle has multiplicity one.  Every nontrivial
zeta zero occurs for some \(L\), and the set of such lengths is countable
and dense in \(\mathbb R_{>0}\).  The sheaf cohomology realization of
S22, Theorems 1.2 and 5.4, now covers all nontrivial zeros, with their
local jet multiplicities.
The Laplacian realization in S22, Proposition 4.2 and Corollary 4.3
(labels \texttt{laplacian}, \texttt{rhequiv}), has spectrum
\[
 \{-1/4-\gamma^2:\zeta(1/2+i\gamma)=0\},
\]
with the multiplicity conventions of that proposition.
\end{proposition}
\begin{proof}
S21, Theorem 6.4 and Corollary 6.5 (labels \texttt{spectralreal},
\texttt{repetitions}), identify the orthogonal complement by precisely
the critical zeros whose characters are periodic on the indicated circle.
The Fourier modes of \(L^2(C_L)\) are one-dimensional, which proves the
single-circle multiplicity statement.  Theorem \(A\) identifies critical
zeros with all nontrivial zeros.  For any positive zero ordinate
\(\gamma\), choose \(L=2\pi n/\gamma\), \(n\geq1\).
Conversely every nonzero \(\mathcal H_L\) has this form for some
\(\gamma,n\).  This is a countable set.  For any \(L_0>0\), take
zero ordinates \(\gamma_j\to\infty\), and choose positive integers
\(n_j\) nearest to \(L_0\gamma_j/(2\pi)\).  Then
\(2\pi n_j/\gamma_j\to L_0\), proving density.
S22, Theorem 5.4, identifies its cohomology with the Schwartz quotient
by the closed ideal generated by \(t\mapsto\zeta(1/2+it)\); its local
description records the zero jets.  Finally substitute
\(\rho=1/2+i\gamma\) into the eigenvalue
\((\rho-1/2)^2-1/4\) of S22, Proposition 4.2.
\end{proof}

\begin{proposition}[AC-S02: the multiplicity obstruction to unitarization]
\label{AC-S02}
At a zero of order \(m\) and real ordinate \(\gamma\), the local jet
representation in S22 is the module
\(J_\gamma=\mathbb C[t]/(t-\gamma)^m\).  On this module, multiplication
by \(t\) is self-adjoint for some positive definite Hermitian form if
and only if \(m=1\).  Equivalently, the scaling representation on this
finite-dimensional block is unitarizable if and only if \(m=1\).
For \(\gamma\ne0\), the same obstruction applies to the Laplacian
\(-1/4-t^2\).  Thus AC-S01 proves a statement about its spectrum;
unitarizing the given jet realization additionally requires simplicity
of the zeros.  Theorem \(A\) as stated does not supply that hypothesis.
\end{proposition}
\begin{proof}
In the basis \(1,t-\gamma,\ldots,(t-\gamma)^{m-1}\), multiplication
by \(t\) is \(\gamma I+J\), with \(J\) a nilpotent Jordan block of
length \(m\).  If \(m>1\), it is not diagonalizable.  Every
self-adjoint operator in a positive finite-dimensional inner product is
diagonalizable, so no such inner product exists.  A unitary one-parameter
group has a self-adjoint infinitesimal generator, giving the equivalent
formulation.  If \(m=1\), the usual Hermitian form works.  The
nilpotent part of \(-1/4-(\gamma I+J)^2\) is
\(-2\gamma J-J^2=-J(2\gamma I+J)\).  If \(\gamma\ne0\), the
second factor is invertible, so for \(m>1\) this has nonzero nilpotent
part and again is not diagonalizable.  There is no zeta zero at ordinate
zero: \(\zeta(1/2)<0\), for example by its alternating-series
continuation.  The argument is an exact obstruction for multiple zeros,
not a claim that a multiple zero exists.
\end{proof}

\subsection{The closed Weil operators on finite windows}

Write \(L=\log\lambda\), and identify
\(L^2([\lambda^{-1},\lambda],d^*u)\) with
\(H_L=L^2([-L,L],dx)\).  All functions on the interval are extended by
zero.  Use
\(\widehat f(t)=\int_{\mathbb R}f(x)e^{-itx}\,dx\).
Let \(q_L\) be the closed lower semibounded Weil form of S25,
\S3.1, and let \(A_L\) be its associated operator.
S25, Proposition 3.4 (\texttt{Hilbert1}), gives the Fourier-polynomial
form core; Theorem 3.6 (\texttt{thmsmallest}) gives compact resolvent.
Write its eigenvalues with multiplicity as
\(a_1(L)\leq a_2(L)\leq\cdots\).

\begin{theorem}[AC-S03: positive gap and exact Galerkin limits]
\label{AC-S03}
Under \(A\), for every finite \(L>0\),
\[
 q_L(f)=\sum_{\gamma}m_\gamma|\widehat f(\gamma)|^2
 \quad(f\in\operatorname{Dom}q_L),\qquad
 0<a_1(L)\leq a_2(L)\leq\cdots\longrightarrow\infty.
\]
Here the sum is over distinct ordinates of all nontrivial zeta zeros.
For the nested Fourier-polynomial spaces \(E_N\) of S25, let
\(a_{r,N}(L)\) denote the \(r\)-th eigenvalue of the restricted form.
For every fixed \(r\), once \(\dim E_N\geq r\),
\[
 a_{r,N}(L)\downarrow a_r(L)>0.
\]
The result is strict positivity of the Weil operator and its finite
matrices; it makes no assertion that \(a_1(L)\) is simple or even.
Strict positivity under RH is already stated in the introduction of
S21.  The proof here supplies the full closed-domain and spectral-gap
details and the higher Galerkin limits.
\end{theorem}
\begin{proof}
The explicit formula and \(A\) give the displayed sum on the
piecewise smooth core.  For \(f\in\operatorname{Dom}q_L\), choose
core vectors \(f_n\to f\) in form norm.  The vectors
\((\sqrt{m_\gamma}\widehat f_n(\gamma))_\gamma\) are Cauchy in
\(\ell^2\), since their squared distance is \(q_L(f_n-f_k)\).
For each fixed \(\gamma\), evaluation of the Fourier transform is
continuous on \(H_L\), by
\(|\widehat f(\gamma)|\leq\sqrt{2L}\|f\|_2\).
Consequently the limit has exactly the displayed coordinates, and the
identity extends to the closed form domain.

If \(q_L(f)=0\), its entire Fourier transform vanishes at every
distinct ordinate.  A nonzero Fourier transform of an \(L^2\) function
supported in \([-L,L]\) has only \(O_f(T)\) zeros in \(|z|\leq T\),
counting multiplicity.  Indeed
\(|\widehat f(z)|\leq\sqrt{2L}\|f\|_2e^{L|\operatorname{Im}z|}\),
and Jensen's formula about any point at which the transform is nonzero,
applied on discs of radii \(T\) and \(2T\), gives that bound.
The distinct zeta ordinates have superlinear counting growth.  One may
use Conrey's unconditional positive-proportion simple-zero theorem
(\emph{Bull. Amer. Math. Soc.} 20 (1989), Theorem 1, p.~79), together
with the Riemann--von Mangoldt formula.  Alternatively, directly under
\(A\), Carneiro--Chandee--Milinovich,
\emph{Bounding \(S(t)\) and \(S_1(t)\) on the Riemann hypothesis},
arXiv:1309.1526, Theorem 2 and equation (1), give
\(m_\gamma=O(\log T/\log\log T)\) on \(T/2\leq\gamma\leq T\)
and \(N(T)-N(T/2)\gg T\log T\).  The number of distinct ordinates in
that interval is therefore \(\gg T\log\log T\).
Both routes contradict the Jensen bound unless \(f=0\).
The compact-resolvent theorem cited above means that the bottom of the
spectrum is an attained eigenvalue.  Its nonnegativity and absence of a
kernel therefore imply \(a_1(L)>0\).

The minimum--maximum principle gives
\(a_{r,N+1}(L)\leq a_{r,N}(L)\) and
\(a_{r,N}(L)\geq a_r(L)\).  Approximate an orthonormal basis of the
span of the first \(r\) eigenvectors in form norm by vectors in a
single \(E_N\).  On this fixed finite-dimensional span the Gram and
form matrices converge, so the largest Rayleigh quotient converges to
\(a_r(L)\).  Applying the same principle gives the reverse limit
inequality.  This proves the claimed convergence.
\end{proof}

\begin{lemma}[AC-S04: a normalized global null function]
\label{AC-S04}
Define
\[
 h(u)=2\pi u^2(2\pi u^2-3)e^{-\pi u^2},\qquad
 K(x)=e^{x/2}\sum_{n\geq1}h(ne^x).
\]
Then \(K\) is a nonzero even smooth function, every derivative decays
faster than every exponential at both ends, and
\[
 \widehat K(t)=\Xi(t):=\xi(1/2+it),\qquad
 \xi(s)=\tfrac12s(s-1)\pi^{-s/2}\Gamma(s/2)\zeta(s).
\]
For every integer \(j\geq0\) and \(u=e^x\geq2\),
\[
 |K^{(j)}(x)|\leq C_j u^{9/2+2j}e^{-\pi u^2}.
\]
With the unnormalized Fourier transform explicitly used in S25,
equation \texttt{mhat}, the function \(h\) printed in equation
\texttt{ku} of that source is one quarter of the one above.  Its
transform is \(\Xi/4\), so the equality as printed in Lemma 7.1
(\texttt{hermfact}) requires this scalar correction.
\end{lemma}
\begin{proof}
Gaussian Fourier differentiation gives
\(\mathcal F(h)=h\) for the Fourier kernel \(e^{2\pi iut}\).
Explicitly, the transforms of \(u^2e^{-\pi u^2}\) and
\(u^4e^{-\pi u^2}\) are respectively
\((1/(2\pi)-t^2)e^{-\pi t^2}\) and
\((t^4-3t^2/\pi+3/(4\pi^2))e^{-\pi t^2}\); substituting the
coefficients of \(h\) proves the assertion.  Also \(h(0)=0\) and
\(\int h=\mathcal Fh(0)=0\).  Poisson summation, pairing positive
and negative integers, therefore gives \(K(x)=K(-x)\).
Termwise differentiation on \(u\geq2\) is justified by Gaussian
domination.  The \(j\)-th derivative is a sum of
\(u^{1/2}P_j(nu)e^{-\pi n^2u^2}\), where \(P_j\) has degree at
most \(4+2j\).  Factoring out \(e^{-\pi u^2}\) bounds the remaining
sum uniformly for \(u\geq2\), and proves the stated estimate and
the decay at both ends.

For \(\operatorname{Re}s>1\), Fubini and the substitution \(v=nu\)
give
\[
 \int_0^\infty u^{1/2}\sum_{n\geq1}h(nu)\,
                  u^{s-1/2}\frac{du}{u}
   =\zeta(s)\int_0^\infty h(v)v^{s-1}\,dv.
\]
The elementary Gaussian integral gives
\[
 \int_0^\infty h(v)v^{s-1}\,dv
   =\tfrac12s(s-1)\pi^{-s/2}\Gamma(s/2).
\]
The rapid decay already proved makes the integral on the left entire
in \(s\).  Analytic continuation, followed by \(s=1/2-it\) and
\(\xi(s)=\xi(1-s)\), proves the Fourier identity.  Since \(\xi\)
is not identically zero, \(K\ne0\).  Repeating the same integral
with the coefficient \(\pi/2\) of S25 instead of \(2\pi\) divides
the result by four, proving the scalar correction.
\end{proof}

\begin{theorem}[AC-S05: every fixed low Weil eigenvalue tends to zero]
\label{AC-S05}
Under \(A\), for each fixed positive integer \(r\) there are
\(C_r<\infty\) and \(\lambda_r\) such that
\[
 0<a_r(\log\lambda)
 \leq C_r\lambda^{4r+5}e^{-2\pi\lambda^2}
 \quad(\lambda\geq\lambda_r).
\]
In particular \(\mu_\lambda=a_1(\log\lambda)\) is positive for every
finite window, decreases to zero, and satisfies
\(\mu_\lambda=O(\lambda^9e^{-2\pi\lambda^2})\).
Together with S25, Corollary 3.8 (\texttt{strange}), this gives
\[
 \mathrm{RH}\quad\Longleftrightarrow\quad
 \lim_{\lambda\to\infty}\mu_\lambda=0.
\]
It does not give the numerical asymptotic
\(\mu_\lambda\asymp e^{-4\pi\lambda^2}\) discussed in S25.
\end{theorem}
\begin{proof}
Put \(L=\log\lambda\), \(K_j=K^{(j)}\), and
\(f_{j,L}=1_{[-L,L]}K_j\), \(0\leq j<r\).
These are piecewise smooth compactly supported functions and belong
to the form domain; this is also explicitly observed immediately after
S25, Corollary 3.7.  They need not lie in the operator domain.
Here is a verification in the full form norm.  For fixed \(L\), multiply
\(K_j\) by a smooth cutoff equal to one on
\([-L+\delta,L-\delta]\) and supported in \((-L,L)\), with derivative
bounded by \(C/\delta\).  The difference \(w_\delta\) from
\(f_{j,L}\) is supported in two intervals of total length \(2\delta\),
is uniformly bounded, and has uniformly bounded variation, including
its endpoint jumps.  Hence
\[
 |\widehat w_\delta(t)|\leq C_{j,L}
                   \min\{\delta,|t|^{-1}\}.
\]
The archimedean symbol is bounded in absolute value by
\(C_L+\log(2+|t|)\); the finite prime and pole terms are bounded
operators on \(H_L\), as in S25, the proof of Theorem 3.6.
Splitting the Fourier integral at \(|t|=\delta^{-1}\) gives
\[
 \|w_\delta\|_2^2+
 \int_{\mathbb R}\log(2+|t|)|\widehat w_\delta(t)|^2\,dt
       =O_{j,L}(\delta\log(1/\delta)).
\]
Thus the smooth cutoffs converge to \(f_{j,L}\) in the closed form
norm.  In particular, membership in the form domain and the spectral
identity used below are justified by form convergence, not just by
\(L^2\) convergence.
The Fourier transforms of the uncut functions are
\((it)^j\Xi(t)\), and vanish at every zero ordinate.
For \(R_{j,L}=K_j-f_{j,L}\), the derivative in the sense of measures
has two boundary atoms as well as the derivatives on the two tails.
The estimates of AC-S04 and
\[
 \int_\lambda^\infty u^a e^{-\pi u^2}\frac{du}{u}
       \leq C_a\lambda^{a-2}e^{-\pi\lambda^2}
 \quad(\lambda\geq\lambda_a)
\]
show that its total variation satisfies
\[
 \operatorname{Var}(R_{j,L})
 \leq C_j\lambda^{9/2+2j}e^{-\pi\lambda^2}.
\]
Integration by parts against that finite measure gives, for \(t\ne0\),
\[
 |\widehat R_{j,L}(t)|
       \leq |t|^{-1}\operatorname{Var}(R_{j,L}).
\]
The Riemann--von Mangoldt estimate implies
\(B:=\sum_\gamma m_\gamma/\gamma^2<\infty\); there is no zero
ordinate at zero.  For \(a=(a_0,\ldots,a_{r-1})\), the spectral
formula of AC-S03 and the vanishing of \(\widehat K_j(\gamma)\)
therefore give
\[
 q_L\left(\sum_{j<r}a_jf_{j,L}\right)
 \leq C_r\|a\|^2\lambda^{4r+5}e^{-2\pi\lambda^2}.
\]
The functions \(K_0,\ldots,K_{r-1}\) are linearly independent:
after Fourier transform a linear relation is
\(P(it)\Xi(t)=0\), which forces the polynomial \(P\) to vanish.
Their Gram matrix is thus positive definite.  The Gram matrices of
the cutoffs converge to it in every entry by dominated convergence.
Consequently, for all sufficiently large \(L\),
\[
 \left\|\sum_{j<r}a_jf_{j,L}\right\|_2^2
           \geq c_r\|a\|^2\quad(c_r>0).
\]
This is an \(r\)-dimensional trial space.  The minimum--maximum
principle proves the upper bound; AC-S03 proves strict positivity.
Monotonicity for \(r=1\) follows by extending test functions by zero
as the interval grows, or from S25, equation \texttt{monotone}.
The converse in the final equivalence is exactly S25, Corollary 3.8.

An explicit constant is available for the lowest eigenvalue.  Put
\(c_0=\int_{-\log2}^{\log2}|K(x)|^2\,dx>0\) and
\(B=\sum_\gamma m_\gamma/\gamma^2\).  For \(u\geq2\), every
summand of \(K(\log u)\) is positive and decreasing: its logarithmic
derivative has the sign of
\(-\pi v^2(8\pi^2v^4-30\pi v^2+15)\), \(v=nu\).
Consequently \(\operatorname{Var}(R_{0,L})=4K(L)\).
Also
\[
 K(\log\lambda)\leq
 4\pi^2\lambda^{9/2}e^{-\pi\lambda^2}
       \sum_{n\geq1}n^4e^{-4\pi(n^2-1)}
 \leq8\pi^2\lambda^{9/2}e^{-\pi\lambda^2}
 \quad(\lambda\geq2).
\]
For the last bound, use \(n^2-1\geq n-1\) and
\(\sum_{n\geq1}n^4q^{n-1}=(1+11q+11q^2+q^3)/(1-q)^5<2\)
at \(q=e^{-4\pi}\).  Therefore the concrete inequality
\[
 \mu_\lambda\leq
 \frac{1024\pi^4B}{c_0}\lambda^9e^{-2\pi\lambda^2}
 \quad(\lambda\geq2)
\]
holds.  Both \(B\) and \(c_0\) are fixed positive constants defined
by absolutely convergent expressions.
\end{proof}

\subsection{The convergence that can and cannot presently be asserted}

\begin{proposition}[AC-S06: fixed-window eigenfunction convergence]
\label{AC-S06}
Fix \(L>0\).  Any sequence of \(L^2\)-normalized lowest Galerkin
eigenvectors has a subsequence converging in \(L^2\) and in form norm
to a lowest eigenvector of \(A_L\).  Their Fourier transforms converge
uniformly on compact subsets of \(\mathbb C\) along that subsequence.
If the lowest eigenvalue of \(A_L\) is simple, phases can be chosen so
that the whole sequence converges.  If that eigenvector is also even,
then for all sufficiently large \(N\) the Galerkin matrix is even-simple,
and the normalized determinant of S25, Theorem 5.10
(\texttt{finmain}), converges locally uniformly to a nonzero scalar
multiple of the Fourier transform of that eigenvector, after removing
the fixed phase factor \(-i\lambda^{-iz}\).
This assertion permits an \(L^2\) normalization.  It does not assert
convergence of the boundary normalization \(\xi_N(\lambda)=1\).
\end{proposition}
\begin{proof}
The eigenvalues converge by AC-S03.  Boundedness of the form norm and
compactness of its embedding in \(H_L\) give a strongly \(L^2\)
convergent subsequence and a weakly form-convergent subsequence with
the same limit.  For \(g\) in the union of the \(E_N\), the Galerkin
eigenvalue identity eventually holds.  Passing to the limit gives
\(q_L(f,g)=a_1(L)\langle f,g\rangle\).  The form-core property
extends this equality to the whole form domain.  Hence \(f\) is a
normalized lowest eigenvector.  Its form norm equals the limiting form
norm of the approximants, so weak convergence in that Hilbert norm is
strong.  Simplicity identifies every subsequential limit up to phase.

For a compact set \(|\operatorname{Im}z|\leq R\),
\[
 |\widehat f_N(z)-\widehat f(z)|
 \leq\left(\int_{-L}^Le^{2R|x|}\,dx\right)^{1/2}\|f_N-f\|_2,
\]
proving local uniform convergence.  If the lowest eigenvector is simple
and even, the spectral gap to the second eigenvalue and convergence
of both Galerkin eigenvalues make the first Galerkin eigenvalue simple
for large \(N\).  Inversion preserves \(E_N\), and its eigenvectors
have parity.  Convergence to the even lowest vector forces even parity.
The boundary functional on that finite-dimensional eigenvector is
nonzero by the argument between Definition 5.3 and Lemma 5.4 of S25.
The determinant formula of Theorem 5.10 is thus applicable.  Replacing
its boundary normalization by the \(L^2\) normalization multiplies
its determinant formula by a nonzero constant; the preceding transform
convergence proves the claim.  Point evaluation at a boundary is not
continuous in \(L^2\), so no boundary-normalization limit was used.
\end{proof}

\begin{proposition}[AC-S07: a precise sufficient gate for the large-window limit]
\label{AC-S07}
Let \(u_L\) be a normalized even simple lowest eigenvector, and let
\(v_L\) be a normalized trial vector in \(\operatorname{Dom}q_L\).
Write \(g_L=a_2(L)-a_1(L)>0\).  After choosing a phase,
\[
 \|u_L-v_L\|_2^2
 \leq 2\frac{q_L(v_L)-a_1(L)}{g_L}.
\]
If trial vectors \(c_Lv_L\), extended by zero, converge to \(K\)
in \(L^1(e^{b|x|}dx)\), and
\[
 |c_L|^2\left(\int_{-L}^Le^{2b|x|}\,dx\right)
       \frac{q_L(v_L)-a_1(L)}{g_L}\longrightarrow0,
\]
then \(c_L\widehat u_L\to\Xi\) uniformly on
\(|\operatorname{Im}z|\leq b\).  The same statement with a compact
strip region follows from its corresponding weaker norm bound.
\end{proposition}
\begin{proof}
Expand \(v_L=\alpha u_L+w\), \(w\perp u_L\), and choose the
phase so that \(\alpha\geq0\).  The spectral theorem for the closed
form gives
\(q_L(v_L)-a_1(L)\geq g_L\|w\|^2\), while
\(\|v_L-u_L\|^2=2(1-\alpha)\leq2(1-\alpha^2)=2\|w\|^2\).
Weighted Cauchy--Schwarz on \([-L,L]\) bounds the weighted \(L^1\)
distance between \(c_Lu_L\) and \(c_Lv_L\) by the square root of
the second displayed expression, times \(\sqrt2\).
The triangle inequality and
\(|\widehat f(z)|\leq\int|f(x)|e^{b|x|}dx\) complete the proof.
\end{proof}

\begin{proposition}[AC-S13: the prolate trial function has a proved strip limit]
\label{AC-S13}
Let \(h\) be the correctly normalized function in AC-S04.  Suppose
\(h_\lambda\) is supported in \([-\lambda,\lambda]\) and
\[
 \sup_{|u|\leq\lambda}|h_\lambda(u)-h(u)|\leq C\lambda^{-2}.
\]
This is the estimate in S25, Lemma 7.2 (\texttt{meixnerlem}), after
the same scalar normalization as AC-S04.  Define
\[
 k_\lambda(x)=1_{[-\log\lambda,\log\lambda]}(x)e^{x/2}
                        \sum_{n\geq1}h_\lambda(ne^x).
\]
For every \(0\leq b<1/2\),
\[
 \|k_\lambda-K\|_{L^1(e^{b|x|}dx)}
      =O_b(\lambda^{b-1/2}),\qquad
 \sup_{|\operatorname{Im}z|\leq b}
          |\widehat k_\lambda(z)-\Xi(z)|
      =O_b(\lambda^{b-1/2}).
\]
This is an unconditional trial-function limit.  It does not assert
that \(k_\lambda\) is a ground eigenvector.
\end{proposition}
\begin{proof}
Write \(u=e^x\), \(L=\log\lambda\), and
\(w_b(u)=\max(u^b,u^{-b})\).  On \(\lambda^{-1}\leq u\leq\lambda\),
the at most \(\lambda/u\) retained terms contribute an error at most
\(C\lambda^{-1}u^{-1/2}\).  Its weighted integral is bounded by
\[
 C\lambda^{-1}\int_{\lambda^{-1}}^\lambda
          u^{-1/2}w_b(u)\frac{du}{u}
       =O_b(\lambda^{b-1/2}).
\]
There is also a tail from the terms of \(h\) with \(nu>\lambda\),
which is not included in that first estimate.  For \(\lambda\geq2\),
\(h\) is positive and decreasing on \([\lambda,\infty)\).
An integral comparison bounds this tail by
\[
 u^{1/2}\sum_{nu>\lambda}|h(nu)|
 \leq u^{1/2}h(\lambda)+u^{-1/2}\int_\lambda^\infty h(v)\,dv.
\]
Here \(h(\lambda)=O(\lambda^4e^{-\pi\lambda^2})\) and
\(\int_\lambda^\infty h(v)\,dv=O(\lambda^3e^{-\pi\lambda^2})\).
Its weighted integral on the window is therefore
\(O_b(\lambda^{b+9/2}e^{-\pi\lambda^2})\).
Finally the weighted integral of \(K\) outside the window is
\(O_b(\lambda^{b+5/2}e^{-\pi\lambda^2})\), by its evenness and
AC-S04.  This proves the first estimate, including both omitted tails.
The Fourier estimate follows from the weighted \(L^1\) bound.
\end{proof}

The two obstacles expressly listed in S25, \S8 (\texttt{missing}),
are simplicity/evenness and comparison with the prolate candidate.
AC-S03 and AC-S05 do not settle them.  In particular, AC-S05 also
makes \(a_2(L),a_3(L),\ldots\) tend to zero; a small Rayleigh quotient
without a comparison to the spectral gap cannot select the ground
state.  AC-S07 makes the missing comparison quantitative.  Positivity
and parity invariance alone do not force an even simple minimum:
in a two-dimensional space with one even and one odd vector,
\(\operatorname{diag}(2,1)\) has an odd minimum and \(I\) has a
multiple minimum.  These are counterexamples to that abstract
inference, not counterexamples for the specific Weil matrices.

S21, \S\S4--5, exhibits numerical agreements for its different
operators \(D(\lambda,k)\); Proposition 4.2 proves a counting
asymptotic, and Theorem 6.4 proves the exact zeta-cycle realization.
Neither is a convergence theorem for \(D(\lambda,k)\) along an
unspecified joint limit.  A proof of that numerical convergence still
needs approximation estimates for the chosen prolate projection,
a specified path \(k=k(\lambda)\), and control of spectral pollution.
For S25, the fixed-window result above is a proved form of convergence;
the further \(\lambda\to\infty\) determinant convergence retains
the stated gates.  This is not a claim of logical independence from RH.
\subsection{A spectral determinant construction with a proved limit}

\begin{theorem}[AC-S08: an explicit Hilbert realization and determinant]
\label{AC-S08}
Under \(A\), form the Hilbert space with orthonormal basis
\(e_{\gamma,j}\), where \(\gamma\) ranges over distinct positive and
negative ordinates and \(1\leq j\leq m_\gamma\), and put
\[
 De_{\gamma,j}=\gamma e_{\gamma,j},\qquad
 \operatorname{Dom}D=
 \left\{(c_{\gamma,j}):\sum_{\gamma,j}\gamma^2|c_{\gamma,j}|^2<\infty\right\}.
\]
Then \(D\) is self-adjoint with compact resolvent,
\(D^{-1}\) is Hilbert--Schmidt, and
\[
 \det{}_2(I-zD^{-1})=\frac{\Xi(z)}{\Xi(0)}.
\]
For the spectral cutoff \(P_T=1_{[-T,T]}(D)\),
\[
 \left\|D^{-1}-P_TD^{-1}P_T\right\|_2^2
       =O(\log T/T),\qquad
 \det(I-zD_T^{-1})\longrightarrow\Xi(z)/\Xi(0),
\]
where \(D_T\) acts on \(P_TH\), and the determinant convergence is
locally uniform with error \(O_R(\log T/T)\) on \(|z|\leq R\).
The diagonal algebra \(c_0(\{(\gamma,j)\})\), or its unitization,
with this \(D\) gives a spectral triple; its summability threshold is
\(p>1\).  This construction uses the zero multiset and does not identify
the operators of S21 or S25 with these diagonal cutoffs.
\end{theorem}
\begin{proof}
The multiplication operator on the stated maximal domain is self-adjoint.
There are finitely many ordinates in any bounded interval, and their
absolute values tend to infinity, so its resolvent is compact.  Partial
summation of \(N(T)=O(T\log T)\) gives
\(\sum_{\gamma>T}m_\gamma/\gamma^2=O(\log T/T)\), proving both
Hilbert--Schmidt assertions.  The order-one Hadamard factorization of
the even entire function \(\Xi\), paired at \(\gamma,-\gamma\),
is
\[
 \Xi(z)/\Xi(0)=\prod_{\gamma>0}(1-z^2/\gamma^2)^{m_\gamma}.
\]
Indeed the paired product converges by the preceding summability;
Hadamard factorization leaves an exponential \(e^{az+b}\), whose
linear term vanishes by evenness and whose constant is fixed at zero.
The defining regularized determinant product
\(\prod_{\gamma,j}(1-z/\gamma)e^{z/\gamma}\) has precisely these
paired factors.  The finite ordinary determinant has the truncated
paired product.  For \(T>2R\), the logarithm of its omitted tail is
bounded by
\(C_R\sum_{\gamma>T}m_\gamma/\gamma^2\).
The truncated products are bounded uniformly on \(|z|\leq R\) by
\(\exp(R^2\sum_{\gamma>0}m_\gamma/\gamma^2)\), which gives the
stated absolute error.  Finite diagonal sequences form a dense
\(*\)-subalgebra with zero commutators with \(D\); compact resolvent
gives the remaining spectral-triple condition.  Finally partial
summation of \(N(T)\sim T\log T/(2\pi)\) proves convergence of
\(\operatorname{Tr}(1+D^2)^{-p/2}\) exactly for \(p>1\).
\end{proof}

\subsection{What the semilocal prolate construction actually supplies}

For a finite set of places \(S\ni\infty\), S23,
Propositions \texttt{groundstate} and \texttt{httransfoS}, gives the
unitary model
\[
 D_S=M_t\ \text{on }L^2(\mathbb R,dm_S),\qquad
 dm_S(t)=c_S\left|\Gamma(1/4+it/2)\right|^2
       \prod_{p\in S\setminus\{\infty\}}
          |1-p^{-1/2+it}|^{-2}\,dt,
\]
where a positive constant \(c_S\) may normalize the cyclic vector.
Let \(N_S\) be the number operator of the associated orthogonal
polynomials.  The formal candidate in Definition \texttt{prolop} is
\(T_{\lambda,S}=-D_S^2+\lambda^2N_S\); the physical archimedean
operator uses the rescaling and additive constant in S23, equation
\texttt{wlambdconj1}.

\begin{proposition}[AC-S09: spectral type and self-adjoint extensions]
\label{AC-S09}
For every finite \(S\), \(D_S\) has purely absolutely continuous
spectrum \(\mathbb R\), of multiplicity one, and no eigenvalues.
The formal \(T_{\lambda,S}\) on the finite span of orthogonal
polynomials is densely defined and symmetric.  It admits self-adjoint
extensions; extensions preserving the even/odd grading can be chosen.
Their existence and real spectra are unconditional.  The construction
does not specify a unique extension or establish a discrete zeta
spectrum for it.
\end{proposition}
\begin{proof}
The displayed density is strictly positive and finite everywhere on
\(\mathbb R\), since the gamma function has no zeros or poles on
this line and \(p^{-1/2}<1\).  Multiplication by the square root of
the density identifies this model with multiplication by \(t\) on
ordinary \(L^2(\mathbb R,dt)\).  This proves the spectral assertion.
Every polynomial has all required moments, and multiplication by
\(t^2\) and \(N_S\) preserve the polynomial span, on which they
are symmetric.  Complex conjugation \(J\) commutes with this
 symmetric operator.  The span is dense: the finite Euler factors
give norms equivalent to the archimedean norm, in which completeness
of these polynomials follows from S23, Proposition \texttt{numberN}
and its unitary transform \texttt{httransfo}.
Complex conjugation \(J\) therefore maps
\(\ker(T_{\lambda,S}^*-i)\) antiunitarily onto
\(\ker(T_{\lambda,S}^*+i)\), so the deficiency indices agree.
Explicitly, choose an orthonormal basis \((e_j)\) of the first
deficiency space and the basis \((Je_j)\) of the second, and let
\(Ue_j=Je_j\) with \(U\) extended linearly.  The extension has
domain
\[
 \operatorname{Dom}\overline T\ \dotplus\
 \{h+Uh:h\in\ker(T^*-i)\},
\]
and acts by \(\overline T f+ih-iUh\).  The standard von Neumann
extension theorem gives self-adjointness.  One may make these choices
separately on the two parity subspaces because \(T\) preserves them.
Theorem \(A\) is not used in any step.
\end{proof}

\begin{proposition}[AC-S10: the semilocal number operator is not a transported one]
\label{AC-S10}
If \(S\) contains a finite prime, the identity map on polynomials
between the archimedean and semilocal weighted spaces does not
intertwine their number operators.  Thus the bounded invertible
\(\eta_S\) of S23 cannot be used to transport the archimedean
prolate spectrum by the claimed intertwining
\(N_S\eta_S=\eta_SN_\infty\).
\end{proposition}
\begin{proof}
If it did intertwine them, every archimedean orthogonal polynomial
of degree \(n\) would also be an eigenvector of the semilocal
\(N_S\) for \(n\).  Since these eigenspaces are one-dimensional,
the two monic orthogonal-polynomial systems would coincide.
In particular the integral of every positive-degree archimedean
orthogonal polynomial against the normalized semilocal measure would
vanish.  Expanding arbitrary polynomials in that basis shows that all
normalized moments of the two measures would agree.  These moments
are determinate: the archimedean Jacobi coefficients are
\(a_n=\sqrt{(n+1)(n+1/2)}\), and
\(\sum a_n^{-1}=\infty\) gives the Carleman determinacy criterion
(S23, Proposition \texttt{fourierfmuall}, and the determinacy argument
in Theorem \texttt{dzero}, part (ii)).
But their density ratio is a constant times
\[
 R_S(t)=\prod_{p\in S\setminus\{\infty\}}
                         |1-p^{-1/2+it}|^{-2}.
\]
This function is not constant: each factor is at most its value at
\(t=0\), and choosing \(t\) for which one of the phases is not one
makes the inequality strict.  Thus the measures are different, a
contradiction.  The conclusion agrees with S23, Remark
\texttt{care}(i); it has been re-proved here to identify exactly
what prevents this transfer.
\end{proof}

\begin{proposition}[AC-S11: correction of the Jacobi metaplectic normalization]
\label{AC-S11}
On finite sequences, let \(N e_n=ne_n\), and let \(B\) be a real
Jacobi matrix with zero diagonal and positive off-diagonal entries
\(b_n\).  Put \(C=[B,N]\).
The assignments \(h\mapsto-iB\), \(k\mapsto N+c\) in S23,
\S5.2, equation \texttt{h}, are incompatible with the standard
relations \([h,e_\pm]=\pm2e_\pm\), \([e_+,e_-]=h\) and
\(k=i(e_--e_+)\): they give
\([\sigma(e_+),\sigma(e_-)]=\sigma(h)/4\).
The corrected assignments
\[
 \sigma(h)=-iB,\quad \sigma(k)=2(N+c),\quad
 \sigma(e_+)=i(N+c)+\tfrac12C,\quad
 \sigma(e_-)=-i(N+c)+\tfrac12C
\]
satisfy the Lie algebra relations if and only if
\[
 b_n^2=(n+1)(n+2c)\quad(n\geq0).
\]
For every \(c>0\) these entries are positive, so the Lie algebra
relations alone do not select \(c=1/4\).  The even metaplectic
realization selects that value by its compact-generator normalization;
it then gives the coefficients printed in S23, equation
\texttt{coefffixed}.  The uniqueness assertion of Theorem
\texttt{dzero}(i) requires this extra normalization.
\end{proposition}
\begin{proof}
Directly from the tridiagonal matrix,
\([[B,N],N]=B\), equivalently \([N,C]=-B\), and
\[
 [B,C]e_n=2(b_{n-1}^2-b_n^2)e_n,\qquad b_{-1}=0.
\]
Solving \(k=i(e_--e_+)\) and
\([h,k]=-2i(e_++e_-)\) gives the source's formulas
\(\sigma(e_+)=\tfrac i2(N+c)+\tfrac14C\),
\(\sigma(e_-)=-\tfrac i2(N+c)+\tfrac14C\).
Their commutator is
\(\tfrac i4[N,C]=-iB/4\), proving the defect independent of the
choice of \(b_n\).
For the corrected assignments the commutator is \(i[N,C]=-iB\).
The other two Lie relations are equivalent to
\([B,C]=-4(N+c)\), or
\(b_n^2-b_{n-1}^2=2(n+c)\).  Summing from zero gives precisely
\((n+1)(n+2c)\); substitution verifies sufficiency.
In the actual even oscillator representation the compact generator
has eigenvalues \(2n+1/2\), by S23, equation \texttt{b-k} and
Proposition \texttt{numberN}.  Thus \(2(n+c)=2n+1/2\), so
\(c=1/4\).  This proves the corrected statement without altering
the independently proved archimedean Jacobi coefficients.
\end{proof}

\subsection{Classical prolate spectrum: established facts and remaining identifications}

The self-adjoint full-line prolate operator is not made self-adjoint
by \(A\): that result was already proved by Connes--Moscovici,
\emph{The UV prolate spectrum matches the zeros of zeta},
\emph{PNAS} 119 (2022), e2123174119, Theorem 1.6.  That theorem gives
a discrete spectrum unbounded in both directions, simple negative
eigenvalues, and the extension commuting with both time and Fourier
cutoffs.  Corollary 1.7 places the negative eigenvectors in the Sonin
space.  These are background theorems.
The finite-dimensional possible excess of positive eigenvalues
mentioned in S23 is no longer an open conjecture in the current
literature: Ramis--Richard-Jung--Thomann,
\emph{New eigenfunctions for the negative part of the Connes--Moscovici
prolate spectrum}, \emph{C. R. Math.} 363 (2025), 1065--1081,
Corollary 15, p.~1076, proves that all nonclassical CM eigenvalues are
negative.  This is an unconditional result, not a consequence of \(A\).

The full Darboux square root \(2\!\not\!D\) in CM, Theorem 5.1,
has both real and imaginary spectrum.  It is consequently neither a
self-adjoint nor a skew-adjoint Hilbert-space operator.  A theorem
about the zeros of zeta does not change those existing spectral values.
Its negative-square spectral restriction is meaningful, but identifying
its exact spectral determinant with \(\Xi\) requires an equality
of spectral data, beyond a counting asymptotic.  The same applies to
the semilocal prolate candidates: the required domain, cutoff
commutation, and trace identification have not been supplied by \(A\).

\begin{proposition}[AC-S12: the imaginary-axis differential expression]
\label{AC-S12}
For the classical differential expression
\[
 W_\lambda=-\partial_x((\lambda^2-x^2)\partial_x)
                                  +(2\pi\lambda x)^2,
\]
put \(x=iy\).  The negative of the resulting real-axis expression
is
\[
 B_\lambda=-\partial_y((\lambda^2+y^2)\partial_y)
                                      +(2\pi\lambda y)^2.
\]
Under \(y=\lambda\sinh t\) and the unitary change
\(v(t)=\sqrt{\lambda\cosh t}\,f(\lambda\sinh t)\), it becomes
\[
 -\partial_t^2+(2\pi\lambda^2)^2\sinh^2t
                    +\frac14+\frac{1}{4\cosh^2t}.
\]
On a half-line, parity at zero gives the Dirichlet or Neumann boundary
condition respectively.  In particular its ordinary self-adjoint
realization is strictly positive and has compact resolvent.
\end{proposition}
\begin{proof}
The substitution \(\partial_x=-i\partial_y\) gives the first
formula.  For \(p(y)=\lambda^2+y^2\), the Liouville change
\(dt/dy=p^{-1/2}\), \(v=p^{1/4}f\), gives additional potential
\(p''/4-(p')^2/(16p)\).
Here this is
\(1/2-y^2/(4(\lambda^2+y^2))
 =1/4+1/(4\cosh^2t)\), proving the displayed expression.
The norm is preserved because \(dy=\lambda\cosh t\,dt\).
At zero the multiplier is even with derivative zero, so odd/even
parity gives the claimed separated boundary conditions.
The potential is at least \(1/4\) and tends to infinity at infinity.
The positive closed form is the sum of derivative energy and potential
energy; on a bounded interval its bounded-energy sets are precompact,
while the diverging potential makes the tails uniformly small.
Thus the form embedding is compact, proving compact resolvent and
strict positivity.  This is a differential-expression calculation;
transporting a particular CM boundary condition through analytic
continuation is the separate theorem in RRJT, Theorem 14.
\end{proof}

\subsection{Algebraic-geometric interpretation and its limits}

The arithmetic L-function here is that of the unit Artin--Tate object
over \(\operatorname{Spec}\mathbb Z\).  Its finite-place local factor
comes from the rank-one constant representation and is
\((1-p^{-s})^{-1}\).  The sheaf in AC-S01 lives on the scaling topos,
not on the small etale site of an ordinary scheme.  Its local jet
algebra is a nonreduced zero-dimensional analytic algebra supported
at \(\gamma\).  Passing to the associated graded of its maximal-ideal
filtration gives \(m_\gamma\) copies of its residue field; this is
the semisimplified multiplicity representation used by AC-S08.
Taking this associated graded preserves the local length (and its
class in the Grothendieck group of finite-length modules) but removes
the Jordan extensions exhibited in AC-S02.  No weight or purity
theorem for finite-field etale cohomology removes these extensions
on the basis of \(A\) alone.

AC-S03--AC-S07 are analytic inequalities and convergence statements
for a closed form; they do not construct an algebraic square or an
algebraic correspondence.  AC-S09--AC-S12 concern real/complex
operator domains and an enveloping-algebra action.  Their rigorous
algebraic content is the explicit \(\mathfrak{sl}_2\) module and its
normalization, not a new pure motive or an etale Frobenius action.
Exact archive cotes or scan numbers for these new analytic statements
have not been supplied or verified; none is invented here.

\section{The arithmetic square, Riemann--Roch, and absolute counting}
\label{sec:sites-twistor}

\subsection{Source inventory and the exact numerical conclusion}

The statements in this subsection distinguish a proved identity on test
functions from an algebro-geometric realization of that identity.  The
attached version of Connes--Consani, \emph{The Riemann--Roch strategy}
(arXiv:1805.10501), \S3.1, equations labelled \texttt{negcrit},
\texttt{negcrit1}, and \texttt{negcrit1bis}, gives the numerical target.
Its \S3.4 lists five steps: construction of a complex lift, regularization
of the relative intersection theory, a surface Riemann--Roch formula,
an existence theorem for sections, and tropical descent.  Construction of
the complex lift and a lift of the scaling correspondences is part of that
paper; the complete surface intersection/Riemann--Roch theorem is not.

Here are the relevant already proved results, with locators in the supplied
sources.  They are background, not findings of this work.
\begin{itemize}
\item \emph{Geometry of the scaling site} (1603.03191), Theorems 3.9
(\texttt{scaltop}), 4.2 (\texttt{structure2}), and 4.3
(\texttt{structure3}): points, stalks, and points over the tropical
semifield; Theorem 5.6 (\texttt{thmjaccp}): the divisor-class description
on the periodic orbit $C_p$; Theorem 5.17 (\texttt{RRperiodic}):
$\operatorname{Dim}_{\mathbb R}H^0(D)-\operatorname{Dim}_{\mathbb R}H^0(-D)
=\deg D$, with the continuous dimension defined as the limit in the
source's \texttt{rr1}.  These statements do not assume RH.
\item \emph{Riemann--Roch for the Arakelov compactification of
$\operatorname{Spec}\mathbb Z$} (2205.01391), Theorem 4.3
(\texttt{rrspzb}):
\[
 \dim_{\mathbb S[\pm1]}H^0(D)-\dim_{\mathbb S[\pm1]}H^1(D)
 =\left\lceil\frac{\deg D+\log2}{\log3}\right\rceil'
 -\mathbf1_L(\deg D).
\]
Here the odd ceiling and the exceptional set $L$ are exactly those in
\texttt{ceilingprime} and \texttt{excL} of that paper.  Theorem 5.3
(\texttt{serredual}) identifies
$H^0(K-D)$ with the appropriate tolerant-module internal Hom from
$H^1(D)$, where $K=-2\{2\}$.  Both statements are unconditional.
\item \emph{Riemann--Roch for the ring $\mathbb Z$} (2306.00456),
Theorem 5.1 (\texttt{rrspzb}):
\[
 \dim_{\mathbb S}H^0(D)-\dim_{\mathbb S}H^1(D)
 =c(\deg D/\log2)+1.
\]
The function $c$ is the source's \emph{right-continuous} extension of
$\lceil x\rceil$ for positive nonintegral $x$ and
$-\lceil-x\rceil$ for negative nonintegral $x$.  It must not be replaced
by the ordinary ceiling at integral arguments; for example $c(0)=1$.
The exact formula, including this convention, is already unconditional.
\end{itemize}
None is a Riemann--Roch theorem for the two-dimensional arithmetic square.
Their existence therefore supplies no missing surface theorem merely by
substitution of A.

Write $d^*u=du/u$ and
$Mf(s)=\int_0^\infty f(u)u^s\,d^*u$.  All test functions below belong to
$\mathcal D=C_c^\infty(\mathbb R_+^*,\mathbb C)$.
Put $\widetilde f(u)=u^{-1}f(u^{-1})$, and let $N$ be the distribution
of \texttt{negcrit1} in the preceding source:
\[
 N(h)=\sum_{n\geq1}\Lambda(n)h(n)
 +\int_1^\infty\frac{u^2h(u)-h(1)}{u^2-1}\,d^*u
 +\tfrac12(\log\pi+\gamma)h(1).
\]
The subtraction makes the integrand bounded at $u=1$ and the tail
integrable.  Thus this expression is defined on $\mathcal D$.
For complex functions the Hermitian extension of the real intersection
form is
\[
 I(f,g)=\tfrac12\left\{N(f*\widetilde{\overline g})+
             \overline{N(g*\widetilde{\overline f})}\right\}.
\]
This symmetrization matters: the one-sided distribution $N$ is not itself
an invariant trace on the whole multiplicative group.

\begin{proposition}[AC-G01: the full numerical Hodge inequality]
\label{AC-G01}
Under A, with the nontrivial zeta zeros written
$\rho=\tfrac12+i\gamma$ and $m_\rho$ their multiplicities,
\begin{equation}\label{eq:G01}
 2I(f,g)=Mf(0)\overline{Mg(1)}+Mf(1)\overline{Mg(0)}
 -\sum_\rho m_\rho Mf(\rho)\overline{Mg(\rho)}.
\end{equation}
The sum converges absolutely.  In particular, on
$\mathcal D_{00}=\{f:Mf(0)=Mf(1)=0\}$ the matrix
$\bigl(I(f_i,f_j)\bigr)_{i,j=1}^r$ is negative semidefinite for every finite
family.  For real $f$, this is precisely the inequality
$N(f*\widetilde f)\leq0$ requested in the Riemann--Roch strategy.
Equality on $\mathcal D_{00}$ holds exactly when $Mf(\rho)=0$ for every
nontrivial zeta zero.
\end{proposition}
\begin{proof}
The usual explicit formula in the normalization of $N$ is
\[
 N(h)+N(\widetilde h)=Mh(0)+Mh(1)-\sum_\rho m_\rho Mh(\rho).
\]
This is the same explicit formula used in the source's \S3.1, and in
CCM (math/0703392), \S7, Proposition \texttt{PosTrRHint} and equation
\texttt{terms01Tr}.  To check its normalization, the Mellin transform
of the distribution on $[1,\infty)$ is
$-\Lambda_\zeta'(s)/\Lambda_\zeta(s)$ for
$\Lambda_\zeta(s)=\pi^{-s/2}\Gamma(s/2)\zeta(s)$; adding its reflected
contour picks up the poles at $0,1$ and the negatives of the zero residues.
Thus there is no extra factor of two in this displayed identity.
Now
$M(f*\widetilde{\overline g})(s)=Mf(s)\overline{Mg(1-\overline s)}$
and
$\widetilde{(f*\widetilde{\overline g})}
 =\overline{g*\widetilde{\overline f}}$.
A makes $1-\overline\rho=\rho$, which gives \eqref{eq:G01}.
For every $B>0$, integration by parts after $u=e^x$ gives
$Mf(\tfrac12+it)=O_{f,B}((1+|t|)^{-B})$.
The zero-counting bound $N_\zeta(T)=O(T\log(T+2))$ recalled in
Connes 1999, \S I, equations (1)--(3) and the following paragraph,
proves absolute convergence.  On $\mathcal D_{00}$,
$-2I(f,f)=\sum m_\rho|Mf(\rho)|^2$.  Apply this identity to
$\sum c_if_i$ for the matrix assertion, and use that each summand is
nonnegative for the equality criterion.  If $f$ is real,
$f*\widetilde f$ is fixed by the tilde operation, so
$I(f,f)=N(f*\widetilde f)$.
\end{proof}

\begin{proposition}[AC-G02: a constructed numerical square]
\label{AC-G02}
Let $\mathcal H$ be the closure of the vectors
$\Phi(f)=(\sqrt{m_\rho}Mf(\rho))_\rho$, $f\in\mathcal D$,
in $\ell^2$ of the distinct zeros.  Under A the Hilbert space
$\mathcal H$ is well-defined.  The space
$\mathcal N=\mathbb C^2\oplus\mathcal H$, with Hermitian form
\[
 B((a,b,v),(a',b',v'))=a\overline{b'}+b\overline{a'}-
 \langle v,v'\rangle,
\]
contains the explicitly defined classes
$D_f=(Mf(0),Mf(1),\Phi(f))$, and
$B(D_f,D_g)=2I(f,g)$.
Its two coordinate classes $E=(1,0,0)$, $F=(0,1,0)$ satisfy
$E^2=F^2=0$, $B(E,F)=1$; on the orthogonal of $E+F$ the form
is negative definite.  Its minimal spectral Hilbert realization is unique
up to the isometry carrying each marked vector $\Phi(f)$ to its counterpart.
\end{proposition}
\begin{proof}
Square summability is AC-G01.  The form and the classes are given by the
displayed formulas, so their compatibility is \eqref{eq:G01}.
Orthogonality to $E+F$ means $a+b=0$; consequently
$B((a,-a,v),(a,-a,v))=-2|a|^2-\|v\|^2$, which is negative
unless the vector is zero.  For uniqueness, two dense families with the
same Gram kernel have an isometry on their finite linear spans:
$\sum c_i\Phi(f_i)\mapsto\sum c_i\Phi'(f_i)$.
Equality of squared norms proves that this map is well-defined, and its
continuous extension is onto by density.
\end{proof}

This constructs a numerical intersection space, not a surface, a category
of divisors, an integral N\'eron--Severi lattice, or a sheaf cohomology.
It forces any proposed trace-compatible geometric realization to have the
same negative form on the marked primitive correspondences.  In particular,
a model that gives one of these correspondences positive square is excluded.
It does not determine a geometric square: even at the numerical level one
may adjoin any extra negative definite space with zero degree and codegree,
leaving all the displayed marked intersections unchanged.  The quotient
by the marked subspace can have arbitrarily prescribed finite dimension.
Thus the missing construction remains the surface intersection/Riemann--Roch
and effectivity input in \S3.4, steps 2--5, of 1805.10501.  This is a
statement about the precise missing input, not a formal independence theorem.

\begin{proposition}[AC-G03: a finite single-Frobenius obstruction]
\label{AC-G03}
No finite collection of finite-dimensional vector spaces with invertible
operators $F_i$ and one fixed real $q>1$ can express the nontrivial-zero
divisor of $\zeta(s)$ as the divisor of
$\prod_i\det(1-q^{-s}F_i)^{n_i}$, $n_i\in\mathbb Z$.
\end{proposition}
\begin{proof}
Each nonzero eigenvalue $\alpha$ contributes only the arithmetic progression
$s=(\log\alpha+2\pi i k)/\log q$, $k\in\mathbb Z$,
with multiplicity bounded by the fixed dimension.  A finite product has
$O(T)$ zeros and poles, counted with multiplicity, in $|\Im s|\leq T$;
cancellation only decreases this bound.  The Riemann--von Mangoldt formula (Connes 1999, \S I, (2)--(3)) gives
$N_\zeta(T)=(T/2\pi)\log(T/2\pi)-T/2\pi+O(\log(T+2))$
for the upper half-plane, a contradiction.  This does not exclude
$\operatorname{Spec}\mathbb Z$ as a scheme or its usual prime-indexed Euler
product.  It excludes the stated finite, single-Frobenius determinant model.
\end{proof}

\subsection{The absolute twistor line and absolute cyclotomy}

In 2609.00299, Proposition 3.2 (\texttt{prop:points\_X\_K}) describes the
unfolded line over a field containing square roots of $-1$; Theorem 5.1
(\texttt{thm:twistor\_space\_scheme}) gives
$\sigma(z)=-1/\overline z$ on the folded complex line.  The absence of real
fixed points follows already from $|z|^2=-1$, with $0,\infty$ interchanged.
Proposition 7.3 (\texttt{prop:global\_sections\_frobenius}) states the
$\mathbb F_{1^2}$ global sections in its specified odd arithmetic topos.
Proposition 7.5 (\texttt{prop:twistor\_frobenius}) gives, for odd $n$,
\[
 f_n(z)=\begin{cases}z^n&n\equiv1\pmod4,\\-z^{-n}&n\equiv3\pmod4.
 \end{cases}
\]
These are source constructions independent of A.  Their direct coordinate
formulas obey $f_mf_n=f_{mn}$ and commute with $\sigma$: for
$h(z)=-z^{-1}$, all odd power maps commute with $h$, $h^2=1$, and
$f_n=h^{(1-\chi_{-4}(n))/2}\circ(z\mapsto z^n)$.
No zero-location assertion is involved in these verifications.

Likewise, 2112.08820, Theorem 2.11 (\texttt{w0s1}), identifies the
endomorphism invariant ring $W_0(\overline{\mathbb S})$ with
$\mathbb Z[\mathbb Q/\mathbb Z]$; Theorem 2.12
(\texttt{rabi10}) identifies the Frobenius/Verschiebung structure with
the integral BC-system.  These are unconditional statements in
endomorphism $K$-theory.  A supplies the zero-location theorem for the
Dirichlet factors that occur in their arithmetic realizations; it neither
changes these rings nor adds an intersection theorem to them.

The phrase ``Gaussian avatar'' has a precise safe arithmetic realization:
the Artin motive of $\operatorname{Spec}\mathbb Q(i)$, or the ordinary
scheme $\operatorname{Spec}\mathbb Z[i]$.  Its zeta is
$\zeta(s)L(s,\chi_{-4})$.  The presence of the same quadratic sign in the
absolute twistor formulas is not by itself a theorem that the zeta of the
entire absolute twistor object equals this Dedekind zeta.  Such an
identification requires a specified counting/trace formalism for that object.

\subsection{Artin--Tate objects, global zeros, and the prime 2}

Let $\varepsilon$ denote the nontrivial one-dimensional rational
representation of $C_2=\operatorname{Gal}(\mathbb Q(i)/\mathbb Q)$, and
let $\mathbb L$ denote the Tate object with good-prime Frobenius $p$.
Finite direct sums of $\mathbb L^j$ and $\varepsilon\mathbb L^j$,
$j\in\mathbb Z$, form the explicitly constructed tensor category of
finite-dimensional, integer-graded rational $C_2$-representations.
Its Grothendieck ring is
\[
 R=\mathbb Z[\mathbb L,\mathbb L^{-1},\varepsilon]/(\varepsilon^2-1).
\]
These are actual Artin--Tate motives over $\mathbb Q$: the sign summand
is cut out of $h^0(\operatorname{Spec}\mathbb Q(i))$ by the projector
$(1-c)/2$, and the other summand by $(1+c)/2$.
Over $U=\operatorname{Spec}\mathbb Z[1/2]$ their Artin part comes from
the finite \emph{\'etale} double cover $\operatorname{Spec}\mathbb Z[i,1/2]$.
The associated $\ell$-adic sheaves are the trivial and sign sheaves,
tensored with $\mathbb Q_\ell(-j)$ on the open set where $\ell$ is
invertible.  This specifies the objects and their realizations, rather
than only a polynomial that might be their trace.

\begin{proposition}[AC-G04: the complete sign--Tate construction]
\label{AC-G04}
For integers $a_j^+,a_j^-$ of finite support, the Artin--Tate class
$V=\sum_j(a_j^++a_j^-\varepsilon)\mathbb L^j$ has complete finite-prime
$L$-function
\begin{equation}\label{eq:G04}
 L(V,s)=\prod_j\zeta(s-j)^{a_j^+}L(s-j,\chi_{-4})^{a_j^-}.
\end{equation}
At an odd prime its $n$th Frobenius trace is
$\sum_j(a_j^++a_j^-\chi_{-4}(p)^n)p^{nj}$.
Under A the divisor contributed by the nontrivial zeros of the primitive
factors in \eqref{eq:G04} is supported on
$\Re s=j+\tfrac12$.  Negative multiplicities can produce poles and
cancellations; this is not a claim that a virtual class has only zeros.
Local purity is unconditional: every constituent in degree $j$ has
weight $2j$, since its good-prime eigenvalue is $p^j$ or $-p^j$.
A direct sum with distinct $j$ is graded by these weights, not pure of
one weight.
\end{proposition}
\begin{proof}
For $\Re s$ sufficiently large, use
$\exp(\sum_{n\geq1}\alpha^n p^{-ns}/n)=(1-\alpha p^{-s})^{-1}$.
At odd $p$ this gives
$\prod_j(1-p^{j-s})^{-a_j^+}
(1-\chi_{-4}(p)p^{j-s})^{-a_j^-}$.
At $2$, the sign representation has no inertia invariants and the
trivial one has one-dimensional invariants, so the factor is
$\prod_j(1-2^{j-s})^{-a_j^+}$.
Multiplication gives \eqref{eq:G04} in an absolutely convergent half-plane
and hence by meromorphic continuation.  The trace formula follows from
the stated eigenvalues.  Under A a nontrivial factor zero is
$s=j+\rho$ with $\Re\rho=1/2$; the divisor of a product is the sum of
its divisors.  Finally $|\pm p^j|=p^{(2j)/2}$ proves local purity without
using A.  Neither this computation nor A asserts one common weight for
several different $j$.
\end{proof}

\begin{proposition}[AC-G05: tensor multiplicativity and its ramified failure]
\label{AC-G05}
The assignment of local rational zeta functions of $V$ is a ring map
from $R$ into the product of rational Witt rings over odd primes,
where direct sum is Witt addition and tensor product is Witt multiplication.
It is injective.  Including the factor at $2$ defined by inertia
invariants destroys tensor multiplicativity.  Specifically,
\[
 L_2(\varepsilon,s)=1,
 \qquad L_2(\varepsilon\otimes\varepsilon,s)
 =(1-2^{-s})^{-1}.
\]
Thus the complete Euler-product assignment cannot be a ring map for this
primewise Witt multiplication.
\end{proposition}
\begin{proof}
At an odd prime the eigenvalues of a direct sum are their multiset union
and those of a tensor product are all pairwise products, which are the
stated Witt operations.  At $p=3$, the distinct eigenvalues $3^j$ and
$-3^j$ give distinct linear factors in $\mathbb Q(T)$.
Unique factorization therefore recovers every $a_j^+$ and $a_j^-$ from
that one local rational function, proving injectivity even for virtual
classes.  At $2$, inertia acts on $\varepsilon$ by $-1$, giving zero
invariants, whereas its square is the trivial representation with one
invariant.  Witt multiplication of the two unit rational functions
$1$ is still $1$, which contradicts the displayed tensor factor.
This is precisely the failure of
$(V\otimes W)^I=V^I\otimes W^I$ in this example.
\end{proof}

Over $U$ the removed factor also matters analytically:
$\zeta^{(2)}(s-j)=(1-2^{j-s})\zeta(s-j)$ has additional elementary
zeros on $\Re s=j$.  These are not nontrivial zeros of its primitive
completion and must not be relabelled violations of A or assigned the
line $j+1/2$.

\begin{proposition}[AC-G06: counts, categorical positivity, and the Tate tower]
\label{AC-G06}
In the category just constructed, a class is represented by an actual
object exactly when all $a_j^\pm\geq0$.  Nonnegative Frobenius trace
counts do not suffice: $1-\varepsilon$ has all good-prime power traces
$1-\chi_{-4}(p)^n\in\{0,2\}$, but is not the class of an object.
For every integer $k\geq0$, the object
$40h((\mathbb P^1)^{2k})(k)$ has zeta
\[
 Z_k(s)=\prod_{a=-k}^{k}\zeta(s-a)^{40\binom{2k}{k+a}}.
\]
A gives its nontrivial critical lines $\Re s=a+1/2$.
This constructs one object per word-count of the stated form; it does
not construct the multiplication or associativity constraints of a
categorification of all GCT structure constants.
\end{proposition}
\begin{proof}
The simple objects of the graded representation category are exactly
$\mathbb L^j,\varepsilon\mathbb L^j$, so their multiplicities in an actual
object are nonnegative and conversely determine such a direct sum.
The virtual class $1-\varepsilon$ has the displayed traces, but its
ordinary dimension is zero and it is nonzero in the free additive basis
of $R$, so it cannot be an actual object.
The projective-line decomposition is $h(\mathbb P^1)=1\oplus\mathbb L$.
The tensor binomial formula gives
$40\mathbb L^{-k}(1+\mathbb L)^{2k}
 =40\sum_{a=-k}^k\binom{2k}{k+a}\mathbb L^a$;
apply AC-G04.  The construction uses finite products for each $k$.
Their cohomological weights range from $-2k$ to $2k$, which grow with
$k$.  No common finite cohomological realization or compatible
multiplication has been supplied by this calculation.
\end{proof}

Here ``actual'' is deliberately used instead of the potentially ambiguous
word ``effective'': allowing $j<0$ allows inverse Tate twists, which are
not effective motives in the customary geometric convention.  Conversely,
an actual scheme can have a virtual compactly supported cohomology class:
$\mathbb G_m$ has count $q-1$ and zeta $\zeta(s-1)/\zeta(s)$.
Thus ``no inverse zeta factors'' characterizes this specified additive
semisimple cone, not all schemes with nonnegative point counts.

\begin{proposition}[AC-G07: adjoining a point is not shifting an ideal norm]
\label{AC-G07}
Disjointly adjoining a section to a scheme over $\operatorname{Spec}\mathbb Z$
adds one to each prime-power point count and multiplies its zeta by
$\zeta(s)$.  In contrast, the norm-shifted series of the supplied note,
already for ideals of $\mathbb Z$, is
\[
 \sum_{n\geq1}(n+1)^{-s}=\zeta(s)-1\qquad(\Re s>1).
\]
It is not the Hasse--Weil zeta of $\mathbb P^1_{\mathbb Z}$,
which equals $\zeta(s)\zeta(s-1)$.
\end{proposition}
\begin{proof}
In the exponential definition of each local zeta, adding one to every
point count adds $\sum_{n\geq1}p^{-ns}/n$ to its logarithm and hence
multiplies the factor by $(1-p^{-s})^{-1}$.
The norm-shifted formula follows by replacing $n+1$ by $m\geq2$.
The projective line has count $p^n+1$, so the same exponential calculation
gives its two factors.  For example, the shifted series has zero
Dirichlet coefficient at $1$, whereas a normalized Euler product and
$\zeta(s)\zeta(s-1)$ have coefficient one.  Thus the constructions differ
before any question about zeros arises.  A's theorem on primitive
Dirichlet functions cannot be substituted for a zero theorem on
$\zeta(s)-1$ by this counting analogy.
\end{proof}

\subsection{A precise split-zero construction and its scope}

For a commutative unital ring $R$, let $G(R)=R\amalg\{\tau\}$.
Extend ring addition by making $\tau$ its identity and ring multiplication
by making $\tau$ absorbing.  This is a commutative semiring, whose zero
is $\tau$ and whose unit is $1_R$; write $e=0_R$ for the supported zero.
The elementary semiring identities are exactly those in the supplied
split-zero notes.  They must not be interchanged with statements about
monoid algebras that forget ring addition.

\begin{proposition}[AC-G08: the sign and the supported zero, corrected]
\label{AC-G08}
Multiplication by $\varepsilon=-1_R$ fixes
$\{\tau\}\cup R[2]$.  If $R$ has no nonzero $2$-torsion, this fixed set
has Boolean operations on $\{\tau,e\}$, but its inclusion in $G(R)$
is not unital when $R\ne0$.  For $R=\mathbb Z$,
$\varepsilon^{2m}=1_R$ is \emph{not} fixed by multiplication by
$\varepsilon$.  Thus the claim in the supplied split-zero note \S8.2
that even powers land in the Boolean fixed locus is false.
There is a surjective semiring map
\[
 \mathbb N[\varepsilon]/(\varepsilon^2-1)\longrightarrow G(\mathbb Z),
 \quad (a,b)\longmapsto
 \begin{cases}\tau&(a,b)=(0,0),\\a-b& a+b>0.
 \end{cases}
\]
The supported zero is the image of all positive multiples of
$1+\varepsilon$, not equality with that class in the representation
semiring.  The genuine Artin object $1\oplus\varepsilon$ has dimension
$2$ and zeta $\zeta_{\mathbb Q(i)}(s)$.
\end{proposition}
\begin{proof}
For a supported element, $-r=r$ is exactly $2r=0$; $\tau$ is fixed.
On $\{\tau,e\}$ addition is maximum and multiplication is minimum after
ordering $\tau<e$.  Its multiplicative unit is $e$, not $1_R$.
In $G(\mathbb Z)$, $\varepsilon\cdot1_R=-1_R\ne1_R$, proving the
refutation.  Addition of two nonzero coefficient pairs remains nonzero;
their product has total coefficient $(a+b)(c+d)>0$ and signed value
$(a-b)(c-d)$.  This proves the map respects both operations, including
zero.  Every integer is attained by a nonzero pair; zero signed value
occurs precisely at $(a,b)=(m,m)$, $m>0$.
In the representation semiring, $(1+\varepsilon)^2=2(1+\varepsilon)$,
whereas $e^2=e$, so the classes cannot be identified there.
The arithmetic assertion follows from the trivial/sign decomposition
of the permutation representation on $\mathbb Q(i)$ and AC-G04.
\end{proof}

Multiplication by $-1$ is an action on the underlying additive object;
it is not in general a unital semiring automorphism.  Its Boolean fixed
set is therefore not a subalgebra over the same unit.  At inert primes,
$1\oplus\varepsilon$ has zero odd Frobenius traces and trace $2$ at even
powers.  At split primes it has trace $2$ at every power.  The inert trace
zero is the relevant cancellation; it is not the zero object.

\begin{proposition}[AC-G09: sheafification on a fixed space]
\label{AC-G09}
Let $\mathcal O$ be a sheaf of unital rings on a topological space $X$.
The sectionwise presheaf $U\mapsto G(\mathcal O(U))$ is generally not a
sheaf.  Its associated sheaf has the explicit description
\[
 \mathcal G(U)=\{(V,r):V\subset U\text{ open and closed},\ r\in\mathcal O(V)\}.
\]
The bottom section has $V=\varnothing$.  Addition uses the union of
supports and the original ring sum on their intersection; multiplication
uses their intersection and the original ring product.
Its stalk at $x$ is $G(\mathcal O_x)$.  If $\mathcal O$ has no
$2$-torsion, the fixed sections for the sign over $U$ are precisely the
open-and-closed subsets of $U$, with Boolean union and intersection.
This construction changes the structure sheaf on $X$, not its underlying
topos or points.
\end{proposition}
\begin{proof}
A sheaf must take the empty set to a singleton, while
$G(\mathcal O(\varnothing))=G(0)$ has two elements.  Even after repairing
that value, disjoint nonempty opens reveal the defect: a supported
section on one component and $\tau$ on the other glue in a sheaf but
are not any section of $G(\mathcal O(U))$.
In the displayed construction, compatible local support indicators glue
to a locally constant indicator $U\to\{0,1\}$, hence an open-and-closed
$V$.  The supported ring sections glue uniquely by the sheaf property of
$\mathcal O$.  This proves descent.  The prescribed algebra operations
are defined on the disjoint open-and-closed pieces and satisfy the
semiring axioms stalkwise.  Every displayed section is locally either
$\tau$ or a section of $\mathcal O$, so it is locally in the image of
the presheaf.  Two presheaf sections with equal images have equal germs.
These facts prove the asserted associated-sheaf property and identify
the stalks.  For the sign, equality requires $2r=0$ on $V$; under the
hypothesis this means $r=0$, leaving just $V$.
\end{proof}

The relevant general topos input is SGA~4, Expos\'e II, Theorem 4.1(1):
the associated-sheaf functor is exact and commutes with colimits.
At the presheaf level, $G(A\times_C B)=G(A)\times_{G(C)}G(B)$: a compatible
pair is either $(\tau,\tau)$ or a supported pair in $A\times_C B$.
Sheafification therefore preserves this connected gluing identity.
The terminal object and disconnected covers require the correction just
proved; preservation of fibre products alone never proved them.

\begin{proposition}[AC-G10: a distinct construction that adds a generic point]
\label{AC-G10}
In the ordinary ideal spectrum of the semiring $G(R)$ the prime ideals
are exactly
\[
 \{\tau\},\qquad \{\tau\}\amalg\mathfrak p\quad
 (\mathfrak p\in\operatorname{Spec}R).
\]
Thus it is $\operatorname{Spec}R$ with a new open generic point $\xi$;
the closed complement is $V(e)=\operatorname{Spec}R$.
More generally, for a ringed space $(X,\mathcal O)$, define $X^+=X\amalg\{\xi\}$
with open sets $\varnothing$ and $U^+=U\cup\{\xi\}$ for $U\subset X$
open.  Set $\mathcal O^+(U^+)=G(\mathcal O(U))$ and give the empty open
the terminal semiring.  This is a sheaf of semirings.  For affine $X$
it is the semiring spectrum just described, with its usual localization
structure sheaf.  The stalk at $\xi$ is $\mathbb B$.
Every unital semiring map $G(R)\to A$ to a ring factors uniquely through
$p:G(R)\to R$, $p(\tau)=0$, $p(r)=r$.
Consequently its ring-valued points are the old ring-valued points;
$|G(\mathbb F_q)|=q+1$ is a cardinality of elements, not the number of
$\mathbb F_q$-points of this affine spectrum.
\end{proposition}
\begin{proof}
An ideal containing any supported $r$ contains $er=e$, and thereafter
its supported part is a ring ideal, since multiplication by $-1_R$
provides additive inverses.  An ideal with no supported element is
$\{\tau\}$.  Products of two supported elements stay supported, even
when their ring product is zero; hence $\{\tau\}$ is prime.  The other
prime conditions are exactly those for the ring ideal $\mathfrak p$.
The asserted topology follows from $D(e)=\{\xi\}$ and
$D(r)=D_R(r)\cup\{\xi\}$ for supported $r$.
For a nonempty cover by $U_i^+$, the overlaps all contain $\xi$ and the
restriction to $\{\xi\}$ records whether a section is supported or
$\tau$.  Compatibility forces that choice to agree on the whole cover.
In the supported case the ring sections glue over the $U_i$; in the
other case they are all $\tau$.  This proves the sheaf assertion.
Localization at a supported $r$ gives $G(R_r)$ by the usual fraction
equivalence relation, including $r=0$, when both sides are $\mathbb B$.
This identifies the affine structure sheaves and gives the stalk at
$\xi$; equivalently, inverting $e$ forces $e=1$ and then forces all
supported elements to equal $1$.
For a map $\phi$ to a ring, $\phi(e)+\phi(e)=\phi(e)$ gives $\phi(e)=0$.
The restriction to $R$ is therefore a ring homomorphism, and
$\phi(\tau)=0$ determines the factorization.  In particular, over
$\mathbb F_q$ there is one such $\mathbb F_q$-valued point when the
supported constants are required to be fixed, not $q+1$.  This is a
condition on the induced map $\mathbb F_q\to\mathbb F_q$ after ring
completion; it does not presume a unital ring map
$\mathbb F_q\to G(\mathbb F_q)$, which does not exist.
\end{proof}

AC-G09 and AC-G10 are two different, fully specified constructions.
The former answers a fixed-topos sheaf problem; the latter answers an
ordinary semiring-spectrum problem.  Neither identifies that spectrum
with the Connes--Consani spherical spectrum without a comparison of
categories and structure sheaves.  In the pointed-monoid version one can
forget addition and adjoin a new basepoint to a spatial monoid.  The new
generic residue is then the two-element pointed monoid $\mathbb F_1$;
calling it the Boolean \emph{semiring} requires the additive structure
used above.

\begin{proposition}[AC-G11: the power-map and Hopf type restrictions]
\label{AC-G11}
On the pointed multiplicative monoid underlying $G(R)$, the maps
$x\mapsto x^n$, with $\tau\mapsto\tau$, are monoid endomorphisms and
preserve a nontrivial base sign $\varepsilon=-1$ exactly when $n$ is odd.
They need not be semiring endomorphisms.  For $R=\mathbb Z$, already
$n=3$ fails additivity.
For an unpointed commutative monoid $M$, its spherical monoid algebra
has the group-like diagonal over $\mathbb F_1$; when $M=(R,\cdot)$ for a nonzero ring $R$, this
has no antipode because the supported zero is not invertible.
Moreover the diagonal on \emph{all} elements cannot descend to an
$\mathbb F_{1^2}$-linear coproduct after identifying the element
$[-1_R]$ with the base sign, unless that sign has collapsed.
\end{proposition}
\begin{proof}
Commutativity gives $(xy)^n=x^ny^n$.  The sign condition is
$(-1)^n=-1$.  In $G(\mathbb Z)$,
$(1+1)^3=8\ne2=1^3+1^3$, proving failure of semiring additivity.
Over $\mathbb F_1$, the coproduct sends $[m]$ to $[m]\otimes[m]$,
and the counit sends every supported monoid element to $1$.
An antipode would give an inverse to the group-like supported zero,
which is impossible in the multiplicative monoid of a nonzero ring:
$0r=0\ne1$.  After the stated signed identification, the same diagonal
would send $\varepsilon=[-1_R]$ to
$\varepsilon\otimes\varepsilon=1$, whereas base linearity requires
$\varepsilon\otimes1=\varepsilon$.  The two are unequal when the sign
is nontrivial.  Thus one must keep a separate spatial monoid, as in the
free signed construction of 2609.00299, \S2.2, or supply a different
coproduct and verify its relations.  The displayed diagonal is not such
a coproduct on the identified signed ring monoid.
\end{proof}

The statement about the diagonal is deliberately restricted to that
specified map; it does not assert that no other bialgebra structure can
exist.  The two spherical objects
$\mathbb S[(R,\cdot)_+]$ and $HG(R)$ differ: the first remembers monoid
multiplication, the second also remembers ring addition.  The source
notes' proposed extension must choose one before Frobenius and Hopf claims
can be transferred.

\begin{proposition}[AC-G12: a global realization of the trace-unimodality example]
\label{AC-G12}
Let $Y=\operatorname{Res}_{\mathbb Q(i)/\mathbb Q}\mathbb P^1$ and
$X=Y^2\times\mathbb P^2$.  The centered motive $h(X)(3)$ has graded
dimensions and inert-prime normalized trace coefficients respectively
\[
 d=(1,5,11,14,11,5,1),\qquad c=(1,1,3,2,3,1,1),
\]
in Tate degrees $j=-3,-2,\ldots,3$.  Its multiplicities are
\[
 a^+=(1,3,7,8,7,3,1),\qquad a^-=(0,2,4,6,4,2,0).
\]
At an inert good prime $p$, its first normalized trace is
$(p+p^{-1})^2(p+1+p^{-1})$.
Its complete motivic $L$-function is AC-G04 with these multiplicities,
so A proves its nontrivial zero lines.  The graded dimensions are unimodal
while the inert trace coefficients are not.
\end{proposition}
\begin{proof}
After the quadratic base extension, $Y$ becomes $\mathbb P^1\times\mathbb P^1$.
Galois acts on the two degree-two classes by exchanging them, so their
representation is $1\oplus\varepsilon$.  The degree-zero and top classes
are invariant.  The projective bundle formula and this explicit swap give
\[
 h(Y)=1\oplus(1\oplus\varepsilon)\mathbb L\oplus\mathbb L^2.
\]
This is an isomorphism of rational Chow motives, not just an equality
of traces: the split projective-bundle correspondences are equivariant
for the exchange of factors and descend by finite Galois descent;
restriction followed by transfer is multiplication by $2$, and the
rational projectors $(1\pm\sigma)/2$ separate the two Artin summands.
K\"unneth and $h(\mathbb P^2)=1\oplus\mathbb L\oplus\mathbb L^2$ now give
\[
 [h(X)(3)]=\mathbb L^{-3}
 (1+(1+\varepsilon)\mathbb L+\mathbb L^2)^2
 (1+\mathbb L+\mathbb L^2).
\]
Evaluation at $\varepsilon=1$ gives $d$; evaluation at
$\varepsilon=-1$ gives $c$; resolving the two characters gives
$a^+=(d+c)/2$ and $a^-=(d-c)/2$, exactly as displayed.
At an inert prime $\varepsilon$ has Frobenius $-1$, so the trace is
$p^{-3}(1+p^2)^2(1+p+p^2)$, which is the asserted expression.
The descent has a smooth model over $\mathbb Z[1/2]$, but at the ramified
prime we mean the inertia-invariant motivic factor of AC-G04, not an
unverified point count of a chosen bad integral model.  The elementary
inequalities in the two displayed coefficient sequences prove the last
claim: $d$ rises to its central coefficient and then falls, whereas
$c$ contains the strict dip $3,2,3$.
\end{proof}

\subsection{What is not settled by these constructions}

The exact remaining geometric gates are the trace-compatible
intersection theory and surface Riemann--Roch theorem on the relevant
square, together with an effectivity/descent theorem in that category.
A proves their numerical target and the constructions above give its
minimal analytic realization.  They give no functor from that Hilbert
space to algebraic cycles, no algebraicity of projectors, and no
polarization on a putative global arithmetic motive.  Accordingly they
do not prove the standard conjectures, a general motivic RH, or the
cubic Kummer window bound.

The supplied counting notes cite rank-three tables, the ``dimension
ring'' certificates, and the GCT runs by IDs F35, F41, F42, F45 and
P-F-E3-01.  The theorem AC-G06 applies to each coefficient family for
which the required Artin--Tate realization is actually given.  It does
not re-certify those absent GCT tables or prove their associativity.
An associative numerical based ring alone is also not a construction of
associators satisfying the pentagon in a tensor category; that extra
categorical realization is a separate gate.  No theorem about it is
inferred from A here.

For Grothendieck's frameworks, the precise established connections are:
AC-G01--G02 supply the numerical Hodge-index form modeled on
\emph{Sur une note de Mattuck--Tate}, J. reine angew. Math. 200 (1958),
208--215, \S1 (Hodge-index theorem and product-of-curves inequality
(1.8)), with the surface Riemann--Roch argument in \S2, (2.1).
AC-G04--G06 are constructions in Artin--Tate motives and their
$\ell$-adic realizations, with exterior powers and hence the usual
representation $\lambda$-ring structure on $R$; inertia invariants
explain AC-G05.  AC-G09 is a sheaf construction under SGA~4,
Expos\'e II, Theorem 4.1; AC-G10 is a semiring analogue and is not an
EGA scheme since its affine coefficient objects are not rings.
The RR duality paper explicitly compares its Theorem 5.3 with the
Grothendieck--Serre letter dated 15 December 1955; that comparison is
not a new theorem identifying the categories.

The Montpellier catalogue verifies the relevant archival subjects:
cote 2 (Riemann--Roch notes), cote 10 (tensor categories), cote 11
(correspondences), cote 14 (Hodge--Tate--Lefschetz--Weil cycles), cote 16
(correspondences and Poincar\'e--Hodge algebras), and cote 19 (topos).
No theorem-level scan from those cotes was successfully retrieved in
this workstream.  These are catalogue mappings only.  An exact
cote-and-scan attribution remains an explicit bibliography gate; it
must not be manufactured from the subjects of the folders.

\section{The relative ray-class rung and the cubic arithmetic gate}
\label{sec:kummer-lambda}

We use the newest supplied ALP\_MACHINE checkpoint, dated
2026-11-21T2304Z. Its AM-H01--AM-H04 identify the cubic representation,
its full Euler factors, the obstruction to a cyclotomic or CM overfield,
and the analytic window implication. They are inputs here, with the
checks recorded separately. In particular, we do not announce those
calculations as new findings. Put
\[
 K=\mathbb Q(\omega),\quad \omega^2+\omega+1=0,
 \quad E_p=K(\sqrt[3]p),\quad F_p=\mathbb Q(\sqrt[3]p),
 \quad L_p(s)=L_K(s,\chi_{p,3}).
\]
The character is understood primitively, with every local factor
defined on inertia invariants. AM-H01 gives
\begin{equation}
 L_p(s)=\zeta_{F_p}(s)/\zeta(s),\qquad
 \zeta_{E_p}(s)=\zeta_K(s)L_p(s)^2.                 \label{eq:kummer-zeta}
\end{equation}
These are identities of meromorphic functions, including the bad primes.

\subsection{The precise meaning of integral Lambda descent}

\paragraph{Known classification, with its actual scope.}
Borger--de Smit, \emph{Galois theory and integral models of
$\Lambda$-rings}, Theorem~1.1 (source label \texttt{thm-A}),
classifies finite \'etale \emph{algebras}, with their commuting
Frobenius lifts and an integral model, by finite continuous sets
for the multiplicative monoid $\widehat{\mathbb Z}$.
The paragraph immediately following that theorem explicitly says
that such an algebra cannot be a field other than $\mathbb Q$.
Theorem~1.2 (\texttt{thm-B}) identifies the maximal integral model
of $\mathbb Q[t]/(t^n-1)$ as $\mathbb Z[t]/(t^n-1)$.
For a number field base, Borger--de Smit, \emph{Lambda actions of
rings of integers}, Theorems~1.1 and~1.2 (\texttt{local-thm},
\texttt{global-thm}), replace the monoid by the Deligne--Ribet
monoid. The source has the typographical phrase ``Then $K$ has an
integral model'' in both theorem statements; their quantified
algebra and proofs make clear that $E$ is intended.

\begin{proposition}[AC-K01: the corrected rational equivalence]
\label{AC-K01}
The following zero-location assertions are equivalent:
\begin{enumerate}
\item input theorem A;
\item RH for every cyclotomic Dedekind zeta function;
\item RH for the product of Dedekind zeta functions of the field
factors of every finite \'etale $\mathbb Q$-algebra admitting an
integral $\Lambda_{\mathbb Z}$-model.
\end{enumerate}
Every abelian number field occurs as a field factor of an algebra
in (3). The assertion that the field itself has such a model is
false unless it is $\mathbb Q$.
\end{proposition}
\begin{proof}
Abelian Artin factorization expresses a cyclotomic Dedekind zeta
function as a product of primitive Dirichlet $L$-functions.
Every factor is holomorphic in $0<\Re s<1$, so an off-line zero
of one factor cannot be canceled by a pole of another. This proves
both (1)$\Rightarrow$(2) and (2)$\Rightarrow$(1): each primitive
Dirichlet character occurs at its conductor. Borger--de Smit's
classification makes each field factor in (3) abelian over
$\mathbb Q$, proving (1)$\Rightarrow$(3) by the same factorization.
For the converse use $\mathbb Q[t]/(t^n-1)$, whose factors include
$\mathbb Q(\mu_n)$; the same no-cancellation argument proves (2).

For the additional assertion, write the abelian field as
$F=\mathbb Q(\mu_n)^H$. The invariant subring
$\bigl(\mathbb Z[t]/(t^n-1)\bigr)^H$ is finite torsion free over
$\mathbb Z$ and is preserved by all maps $t\mapsto t^p$.
The Frobenius congruence remains valid in the invariant ring:
if an invariant element is $p$ times an element of the larger
torsion-free ring, that quotient is invariant too. Its generic
algebra has $F$ as the factor from primitive $n$th roots. The final
claim is exactly the cited paragraph after Theorem~1.1; the
relative elementary obstruction below also proves it.
\end{proof}

\begin{proposition}[AC-K02: a field obstruction at ramification]
\label{AC-K02}
Let $L/K$ be a finite field extension. If $L$ admits an integral
$\Lambda_{\mathcal O_K}$-model, it is unramified at every finite
prime of $K$. In particular no $E_p$ above admits such a model as
a field, even though $E_p/K$ is cyclic. This refutes the literal
field version of item~4 of the initial first-pass note.
\end{proposition}
\begin{proof}
Let $R\subset L$ be the finite flat integral model. Every
$K$-endomorphism $\psi_{\mathfrak p}$ of $L$ is an automorphism
of finite order. Since it preserves $R$, its inverse, a positive
power of itself, preserves $R$ too. Thus its reduction modulo
$\mathfrak p$ is an automorphism. By the defining congruence that
reduction is $x\mapsto x^{N\mathfrak p}$. A finite-dimensional
algebra over a finite field has injective Frobenius only if it
is reduced: a nonzero nilpotent would be killed by an iterate.
A reduced finite algebra over a perfect field is \'etale.
Hence $R$ is finite flat with \'etale geometric fibres, so it is
finite \'etale over $\mathcal O_K$. Consequently its fraction
field is unramified at every finite place.

For $p\ne3$, the valuation of $p$ at each prime of $K$ above $p$
is one; $X^3-p$ is Eisenstein there and $E_p/K$ is ramified.
For $p=3$, at $\mathfrak l=(1-\omega)$ the valuation of $3$ is
two, so the valuation of its cube root is $2/3$; the ramification
index is divisible by three. This also excludes that case.
\end{proof}

\paragraph{Geometric translation.}
AC-K01 concerns generic fibres of finite flat schemes over
$\operatorname{Spec}\mathbb Z$ with commuting Frobenius lifts.
The field factors need not be stable under these lifts.
AC-K02 is the distinction between a ramified cover and a finite
\'etale cover of $\operatorname{Spec}\mathcal O_K$, not an
obstruction to a rank-one lisse sheaf on an open subscheme.
The Galois-category framework is SGA~1, expos\'e~V, Theorem~4.1;
the arithmetic classification is the cited Borger--de Smit theorem,
not a theorem stated in SGA~1.

\subsection{An actual finite host, and its integral construction}

\begin{proposition}[AC-K03: the minimal-conductor relative host]
\label{AC-K03}
Let $\mathfrak f$ be the conductor of $\chi_{p,3}$ and let
$d(\mathfrak f)$ denote the number of integral-ideal divisors
of $\mathfrak f$. There is a finite \'etale $K$-algebra with
integral $\Lambda_{\mathcal O_K}$-model whose underlying algebra is
\begin{equation}
 B_{p}=E_p\times K^{d(\mathfrak f)-1}.              \label{eq:minhost}
\end{equation}
Its Frobenius lifts and an algorithm terminating at its maximal
integral model are constructed below. Its Dedekind zeta product is
\begin{equation}
 Z(B_p,s)=\zeta_K(s)^{d(\mathfrak f)}L_p(s)^2.       \label{eq:minhost-zeta}
\end{equation}
Under A, RH for $B_p$ is equivalent to RH for $L_p$.
Thus existence of this relative model does not remove the cubic
zero-location question.
\end{proposition}
\begin{proof}
Write $G=\operatorname{Cl}_{\mathfrak f}(K)$ and let
$H=\ker(G\twoheadrightarrow\operatorname{Gal}(E_p/K))$.
The finite set $S=\operatorname{DR}_{\mathfrak f}/H$ is a set
of orbits for multiplication by $H$. Commutativity makes
multiplication by every ideal a well-defined action on $S$.
It commutes with the Galois action through Artin reciprocity.
Set
\[
 B_p=\{f:S\longrightarrow\overline K:
                 f(gs)=g f(s)\text{ for }g\in G_K\},\qquad
 (\psi_{\mathfrak a}f)(s)=f(\mathfrak a s).
\]
This is an explicit finite \'etale algebra with commuting
$K$-algebra endomorphisms. The Deligne--Ribet factorization
criterion gives its integral model.

The decomposition
$\operatorname{DR}_{\mathfrak f}=
 \coprod_{\mathfrak d\mid\mathfrak f}
 \operatorname{Cl}_{\mathfrak f/\mathfrak d}(K)$
is stated immediately before Borger--de Smit's global theorem;
Yalkinoglu's equation \texttt{arithalg} gives the corresponding
product of ray class fields. After quotienting by $H$, the field
at $\mathfrak d$ is
$E_p\cap K_{\mathfrak f/\mathfrak d}$, by ordinary Galois
correspondence and the surjection of ray class groups.
For $\mathfrak d=1$ it is $E_p$. For $\mathfrak d\ne1$ it is
$K$: a degree-three extension has no proper intermediate field,
and containment in the smaller ray class field would contradict
minimality of its conductor. This proves \eqref{eq:minhost}.
Equation \eqref{eq:kummer-zeta} proves \eqref{eq:minhost-zeta}.
All factors are holomorphic in the open critical strip, so A and
the same no-cancellation argument prove the last assertion.

Here is a terminating integral construction, to avoid leaving
``has a model'' as merely an existential instruction. Put
$R_0=\mathcal O_{E_p}\times\mathcal O_K^{d(\mathfrak f)-1}$,
and let $T$ be the finite set of primes dividing $\mathfrak f$.
Recursively compute the subring
\begin{equation}
 R_{n+1}=\{x\in R_n:
 \psi_{\mathfrak p}(x)-x^{N\mathfrak p}
       \in\mathfrak p R_n\text{ for all }\mathfrak p\in T\}.
                                                               \label{eq:integral-iteration}
\end{equation}
It is an $\mathcal O_K$-subalgebra: modulo $\mathfrak p$ the
power map is additive, multiplicative and $\mathcal O_K/\mathfrak p$-linear.
All $\psi$ preserve $R_0$, because they preserve integrality.
Their commutativity proves by induction that they preserve every
$R_n$: apply $\psi_{\mathfrak q}$ to each congruence in
\eqref{eq:integral-iteration}. Each step is a finite computation
of kernels in $R_n/\mathfrak p R_n$, hence of a finite-index
submodule. Outside $T$ the original Frobenius congruences hold
in $R_0$, and each index $[R_0:R_n]$ is supported on $T$.
Consequently they also hold in $R_n$, since localization at
$\mathfrak q\notin T$ leaves $R_n$ unchanged.

By the classification there is some integral model $R$ of full
rank inside $R_0$. Its Frobenius congruences show inductively
that $R\subset R_n$ for every $n$. There are only finitely many
submodules between $R$ and $R_0$. The descending sequence
therefore reaches a fixed point. At that point all congruences
hold at all primes, so it is an integral model. The same argument
for any integral model shows it is the largest one. This proves
termination without asserting a uniform bound on its running
time or its index. Precisely such quantitative bounds would be
additional information in an arithmetic window argument.
\end{proof}

\begin{proposition}[AC-K04: the explicit host for $p=269$]
\label{AC-K04}
For $p=269$, the conductor of $\chi_{269,3}$ is the prime
$\mathfrak p=(269)$ of $\mathcal O_K$, of norm $269^2$.
Let $\mathfrak P$ be the unique prime of $E=E_{269}$ above
$\mathfrak p$. An explicit maximal integral model of
$B=E\times K$ is
\begin{equation}
 R=\{(x,y)\in\mathcal O_E\times\mathcal O_K:
                  x\bmod\mathfrak P=y\bmod\mathfrak p\}.
                                                               \label{eq:269-model}
\end{equation}
Its lifts are
\[
 \psi_{\mathfrak p}(x,y)=(y,y),\qquad
 \psi_{\mathfrak q}(x,y)=
       (\operatorname{Frob}_{\mathfrak q}x,y)
                     \quad(\mathfrak q\ne\mathfrak p),
\]
where arithmetic Frobenius is used. The index
$[\mathcal O_E\times\mathcal O_K:R]$ is $269^2$.
\end{proposition}
\begin{proof}
The prime is inert in $K$ since $269\equiv2\pmod3$.
The polynomial $X^3-269$ is Eisenstein at $\mathfrak p$, so
$E/K$ is totally and tamely ramified of degree three there.
All nontrivial characters of this tame inertia group have
conductor exponent one. The polynomial discriminant restricts
other possible ramification to the prime above three.

Since $269=-1+9\cdot30$, the equation
$t-3t^2+3t^3=30$ has a solution in $\mathbb Z_3$:
its derivative is one modulo three, and Hensel lifting applies.
Then $(-1+3t)^3=269$. As $\mu_3\subset K_3$, all three roots
are in $K_3$, so $E/K$ splits at the prime above three.
This proves the conductor assertion. The ray-class exact
sequence and class number one of $\mathbb Z[\omega]$ give
\[
 \operatorname{Cl}_{\mathfrak p}(K)
   \simeq\mathbb F_{269^2}^{\times}/\mu_6,
       \qquad |\operatorname{Cl}_{\mathfrak p}(K)|=12060.
\]
The group is cyclic and has one quotient of order three. Thus
the quotient Deligne--Ribet set in AC-K03 is a free $C_3$-orbit
and one fixed point. A prime away from $\mathfrak p$ acts on
the orbit by its Artin symbol; $\mathfrak p$ sends every point
to the fixed point. This gives exactly the displayed lifts.

The congruence condition in \eqref{eq:269-model} is stable under
these lifts, and the lifts commute. Away from $\mathfrak p$,
the algebra of integers is \'etale and the specified Frobenius
automorphism lifts residue Frobenius. Since the index of $R$
is supported at $\mathfrak p$, the same congruence holds in $R$.
At $\mathfrak p$ put $q=269^2$ and $x=y+t$ with
$t\in\mathfrak P$. We have $\mathfrak p\mathcal O_E=\mathfrak P^3$.
Every intermediate binomial coefficient of $(y+t)^q$ is
divisible by $269$, and $q\ge4$, so
\[
                 x^q-y^q\in\mathfrak P^4
                                      =\mathfrak p\mathfrak P.
\]
Also $y-y^q\in\mathfrak p$. Therefore
\[
 \frac{\psi_{\mathfrak p}(x,y)-(x,y)^q}{269}
 =\left(\frac{y-x^q}{269},\frac{y-y^q}{269}\right)\in R,
\]
because the difference of its coordinates belongs to
$\mathfrak P$. This is the required Frobenius congruence.
The residue-difference map from the maximal order onto
$\mathbb F_{269^2}$ has kernel $R$, proving its index.

Finally any integral model lies in the maximal order. Its
Frobenius congruence at $\mathfrak p$, reduced in the first
coordinate modulo $\mathfrak P$, forces $y=x^q=x$ in
$\mathbb F_{269^2}$. Hence it is contained in $R$, which proves
maximality. No bound on a Selmer group or on H04.4 has been
deduced from this exact index calculation.
\end{proof}

\subsection{Which Hecke sectors A actually supplies}

\begin{proposition}[AC-K05: norm sectors and infinity types]
\label{AC-K05}
For a primitive Dirichlet character $\eta$, the Hecke character
$\eta\circ N_{K/\mathbb Q}$ satisfies
\begin{equation}
 L_K(s,\eta\circ N)=L_{\mathbb Q}(s,\eta)
                    L_{\mathbb Q}(s,\eta\chi_{-3}),             \label{eq:norm-sectors}
\end{equation}
with primitive Artin local factors on both sides. A therefore
proves its critical-line assertion, and that of each unitary
norm twist. Every finite-order character of $G_K$ invariant
under complex conjugation is of this form.

A unitary Hecke character at the complex place has the form
$(z/|z|)^m|z|^{2it}$, with $m\in\mathbb Z$. Its finite component
on $\mu_6$ must be $\epsilon\mapsto\epsilon^{-m}$.
If the finite component is unramified everywhere, $6\mid m$;
ramification allows other $m$. The algebraic Hecke characters
of CM elliptic curves have nonzero angular type after unitary
normalization and are not norm sectors in \eqref{eq:norm-sectors}.
\end{proposition}
\begin{proof}
Induction of the restriction of the one-dimensional
$G_{\mathbb Q}$-representation $\eta$ is
$\eta\otimes(1\oplus\chi_{-3})$. Artin induction, including
inertia invariants, gives \eqref{eq:norm-sectors}. Tensoring by
the unitary norm character shifts $s$ vertically and leaves
the real critical line unchanged. If a finite-order character
$\xi$ of $G_K$ is conjugation invariant, the exact sequence
$G_K\subset G_{\mathbb Q}$ is split by an actual complex
conjugation $c$, with $c^2=1$. The formula
$\widetilde\xi(gc^e)=\xi(g)$, $e=0,1$, defines a continuous
finite-order character of $G_{\mathbb Q}$. By Kronecker--Weber
it is a Dirichlet character, and reciprocity identifies its
restriction with the norm pullback. Conversely a restriction
is conjugation invariant.

The continuous unitary characters of $\mathbb C^\times$
are exactly the displayed characters, by
$\mathbb C^\times\simeq\mathbb R_{>0}\times S^1$.
Triviality on principal global units gives the assertion on
$\mu_6$. A finite unramified component is trivial on all
finite local units, whence $6\mid m$. Finally complex norm
has no angular dependence, whereas a CM elliptic $H^1$
has types $(1,0)$ and $(0,1)$; its unitary characters have
angular exponents $1$ and $-1$ up to the reciprocity convention.
Thus they cannot be obtained by the specified norm pullback.
\end{proof}

\begin{proposition}[AC-K06: CM base change does not repair the slot]
\label{AC-K06}
Let $M$ be any CM number field containing $K$. The restriction
of $\chi_{p,3}$ to $G_M$ is nontrivial and cannot become the
restriction of a Dirichlet character of $G_{\mathbb Q}$.
Twisting by such a restriction does not change this obstruction.
\end{proposition}
\begin{proof}
AM-H03 proves that no CM field contains $E_p$.
Since $[E_p:K]=3$, its intersection with $M$ is $K$, so the
restricted character cuts out $E_pM/M$ and remains nontrivial.
If it were the restriction of a Dirichlet character of conductor
$n$, this field would lie in $M(\mu_n)$. A compositum of CM
fields is CM: in a common normal closure, complex conjugation
is central and restricts to the CM involution on each factor;
its fixed field is totally real and its nontrivial extension
has degree two. Thus $M(\mu_n)$ would be a CM field containing
$E_p$, contrary to AM-H03. A putative equality after a
Dirichlet twist can be divided by that twist and gives the
same contradiction.
\end{proof}

\subsection{The Bost--Connes realization and the missing positivity}

\paragraph{Known constructions.}
Laca--Larsen--Neshveyev, Theorem~2.1 (\texttt{mainthm}),
constructs and classifies the KMS states of the number-field
system, with the partition function $\zeta_K$ and the explicit
state formula \texttt{eKMSext}. Yalkinoglu, Theorems~A and~B
(\texttt{ThmA}, \texttt{ThmB}), identifies its measured
endomotive and constructs its arithmetic subalgebra for every
number field. These were unconditional. For our $K$, A makes
the nontrivial zeros of its partition function
$\zeta_K=\zeta L(\,\cdot\,,\chi_{-3})$ lie on the critical
line. It does not newly construct this dynamical system.

\begin{proposition}[AC-K07: a character observable is not a positive trace]
\label{AC-K07}
On $\ell^2(I_K)$, where $I_K$ is the set of nonzero integral
ideals, set $H e_{\mathfrak a}=\log N\mathfrak a\,e_{\mathfrak a}$.
For a finite-order Hecke character $\chi$ of conductor
$\mathfrak f$, extend $\chi(\mathfrak a)$ by zero when
$(\mathfrak a,\mathfrak f)\ne1$ and define
$D_\chi e_{\mathfrak a}=\chi(\mathfrak a)e_{\mathfrak a}$.
For $\Re s>1$,
\[
 \operatorname{Tr}(e^{-sH})=\zeta_K(s),\qquad
 \operatorname{Tr}(D_\chi e^{-sH})=L_K(s,\chi),\qquad
 |L_K(s,\chi)|\le\zeta_K(\Re s).
\]
For nontrivial cubic $\chi$, $D_\chi$ is not positive and
$\Re D_\chi$ has eigenvalue $-1/2$. Positivity of the
Bost--Connes Gibbs states therefore does not, by this trace
identity, supply the positivity of the cubic Weil form.
\end{proposition}
\begin{proof}
For $\Re s>1$ the sum of the absolute diagonal entries is
bounded by the absolutely convergent ideal zeta series. This
proves trace class and the identities by the definition of
the primitive Hecke series. The bound is the triangle
inequality. A faithful cubic ray-class character takes the
values $1,\omega,\omega^2$ on ideals prime to its conductor,
because every ray class is represented by such an ideal.
Its diagonal operator consequently has nonreal eigenvalues,
and its real part has eigenvalue $\Re\omega=-1/2$.
The assertion excludes this specific positivity transfer;
it is not a counterexample to cubic RH.
\end{proof}

\begin{proposition}[AC-K08: the cubic tensor bootstrap retains its target]
\label{AC-K08}
Let $V$ be the two-dimensional standard representation of
$S_3$ in AM-H01, and let $\varepsilon$ be its sign character.
For every integer $m\ge1$,
\[
 V^{\otimes m}\simeq
    1^{\oplus a_m}\oplus\varepsilon^{\oplus a_m}
                           \oplus V^{\oplus c_m},\qquad
 a_m=\frac{2^m+2(-1)^m}{6},\quad
 c_m=\frac{2^m-(-1)^m}{3}.
\]
In particular the coefficient of $V$ is positive for all $m$.
For even $m$, all Frobenius traces of this tensor power are
nonnegative, but its Artin $L$-function is still
\[
 L(s,V^{\otimes m})=\zeta_K(s)^{a_m}L_p(s)^{c_m}.
\]
Under A, its critical-line assertion is exactly equivalent
to that of $L_p$. Passing to an even tensor power therefore
does not eliminate the outstanding sector.
\end{proposition}
\begin{proof}
The three conjugacy classes have sizes $1,3,2$, and $V$ has
character values $2,0,-1$. Taking inner products of
$(2^m,0,(-1)^m)$ with the three irreducible characters gives
the stated multiplicities. They are nonnegative integers
(for $m=1$, $a_1=0$), and $c_m>0$. At a ramified prime the
same representation identity holds before taking inertia
invariants, so Artin multiplicativity gives the displayed
full $L$-function identity there too. All its factors have
no poles in the open critical strip, which proves the last
assertion. At a ramified prime, with inertia group $I$ and a
Frobenius lift $g$, the trace on invariants is
$|I|^{-1}\sum_{h\in I}\operatorname{tr}V(gh)^{m}$, which is
nonnegative for even $m$ as well. Nonnegative local traces of even powers are
therefore not the missing global one-sided window estimate.
\end{proof}

\paragraph{The commuting phase in geometric language.}
Over $U=\operatorname{Spec}\mathcal O_K[1/(3p)]$, the equation
$z^3=p$ is a finite \'etale $\mu_3$-torsor. Its two nontrivial
characters give rank-one lisse Artin sheaves, pure of local
weight zero. The action of $\mu_3$ commutes with Frobenius
over $K$, so the individual character projectors exist over
$\mathbb Q(\omega)$ or $\mathbb Q_\ell(\omega)$. Over
$\mathbb Q$ only the sum of the two nontrivial projectors is rational.
The integral failure of these projectors at three is already
computed in AM-H06--AM-H08; it is not repaired by a change
of vocabulary. For the CM Fermat cubic $x^3+y^3=1$, the
automorphism $(x,y)\mapsto(\omega x,y)$ instead acts on
$H^1$, giving weight-one Hecke realizations. Their local
weight theorem and their angular types are different data
from RH of their number-field Euler products. The relevant
frameworks are SGA~1, expos\'e~XI, \S6 (Kummer torsors),
and SGA~6, expos\'e~V, 2.2(b) and Lemma~2.2.1 (exterior
powers and representation rings). No inequality in these
framework statements implies H04.4 merely from the phase.

\subsection{An unconditional cubic window reached by the operator route}

\begin{proposition}[AC-K09: strict cubic positivity through $\lambda=4$]
\label{AC-K09}
For $p=269$, let $W_{\chi,\lambda}$ be the Weil quadratic form
of the primitive completed $L_K(s,\chi_{269,3})$ on logarithmic
test functions supported in $[-\log\lambda,\log\lambda]$.
With Fourier normalization as in AC-T05, for every
$1<\lambda\le4$ and every vector in its closed form domain,
\[
                     W_{\chi,\lambda}(v)\ge\|v\|_2^2.
\]
This assertion uses no RH for the cubic character and no input A.
It settles this finite range of windows, not all windows.
\end{proposition}
\begin{proof}
The primitive conductor in AC-K04 gives
$Q=|D_K|N\mathfrak f=217083$ and completed factor
$Q^{s/2}\Gamma_{\mathbb C}(s)L_p(s)$, where
$\Gamma_{\mathbb C}(s)=2(2\pi)^{-s}\Gamma(s)$.
The explicit formula, proved by the same logarithmic-derivative
contour calculation as AC-T05, has archimedean symbol
\[
 a(t)=\log Q-2\log(2\pi)+2\Re\psi(1/2+it).
\]
There is no pole term, since the character is nontrivial.
The convergent digamma difference series gives
$\Re\psi(1/2+it)\ge\psi(1/2)=-\gamma-2\log2$.
Thus $a(t)\ge c_0=\log Q-2\log(8\pi)-2\gamma$.

For the largest interval, $\lambda=4$, the only nonzero
prime-power coefficients with $n<16$ are
\[
\begin{array}{c|ccccc}
n&3&4&7&9&13\\ \hline
b_p(n)&\log3&2\log2&-\log7&\log3&-\log13.
\end{array}
\]
The prime above three splits by AC-K04. At inert rational
primes only even powers contribute, as in AM-H02. At seven
and thirteen, $269^{(\ell-1)/3}\not\equiv1\pmod\ell$, so
AM-H02 gives the two negative entries. These assertions
also follow by direct residue computations ($3^2\not\equiv1$
modulo seven and $9^4\not\equiv1$ modulo thirteen).
The translation by $\log16$ has zero compression at the
endpoint and contributes nothing.

For a shift $a$ in an interval of length $T$, AC-T14's
finite-chain decomposition gives
$\|T_a+T_{-a}\|=2\cos(\pi/(\lceil T/a\rceil+1))$
when $a<T$, and zero when $a\ge T$. At $T=\log16$, the
maximal chain lengths are respectively $3,2,2,2,2$ for
the five displayed entries. Consequently the norm of the
whole prime operator is at most
\[
 B_4=\sqrt2\,\frac{\log3}{\sqrt3}+\log2
       +\frac{\log7}{\sqrt7}+\frac{\log3}{3}
       +\frac{\log13}{\sqrt{13}}.
\]
The exact rational-interval certificate
\path{scripts/check_cubic_269.py} proves $c_0-B_4>1$.
Its elementary enclosures use the positive arctanh series
for logarithms with its geometric remainder, the alternating
Machin arctangent series for $\pi$, integer square-root
bounds, and $H_{100}-\log101<\gamma<H_{100}-\log100$.
All comparisons in that certificate are rational, not
floating-point sign decisions. The explicit formula therefore
gives the stated inequality first on smooth test functions.
It extends by closure to the full logarithmic-symbol form
domain. Restriction to a smaller support interval preserves
both the form and the norm, proving all $\lambda\le4$.
\end{proof}

\subsection{The exact gate retained from the newest checkpoint}

Write $b_p(n)$ for the full coefficients of $-L_p'/L_p$ and
$D_p=L_p'/L_p$. AM-H04 defines
\[
 G_p(y)=\sum_{n\le y}b_p(n)(y/n-1)+yD_p(1)
\]
and verifies every input of the generic theorem
\texttt{Deduction.landau\_window} in
\path{Development/AutomorphicArithmetic/DedekindLFunctions/LandauTheorems/landauWindow.lean}.
The recovered source of that theorem is present inside the
newest checkpoint's inherited source directory and was read
in this run; its dependencies were not rebuilt in Lean.
The remaining premise is, for each $\theta>1/2$, a constant
$C_\theta$ and arbitrarily large $Y$ such that
\begin{equation}
 G_p(y)\le C_\theta y^\theta
             \quad(Y\le y\le Y^{1+\delta}),\qquad\delta>0.
                                                               \label{eq:H044}
\end{equation}
None of AC-K01--AC-K09 proves \eqref{eq:H044}; the window in
AC-K09 is a bounded multiplicative test window and must not
be confused with an unbounded arithmetic ramp window in $y$.

The minimal-conductor host above removes ambiguity about the
finite algebra and its ramified Frobenius lifts. It supplies
no inequality on the individual twisted cohomology lengths
or the growing-host costs in ALP\_MACHINE's final Q3 gate.
The stronger full trace formula over $K$ would require the
Weil forms for \emph{all} its unitary Hecke characters to be
positive. AC-K05 proves the norm sectors. The cubic sector
and the nonzero angular sectors remain separate analytic
questions. Their failure to follow from the constructions
listed here is a missing proof, not a theorem of logical
independence from A.

\section{BCM systems, CM ground states, and exact neutrality}
\label{sec:bcm-neutrality}

This section incorporates Follow-up~1. HP denotes the supplied version
of Ha--Paugam, \emph{Bost--Connes--Marcolli systems for Shimura varieties I},
arXiv:math/0507101. CMR denotes Connes--Marcolli--Ramachandran,
\emph{KMS states and complex multiplication}, arXiv:math/0501424.
LLN denotes Laca--Larsen--Neshveyev, \emph{On Bost--Connes type systems
for number fields}, arXiv:0710.3452. Y denotes Yalkinoglu,
\emph{On arithmetic models and functoriality of Bost--Connes systems},
arXiv:1105.5022, including Neshveyev's appendix.
The section numbers for HP below are those of the supplied source:
its Dirichlet construction is \S8, not \S6.

\subsection{Positive partition functions and character insertions}

\paragraph{The source statements.}
HP \S7.3, first and second lemmas, gives the number-field Hamiltonian
$H e_{\mathfrak a}=\log N\mathfrak a\,e_{\mathfrak a}$ and partition
function $\zeta_F$. HP \S8.2 explicitly supposes that the character
is trivial at infinity. Its second and third definitions introduce
\emph{twisted} trace and partition function, and its lemma identifies
that twisted partition function with $L_F(s,\chi)$. Its final remark
leaves construction for nontrivial infinite components unresolved.
Thus the displayed construction covers even rational Dirichlet
characters, and finite-order characters over imaginary quadratic
fields; its own text does not supply the odd rational construction.
The zeros of odd Dirichlet functions are nevertheless covered by A.

\begin{proposition}[AC-B01: the exact analytic and thermodynamic consequences]
\label{AC-B01}
For a finite-order Hecke character $\chi$ with conductor $\mathfrak f$,
put $D_\chi e_{\mathfrak a}=\chi(\mathfrak a)e_{\mathfrak a}$, with
the value zero on ideals not prime to $\mathfrak f$. Then
\[
 \operatorname{Tr}(D_\chi e^{-sH})=L_F(s,\chi),\qquad \Re s>1.
\]
On the conductor-restricted Hilbert space used for the twisted trace,
the \emph{ordinary} positive heat trace is instead
\[
 Z_{F,\mathfrak f}(s)=
 \zeta_F(s)\prod_{\mathfrak p\mid\mathfrak f}
                      (1-N\mathfrak p^{-s}).                 \tag{B1.1}
\]
For $F=\mathbb Q$, A puts the nontrivial zeros of the primitive
twisted function on $\Re s=1/2$. For abelian $F/\mathbb Q$, A puts
the nontrivial zeros of $\zeta_F$ on that line. Finite Euler factors
deleted in (B1.1) have zeros on $\Re s=0$, and do not alter this
critical-strip conclusion. Continuation of these functions is a
classical input; the additional conclusion from A is their zero
location, not a new construction of the continuation.

The positive Gibbs state for the ordinary number-field system, when
$\beta>1$, is
\[
 \Phi_{\beta,w}(a)=\zeta_F(\beta)^{-1}
                 \operatorname{Tr}(\pi_w(a)e^{-\beta H}).
\]
Its established phase diagram is unchanged: there is one KMS state
for $0<\beta\le1$ and the extremal states for $\beta>1$, and at
zero temperature, are parametrized by
$\operatorname{Gal}(F^{\rm ab}/F)$. This is LLN Theorem~2.1
(\texttt{mainthm}), unconditional for every number field.
Meromorphic continuation does not make the heat operator trace
class for $0<\beta\le1$.
\end{proposition}
\begin{proof}
In the ideal basis, both operators in the first trace are diagonal.
The sum of absolute values is at most $\zeta_F(\Re s)$, which
converges for $\Re s>1$. Taking the trace therefore gives precisely
the ideal Dirichlet series with its primitive bad-prime convention.
Removing the zero entries is equivalent, for this trace, to
restricting to ideals prime to $\mathfrak f$. Without the insertion
the same basis gives the Euler product (B1.1).

For $F=\mathbb Q$ the first zero assertion is A. For abelian $F$,
Artin factorization expresses $\zeta_F$ as the product of the
primitive Dirichlet functions of the characters of
$\operatorname{Gal}(F/\mathbb Q)$. Each is holomorphic in the open
critical strip, so their zeros, counted with multiplicity, are the
union of those of the factors. A applies to all factors. A zero
of $1-N\mathfrak p^{-s}$ has $|N\mathfrak p^{-s}|=1$, hence
$\Re s=0$.

The state and its phase diagram are the cited LLN theorem; CMR
Theorem~5.1 (\texttt{KMSCM}) gives the imaginary-quadratic instance,
including its arithmetic action. These theorems are used, not
claimed as new consequences of A. Finally the positive trace
$\sum_{\mathfrak a}N\mathfrak a^{-\beta}$ diverges for
$0<\beta\le1$; the conductor-restricted series has the same
abscissa and divergence, by its nonzero finite Euler factor at one
and its positive coefficients. An analytic continuation changes
none of the positive eigenvalues $N\mathfrak a^{-\beta}$ or their
sum. It therefore cannot provide a normalized Gibbs density there.
\end{proof}

\begin{proposition}[AC-B02: a twisted trace is generally not a state]
\label{AC-B02}
For a nontrivial finite-order $\chi$, normalizing
$\operatorname{Tr}(D_\chi\,\cdot\,e^{-\beta H})$ by
$L_F(\beta,\chi)$ does not, in general, define a positive state,
even when that denominator is nonzero. In particular this failure
occurs for the even primitive character $\chi_5$ and for a cubic
character whose ideal values include $1,\omega,\omega^2$.
Inserting $D_\chi$ does not change the eigenvalues $\log N\mathfrak a$
of $H$. Hence the nontrivial zeros of the twisted partition
function are not thereby eigenvalues of this Hamiltonian.
\end{proposition}
\begin{proof}
On the positive rank-one projection onto $e_{\mathfrak a}$, the
normalized functional takes value
\[
 \frac{\chi(\mathfrak a)N\mathfrak a^{-\beta}}
      {L_F(\beta,\chi)}.                                    \tag{B2.1}
\]
For $F=\mathbb Q$, $\chi_5(2)=-1$, while $\chi_5(1)=1$ and
$L(\beta,\chi_5)>0$ for $\beta>1$: each real Euler factor is
positive. Formula (B2.1) is negative at $\mathfrak a=(2)$.
For a cubic character, no fixed nonzero complex scalar can send
all three numbers $1,\omega,\omega^2$ to nonnegative reals.
Thus positivity already fails on diagonal projections of
$B(\ell^2)$, the domain on which HP defines the twisted trace.
Equivalently, a positive normal functional on this algebra must
have a positive density operator, whereas the density displayed
here has these different phases. This proves the stated general
failure without asserting that every restriction to every smaller
subalgebra must fail. The diagonal formulas for $H$ and $D_\chi$
show the final assertion directly. A separate zero spectral
realization, such as AC-T04, uses different operators and a
different Hilbert space.
\end{proof}

For arbitrary Shimura BCM pairs, HP \S6.2,
Proposition \texttt{prop-construction-KMS}, constructs extreme
type-I Gibbs states in the convergence range. The following
Question \texttt{classif-etats-KMS-general} asks whether all
low-temperature extreme states arise from the Shimura variety.
The abelian number-field factorization above supplies no proof
of that general classification. The actual positive algebra,
dynamics, and measure classification remain necessary data.
There is also a harmless typographical omission of $\beta$ in
the first displayed cyclic vector in that proposition's proof:
its coefficient must be $\det(\phi(h))^{-\beta/2}$, as the
next computation and the Gibbs formula themselves require.

\subsection{Constructing the cubic ground-state observable}

\begin{proposition}[AC-B03: the radical is an actual arithmetic observable]
\label{AC-B03}
Let $K=\mathbb Q(\omega)$, $\alpha=\sqrt[3]p$, and let
$\mathfrak f$ be the conductor of $E_p/K$. In the arithmetic
endomotive of Y, one can construct a bounded finite-level
observable $a$ and a boundary idempotent $e$ such that
\[
                         a^3=p e.
\]
On a ground state belonging to the identity unit stratum,
$a$ has value $\alpha$; its three conjugate ground-state values
are $\alpha,\omega\alpha,\omega^2\alpha$.
The CMR CM arithmetic algebra likewise has ground-state values
generating $E_p$ and indeed all of $K^{\rm ab}$. This gives an
arithmetic realization of the radical, without an estimate of
the H04.4 ramp.
\end{proposition}
\begin{proof}
Y's finite-level algebra, equation \texttt{arithalg} in \S6,
is $E_{\mathfrak f}=\prod_{\mathfrak d\mid\mathfrak f}
K_{\mathfrak f/\mathfrak d}$. The quotient of finite
Deligne--Ribet sets constructed in AC-K03 induces an injection
\[
 B_p=E_p\times K^{d(\mathfrak f)-1}\lhook\joinrel\longrightarrow
 E_{\mathfrak f}.
\]
Take $a=(\alpha,0,\ldots,0)$ and $e=(1,0,\ldots,0)$ in $B_p$.
In the function description, $a(gH)=g\alpha$ on the unit
stratum and is zero on every other stratum. This is a
$G_K$-equivariant function, hence precisely an element of that
finite algebra. Pullback embeds it in Y's arithmetic subalgebra;
it is bounded because the finite-level function has finitely
many values. Its cubic identity is componentwise. Ground states
are evaluations on the invertible boundary, by LLN Theorem~2.1(iv)
and Y \S1.2, subsection \texttt{kmsinfty}. Evaluation at the
identity and Galois equivariance give the claimed values.

For CMR, Theorem~5.1 and the proof in \S5 under
``KMS states at zero temperature and Galois action'' give
$\varphi_{\infty,L}(f)=f(L,L)$, equation \texttt{Kground},
and the modular-value identity $F_\tau=K^{\rm ab}$,
equation \texttt{FtauKab}. In fact the values of functions
regular at $\tau$ form a field: if $f(\tau)\ne0$, then
$1/f$ is regular at $\tau$ as well. Thus the generated field
statement realizes $\alpha$ as the value of a rational
expression in finitely many such functions, regular at
$\tau$. This is an existence conclusion from the CM theorem,
not an explicit formula for such a modular function and not
an unproved identification of two arithmetic subalgebras.
None of these identities contains the logarithmic-derivative
coefficients or a bound on their sparse arithmetic windows.
\end{proof}

CMR's paragraph immediately after Theorem~5.1 supplies a further
necessary distinction. Inclusion of a subgroupoid in the
$\mathrm{GL}_2$ groupoid, even compatible with time evolution,
does not transport its KMS states to those of the ambient
system. Its explicit example is a finite full pair groupoid
inside an infinite one, with trivial dynamics: the first has
a normalized trace, the compact-operator algebra of the
second has none. Accordingly, CM ground states cannot be
transported to an ambient analytic zero theorem merely by
this inclusion.

\subsection{Neutrality at the units, at three, and at infinity}

\begin{proposition}[AC-B04: the exact finite and infinite alternatives]
\label{AC-B04}
For a unitary Hecke character of $K=\mathbb Q(\omega)$ with
infinite component
$\chi_\infty(z)=(z/|z|)^{-k}|z|^{2it}$, one has
\[
             \prod_{v<\infty}\chi_v(\epsilon)=\epsilon^k
                         \qquad(\epsilon\in\mu_6).          \tag{B4.1}
\]
A finite-order character necessarily has $k=t=0$.
Let $\mathfrak p=(\pi)\nmid3$ and let $r_\pi$ be the
nontrivial cubic residue character on
$(\mathcal O_K/\mathfrak p)^\times$, with values in $\mu_3$.
Then
\[
       r_\pi(\omega)=\omega^{(N\mathfrak p-1)/3}.            \tag{B4.2}
\]
If this value is $1$, $r_\pi$ descends to a finite-order ray
character of conductor $\mathfrak p$. Otherwise it can be
corrected by a uniquely specified cubic character modulo
$(3)=(1-\omega)^2$, yielding conductor exactly
$(3)\mathfrak p$. Correcting the same nontrivial unit colour
by an infinity type $k\equiv2$ or $4\pmod6$ instead yields
an infinite-order Hecke character.
\end{proposition}
\begin{proof}
A Hecke character is trivial on diagonal $K^\times$.
Applying this to $\epsilon\in\mu_6$, whose absolute value
is one, proves (B4.1). A continuous finite-image character
of the connected group $\mathbb C^\times$ is trivial.
The angular character and the radial character in the
displayed infinity type are both trivial only if $k=t=0$.
This also proves that a nonzero angular compensation cannot
remain finite order. For $\mathbb Q$, the same argument
at $-1$ gives $\chi_\infty(-1)=\chi_f(-1)$, since a sign
is its own inverse.

Formula (B4.2) is the definition of the cubic residue
character, applied to $\omega$. Also $r_\pi(-1)=1$:
its order divides both two and three. Since $K$ has class
number one, the ray class group for a finite modulus is
the group of residue units modulo the image of $\mu_6$.
Thus $r_\pi$ descends exactly when it is trivial on $\mu_6$,
equivalently when (B4.2) is one, or
$N\mathfrak p\equiv1\pmod9$.

The reduction map
\[
                 \mu_6\longrightarrow(\mathcal O_K/(3))^\times
\]
is an isomorphism. Its target has six elements. The images
of $-1$ and $\omega$ have orders two and three respectively:
$\omega-1$ is divisible by $1-\omega$ but not by its square.
They therefore generate six distinct units. Define $c$ on
the target by $c(\epsilon)=r_\pi(\epsilon)^{-1}$.
The Chinese remainder product $c\,r_\pi$ modulo
$(3)\mathfrak p$ is trivial on all global units, so descends
to the ray class group. If (B4.2) is nontrivial, $c$ has
order three and is nontrivial on the principal units
modulo $(1-\omega)^2$; it cannot factor through the
two-element unit group modulo $1-\omega$. Its conductor
at that prime is exactly the square. At $\mathfrak p$,
$r_\pi$ is nontrivial and has conductor $\mathfrak p$.
This proves the exact finite conductor.

For completeness, the infinity-type alternative is also
well-defined, not just a necessary unit test. Choose
$k$ with $r_\pi(\epsilon)=\epsilon^k$ for every unit.
On principal ideals prime to $\mathfrak p$ define
$r_\pi(\alpha)^{-1}(\alpha/|\alpha|)^k$ for a generator
$\alpha$. Changing the generator to $\epsilon\alpha$
does not change this value. Its finite local component on
residue units is $r_\pi$: multiplying the principal-ideal
value by $r_\pi(\alpha)$ and by the infinity component
$(\alpha/|\alpha|)^{-k}$ gives one. Thus the resulting
ideal character, with this infinity component and local
conductor, is a unitary Hecke character. Its angular
component is nontrivial when $k\equiv2$ or $4\pmod6$,
so it has infinite order by connectedness as above.
\end{proof}

This corrects the finite-order clause in the updated first-pass
note, \S2b item~10. Neutralization at infinity and finite-order
neutralization at the prime over three are different constructions.
For $p=269$, $N(p)=269^2\equiv1\pmod9$, consistent with
the conductor $(269)$ proved in AC-K04.

\subsection{The divisor-pair statement and its positivity limit}

\begin{proposition}[AC-B05: exact norm congruences and Fourier counting]
\label{AC-B05}
Let $R\ge2$ and let $d,e$ be integers prime to $R$. Then
\[
\begin{split}
 N_K(d+e\omega)&=d^2-de+e^2,\\
 N_K(d-e\omega)&=d^2+de+e^2.
\end{split}                                                  \tag{B5.1}
\]
The two respective divisibility tests by $R$ are equivalent
to $d/e$ being a root modulo $R$ of
$\Phi_6(T)=T^2-T+1$ or $\Phi_3(T)=T^2+T+1$.
For either polynomial, roots exist precisely when
\[
 v_3(R)\le1,\qquad
 \ell\equiv1\pmod3\quad\hbox{for every prime }\ell\mid R,
                                                 \ \ell\ne3. \tag{B5.2}
\]
When they exist, their number is
$2^{\#\{\ell\mid R:\ell\ne3\}}$.

For any finite multiset $S$ of integers prime to $R$ and
$\zeta\in(\mathbb Z/R\mathbb Z)^\times$, the ordered-pair
count is exactly
\[
 \#\{(d,e)\in S^2:d\equiv\zeta e\pmod R\}
 =\frac1{\varphi(R)}\sum_{\chi\bmod R}
          \overline{\chi(\zeta)}
          \left|\sum_{d\in S}\chi(d)\right|^2.              \tag{B5.3}
\]
Summing (B5.3) over the roots of the appropriate polynomial
gives the exact norm-divisor-pair count. It does not establish
that a prescribed Erd\H{o}s--Straus divisor set occupies any
of these congruence classes, and it does not give positivity
of the phase-weighted form on all complex test vectors.
\end{proposition}
\begin{proof}
Multiply each element of (B5.1) by its conjugate, using
$\omega+\omega^2=-1$ and $\omega^3=1$. Division by the
unit $e^2$ gives the polynomial criteria. For
$\ell\ne2,3$, the roots of $\Phi_3$ are the two primitive
cube roots in $\mathbb F_\ell^\times$, so exist exactly
when $3\mid\ell-1$. The roots of $\Phi_6$ are their
negatives. In either case the roots are simple and lift
uniquely to every $\ell^a$ by the elementary simple-root
lifting argument. There are no roots at two. At three
there is one root, namely $1$ for $\Phi_3$ and $-1$ for
$\Phi_6$. Substitution $1+3u$ or $-1+3u$ gives a value
of exact 3-adic valuation one, so no root exists modulo
nine. The Chinese remainder theorem proves (B5.2) and
the count. At a composite modulus the local choices can
be made independently; they need not come from one
preselected choice of a global Eisenstein phase.

Character orthogonality on $(\mathbb Z/R\mathbb Z)^\times$
gives
$\mathbf1_{d=\zeta e}=\varphi(R)^{-1}
\sum_\chi\chi(d)\overline{\chi(\zeta e)}$.
Summing over the two copies of $S$ proves (B5.3), including
multiplicities. The distinct roots specify disjoint
congruence conditions for $d/e$, so summing gives the
norm condition exactly.

There is a direct obstruction to promoting this counting
identity to arbitrary-vector positivity. On the cyclic
three-point set let $U$ be cyclic translation. The form
for the union of its two nonidentity phases is represented
by $U+U^*$, whose matrix has zeros on the diagonal and
ones elsewhere. Its eigenvalues are $2,-1,-1$; the vector
$(1,\omega,\omega^2)$ has eigenvalue $-1$. Its values on
nonnegative counting vectors are nonnegative, but it is
not a positive Hermitian form. Finally, if either $d$ or
$e$ is not a unit, the division argument is invalid;
for instance $d=e=R$ satisfies both norm divisibilities
without specifying any unit ratio. These qualifications
exhibit the precise limits of the proposed phase bridge.
\end{proof}

\subsection{The common cyclotomic field and the surviving cubic factors}

\begin{proposition}[AC-B06: the biquadratic base change is exact but not a closure]
\label{AC-B06}
Put $M=\mathbb Q(\zeta_{12})=\mathbb Q(i,\omega)$.
On its Galois group the three nontrivial quadratic
characters satisfy $\chi_{-4}\chi_{-3}\chi_{12}=1$, and
\[
 \zeta_M(s)=\zeta(s)L(s,\chi_{-4})L(s,\chi_{-3})L(s,\chi_{12}).
                                                               \tag{B6.1}
\]
Thus A proves RH for $\zeta_M$. Moreover
\[
 N_{M/\mathbb Q}(r-s\zeta_{12})=r^4-r^2s^2+s^4.                \tag{B6.2}
\]
For a rational prime $p$, let $E'=M(\sqrt[3]p)$ and
let $\eta$ be the quadratic character of $M/K$, namely
the restriction of $\chi_{-4}$, or equivalently its
norm pullback to $K$. Then
\[
\begin{split}
 \operatorname{Gal}(E'/\mathbb Q)&\simeq S_3\times C_2,\\
 L_M(s,\chi_{p,3}|_{G_M})
       &=L_K(s,\chi_{p,3})L_K(s,\chi_{p,3}\eta),\\
 \zeta_{E'}(s)&=\zeta_M(s)
          L_K(s,\chi_{p,3})^2L_K(s,\chi_{p,3}\eta)^2.
\end{split}                                                \tag{B6.3}
\]
All equalities include ramified Euler factors. Consequently
the nonabelian factors have not been eliminated by this
cyclotomic base change or by the product of the three signs.
Under A, RH for $\zeta_{E'}$ is equivalent to RH for both
cubic factors displayed in (B6.3).
\end{proposition}
\begin{proof}
The unit classes $1,5,7,11$ modulo twelve give the table
\[
\begin{array}{c|rrrr}
 &1&5&7&11\\\hline
 \chi_{-4}&1&1&-1&-1\\
 \chi_{-3}&1&-1&1&-1\\
 \chi_{12}&1&-1&-1&1.
\end{array}
\]
The product assertion and (B6.1) follow from this table
and abelian Artin factorization. The product is an
identity of Galois characters, or of characters on
units, not the constant-one identity for their
zero-extended values at all integers. Formula (B6.2)
is $s^4\Phi_{12}(r/s)$, with the case $s=0$ obtained
by the polynomial identity.

The splitting field $E_p$ of $X^3-p$ has group $S_3$
and unique quadratic subfield $K$. Hence
$E_p\cap\mathbb Q(i)=\mathbb Q$, and the compositum
$E'=E_p\mathbb Q(i)$ has group $S_3\times C_2$.
Restriction followed by induction for the index-two
subgroup $G_M\subset G_K$ gives
\[
 \operatorname{Ind}_{G_M}^{G_K}
      (\chi_{p,3}|_{G_M})
              =\chi_{p,3}\oplus\chi_{p,3}\eta.
\]
Artin induction of local inertia invariants proves
the middle identity in (B6.3) at every prime, not only
the unramified ones. The two nontrivial characters
of the cubic extension $E'/M$ are interchanged by
complex conjugation, which preserves ideal norms
and the quadratic character $\eta$. Their global
$L$-functions are therefore equal. Applying the
cyclic Artin factorization over $M$ proves the last
identity. The Hecke functions there are entire,
because both characters are nontrivial. No poles
can cancel their zeros in the open critical strip.
Together with A for $\zeta_M$, this proves the final
equivalence. The field $E'$ is not CM, since it
contains the non-CM-overfield obstruction $E_p$
of AM-H03; adjoining $i$ has not removed that
obstruction.
\end{proof}

\paragraph{Geometric interpretation and an attribution gate.}
Over $\operatorname{Spec}\mathbb Z[1/6]$, the three
quadratic characters are rank-one lisse Artin sheaves
whose tensor product is the unit object; (B6.2) is the
norm of the finite \'etale cyclotomic algebra. Over
$\operatorname{Spec}\mathcal O_M[1/(3p)]$, the cubic
cover is a finite \'etale $\mu_3$-torsor. Its character
summands exist over a coefficient field containing
$\mu_3$, are locally pure of weight zero, and still
have the global functions in (B6.3). These are precise
uses of Galois categories, Kummer torsors and Artin
motives; local weight zero is not a proof of their
global critical-line assertion. Relevant verified
framework locators are recorded in the source appendix.
The updated note attributes $\Delta=\chi_{12}\circ u$
to ``DELTA23 intake L1''. That intake and the map $u$
have not been located among the supplied files in this
run. The elementary character table above is checked;
the specific Grothendieck attribution remains a
source gate and is not asserted as a verified fact.

\InputIfFileExists{modules/trace_followup.tex}{}{}
\InputIfFileExists{modules/trace_cutoff_bridge.tex}{}{}
\section{A local upper bound for the actual prolate defect}

Throughout this section \(\chi\) is a fixed nonprincipal primitive
Dirichlet character, \(R=\Lambda/q\), \(c=\Lambda R\), and
\(\Pi_\Lambda\) is the original full selected raw prolate projection.
The target projection \(Q_\Lambda\), the sharp log window \(S_\Lambda\),
and \(D_\Lambda=S_\Lambda-S_\Lambda Q_\Lambda S_\Lambda\)
are those of AC-T30--T31.
Raw Mellin coordinates are given by the unitary
\((U\psi)(t)=e^{t/2}\psi(e^t)\).
The Sobolev norm \(H_D^s\) means the ordinary \(H^s\) norm
after this unitary.  For a symbol \(a\), \(C_a\) denotes the
logarithmic Fourier multiplier with symbol \(a(\tau)\), using
the positive-sign Fourier convention of AC-T31.

\begin{lemma}[AC-T38: graph approximation with a fixed cutoff margin]
Choose an even \(p\in C_c^\infty((-1,1))\), with \(0\leq p\leq1\),
which equals one on a neighborhood of zero.  Put
\[
 P_\Lambda^p=p(u/\Lambda),\qquad
 Q_R^p=\mathcal F_\epsilon^{-1}p(\xi/R)\mathcal F_\epsilon.
 \tag{TU.1}
\]
Let \(T_\Lambda\) be a Hilbert--Schmidt operator into the raw
parity space with
\(\|T_\Lambda\|_2=O((1+\log\Lambda)^d)\) for some fixed \(d\).
Then, for every \(A>0\),
\[
 \left\|E_\chi(1-\Pi_\Lambda)
             P_\Lambda^pQ_R^pP_\Lambda^pT_\Lambda\right\|_2
                    =O_A(c^{-A}).
 \tag{TU.2}
\]
No form of RH is needed for this lemma.
The same graph approximation, at every polynomial rate, holds
for any fixed finite family of smooth parity profiles compactly
supported strictly inside the raw time interval, after dilation
to radius \(\Lambda\), or compactly supported strictly inside
the raw Fourier interval, after dilation to radius \(R\).
\end{lemma}
\begin{proof}
Let \(b<1\) contain the support of \(p\).
The Fourier kernel of multiplication by \(p(u/\Lambda)\) is
\(\Lambda\widehat p(\Lambda(\xi-\eta))\).
For \(|\eta|\leq bR\) and \(|\xi|>R\), their separation
is at least \((1-b)R\).
Schur's test and rapid decay of \(\widehat p\) consequently give
\[
 \|1_{|\xi|>R}\mathcal F_\epsilon
                         P_\Lambda^pQ_R^p\|
                  \leq C_M c^{-M}\quad(M>0).
 \tag{TU.3}
\]
The constants depend on the fixed \(p,M\), not on the cutoffs.
The full-cutoff prolate inequality (TB.82) applies to the
columns \(f\) of \(P_\Lambda^pQ_R^pP_\Lambda^pT_\Lambda\);
their raw time tail is zero.
Thus their summed squared projection error is bounded by
eight times the squared Fourier-tail norm in (TU.3).
It follows that the raw Hilbert--Schmidt error is
\(O_M(c^{-M}(1+\log\Lambda)^d)\).

For \(0<u<1/2\), the sharp-endpoint Sobolev estimate
(TB.33)--(TB.34), with the full retained eigenvalues bounded
below, gives
\[
 \|\Pi_\Lambda f\|_{H_D^u}
                  \leq C_u(1+c)^u\|f\|_2.
 \tag{TU.4}
\]
For the smooth-cutoff vector itself, differentiate
\(u^{1/2}p(u/\Lambda)(Q_R^pP_\Lambda^p f)(u)\).
The term from \(p'\) is bounded because
the raw support has \(u/\Lambda\) in a fixed compact set;
the band derivative is bounded by \(2\pi R\), multiplied
by at most \(\Lambda\).
Hence its \(H_D^1\) norm is at most \(C(1+c)\|f\|_2\).
Interpolation gives the same estimate (TU.4) for this
smooth-cutoff vector.  Applying these inequalities to an
orthonormal input basis gives their Hilbert--Schmidt versions.

Use the unconditional multiplier bound already used in
AC-T25,
\[
 \|E_\chi f\|_2\leq C_\chi\|f\|_{H_D^{5/16}},
 \tag{TU.5}
\]
and interpolate the small raw \(L^2\) error with the coarse
\(H_D^{3/8}\) bound.  Since
\((5/16)/(3/8)=5/6\), the resulting bound is
\(O_M(c^{-M/6+5/16}(1+\log\Lambda)^d)\).
Choosing \(M\) arbitrarily large proves (TU.2).
For a fixed smooth compact time profile, its Fourier tail
has every polynomial decay rate, while (TU.4) remains
available.  This proves the first profile assertion.
For a compact smooth Fourier profile the time tail has
every such rate and its Fourier tail is zero.
Fourier conjugation changes the dilation generator to its
negative, so its fixed-profile \(H_D^{3/8}\) norm is bounded
after the stated unitary scaling.  The same interpolation
proves the second assertion.
\end{proof}

\begin{theorem}[AC-T39: local convergence of the actual positive defect]
Assume \(\mathsf A\).  For every \(a\in C_c^\infty(\mathbb R)\),
\[
 \lim_{\Lambda\to\infty}
    \operatorname{Tr}\bigl(D_\Lambda C_aC_a^*\bigr)
                =\sum_\gamma m_{\chi,\gamma}|a(\gamma)|^2.
 \tag{TU.6}
\]
Equivalently, the positive measures \(\mu_\Lambda\)
defined in AC-T31 converge vaguely to
\(\sum_\gamma m_{\chi,\gamma}\delta_\gamma\).
This statement concerns compact Mellin-frequency tests.
It does not assert the uniform high-frequency tightness
needed to pass to every original compact norm test.
\end{theorem}
\begin{proof}
Write \(L=\log\Lambda\), \(S_L=1_{[-L,L]}\).
Choose a fixed compact interval containing the support of
\(a\) in its interior, whose endpoints are not zero ordinates,
and list the finitely many distinct ordinates in that interval.
Let \(Z_L\) be the span in \(S_LL^2(\mathbb R)\) of
\[
       1_{[-L,L]}(t)e^{-i\gamma t}t^j,
                         \qquad 0\leq j<m_{\chi,\gamma},
 \tag{TU.7}
\]
and let \(P_{Z,L}\) be its orthogonal projection.
Set \(W_L=S_L-P_{Z,L}\).
The normalized-window Gram calculation in AC-T31 shows that
its dimension is constant for large \(L\), that its normalized
basis coefficients are uniformly bounded, and that
\[
       \|S_LC_aP_{Z,L}\|_2^2
              \longrightarrow
                  \sum_\gamma m_{\chi,\gamma}|a(\gamma)|^2.
 \tag{TU.8}
\]
Indeed each normalized exponential-polynomial jet has Fourier
mass concentrating at its ordinate, so multiplication by the
bounded continuous \(a\) is asymptotically multiplication by
\(a(\gamma)\).  The distinct-ordinate Gram blocks tend to zero,
and each same-ordinate block is the invertible Hilbert matrix
already computed in AC-T31.

\emph{A finite-pole inverse with both exact boundary descriptions.}
On the compact frequency interval, \(a(\tau)/
L(1/2+i\tau,\chi)\) has only the finitely many zero poles
in (TU.7).  Give each pole the boundary prescription
\((\tau-\gamma-i0)^{-j}\).
This defines a compactly supported distribution \(b\).
Its inverse Fourier transform \(\kappa\) is smooth and has
the following two properties, with all derivatives:
\[
 \begin{cases}
 \kappa(t)\text{ decreases faster than every power as }t\to+\infty,\\
 \kappa(t)=\displaystyle\sum_{\gamma,j}d_{\gamma j}
                  t^je^{-i\gamma t}+O_N((1+|t|)^{-N}),
                                      &t\to-\infty .
 \end{cases}
 \tag{TU.9}
\]
Here the polynomial degrees are less than the corresponding
zero multiplicities.  To verify (TU.9), subtract the finitely
many principal parts using smooth cutoffs supported near each
pole.  The remainder is smooth compactly supported.
The inverse transform of each prescribed pole is a constant
multiple of \(1_{t<0}t^{j-1}e^{-i\gamma t}\).
Convolution with the Schwartz transform of its smooth cutoff
gives its full polynomial asymptote on the negative half-line
and a Schwartz remainder at both ends.  This also proves the
assertions for every derivative.

Define
\[
                T_L=U^{-1}C_bW_L.
 \tag{TU.10}
\]
This is an ordinary Hilbert--Schmidt operator, despite the
distribution notation for \(b\).
The exact zero-jet conditions on \(W_Lf\) cancel every
polynomial term of (TU.9) when the output coordinate is
below \(-L-1\).  Above \(L+1\), the positive tail in
(TU.9) is Schwartz.  In the remaining strip, the kernel
and each fixed derivative are bounded by a fixed power of
\(1+L\).  Integrating its squared kernel therefore proves
\[
       \|UT_L\|_{\mathcal S_2(L^2,H^j)}
                    \leq C_j(1+L)^{d_j}\quad(j\geq0).
 \tag{TU.11}
\]
The same conclusion for the lower tail follows by writing
the canceled polynomial convolution first and estimating
the remaining Schwartz kernel; right multiplication by
the contraction \(W_L\) cannot increase that estimate.
Multiplication of distributions by the analytic \(L\)-factor
is legitimate and cancels every prescribed pole.  Thus the
closed multiplier realization gives the exact identity
\[
                         E_\chi T_L=C_aW_L .
 \tag{TU.12}
\]

We also need its ordinary Fourier transform in raw space.
On a fixed parity sector, the unitary raw Fourier transform
in Mellin coordinates is reflection \(\tau\mapsto-\tau\)
followed by a smooth unimodular multiplier on the real line.
For example these multipliers are
\[
 2(2\pi)^{-s}\Gamma(s)\cos(\pi s/2),\qquad
 -2i(2\pi)^{-s}\Gamma(s)\sin(\pi s/2),
                   \quad s=\tfrac12+i\tau ,
 \tag{TU.13}
\]
in the even and odd conventions, respectively.
The Mellin integral of cosine or sine first proves the
formula in its convergence strip; unitarity proves its
extension to the critical line.
Multiply \(b(-\tau)\) by this smooth multiplier.
Changing its finite pole prescriptions back to those used
in (TU.9) adds only delta derivatives at the reflected
zero ordinates.  These vanish on the reflected input
\(JW_L\), where \(Jf(t)=f(-t)\).
Consequently
\[
       \mathcal F_\epsilon T_L=U^{-1}C_{b^\#}JW_L
 \tag{TU.14}
\]
for a compactly supported distribution \(b^\#\) whose
kernel has precisely the two properties (TU.9).
This is an exact identity, not an approximation of a pole.

\emph{The two boundary families are compact in the required graph norms.}
Use the fixed smooth \(P=P_\Lambda^p\), \(Q=Q_R^p\) of
AC-T38.  The raw upper tail \((1-P)T_L\), after unitary
raw dilation by \(\Lambda\) and input translation by \(L\),
has, before the finite jet correction, the log kernel
\[
        (1-p(e^t))\,\kappa(t-s)\,
                              1_{[-2L,0]}(s).
 \tag{TU.15}
\]
Its output coordinate is bounded below by a fixed constant.
Since \(s\leq0\), (TU.9) makes this a Schwartz kernel in
the quadrant towards infinity.
It converges in Hilbert--Schmidt norm, with every fixed
output log derivative and every ordinary raw derivative,
to the same kernel with \(1_{(-\infty,0]}(s)\).
For ordinary raw derivatives use
\(\partial_u=u^{-1}\partial_{\log u}\), and recall
that the output raw coordinate is bounded away from zero.
All the resulting kernels are square integrable uniformly.

The finite jet correction tends to zero in these same norms.
In the translated input coordinate a normalized jet is
\(O(L^{-1/2})(1+|s|/L)^j\) on \([-2L,0]\), up to its
unit-modulus phase and uniformly bounded Gram coefficients.
The limiting kernel in (TU.15) has integrable moments of
every order in \(s\).
Cauchy--Schwarz, or direct integration of those moments,
therefore bounds each corrected column by \(O(L^{-1/2})\).
There are only finitely many columns.

The Fourier upper tail
\((1-p(\xi/R))\mathcal F_\epsilon T_L\) has the identical
description using (TU.14).
Its raw endpoint is \(\log R=L-\log q\), so only a fixed
translation by \(\log q\) occurs in its kernel.
It has the same Hilbert--Schmidt convergence and all the
same derivative bounds.

These facts also prove the two cross estimates
\[
 \|(1-Q)(1-P)T_L\|_{\mathcal S_2(H,H_D^j)}
 +\|(1-P)(1-Q)T_L\|_{\mathcal S_2(H,H_D^j)}
                   =O_{j,M}(c^{-M})\quad(j,M\geq0).
 \tag{TU.16}
\]
For the first estimate, take the Fourier tail of the
rescaled kernel (TU.15).  Every ordinary output derivative
is uniformly Hilbert--Schmidt, so its Fourier tail above
frequency \(c\) is \(O_M(c^{-M})\).
Applying a log derivative commutes with Fourier transform
by changing the sign of the dilation generator; derivatives
of the smooth cutoff are uniformly bounded in log coordinates.
The same reasoning applied to each derivative proves the
stated \(H_D^j\) bound.
For the second estimate use the Fourier-dual boundary family
and invert Fourier transform.  This proves (TU.16) directly
with the actual conductor scale.

\emph{Exact decomposition of the target residual.}
Let \(S_+=1_{(-\infty,L]}\) and \(S_-=1_{[-L,\infty)}\),
so \(S_L=S_+S_-\).
Raw time support gives \(E_\chi P\) target support at most
\(\Lambda\).  Raw Fourier support and primitive Poisson
give \(E_\chi Q\) target support at least
\(1/(qR)=1/\Lambda\).
All raw terms here belong to \(H_D^1\), by (TU.11) and
boundedness of the smooth log cutoffs on this space.
The support assertions first hold on the corresponding
time- or band-supported smooth cores.  Cutting off near
raw zero and then smoothing within the stated strict margin
approximates in \(H_D^{5/16}\); Fourier conjugation preserves
that norm.  Thus (TU.5) extends the support assertions to
the actual vectors in the following identity.
Consequently, with no commutation of \(P,Q\),
\[
 \begin{split}
 S_LE_\chi T_L-E_\chi PQPT_L
  ={}&S_LE_\chi(1-P)T_L\\
    &+S_LE_\chi(1-Q)PT_L
      +S_+E_\chi(1-P)QPT_L .
 \end{split}
 \tag{TU.17}
\]
For example, the last step of the expansion is
\((S_L-1)E_\chi QPT_L+E_\chi(1-P)QPT_L\);
its lower part is zero by the just stated Fourier gap,
and its upper part cancels because \(E_\chi P\) has
the upper support bound.  This verifies (TU.17).

Equation (TU.16), the boundedness of smooth log cutoffs
on \(H_D^1\), and (TU.5) permit replacing \(PT_L\) by
\(T_L\) in the second term, and replacing \(QPT_L\) by
\(PT_L\) in the third term, with Hilbert--Schmidt error
tending to zero faster than every power of \(c\).
For the latter replacement the exact raw error is
\[
 \begin{split}
 (1-P)(Q-1)PT_L
   ={}&-(1-P)(1-Q)T_L\\
      &+(1-P)(1-Q)(1-P)T_L.
 \end{split}
 \tag{TU.16a}
\]
The first term is the second cross estimate in (TU.16);
the second is the first cross estimate followed by the
uniformly \(H_D^1\)-bounded multiplier \(1-P\).
By (TU.15), the first and third terms, translated to the
upper target endpoint \(L\), converge in Hilbert--Schmidt
norm to fixed operators with range in \(L^2((-\infty,0])\).
The far lower target cutoff tends strongly to the identity
on their fixed Hilbert--Schmidt limits.
For the second term use the exact theta--Fourier identity
from AC-T25/AC-T32:
\(E_\chi\mathcal F_\epsilon^{-1}\) is a unit phase times
\(J_qE_{\bar\chi}\), with \(J_qF(x)=F(1/(qx))\).
The raw Fourier endpoint \(R\) then becomes the target
endpoint \(1/(qR)=1/\Lambda\).
Thus after translation to \(-L\) this second term converges
in Hilbert--Schmidt norm to a fixed operator with range
in \(L^2([0,\infty))\).
The graph-norm convergence of the raw boundary kernels
justifies applying these closed theta operators before
taking the limits.

\emph{Absorption of a fixed boundary operator by actual finite images.}
Under \(\mathsf A\), AC-T24, using Burnol's actual
nonprincipal sector with multiplier \(V=1\), says that
\[
 \overline{\{E_\chi\phi:
          \phi\text{ smooth of the required parity},
                    \operatorname{supp}\phi\Subset(-1,1)\}}
                  =L^2((-\infty,0])
 \tag{TU.18}
\]
in target log coordinates.
The same assertion holds for the dual character.
Smooth profiles may be chosen to vanish near raw zero:
cutting a smooth profile there changes its \(H_D^s\) norm
by a quantity tending to zero for every fixed \(s<1/2\),
as follows directly from its \(O(e^{t/2})\) log tail.
Equation (TU.5) then preserves density.

If a fixed Hilbert--Schmidt operator has range in the
right side of (TU.18), first approximate it by a finite-rank
operator, and then approximate each of its finitely many
output vectors by \(E_\chi\phi\) in (TU.18).
This approximates the operator in Hilbert--Schmidt norm.
Dilate those finitely many raw profiles to the time
endpoint \(\Lambda\).
AC-T38 shows that replacing each by its actual
\(\Pi_\Lambda\) projection changes its theta image by
a quantity tending to zero.
Therefore, for every upper boundary family just obtained,
\[
                 \|(1-Q_\Lambda)B_{\Lambda,+}\|_2\to0.
 \tag{TU.19}
\]
For a lower boundary family use the dual version of
(TU.18), Fourier transform its profiles back to the raw
space, and dilate their Fourier support to radius \(R\).
The second profile assertion of AC-T38 applies.
Primitive Poisson supplies exactly the lower target endpoint
\(1/(qR)\), so the same conclusion holds for that family.
All limits here keep the finite frequency interval and its
finite zero set fixed.  No completed span of all zero jets
in a fixed window is used.

By (TU.11) and AC-T38, the bulk \(E_\chi PQPT_L\) itself
is within \(o(1)\), in Hilbert--Schmidt norm, of actual
columns in \(E_\chi\Pi_\Lambda\).
Equations (TU.12), (TU.17)--(TU.19) now give
\[
                    \|(1-Q_\Lambda)S_LC_aW_L\|_2\to0.
 \tag{TU.20}
\]
There is one additional input boundary:
\(S_LC_a(1-S_L)\).
Its convolution kernel is Schwartz because \(a\) is
smooth compactly supported.
Splitting its input above \(L\) and below \(-L\),
and translating each endpoint to zero, gives two fixed
quadrant Hilbert--Schmidt limits, with outputs on the
appropriate target half-lines.
The same boundary absorption argument proves
\[
                  \|(1-Q_\Lambda)S_LC_a(1-S_L)\|_2\to0.
 \tag{TU.21}
\]

\emph{The matching upper bound.}
The positive defect has the exact distance formula
\[
       \operatorname{Tr}(D_\Lambda C_aC_a^*)
                    =\|(1-Q_\Lambda)S_LC_a\|_2^2 .
 \tag{TU.22}
\]
Decompose its input orthogonally into \(Z_L\) and its
orthogonal complement.
Equations (TU.20)--(TU.21) show that the second summand
tends to zero.  The first is bounded above by
\(\|S_LC_aP_{Z,L}\|_2^2\), whose limit is (TU.8).
This proves the required upper limit.
The lower limit is the finite-jet argument of AC-T31:
that proof uses only boundedness and continuity of the
multiplier at the finitely many ordinates, so it applies
to the present \(a\) verbatim.
This proves (TU.6).

Finally polarization gives convergence against every
smooth compact frequency test: write such a test as
\(a\overline b\), choosing \(a=1\) on its support.
For each compact set choose a smooth \(a\) equal to one
there.  Equation (TU.6) bounds the local masses uniformly
for large \(\Lambda\).
Approximation of continuous compact tests by smooth ones
then gives the stated vague convergence.
\end{proof}

\begin{remark}[What AC-T39 leaves in the original trace problem]
The local upper bound rules out an excess positive measure
on every fixed compact Mellin-frequency interval.
The original cutoff trace tests have rapidly decreasing,
but generally noncompact, Mellin transforms.
Their full limit still requires a uniform weighted tail
estimate for \(\mu_\Lambda\), or an equivalent
justified passage from these compact frequency tests.
Local convergence alone does not supply that passage.
The exact actual-adelic cutoff dictionary and the other
scope qualifications of AC-T11 and AC-T30 remain in force.
All approximations in this section concern actual
Hilbert--Schmidt column operators and orthogonal projections.
Burnol's Theorem 3.2 is used only through the already proved
half-line closure (TU.18), with its nonprincipal multiplier
identically one.  The finite-pole inverse, strict cutoff
margin, and boundary approximation supply the further
finite-prolate comparison here.
No lifting of a Gram inverse to an LF representation,
Grothendieck nuclear trace identification, or new adelic
cutoff equivalence is asserted by this argument.
\end{remark}

\InputIfFileExists{modules/spectral_followup.tex}{}{}
\InputIfFileExists{modules/spectral_boundary.tex}{}{}
\section{Smooth approximation in the actual theta graph norm}
\label{sec:spectral-prolate-graph}

The following estimate concerns the original sharp prolate
projection and the closed arithmetic theta map.  For smooth
input it needs only the unconditional convexity bound, not
theorem~A.  It is also an estimate for families of columns,
which is useful before any arithmetic Gram normalization.

\begin{theorem}[AC-S29: quantitative theta graph approximation]
\label{AC-S29}
Fix a nonprincipal primitive Dirichlet character $\chi$,
with parity $\epsilon$.  Use the Fourier convention
$\widehat f(\xi)=\int_{\mathbb R}f(u)e^{-2\pi i u\xi}\,du$.
On the raw parity space, with norm
$\|f\|_{\mathcal H}^2=\int_0^\infty|f(u)|^2du$, put
\[
 P=1_{[-1,1]},\qquad
 Q_c=\mathcal F^{-1}1_{[-c,c]}\mathcal F,
             \qquad S_c=PQ_cP.                     \tag{SG.1}
\]
Let $\Pi_c$ retain an initial segment of the prolate
eigenvectors in this parity, and suppose that, uniformly
for the values of $c\geq1$ under consideration,
\[
 \min_{\rm retained}\nu_j\geq\eta,
       \qquad\max_{\rm discarded}\nu_j\leq1-\eta
       \quad\text{for some }\eta>0.                 \tag{SG.2}
\]
The original full cutoff satisfies these inequalities for
all sufficiently large $c$, by the concentration bounds
already used in AC-T29 and AC-T33.

Define the dilation Sobolev norm by
\[
 (\mathcal Uf)(t)=e^{t/2}f(e^t),\qquad
       \|f\|_{H_D^s}=\|\mathcal Uf\|_{H^s(\mathbb R)}.
                                                               \tag{SG.3}
\]
The closed theta operator $\mathcal E_\chi$ is the
Mellin multiplier of AC-T21 and AC-T25, agreeing on
its smooth core with
$E_\chi f(x)=\sqrt x\sum_{n\geq1}\chi(n)f(nx)$.
For $u=3/8$, $\vartheta=5/16$ and $f=Pf\in H_D^u$,
\[
 \begin{split}
 \|\mathcal E_\chi(1-\Pi_c)f\|_2
 \leq{}&C_{\chi,\eta}
        \|(1-Q_c)f\|_{\mathcal H}^{1/6}\\
 &\quad\cdot
   \bigl(\|f\|_{H_D^u}+(1+c)^u\|f\|_{\mathcal H}\bigr)^{5/6}.
 \end{split}                                        \tag{SG.4}
\]
All vectors on the left belong to the closed theta domain.

There is a corresponding column estimate.  If
$T:\mathcal A\to\mathcal H$ is Hilbert--Schmidt,
$PT=T$, and $T$ is Hilbert--Schmidt into $H_D^u$, then
\[
 \begin{split}
 \|\mathcal E_\chi(1-\Pi_c)T\|_2
 \leq{}&C_{\chi,\eta}
              \|(1-Q_c)T\|_2^{1/6}\\
 &\quad\cdot
 \bigl(\|T\|_{2,\mathcal A\to H_D^u}
                     +(1+c)^u\|T\|_2\bigr)^{5/6}.
 \end{split}                                        \tag{SG.5}
\]
Here the subscript $2$ on operators denotes the
Hilbert--Schmidt norm, while the left side of (SG.4)
is the target $L^2(dx/x)$ norm.
In particular, if a family $T_c$ satisfies for every $m$
\[
 \|(1-Q_c)T_c\|_2\leq C_m c^{-m}A_c,
 \qquad
 \|T_c\|_{2,\mathcal A\to H_D^u}
                 +c^u\|T_c\|_2\leq Cc^u A_c,
                                                               \tag{SG.6}
\]
then $\|\mathcal E_\chi(1-\Pi_c)T_c\|_2
\leq C_Nc^{-N}A_c$ for every prescribed $N$.
The constants in (SG.6) must be uniform for the family;
no such uniformity is inferred from pointwise smoothness.

For every integer $m\geq1$, let $X_m$ be the parity
subspace of $H^m(\mathbb R)$ supported in $[-1,1]$,
with its ordinary spatial Sobolev norm.  Then
\[
 \|\mathcal E_\chi(1-\Pi_c)\|_{X_m\to L^2(dx/x)}
                \leq C_{\chi,\eta,m}
                         c^{-m/6+5/16}.             \tag{SG.7}
\]
Consequently a fixed $f\in C_c^\infty((-1,1))$ of the
prescribed parity is approximated by $\Pi_cf$ in
both raw $L^2$ norm and theta graph norm faster than
any prescribed negative power of $c$.
If $T:\mathcal A\to X_m$ belongs to a Schatten class
$\mathcal S_p$, $1\leq p\leq\infty$, composition in
(SG.7) gives the same bound multiplied by
$\|T\|_{p,\mathcal A\to X_m}$.
Thus fixed smooth compact kernel operators whose output
has the prescribed parity and is supported strictly inside
$(-1,1)$ have the same arbitrary-power convergence in
trace norm and Hilbert--Schmidt norm.
\end{theorem}
\begin{proof}
\proofstep{1. Raw concentration and the sharp endpoint.}
Expand $f=Pf$ in the time-supported prolate basis.
For every discarded index, $1-\nu_j\geq\eta$.
The spectral expansion and the projection identity give
\[
 \begin{split}
 \|(1-\Pi_c)f\|_{\mathcal H}^2
 &\leq\eta^{-1}\langle f,(1-S_c)f\rangle\\
 &=\eta^{-1}\|(1-Q_c)f\|_{\mathcal H}^2.
 \end{split}                                        \tag{SG.8}
\]
The same inequality holds after summing over an
orthonormal basis of a Hilbert--Schmidt input space.

For $\psi\in\operatorname{ran}\Pi_c$, put
$\phi=S_c^{-1}\psi$ on the retained finite-dimensional
space.  Then $\|\phi\|\leq\eta^{-1}\|\psi\|$ and
$\psi=Q_c\phi$ inside $(-1,1)$.
Band limitation and Plancherel therefore imply
\[
       \|\psi'\|_{L^2(0,1)}
             \leq2\pi c\eta^{-1}\|\psi\|_{\mathcal H}.
                                                               \tag{SG.9}
\]
The positive-half-line normalization follows from parity.
Set $v(t)=e^{t/2}\psi(e^t)$ for $t<0$.
Then $\|v\|_2=\|\psi\|$ and
$\|v'\|_2\leq C_\eta(1+c)\|\psi\|$.
The zero extension of $v$ above zero has a possible jump.
For $0<s<1/2$, the explicit zero-extension estimate
AC-T28, equation (TB.33), applies with this endpoint:
\[
 \|\psi\|_{H_D^s}
   \leq C_{\eta,s}(1+c)^s\|\psi\|_{\mathcal H}.
                                                               \tag{SG.10}
\]
To recall its endpoint justification, integration by
parts gives $i\tau\widehat v$ as the transform of
$v'$ together with one term $v(0)$, and
$|v(0)|^2\leq2\|v\|_2\|v'\|_2$.
Splitting at $1+\|v'\|_2/\|v\|_2$ proves the estimate;
the boundary tail is integrable exactly when $s<1/2$.
Thus (SG.10) includes the actual sharp time cutoff.

\proofstep{2. The theta graph estimate.}
The unconditional bound already proved in AC-T25,
equation (TB.16), is
\[
 |L(1/2+i\tau,\chi)|
                    \leq C_\chi(1+|\tau|)^{5/16}.
                                                               \tag{SG.11}
\]
It follows from the absolutely convergent line
$\operatorname{Re}s=17/16$, the functional equation
on $\operatorname{Re}s=-1/16$, and strip convexity;
it uses no zero-location assumption.
By the closed Mellin realization,
\[
       \|\mathcal E_\chi h\|_2
                      \leq C_\chi\|h\|_{H_D^{\vartheta}},
                  \qquad\vartheta=5/16.             \tag{SG.12}
\]
In particular every vector in that Sobolev space
belongs to the closed multiplier domain.

Let $h=(1-\Pi_c)f$ and $u=3/8$.
Fourier H\"older interpolation gives
\[
 \|h\|_{H_D^{\vartheta}}
       \leq\|h\|_{\mathcal H}^{1-\vartheta/u}
                           \|h\|_{H_D^u}^{\vartheta/u}.
                                                               \tag{SG.13}
\]
Here $\vartheta/u=5/6$.
Equations (SG.8) and (SG.10), and
$\|\Pi_cf\|\leq\|f\|$, prove (SG.4).
For (SG.5), perform exactly the same H\"older inequality
on the product of Fourier measure and counting measure
over an orthonormal basis of $\mathcal A$.
This gives
\[
 \|h\|_{2,\mathcal A\to H_D^{\vartheta}}
 \leq\|h\|_2^{1-\vartheta/u}
                     \|h\|_{2,\mathcal A\to H_D^u}^{\vartheta/u}
\]
for a column operator $h$.
The summed forms of (SG.8) and (SG.10) prove the
stated bound.  Finite-column approximation extends it
to all the stated Hilbert--Schmidt maps, using
closedness of the theta multiplier.
Substitution of (SG.6) gives
$C_m c^{-m/6+5/16}A_c$; choosing $m$ sufficiently
large proves its consequence.

\proofstep{3. Spatial smoothness and Schatten composition.}
For $f\in X_m$, Plancherel bounds its Fourier tail by
\[
       \|(1-Q_c)f\|\leq C_m c^{-m}\|f\|_{H^m}.
\]
Also $\|f\|_{H_D^u}\leq C\|f\|_{H^1}$.
Indeed differentiating $e^{t/2}f(e^t)$ gives the
two terms $e^{t/2}f(e^t)/2$ and $e^{3t/2}f'(e^t)$;
the latter has squared integral
$\int_0^1u^2|f'(u)|^2du\leq\|f'\|_2^2$.
The trace at $u=1$ is zero for the full-line
$H^1$ zero extension, so no boundary distribution
appears in this assertion.
Consequently the second factor in (SG.4) is bounded
by $C_m c^u\|f\|_{H^m}$ for $c\geq1$.
This proves (SG.7), including its exponent
$-m/6+(5/6)(3/8)=-m/6+5/16$.
The raw norm estimate follows separately from (SG.8).

The Schatten conclusion is the operator-ideal inequality
$\|AT\|_p\leq\|A\|\|T\|_p$ applied to the bounded map
in (SG.7).
For a fixed smooth kernel compactly supported in both
variables, a Fourier expansion in its input variable
has coefficients rapidly decreasing in every output
$H^m$ norm, by integration by parts.
The resulting rank-one series is absolutely summable
as a map into $X_m$ whenever its output support is
contained in $(-1,1)$.
Thus it is trace class into every such $X_m$, proving
the last assertion.
\end{proof}

\paragraph{Exact physical scaling and remaining scope.}
For the actual time radius $\Lambda$ and frequency radius
$\Lambda/q$, one has $c=\Lambda^2/q$ after the unitary
dilation $D_\Lambda f(u)=\Lambda^{-1/2}f(u/\Lambda)$.
This dilation preserves all the $H_D^s$ norms, and
$\mathcal E_\chi D_\Lambda=V_\Lambda\mathcal E_\chi$
with $V_\Lambda F(x)=F(x/\Lambda)$ unitary on the target.
Therefore all estimates above apply to the actual
conductor-normalized finite prolate family, with that
exact substitution.  They control smooth profiles or
families satisfying the displayed norm hypotheses.
They do not assert uniform boundedness of $\Pi_c$ on
$H_D^s$, convergence of every varying boundary family,
or the needed comparison of arithmetic polar projections.

\InputIfFileExists{modules/sites_followup.tex}{}{}
\InputIfFileExists{modules/cm_twisted_selmer.tex}{}{}
\section{The exact higher K-theory quantity behind the strict Selmer length}
\label{sec:selmer-k-theory}

The finiteness in AC-G22 can be made into an exact order identity.
This identifies a concrete arithmetic quantity to estimate; it does
not replace an estimate by finiteness.  Throughout this section
$e\geq1$, and $F$ is a totally imaginary number field.  Put
\[
 A_{\ell,e}=\mathbb Q_\ell/\mathbb Z_\ell(-e),\qquad
 S_\ell(F,e)=\ker\left(H^1(G_{F,S_\ell},A_{\ell,e})
       \longrightarrow\bigoplus_{v\mid\ell}
                         H^1(F_v,A_{\ell,e})\right),
\]
where $S_\ell$ comprises the primes above $\ell$ and the complex
places.  Complex places contribute zero modified local cohomology.
This is the untwisted strict group used in AM-H06--AM-H08.
Write $B^\vee=\operatorname{Hom}_{\rm cont}(B,\mathbb Q_\ell/\mathbb Z_\ell)$
for Pontryagin duality on the relevant $\ell$-primary groups.

\begin{proposition}[AC-K10: strict order and the local correction]
\label{AC-K10}
There is an exact sequence of finite groups
\[
 0\longrightarrow S_\ell(F,e)^\vee
 \longrightarrow K_{2e}(\mathcal O_F)_{(\ell)}
 \longrightarrow\bigoplus_{v\mid\ell}
                     H^0(F_v,A_{\ell,e})^\vee
 \longrightarrow H^0(F,A_{\ell,e})^\vee
 \longrightarrow0.                                      \tag{K10.1}
\]
The middle identification is the standard \'etale Chern-character
identification; the last arrow is dual to restriction of invariants.
In particular
\[
 \#S_\ell(F,e)=
 \#K_{2e}(\mathcal O_F)_{(\ell)}\,
 \frac{\#H^0(F,A_{\ell,e})}
      {\prod_{v\mid\ell}\#H^0(F_v,A_{\ell,e})}.          \tag{K10.2}
\]
All these groups are finite, including when $\ell=2$.  For fixed
$F,e$, only finitely many $S_\ell(F,e)$ are nonzero, and
\[
 \sum_\ell \log\ell\,
       \operatorname{length}_{\mathbb Z_\ell}S_\ell(F,e)
 \leq\log\#K_{2e}(\mathcal O_F).                         \tag{K10.3}
\]
More generally, for every integer $t\geq2$, the left side summed
over $1\leq e<t$ is at most
$\sum_{1\leq e<t}\log\#K_{2e}(\mathcal O_F)$.
\end{proposition}
\begin{proof}
Take the compact representation $T=\mathbb Z_\ell(e+1)$.
Its Tate dual $T^\vee(1)$ is exactly $A_{\ell,e}$.
The compact form of Poitou--Tate duality identifies the kernel
of localization in $H^2(G_{F,S_\ell},T)$ with
$S_\ell(F,e)^\vee$.  The final part of the Poitou--Tate sequence is
\[
 H^2(G_{F,S_\ell},T)\longrightarrow
 \bigoplus_{v\mid\ell}H^2(F_v,T)\longrightarrow
 H^0(F,A_{\ell,e})^\vee\longrightarrow0.                 \tag{K10.4}
\]
Local Tate duality identifies $H^2(F_v,T)$ with
$H^0(F_v,A_{\ell,e})^\vee$, with the maps in (K10.4) dual
to the usual restrictions.  The exact statements used are
R.~Sharifi, \emph{Iwasawa Theory}, Appendix~A, Theorems
A.0.3, A.0.6 and A.0.8, with the compact-coefficient extensions
in Remarks A.0.4 and A.0.9.  Their use here includes continuous
compact cohomology, not an unjustified replacement of finite
coefficients in a finite-level sequence.  Because $F$ has no
real places, there is no real-place modification at $\ell=2$.

The integral \'etale comparison gives
\[
 H^2(G_{F,S_\ell},\mathbb Z_\ell(e+1))
 =H^2_{\rm et}(\mathcal O_F[1/\ell],\mathbb Z_\ell(e+1))
 \simeq K_{2e}(\mathcal O_F)_{(\ell)}.                  \tag{K10.5}
\]
Here the first equality is the usual arithmetic Galois-cohomology
description.  The precise K-theory result is C.~Weibel,
\emph{The K-book}, Chapter~VI, Theorem~8.2, with its totally
imaginary formulation in Theorem~8.4 (author's online chapter,
pp.~47--48).  It applies for $2e\geq2$ and includes $\ell=2$
under the present hypothesis.  These are cited known theorems;
their proofs are not claimed as new results of this run.
The even group $K_{2e}(\mathcal O_F)$ is finite.  Substitution
in (K10.4), together with the kernel identification, gives
(K10.1).

For explicit finiteness of the invariant terms, the cyclotomic
image of an absolute number-field or $\ell$-adic local Galois
group is an open subgroup of $\mathbb Z_\ell^\times$.
It contains an element $u$ of infinite order.  Since $e>0$,
$u^{-e}-1\ne0$, and the invariant subgroup of
$\mathbb Q_\ell/\mathbb Z_\ell(-e)$ is killed by that nonzero
scalar.  It is consequently finite.  Taking alternating orders
in (K10.1) proves (K10.2).

Restriction injects global invariants into the invariants at
each chosen place: both are subgroups of the same coefficient
module, and the global Galois group contains its decomposition
group.  In particular the quotient of invariant orders in
(K10.2) is at most one.  Thus the order of $S_\ell(F,e)$ divides
the order of the indicated $\ell$-primary K-group; the injection
in (K10.1) gives the same conclusion.  All but finitely many
such K-groups are zero.  Taking logarithms and summing proves
(K10.3) and its finite-weight version.
\end{proof}

\begin{proposition}[AC-K11: the exact local subtraction for the cubic field]
\label{AC-K11}
Let $K=\mathbb Q(\omega)$ and $E=K(\sqrt[3]{269})$.  Define
\[
 \kappa_e=v_3\bigl(\#K_{2e}(\mathcal O_E)\bigr),\qquad
 k_e=1+v_3(e),\qquad
 a_E(e)=\operatorname{length}_{\mathbb Z_3}S_3(E,e).
\]
Then
\[
                a_E(e)=\kappa_e-2k_e.                  \tag{K11.1}
\]
The corresponding strict group over $K$ is zero.  Hence the
finite cyclic exponent $m_e$ of AC-G22 satisfies
\[
 \kappa_e-2k_e-4\ \leq\ m_e\ \leq\
                       \kappa_e-2k_e+9.               \tag{K11.2}
\]
The lower bound may of course be strengthened to $m_e\geq1$.
These identities and inequalities use no RH hypothesis.
They reduce the positive-weight coefficient-prime-$3$ estimate
to the growth of a definite sequence of even K-groups, up to
an explicit logarithmic local term and a bounded integral
character-comparison error.  They do not supply a growth bound
for $\kappa_e$ or for the fields appearing in the source's
growing-host family.
\end{proposition}
\begin{proof}
The Hensel argument in AC-G22 proves that $269$ is a cube in
$\mathbb Q_3$.  The three roots of $X^3-269$ are then in
$K_3=\mathbb Q_3(\omega)$, so the unique prime of $K$ over $3$
splits into three primes of $E$, each with completion $K_3$.
At the prime $(269)$ of $K$ the cubic extension $E/K$ is
totally ramified, while the cyclotomic $3$-power tower is
unramified there.  Since $[E:K]=3$, its intersection with
$\mathbb Q(\mu_{3^\infty})$ is therefore exactly $K$.
It follows that the $3$-adic cyclotomic images for $G_E$ and
for all three completions are the same group
$1+3\mathbb Z_3$.

For that group the minimum of
$v_3(u^{-e}-1)$ is $k_e=1+v_3(e)$, by the binomial
lifting-of-exponent calculation; $u=4$ attains it.
Thus each of the four invariant groups in (K10.2), one global
and three local, has order $3^{k_e}$.  Formula (K10.2)
therefore gives
$3^{a_E(e)}=3^{\kappa_e}3^{k_e}/3^{3k_e}$,
which is (K11.1).

For $K$, one can use either the direct residual strict-kernel
calculation accompanying AC-G22 or the following known
K-theory computation.  The prime $3$ is regular because
$\mathbb Z[\omega]$ has class number one.  Weibel, Chapter~VI,
Example~8.3.2, gives
$K_{2e}(\mathbb Z[\omega])_{(3)}=0$ for all $e\geq1$.
The injection in (K10.1) proves $S_3(K,e)=0$.

The existing integral comparison AM-H08, (H08.2)--(H08.3),
for $r=3,p=269$, reads
\[
       -9\leq a_E(e)-a_K(e)-m_e\leq4.                 \tag{K11.3}
\]
All its finiteness hypotheses now hold by AC-K10 and AC-G22.
Its $S$ contains the primes above $3$ and $269$; its strict
inertia conditions identify the untwisted groups with
$S_3(E,e)$ and $S_3(K,e)$, as proved in AM-H06.  Thus the
comparison concerns the very groups used here, rather than
groups with a ramified local condition omitted.  Substitution
of $a_K(e)=0$ and (K11.1), followed by rearrangement, gives
(K11.2).
\end{proof}

\begin{proposition}[AC-K12: the single cyclic K-module at three]
\label{AC-K12}
For $E=\mathbb Q(\omega,\sqrt[3]{269})$,
$\operatorname{Cl}(E)_{(3)}=0$.
Let $\Gamma=\operatorname{Gal}(E/K)=\langle\sigma\rangle$.
Via $\sigma\mapsto\omega$, the action on
$H_e=K_{2e}(\mathcal O_E)_{(3)}$ factors through
$\mathcal O=\mathbb Z_3[\omega]$, and
\[
 H_e\simeq\mathcal O/\lambda^{\kappa_e}\mathcal O,
 \qquad
 S_3(E,e)^\vee\simeq
       \mathcal O/\lambda^{\kappa_e-2k_e}\mathcal O.     \tag{K12.1}
\]
The second isomorphism uses the induced $\Gamma$ action under
Poitou--Tate duality; length zero means the zero module.
In particular $\kappa_e\geq2k_e$.
This identifies the module structure, but gives no upper bound
on its exponent.
\end{proposition}
\begin{proof}
The only finite prime ramified in $E/K$ is $(269)$, with
ramification index three, by AC-K04; the primes over three
split by AC-G22.  There are no real places.
Chevalley's ambiguous class number formula gives
\[
 \#\operatorname{Cl}(E)^\Gamma
 =\frac{h_K\prod_v e_v}
 {[E:K]\,[\mathcal O_K^\times:
                 \mathcal O_K^\times\cap N_{E/K}E^\times]}
 =\frac1{[\mathcal O_K^\times:
                 \mathcal O_K^\times\cap N_{E/K}E^\times]}.
                                                               \tag{K12.2}
\]
The formula used is C.-F.~Yu, \emph{Ambiguous class number
formulas}, MPIM Preprint 2014(68), Theorem~1.1, equation~(1.2),
printed p.~1, with the ordinary-class interpretation on p.~2.
Its left side is a positive integer.  Thus the unit index and
the ambiguous class group order are both one.  If the finite
$3$-group $\operatorname{Cl}(E)_{(3)}$ were nonzero, its
nonzero $3$-torsion vector space would have a nonzero fixed
vector: on a vector space in characteristic three,
$(\sigma-1)^3=\sigma^3-1=0$ has nontrivial kernel.
This contradicts (K12.2), proving the first assertion.

We next justify the action on the K-group, including the
ramified prime.  Transfer along the finite free extension
$\mathcal O_K\subset\mathcal O_E$ sends $H_e$ to
$K_{2e}(\mathcal O_K)_{(3)}=0$, as used in AC-K11.
Transfer followed by restriction on the fraction fields is
$1+\sigma+\sigma^2$.  The localization sequence injects
$K_{2e}(\mathcal O_E)$ into $K_{2e}(E)$, since the groups
$K_{2e}$ of its finite residue fields vanish for $e>0$.
Consequently the same norm operator is zero on $H_e$.
This proves that its $\mathbb Z_3[\Gamma]$ action factors
through $\mathbb Z_3[\Gamma]/(1+\sigma+\sigma^2)
\simeq\mathcal O$.  This transfer argument does not assume
that $\operatorname{Spec}\mathcal O_E\to\operatorname{Spec}\mathcal O_K$
is \'etale at $269$.

The mod-three comparison of Weibel, Chapter~VI, Corollary~8.3
and Example~8.3.1, gives
\[
 \dim_{\mathbb F_3}(H_e/3H_e)
 =\dim_{\mathbb F_3}\bigl(\operatorname{Cl}
                         (\mathcal O_E[1/3])/3\bigr)
            +\#\{v\mid3\}-1=2.                       \tag{K12.3}
\]
The localized class group has no three-part, being a quotient
of $\operatorname{Cl}(E)$.  There are exactly three places
over three.

The local quotient in (K10.1) can be identified equivariantly.
Choose a lift of $\sigma$ to $G_K$ with trivial $3$-adic
cyclotomic character; it exists because
$E\cap\mathbb Q(\mu_{3^\infty})=K$.
It cyclically permutes the three places above three and acts
trivially on their common invariant coefficients.  The last
map in (K10.1) is consequently the sum
\[
 (\mathbb Z/3^{k_e})^3\longrightarrow\mathbb Z/3^{k_e}.
\]
Its kernel is the augmentation module for $\Gamma$ over
$\mathbb Z/3^{k_e}$.  The map $a\mapsto(\sigma-1)a$ induces
an isomorphism from $\mathcal O=\mathbb Z_3[\Gamma]/(1+\sigma+\sigma^2)$
onto the integral augmentation ideal; reduction modulo $3^{k_e}$ therefore
identifies this kernel with $\mathcal O/3^{k_e}\mathcal O$.
Thus (K10.1) gives a surjection
\[
              H_e\longrightarrow
                 \mathcal O/\lambda^{2k_e}\mathcal O.  \tag{K12.4}
\]

Write the finite $\mathcal O$-module as
$H_e=\bigoplus_i\mathcal O/\lambda^{n_i}$ with $n_i>0$.
Equation (K12.3) says $\sum_i\min(2,n_i)=2$.
The quotient (K12.4), with $2k_e\geq2$, forces at least
one $n_i\geq2k_e$.  Its summand already contributes two
to that sum, so there are no other summands.
Its length is $\kappa_e$, since the residue field has order
three.  The kernel of the surjection (K12.4) from this
cyclic module is cyclic of length $\kappa_e-2k_e$.
By (K10.1) it is precisely $S_3(E,e)^\vee$, proving (K12.1).
\end{proof}

\paragraph{Exact remaining size statement.}
For any totally imaginary host field $F$ and integer $t$, AC-K10
supplies the concrete sufficient estimate
\[
 \sum_{1\leq e<t}\log\#K_{2e}(\mathcal O_F)\leq B
 \quad\Longrightarrow\quad
 \sum_\ell\log\ell\sum_{1\leq e<t}
       \operatorname{length}_{\mathbb Z_\ell}S_\ell(F,e)
       \leq B.                                        \tag{K11.4}
\]
The quantifiers over the fields, weights and allowed bound in
the source predicate \texttt{VVle} must still be satisfied.
An unspecified finite bound for each fixed $F,t$ is not a
uniform bound with the source's prescribed $B$.
Nor is an estimate over the fixed field $E$ automatically an
estimate over the varying CM hosts in \texttt{SupplyChi}.
This is the exact limitation of the K-theory route at this stage.

\paragraph{Algebraic-geometric dictionary.}
The Kummer character is the lisse rank-one coefficient system
of the finite \'etale torsor over
$\operatorname{Spec}\mathcal O_K[1/(3\cdot269)]$ (SGA~1,
XI~\S6, Propositions~1 and~4 and Corollary~5).  AC-K10 places
its untwisted restriction in localization and duality on the
arithmetic scheme $\operatorname{Spec}\mathcal O_E[1/\ell]$;
AC-K11 gives the numerical correction for its three local
fibers over $3$.  The passage from algebraic K-theory to these
cohomology groups is the modern \'etale comparison just cited,
not a theorem attributed to a Grothendieck archive scan.
The exterior-power and coefficient-tensor language belongs to
SGA~6, V~2.2(b), whereas the order estimate in (K11.4) is an
additional arithmetic requirement.  No previously inaccessible
Grothendieck scan is represented as reread by this argument.

\InputIfFileExists{modules/cm_twist_transport.tex}{}{}
\InputIfFileExists{modules/exact_cubic_descent.tex}{}{}
\section{A certified logarithmic invariant and exact cubic weights}
\label{sec:cubic-logarithmic}

This section retains $K=\mathbb Q(\omega)$,
$E=K(\alpha)$ with $\alpha^3=269$,
$\mathcal O=\mathbb Z_3[\omega]$ and $\lambda=\omega-1$.
The degree-zero logarithmic class group at three is denoted by
$\widetilde{\operatorname{Cl}}_3(E)$.  Degree zero is part of this
notation: the full logarithmic divisor-class group has an additional
free degree direction.  Write
\[
 W_e(E)=\ker\left(H^2(G_{E,S_3},\mathbb Z_3(e+1))
       \longrightarrow\bigoplus_{v\mid3}
                   H^2(E_v,\mathbb Z_3(e+1))\right) \quad(e\geq1).
\]
By AC-K10 and AC-K12 this is the wild \'etale kernel
$S_3(E,e)^\vee$, and
\[
 W_e(E)\simeq\mathcal O/\lambda^{a_e},\qquad
 a_e=\kappa_e-2(1+v_3(e))=m_e-1.                    \tag{KL.1}
\]
We also use $W_0(E)=\widetilde{\operatorname{Cl}}_3(E)$, the
weight-zero convention in the Iwasawa description below.

\begin{proposition}[AC-K13: the residual logarithmic invariant]
\label{AC-K13}
For every $e\geq1$,
\[
 \min\{2,m_e-1\}
 =\dim_{\mathbb F_3}\bigl(
       \widetilde{\operatorname{Cl}}_3(E)/3\bigr).       \tag{KL.2}
\]
Thus a residual logarithmic dimension zero would give $m_e=1$
for every $e$, dimension one would give $m_e=2$ for every $e$,
and dimension two gives $m_e\geq3$ for every $e$.
For the present field, the certified calculation described below gives
\[
 \operatorname{Cl}(E)\simeq\mathbb Z/7\oplus\mathbb Z/7,
 \qquad
 \widetilde{\operatorname{Cl}}_3(E)
       \simeq\mathbb Z/9\oplus\mathbb Z/3.             \tag{KL.3}
\]
Consequently the third alternative applies.  Neither theorem A
nor a conjecture about logarithmic class groups is used.
The group isomorphisms in (KL.3) are computer-assisted arithmetic
statements, with an unconditional certification command and its
returned result retained in the evidence files.
\end{proposition}
\begin{proof}
The precise known result used is J.-F.~Jaulent and A.~Michel,
\emph{Approche logarithmique des noyaux \'etales sauvages des corps
 de nombres}, J.~Number Theory \textbf{120} (2006), 72--91,
Theorem~3 and its complete proof on p.~7 of the author's 15-page
version:
\url{https://www.math.u-bordeaux.fr/~jjaulent/Articles/J44.pdf}.
Apply it with the auxiliary prime-to-three extension equal to the
identity extension of $E$ and with $r=1$.  Its cyclotomic
procyclicity hypothesis holds for the odd prime three, and
$\mu_3\subset E$.  It gives
\[
 W_e(E)/3\simeq\mu_3^{\otimes e}\otimes_{\mathbb Z_3}
                   \widetilde{\operatorname{Cl}}_3(E). \tag{KL.4}
\]
In particular their dimensions over $\mathbb F_3$ agree.
Since $3$ is a unit times $\lambda^2$ and the residue field of
$\mathcal O$ has three elements, (KL.1) gives
$\dim_{\mathbb F_3}(W_e/3)=\min(2,a_e)$.  AC-G26 now proves
(KL.2) and its three alternatives.  The source explicitly does
not assume Gross--Kuz'min, Schneider or Leopoldt conjectures
for its quotient isomorphism; no finiteness of a weight-zero
class group was inserted into this deduction.

Here are the exact input and certification for (KL.3).  With
$z=\omega-\alpha$, the resultant of
$Y^3-269$ and $(z+Y)^2+(z+Y)+1$ is
\[
 P(z)=z^6+3z^5+6z^4+545z^3+813z^2-804z+71824.
                                                               \tag{KL.5}
\]
The element $z$ generates $E$: indeed
\[
 \alpha=\frac{269-(2z+1)(z^2+z+1)}{3z(z+1)},
 \qquad \omega=z+\alpha,
\]
and $z(z+1)\ne0$.  Thus there is no unverified field
identification between the computational polynomial and $E$.
PARI/GP~2.19.1, built from the official source release, was run
on this exact polynomial with \texttt{bnfinit(P,1)}.
It returned discriminant $-141375086667$, class group
\texttt{[7,7]}, and class number $49$.
The complete command \texttt{bnfcertify(b)} returned \texttt{1};
no partial-certification flag was used.  Then
\texttt{bnflog(b,3)} returned
\texttt{[[9,3],[9,3],[]]}.
The first component is precisely the degree-zero logarithmic
class group.  The other components are its subgroup generated
at three and the ordinary $S_3$-class group, respectively.
The meaning and proof of this algorithm are K.~Belabas and
J.-F.~Jaulent, \emph{The logarithmic class group package in
PARI/GP}, Publ.~Math.~Besan\c con (2016), 5--18,
author version \S\S3.2--3.4, pp.~9--11, and \S4.1, p.~12:
\url{https://www.math.u-bordeaux.fr/~kbelabas/research/bnflog-2016.pdf}.
Its finite stopping criterion is proved in Lemma~3.8; the
initial class-group/unit data were unconditionally certified
here.  The reproducible script, complete raw output and binary
arithmetic data are
\texttt{evidence/cubic\_logarithmic/certify\_logclass.gp},
\texttt{certify\_logclass.out}, and
\texttt{cubic\_bnf.bin}.  The runtime build receipt is
\texttt{evidence/pari\_build.json}.  This proves the asserted
computer-assisted values and completes the residual conclusion.
\end{proof}

\begin{theorem}[AC-K14: finite-twist congruence and an exact infinite family]
\label{AC-K14}
Let $e,e'\geq0$ and $n\geq1$.  If
$e\equiv e'\pmod{3^{n-1}}$, then
\[
                   W_e(E)/3^n\simeq W_{e'}(E)/3^n.    \tag{KL.6}
\]
In particular, for every positive integer $e$ divisible by three,
\[
 \boxed{m_e=4,\qquad a_e=3,\qquad
        \kappa_e=2(1+v_3(e))+3.}                     \tag{KL.7}
\]
More generally, an exact value $a_e=a$ at one positive weight
fixes $a_{e'}=a$ and $m_{e'}=a+1$ at every positive weight
\[
                    e'\equiv e\pmod{3^{\lfloor a/2\rfloor}}.
                                                               \tag{KL.8}
\]
Thus the remaining weights can be investigated by finite
certificates on residue classes.  This theorem alone does not
give an upper bound on the two nonzero classes modulo three.
\end{theorem}
\begin{proof}
We spell out why (KL.6) uses a weaker condition than requiring
$\mu_{3^n}\subset E$.  Let $E_\infty/E$ be the cyclotomic
$\mathbb Z_3$-extension, $\Gamma_\infty=\langle\gamma\rangle$,
and let $X'$ be the Galois group of the maximal abelian pro-three
extension of $E_\infty$ unramified everywhere and split at
its primes above three.  The Schneider description used and
stated in the proof of Jaulent--Michel Theorem~3 is
\[
 W_e(E)=\bigl(\mathbb Z_3(e)\otimes_{\mathbb Z_3}X'\bigr)
                   _{\Gamma_\infty},
 \qquad
 W_0(E)=X'_{\Gamma_\infty}
             =\widetilde{\operatorname{Cl}}_3(E).    \tag{KL.9}
\]
The quotients are compact coinvariants.  At positive weights
this is the usual wild \'etale kernel occurring in (KL.1).
These known identifications, not a conjectural characteristic
ideal formula, are the only Iwasawa-theoretic input here.

Choose a basis of the Tate module and put
$u=\varepsilon_3(\gamma)\in1+3\mathbb Z_3$.
The quotient in (KL.9) modulo $3^n$ is the quotient of
$X'/3^nX'$ by the image of $u^e\gamma-1$.
Lifting the exponent in $1+3\mathbb Z_3$ shows
\[
 e\equiv e'\pmod{3^{n-1}}
       \quad\Longrightarrow\quad
                    u^e\equiv u^{e'}\pmod{3^n}.
\]
The two endomorphisms on $X'/3^nX'$ therefore agree, proving
(KL.6).  Images and quotients can equally be taken with closures:
they are continuous images of compact groups.  Changing the
chosen Tate basis only changes the displayed identification by
a unit.  If desired the isomorphism is compatible with the
cubic action by choosing its lift to fix all cyclotomic roots,
which is possible since $E\cap\mathbb Q(\mu_{3^\infty})=K$.

For $3\mid e$, take $e'=0$ and $n=2$.
By (KL.3), the right side of (KL.6) has order $27$.
By (KL.1), the left side has order
$3^{\min(a_e,4)}$.  Hence
$\min(a_e,4)=3$, which forces $a_e=3$.
Equation (KL.1) gives all of (KL.7).

Finally suppose $e,e'>0$ and $a_e=a$.
Take $n=\lfloor a/2\rfloor+1$, so that $2n>a$.
Under (KL.8), (KL.6) and the cyclic structure (KL.1) imply
\[
            \min(a_{e'},2n)=\min(a,2n)=a<2n.
\]
Consequently $a_{e'}=a$, and AC-G26 gives the equality of $m$.
This proves the general congruence assertion.
\end{proof}

\begin{remark}[The exact small computational gate]
The argument does not identify logarithmic-class vanishing
with ordinary three-rationality.  In fact the ordinary class
group has no three-part while the logarithmic class group in
(KL.3) is nonzero.  The three local places over three account
for two residual Brauer directions in AC-K12.
For the remaining weights, certifying $a_1$ and $a_2$ at values
less than four would settle both nonzero classes modulo three
by (KL.8).  A larger value requires finer residue classes.
Neither an uncertified class-group calculation nor a claim
about a regulator or an Iwasawa main conjecture is used as
such a certificate.
\end{remark}

\section{A bounded-rank Iwasawa module for the cubic field}
\label{sec:cubic-iwasawa-tower}

The fixed coefficient-prime-three sequence admits a finite-dimensional
arithmetic description.  This description does not estimate a
specialization by declaring its characteristic series nonzero.
Retain $E=\mathbb Q(\omega,\sqrt[3]{269})$, and set
\[
 K_n=\mathbb Q(\mu_{3^{n+1}}),\quad E_n=EK_n,\quad
 A_n=\operatorname{Cl}(E_n)_{(3)},\quad
 X=\varprojlim_n A_n,
\]
where transition maps are norms.  Let $A'_n$ be the quotient of $A_n$
by the classes of primes above three, and put
$X'=\varprojlim_n A'_n$.  The latter is the split unramified module
in (KL.9).  Write $\sigma$ for a generator of
$\operatorname{Gal}(E_n/K_n)$, chosen compatibly, and let
$\mathcal O=\mathbb Z_3[\omega]$ act by $\omega=\sigma$.

\begin{theorem}[AC-K16: seven generators suffice in the whole tower]
\label{AC-K16}
For every $n\geq0$,
\[
 \operatorname{Cl}(K_n)_{(3)}=0,\qquad
 \#\{\mathfrak p\subset\mathcal O_{K_n}:\mathfrak p\mid269\}
       =3^{\min(n,2)}.                                  \tag{KI.1}
\]
The norm $1+\sigma+\sigma^2$ annihilates $A_n$, and
\[
 \dim_{\mathbb F_3}(A_n/\lambda A_n)
 =
 \begin{cases}
 0,&n=0,\\
 2,&n=1,\\
 7,&n\geq2,
 \end{cases}
 \qquad \lambda=\omega-1.                              \tag{KI.2}
\]
All transition norms $A_{n+1}\to A_n$ are surjective.
Both $X$ and $X'$ are finitely generated $\mathcal O$-modules.
The minimal number of generators of $X$ is exactly seven,
and that of $X'$ is either six or seven.  In particular they are finitely
generated over $\mathbb Z_3$, of rank at most fourteen; their
Iwasawa $\mu$-invariants vanish.  No assumption of A or of a
main conjecture enters these assertions.
\end{theorem}
\begin{proof}
\proofstep{1. The cyclotomic base and the ramified primes.}
The class number of $K_0=\mathbb Q(\omega)$ is one.
For $n\geq1$, the extension $K_n/K_{n-1}$ is cyclic of degree
three, with just one ramified finite prime, the unique prime
above three, totally ramified.  There are no real places.
Suppose inductively that the three-primary class group of
$K_{n-1}$ is zero.  The three-primary part of Chevalley's
ambiguous class-number formula bounds the invariant three-class
group of $K_n$ by $3/3=1$, since its remaining unit index is
an integer at least one.  A nonzero finite abelian three-group
acted on by a group of order three has nonzero fixed points:
on its nonzero three-torsion vector space, $(\sigma-1)^3=0$.
Thus the entire three-class group of $K_n$ is zero.  This proves
the first assertion in (KI.1) by induction.  The exact known
formula used is C.-F.~Yu, \emph{Ambiguous class number formulas},
MPIM preprint 14--68 (2014), Theorem~1.1, equation~(1.2),
printed p.~1, with the ordinary-class interpretation on p.~2.
Its application here is the same ordinary, totally imaginary
case used in AC-K12.

Since $269\equiv-1\pmod{27}$ and
$v_3(269^2-1)=3$, lifting the exponent gives
\[
 \operatorname{ord}_{3^{n+1}}(269)
       =2\,3^{\max(n-2,0)}.                             \tag{KI.3}
\]
For completeness, an odd power of $269$ is $-1$ modulo three;
for an even exponent $2b$ its distance from one has valuation
$3+v_3(b)$.  The smallest exponent producing one modulo
$3^{n+1}$ is exactly (KI.3).  Cyclotomic Frobenius is exponentiation
by $269$, so dividing $[K_n:\mathbb Q]=2\,3^n$ by (KI.3)
proves the second assertion in (KI.1).

The intersection $E\cap K_\infty$ is $K_0$: the cubic extension
$E/K_0$ ramifies at $269$, whereas $K_\infty/K_0$ does not.
Thus every $E_n/K_n$ has degree three.  At each prime over
$269$, the valuation of $269$ is one, and $X^3-269$ is
Eisenstein; these are totally ramified primes of index three.
The Hensel calculation in AC-G22 puts all three roots in
$\mathbb Q_3(\omega)$, so every prime over three splits in
$E_n/K_n$.  The discriminant of the defining cubic shows there
is no further ramification.  Consequently $E_n/K_n$ has exactly
$t_n=3^{\min(n,2)}$ ramified primes, all tame and finite.

\proofstep{2. A uniform genus bound, with the root-of-unity saving.}
For an ideal class in $A_n$, extension after norm to $K_n$ is
the action of $1+\sigma+\sigma^2$.  The norm class has three-power
order and the base three-class group is zero, so this action
vanishes.  The action therefore factors through
\[
 \mathbb Z_3[C_3]/(1+\sigma+\sigma^2)\simeq\mathcal O.
\]
Chevalley's formula gives
\[
 \#A_n^{\langle\sigma\rangle}
 =\frac{3^{t_n-1}}
 {[\mathcal O_{K_n}^{\times}:
   \mathcal O_{K_n}^{\times}\cap
              N_{E_n/K_n}(E_n^{\times})]_{(3)}}.        \tag{KI.4}
\]
Here $[\cdot]_{(3)}$ denotes the three-primary part of the
finite index.  In fact this unit quotient has exponent dividing
three, since the cube of any base unit is a norm.
Without using the denominator, (KI.4) gives the bounds zero
and two on the logarithms for $n=0,1$.

For $n\geq2$ the denominator is at least three.  Indeed, choose
a primitive root $\zeta_{3^{n+1}}\in K_n$.  At a prime over
$269$ the residue-field size is
$q=269^{2\,3^{n-2}}$, and (KI.3) gives
$v_3(q-1)=n+1$.  Reduction preserves the exact order
$3^{n+1}$ of this root of unity, because that order is prime
to $269$.  Hence its residue is not a cube in the residue-field
unit group: the three-part of that group has precisely order
$3^{n+1}$.  In a totally ramified tame extension of degree
three, reduction of the norm of a unit is its residue cubed.
Also a norm which is a unit must come from a unit, since
$v_{K_n}(Nx)=f\,v_{E_n}(x)$ and $f=1$.
Thus $\zeta_{3^{n+1}}$ is not even a local norm, and so cannot
be a global norm.  The unit quotient in (KI.4) is nontrivial
and has order at least three.  This proves
$\#A_n^{\langle\sigma\rangle}\leq3^7$.

The invariant group is $A_n[\lambda]$.  For a finite
$\mathcal O$-module, multiplication by $\lambda$ has kernel
and cokernel of the same cardinality.  Both are vector spaces
over $\mathcal O/\lambda=\mathbb F_3$.
The upper bounds in (KI.2) now follow from (KI.4).

We next check that the unit index in (KI.4) is exactly three
for $n\geq2$ and exactly one for $n=0,1$.
For any unit $u$ in the CM field $K_n$, $u/\overline u$
is a root of unity.  Indeed it is an algebraic integer unit
all of whose conjugates have modulus one; the elementary
Kronecker lemma follows by bounding the integral coefficients
of the minimal polynomials of its powers, so that two powers
coincide.  Thus $u^2$ is a root of unity times the real unit
$u\overline u$.  In a quotient of exponent three, squaring
is bijective.  It is therefore enough to examine roots of
unity and real units.

At a prime over $269$, complex conjugation is the unique
order-two residue Frobenius: it is exponentiation by
$269^{3^{\max(n-2,0)}}$.  A real unit reduces to the
subfield of that cardinality.  Its multiplicative group has
order prime to three, since this odd power of $269$ is
$-1$ modulo three.  Hence every real-unit residue is a cube.
In this tame local extension that is sufficient for being
a unit norm: principal units are cubes, and lifting a
cubic residue reduces to this assertion.  At every other
finite place the extension is unramified or split, so units
are local norms.  There is no condition at a complex place.
The cyclic Hasse norm theorem now makes every real unit
a global norm.  Its exact locator is Milne, \emph{Class Field
Theory}, version~4.03, VIII, Theorem~3.1, printed p.~234.
For unramified local unit norms one may use V,
Proposition~4.12(c), printed p.~177.

The roots of unity of $K_n$ have order $2\,3^{n+1}$.
Their norm obstruction has order at most three, and the
primitive root just tested shows it has order three for
$n\geq2$.  At $n=0,1$, the residue field has cardinality
$269^2$ and its cube subgroup already contains all the
$3^{n+1}$-th roots.  These roots are local norms everywhere
and hence global norms as well.  Thus the unit index is
one in those two cases.  Substitution in (KI.4) proves
all the equalities in (KI.2).

\proofstep{3. Surjective norms and the compact limit.}
The extension $E_{n+1}/E_n$ is cyclic of degree three and
is totally ramified at each of its primes above three: its
local extension there is the cyclotomic step
$\mathbb Q_3(\mu_{3^{n+2}})/\mathbb Q_3(\mu_{3^{n+1}})$.
Let $H_n/E_n$ be the Hilbert three-class field.  Since $H_n/E_n$
is unramified, $H_n\cap E_{n+1}=E_n$.
Its base change $H_nE_{n+1}/E_{n+1}$ remains unramified and
abelian, and has Galois group $\operatorname{Gal}(H_n/E_n)$.
The Hilbert three-class field of $E_{n+1}$ consequently
surjects onto this group.  Under Artin reciprocity this
restriction is exactly the ideal-class norm; hence
$A_{n+1}\to A_n$ is onto.  The known reciprocity and
functoriality statements used are J.~S.~Milne,
\emph{Class Field Theory}, version~4.03 (2020), Chapter~V,
Proposition~3.3 and Theorems~3.5--3.6, printed pp.~157--158.
One can check the norm compatibility directly on an unramified
prime using the power of Frobenius in Proposition~3.3.

The maps $A_{n+1}/\lambda A_{n+1}\to A_n/\lambda A_n$
are therefore onto, and their dimensions form a nondecreasing
sequence bounded by seven.  They are exactly seven from level two onward.
There is a canonical isomorphism
\[
 X/\lambda X\simeq\varprojlim_n A_n/\lambda A_n.       \tag{KI.5}
\]
Here is a compactness verification to make the passage exact.
The inverse limit of any exact sequence of finite abelian
systems is exact: for a coherent element of the quotient,
its possible lifts for each finite initial segment form
nonempty finite sets, and the closed finite-constraint sets
in the product have the finite intersection property.
Similarly, if $y_n\in\lambda A_n$ is a coherent sequence,
solutions of $\lambda x_n=y_n$ over a finite initial segment
are nonempty (choose a solution at its largest level and
apply norms).  Compactness gives compatible solutions at all
levels.  Thus $\varprojlim\lambda A_n=\lambda X$, which
proves (KI.5).

Lift a basis of the finite vector space (KI.5) to at most
seven elements of $X$, and let $B$ be their $\mathcal O$-span.
It is closed, being a continuous image of a finite product
of the compact ring $\mathcal O$.  The compact pro-three
quotient $C=X/B$ satisfies $C=\lambda C$.  Since
$\lambda^2=-3\omega$, it also satisfies $C=3C$.
Every finite continuous quotient of a compact pro-three
abelian group is a finite three-group; multiplication by
three cannot be surjective on such a group unless it is
zero.  Hence all finite quotients of $C$ are zero, which
forces $C=0$.  This proves finite generation and the bound.
The finite-system compactness argument also makes
$X\to\varprojlim A'_n=X'$ surjective.  Its kernel is
cyclic over $\mathcal O$.  To see this, choose one compatible
sequence of primes $\mathfrak p_n$ above three.  Total
ramification in each tower step gives
$N[\mathfrak p_{n+1}]=[\mathfrak p_n]$.  Their three-primary
classes therefore define $p\in X$.  At every level the
three prime classes are its $\sigma$-translates; their sum
is the class of the principal extended cyclotomic prime,
so their span is $\mathcal O p_n$.  For any coherent sequence
in these spans, the equations $b p_n=y_n$ have a simultaneous
solution $b\in\mathcal O$ by the finite-intersection property
in the compact ring $\mathcal O$.  Consequently the kernel
is exactly $\mathcal O p$.  Since (KI.5) has dimension seven,
quotienting by this cyclic module leaves a reduction modulo
$\lambda$ of dimension six or seven.  This proves the
asserted minimal generator counts.

Finally a finitely generated $\mathbb Z_3$-module with a
continuous cyclotomic action is a torsion Iwasawa module
without a factor supported at the prime $(3)$; equivalently
its $\mu$-invariant is zero.  Directly, its quotient modulo
three is finite-dimensional over $\mathbb F_3$, whereas any
nonzero $\Lambda/(3^a)$ elementary factor has infinite-dimensional
reduction modulo three.  This proves the last assertion.
\end{proof}

\begin{corollary}[AC-K16: one matrix and one finite correction]
\label{AC-K16-matrix}
Let $T=X'_{\rm tors}$, with torsion taken over $\mathcal O$,
and write $X'/T\simeq\mathcal O^r$, where $0\leq r\leq7$.
Choose a cyclotomic generator $\gamma$ with
$\varepsilon_3(\gamma)=4$, and let
$G\in\operatorname{GL}_r(\mathcal O)$ be its action on this
free quotient.  For each $e\geq1$,
\[
 \begin{split}
 a_e=m_e-1
  &=\operatorname{length}_{\mathcal O}
          \bigl(T/(4^e\gamma-1)T\bigr)
      +v_\lambda\det(4^eG-I_r),\\
 0&\leq\operatorname{length}_{\mathcal O}
          \bigl(T/(4^e\gamma-1)T\bigr)
      \leq\operatorname{length}_{\mathcal O}T.
 \end{split}                                             \tag{KI.6}
\]
The determinant is nonzero.  For $r=0$ it is interpreted as
one, and the sequence $m_e$ is uniformly bounded.
For positive $r$ this is a reduction to the evaluations of
one degree-at-most-seven polynomial, with fixed finite error.
No lower bound for those nonzero evaluations is asserted.
\end{corollary}
\begin{proof}
The torsion submodule of a finitely generated module over
$\mathcal O$ is finite and stable under $\gamma$.
A generator with cyclotomic value four exists since the
cyclotomic image is $1+3\mathbb Z_3$ and four generates that
group.  Apply the endomorphism $b_e=4^e\gamma-1$ to
$0\to T\to X'\to\mathcal O^r\to0$.
By (KL.9) its cokernel on $X'$ is the finite group $W_e(E)$.
Its cokernel on the free quotient is therefore finite, so
$\det(4^eG-I_r)\ne0$ and its kernel on that free quotient
is zero.  The snake lemma consequently gives the exact sequence
\[
 0\longrightarrow T/b_eT\longrightarrow W_e(E)
 \longrightarrow\mathcal O^r/(4^eG-I_r)\mathcal O^r
 \longrightarrow0.                                    \tag{KI.7}
\]
Smith normal form computes the length of the last group as
the determinant valuation.  Taking lengths, using (KL.1),
proves (KI.6).  The asserted error bound is simply that a
quotient of the fixed finite module $T$ has no greater length.
\end{proof}

\paragraph{Geometric meaning and the remaining estimate.}
At finite level, $A_n$ describes the maximal abelian
three-primary finite \'etale covering of
$\operatorname{Spec}\mathcal O_{E_n}$; quotienting by the
three-prime classes requires those primes to split.
This is the Galois-category setting of SGA~1, V~\S4 and
\S7, with arithmetic class field theory specifying its
abelian quotient.  The inverse limit is a compact
$\mathcal O$-module, not a finite-dimensional pure motive
or an arithmetic surface.  The Tate-twist specialization
is the wild-kernel comparison of (KL.9), related to higher
K-theory by AC-K10.  No Grothendieck archive scan asserts
the new generator bound or supplies the missing size estimate.
The root-of-unity norm argument is an integral ramification
calculation, not a positivity statement about the cubic
Weil distribution.  H04.4 and the comparison with Alp\"oge's
growing hosts remain open after (KI.6).

\section{The remaining three-primary growth gate has rank one}
\label{sec:cubic-rank-gate}

The finite logarithmic certificate has a strong consequence when
combined with the rational complex-conjugation symmetry of this
particular tower.  We use $X'$, $T$, $r$, $G$ and $\gamma$ from
AC-K16, with $\varepsilon_3(\gamma)=4$.

\begin{theorem}[AC-K17: rank alternatives and a finite arithmetic test]
\label{AC-K17}
Let $r=\operatorname{rank}_{\mathcal O}X'$, so $0\leq r\leq7$.
If $r=0$, the sequence $(m_e)_{e\geq1}$ is uniformly bounded.
If $r\geq2$, then
\[
                         m_e=4\quad(e\geq1).           \tag{KR.1}
\]
If $r=1$, there are $g\in1+3\mathbb Z_3$ with
$v_3(g-1)=1$ and $k\in\{5,6\}$ such that
\[
 T\simeq\mathbb F_3[B]/(B^k),\quad B=\gamma-1,
 \qquad
 m_e=2+2v_3(4^eg-1)\quad(e\geq1).                     \tag{KR.2}
\]
In this last case the finite weight-zero invariant of the
first cyclotomic layer necessarily satisfies
\[
 \#\widetilde{\operatorname{Cl}}_3(E_1)=3^7,\qquad
 E_1=\mathbb Q(\mu_9,\sqrt[3]{269}).                    \tag{KR.3}
\]
Consequently an unconditional upper bound strictly less than
$3^7$ for that logarithmic class group would prove that
$(m_e)$ is uniformly bounded.  The tentative arithmetic output
for $E_1$ is not used in this theorem.
\end{theorem}
\begin{proof}
\proofstep{1. The rational characteristic polynomial.}
Complex conjugation $c$ acts on the whole field $E_\infty$.
It conjugates $\sigma$ to $\sigma^{-1}$ and commutes with
$\gamma$.  These identities follow from
\[
 \operatorname{Gal}(E_\infty/\mathbb Q)
       \simeq S_3\times\mathbb Z_3:
\]
the fiber product of the $S_3$-extension $E/\mathbb Q$
and the abelian cyclotomic extension over their common quadratic
field identifies the cyclotomic pro-three factor as central.
The induced action of $c$ on $X'$ is therefore semilinear for
$\omega\mapsto\omega^2$, and commutes with $\gamma$.
It preserves the torsion module $T$ and hence acts on the
free quotient $\mathcal O^r$.
If this semilinear map has matrix $H$, then commutation gives
$G=H\overline G H^{-1}$.  Thus
\[
 P(z)=\det(zI_r-G)\in\mathbb Z_3[z].                   \tag{KR.4}
\]
The finite image of the pro-three group generated by $\gamma$
on the free quotient modulo $\lambda$ is a three-group.
Every eigenvalue of this image is one in characteristic three:
if $M^{3^s}=I$, then $(M-I)^{3^s}=0$.
Consequently
\[
                    P(z)\equiv(z-1)^r\pmod3,          \tag{KR.5}
\]
since reduction of coefficients in $\mathbb Z_3$ modulo
$\lambda$ is the same as reduction modulo three.

The certified base logarithmic group in AC-K13 has order
$27$.  Applying the snake-lemma proof of (KI.7) also at
$e=0$ is legitimate: $X'/(\gamma-1)X'$ is this finite
logarithmic group, so the induced endomorphism $G-I_r$
on the free quotient is injective with finite cokernel.
It gives
\[
 3=t_0+2v_3(P(1)),\qquad
 t_0=\operatorname{length}_{\mathcal O}T/(\gamma-1)T.
                                                               \tag{KR.6}
\]
For $r>0$, (KR.5) implies $v_3(P(1))\geq1$.
The integer $t_0$ is nonnegative, so (KR.6) forces
\[
                     v_3(P(1))=1,\qquad t_0=1.        \tag{KR.7}
\]
In particular $T$ is then nonzero.

\proofstep{2. The finite torsion is killed by $\lambda$.}
Assume $r>0$ and write $B=\gamma-1$ on $T$.
Since $T$ is finite and the cyclotomic group is pro-three,
$B$ is nilpotent modulo three; because a power of three
annihilates $T$, some power of $B$ annihilates $T$ as well.
The finite module $T$ is thus a module over the local ring
$\mathbb Z_3[[B]]$ with maximal ideal $(3,B)$.
By (KR.7), $T/BT$ has order three, and
$T/(3,B)T=\mathbb F_3$.  Nakayama's lemma shows that $T$
is cyclic over this ring.

Conjugation $c$ commutes with $B$ and with $\mathbb Z_3$.
It acts by one sign $\epsilon\in\{1,-1\}$ on $T/(3,B)T$.
If $v$ is a generator, then
$v_\epsilon=(v+\epsilon c(v))/2$ has the same nonzero image
in this quotient.  It is again a generator and satisfies
$c(v_\epsilon)=\epsilon v_\epsilon$.
Since $c$ commutes with every coefficient and with $B$,
it acts by that scalar on all of $T$.
On the other hand $c\sigma c^{-1}=\sigma^{-1}$.
A scalar $c$ commutes with $\sigma$, so
$\sigma=\sigma^{-1}$ on $T$; together with $\sigma^3=1$
this gives $\sigma=1$.  Hence $\lambda T=0$.
Cyclicity and nilpotence now give
\[
                 T\simeq\mathbb F_3[B]/(B^k)
                         \quad\text{for some }k\geq1. \tag{KR.8}
\]
Since $4^e=1$ on this characteristic-three module, its
coinvariant length is exactly one for every twist $e$.
Moreover the $\mathcal O$-module sequence
$0\to T\to X'\to\mathcal O^r\to0$ splits as modules,
so the minimal number of generators of $X'$ is $r+k$.
AC-K16 therefore gives
\[
                              r+k\in\{6,7\}.           \tag{KR.9}
\]

\proofstep{3. All free ranks at least two.}
Suppose $r\geq2$.  Equation (KR.5) gives $P'(1)\equiv0\pmod3$.
For every $z\in1+3\mathbb Z_3$, polynomial Taylor expansion
with integral coefficients gives
\[
                         P(z)\equiv P(1)\pmod9.
\]
Indeed its linear term is a product of two multiples of three,
and its terms of degree at least two contain $(z-1)^2$.
By (KR.7), $v_3(P(z))=1$ throughout that set.
Taking $z=4^{-e}$, the determinant in (KI.6) differs from
$P(4^{-e})$ only by a unit.  Its $\lambda$-valuation is two;
the torsion contribution is one by (KR.8).  Thus $a_e=3$
and $m_e=4$, proving (KR.1).

\proofstep{4. The single remaining free rank.}
If $r=1$, equation (KR.4) says that $G$ is the scalar
$g\in\mathbb Z_3^\times$.  Equation (KR.5) gives
$g\in1+3\mathbb Z_3$, and (KR.7) gives $v_3(g-1)=1$.
Equations (KI.6) and (KR.8) then yield (KR.2).
The determinant cannot vanish at any positive integer $e$,
by the already proved finiteness of $W_e$.
Equation (KR.9) says that $k=5$ or $6$.

The cyclotomic extension above $E_1$ has generator $\gamma^3$
and the same split unramified module $X'$.
Thus its logarithmic class group is
$X'/(\gamma^3-1)X'$, by (KL.9) applied over $E_1$.
On the finite module (KR.8),
$(1+B)^3-1=B^3$, so the cokernel has length
$\min(k,3)=3$.  On the free rank-one quotient, lifting the
exponent gives $v_3(g^3-1)=2$, hence cokernel length four
over $\mathcal O$.  The same exact snake-lemma sequence
as (KI.7) therefore proves total length seven.
This proves (KR.3), including finiteness in this case.

Finally, if $r=0$, $X'$ is itself finite.  Every $W_e$ is
a quotient of that fixed module, so
$m_e\leq1+\operatorname{length}_{\mathcal O}X'$ uniformly.
The final implication follows by eliminating $r=1$ with
(KR.3), then applying either this bound or (KR.1).
\end{proof}

\paragraph{What this does and does not settle.}
The only free-rank case that can still carry unbounded
three-primary lengths is the single scalar specialization
in (KR.2).  Writing $g=4^{-\beta}$ with $\beta\in\mathbb Z_3$
would rewrite its valuation in terms of $e-\beta$; finiteness
at every positive integer does not bound that approximation.
The finite logarithmic test (KR.3) is therefore a strictly
smaller arithmetic gate than an unspecified infinite list
of exponents.  No uncertified class-group output is a proof
of its failure.  Even a uniform fixed-field three-primary
bound would still leave the other coefficient primes and
the source's growing CM hosts in H04.4.

In geometric language, conjugation acts on the compact
abelianized covering module from SGA~1's Galois category;
it is semilinear over the integral cubic character order.
The determinant (KR.4) is the characteristic polynomial
of that arithmetic module, not a motive's archimedean
zero operator.  The argument uses the exact K-theoretic
and logarithmic-class comparisons already cited in
AC-K10 and AC-K13--K14, and asserts no new scan-level
attribution in Grothendieck's archive.

\section{A finite first-layer certificate and a uniform fixed cubic bound}
\label{sec:cubic-layer-certificate}

Let $E_1=\mathbb Q(\mu_9,\sqrt[3]{269})$.  The exhaustive
class-group certification attempted for this degree-eighteen field
was stopped: its unconditional prime-search bound was
$9173231824586$.  Its preliminary class-group and logarithmic-class
outputs are not used as certified answers.  Instead, the following
finite certificate uses proposed algebraic objects from that
calculation and checks their required properties directly.
The verification scripts and complete outputs are retained in
\texttt{evidence/cubic\_logarithmic/}.

\begin{theorem}[AC-K15: a certified first-layer upper bound]
\label{AC-K15}
The actual degree-zero logarithmic class group satisfies
\[
       \#\widetilde{\operatorname{Cl}}_3(E_1)\leq3^6.    \tag{KC.1}
\]
Together with AC-K13 and AC-K16--AC-K17, this proves that the
split unramified Iwasawa module $X'$ of the fixed field $E$ is
finite.  Consequently there is a fixed finite integer
\[
 C_E=1+\operatorname{length}_{\mathcal O}X'
\]
such that, for every $e\geq1$ and every integer $t\geq2$,
\[
 3\leq m_e\leq C_E,\qquad
 \sum_{1\leq e<t}m_e\leq C_E(t-1).                   \tag{KC.2}
\]
The sharper equality $m_e=4$ for $3\mid e$ remains valid.
All assertions are unconditional.  The proof certifies the
upper bound (KC.1), rather than asserting exact first-layer
class invariants from an uncertified class-group routine.
No numerical value for $C_E$ is claimed.
\end{theorem}
\begin{proof}
\proofstep{1. Four explicit unramified cubic characters.}
The field used in the computation is the exact compositum of
$Y^3-269$ and $Y^6+Y^3+1$.  The compositum output retains its
embeddings; its chosen primitive element is
$\zeta_9-\sqrt[3]{269}$.  The degree-eighteen polynomial and
embeddings are in \texttt{certify\_layer1.out}, and the polynomial
is repeated in \path{layer1_unramified_certificate.out}.
The test \texttt{nfcertify(b)} returned the empty vector,
certifying the maximal order used in subsequent ideal operations.
This is a maximal-order check, not a class-group certification.

The file \texttt{layer1\_classfields.bin} retains a number-field
structure, a provisional ray-class structure, and four explicit
polynomials
\[
                         X^3-a_i\quad(1\leq i\leq4).
\]
Their algebraic coefficients $a_i\in E_1$ are printed completely
in \path{layer1_unramified_certificate.out}.
The routine proposing these polynomials is not used as a proof
of their class-field interpretation.  Each polynomial was
independently factored over the actual number field: the exact
factorization returned one factor of degree three.
For each, \texttt{rnfdisc(b,polynomial)[1]} returned the unit
ideal.  Thus each defines a degree-three extension unramified
at every finite prime.  The base field is totally imaginary,
so there are no additional real-place conditions.
The polynomials are pure cubics and $\mu_3\subset E_1$;
therefore each extension is cyclic.  This last argument is
specific to $X^3-a_i$, not to an arbitrary irreducible cubic.
They consequently define four genuine unramified Artin
characters on the ideal class group, with values in
$\mu_3\simeq\mathbb F_3$.

\proofstep{2. A certified upper bound for the ordinary three-class group.}
Let $g_1,g_2,g_3,g_4$ be the four explicit ideal matrices retained
in the field data, and let $\mathfrak q$ be the first prime
above seven in the retained prime decomposition.  Its norm is
$343$.  Cubic residue symbols of the $a_i$, evaluated at the
prime factors of $g_1,g_2,g_3,\mathfrak q$, give the matrix
\[
 \begin{pmatrix}
 0&2&0&0\\
 0&0&2&0\\
 2&0&0&0\\
 0&0&0&2
 \end{pmatrix},\qquad \det=1\quad\hbox{in }\mathbb F_3.
                                                               \tag{KC.3}
\]
All these test primes avoid the radicals' zero divisors.
Explicitly, at a prime $\mathfrak l$ of norm $N\mathfrak l$,
the script compares
$a_i^{(N\mathfrak l-1)/3}\bmod\mathfrak l$
with $1,\zeta_3,\zeta_3^2$; for an ideal product it sums
these exponents with the ideal-factorization multiplicities.
This is the Frobenius action on a cube root of $a_i$.
The root $\zeta_3$ is the sixth power of the displayed
order-eighteen root of unity.  Thus the computation is a
finite-field calculation of the actual Artin characters.

The same script supplies algebraic elements, in compact
factored form, and verifies by exact ideal arithmetic
\[
 \mathfrak q=(b)g_1^{21}g_4^2,
 \qquad
 g_1^{63}=(b_1),\quad g_2^{21}=(b_2),\quad
             g_3^3=(b_3),\quad g_4^3=(b_4).            \tag{KC.4}
\]
Every equality in (KC.4) returned verification value one.
Their complete elements and ideal matrices are in
\path{layer1_generation_certificate.out}; the reproducible
script has the same basename and suffix \texttt{.gp}.
These are direct ideal equalities, not assertions that the
provisional class-group generators form a basis.

Put $A_1=\operatorname{Cl}(E_1)_{(3)}$.
AC-K16 gives $\dim_{\mathbb F_3}A_1/\lambda A_1=2$.
Since $3$ is a unit times $\lambda^2$, the filtration by
$\lambda$ shows $\dim_{\mathbb F_3}A_1/3A_1\leq4$.
The rank-four pairing (KC.3) forces the classes of the four
tested ideals to span this quotient.  By the first relation
in (KC.4), they belong to the subgroup generated by the
$g_i$.  Nakayama's lemma for the finite three-group therefore
shows that the three-primary projections of the $g_i$
generate all of $A_1$.  The remaining relations in (KC.4)
now give the unconditional bound
\[
                       \#A_1\leq3^{2+1+1+1}=3^5.     \tag{KC.5}
\]
No claim that its order is exactly $3^5$ is needed.

\proofstep{3. Removing the classes at three saves a factor of three.}
Let $S$ be the three primes of $E_1$ over three.
For the radicals $a_1,a_2,a_4$, the exact test
\texttt{nfislocalpower(b,pr,a,3)} returned zero at every
$\mathfrak p\in S$; the third radical returned one.
The complete $4\times3$ result is
\[
 \begin{pmatrix}0&0&0\\0&0&0\\1&1&1\\0&0&0\end{pmatrix}.
                                                               \tag{KC.6}
\]
This tests local cubeness, independently of any class-group
or unit-group completeness assumption.  The implementation
in PARI/GP~2.19.1, \texttt{src/basemath/bnflog.c}, function
\texttt{nfislocalpower\_i}, strips the local valuation and
uses the finite principal-unit group.  Here the local
ramification index is six and its test level is
$6\cdot3/(3-1)+1=10$.  The public command semantics are
recorded in the same release's manual under
\texttt{nfislocalpower}.  Thus zero in (KC.6) is a certified
local noncube test, not a floating-point approximation.

Choose the first extension and any prime in $S$.
It is unramified by Step~1 and nonsplit by (KC.6), so its
Frobenius has order three.  The corresponding prime class
has a nontrivial image in an unramified character, and its
three-primary projection has order at least three.
It follows from (KC.5) that the actual localized class group
satisfies
\[
 \#\operatorname{Cl}(\mathcal O_{E_1}[1/3])_{(3)}
                         \leq3^4.                    \tag{KC.7}
\]

\proofstep{4. Two actual units bound the logarithmic kernel.}
Each completion of $E_1$ at three is $\mathbb Q_3(\mu_9)$.
Its logarithmic ramification and residue indices are
$(\widetilde e,\widetilde f)=(2,3)$, also independently
returned for all three primes by \texttt{bnflogef}.
Choose logarithmic degree $d=3\log_3(4)$ at each of these
three places.  This is the normalization allowed by
Belabas--Jaulent, Corollary~3.5 and Remark~3.6, author's
pp.~8--9.  For a genuine $S$-unit $u$, put
\[
  \ell_i(u)=-\frac{\log_3 N_{(E_1)_{\mathfrak p_i}/\mathbb Q_3}(u)}
                         {3\log_3(4)}\in\mathbb Z_3.
\]
The product formula gives $\sum_i\ell_i(u)=0$ exactly.

The script \path{layer1_logkernel_certificate.gp} retains
eleven explicit proposed $S$-units.  It verifies, by exact
ideal factorization, that every prime in every one's divisor
lies above three.  This only proves they are actual $S$-units;
it does not assume they are a full basis.  For units numbered
five and ten, the first two rows of their logarithmic
valuation matrix are
\[
       \begin{pmatrix}39&51\\21&30\end{pmatrix}
                  \pmod{81}.                         \tag{KC.8}
\]
Its determinant is $18\pmod{81}$, and hence has three-adic
valuation exactly two.  The script computes the local norms
as resultants with the three degree-six factors of the field
polynomial over $\mathbb Q_3$.  It starts with factor precision
sixty; the reported matrix precision is at least fifty,
and the determinant precision is fifty-one.  These are
absolute three-adic precision bounds.  In particular the
four digits required for the displayed reduction modulo
$81$ are available with a large margin.  Independent
replays at precisions forty and seventy give determinant
valuation two again.  Permuting the local factors can
change the displayed minor by a sign; it does not change
its valuation or the lattice index.

The lattice of degree-zero divisors supported in $S$ is
\[
 D_S^0=\{(z_1,z_2,z_3)\in\mathbb Z_3^3:
                                z_1+z_2+z_3=0\}.
\]
Its first two coordinates identify it with $\mathbb Z_3^2$.
By (KC.8), the logarithmic divisors of just these two actual
$S$-units span a sublattice of index nine in $D_S^0$.
Quotienting by all principal logarithmic divisors can only
reduce this order.  Therefore the degree-zero subgroup of
logarithmic classes generated at three has order at most
nine.

\proofstep{5. The logarithmic-class upper bound.}
The logarithmic divisor-class presentation has the exact
sequence
\[
 0\longrightarrow\widetilde{\operatorname{Cl}}_3(E_1)(3)
 \longrightarrow\widetilde{\operatorname{Cl}}_3^{\rm all}(E_1)
 \longrightarrow\operatorname{Cl}(\mathcal O_{E_1}[1/3])_{(3)}
 \longrightarrow0,                                    \tag{KC.9}
\]
where the middle group here includes the degree direction
and the left group is generated by the primes over three.
The source is Belabas--Jaulent, beginning of \S3, author's
p.~6, with the $S$-unit relation presentation in \S3.3,
p.~10.  Upon restriction to degree zero, surjectivity at
right need not persist; the image is a subgroup of the
last group.  Its order is at most $3^4$ by (KC.7), while
the kernel has order at most $3^2$ by Step~4.
Thus the degree-zero middle group is finite and has order
at most $3^6$, proving (KC.1).  This argument explains why
neither a saturated unit basis nor a full certified class
presentation is needed: all extra true relations only
make the group smaller.

\proofstep{6. All positive free ranks are excluded.}
Use $T=X'_{\rm tors}$, $X'/T\simeq\mathcal O^r$ and $G$
from AC-K16, and put $P(z)=\det(zI-G)$.
Suppose $r>0$.  AC-K17 proves
\[
 P\in\mathbb Z_3[z],\quad
 P(z)\equiv(z-1)^r\pmod3,\quad v_\lambda P(1)=2,
 \quad T\simeq\mathbb F_3[B]/B^k,
                                                               \tag{KC.10}
\]
where $B=\gamma-1$, $k\geq1$, and $r+k\in\{6,7\}$.
For clarity, the torsion assertion follows by applying
Nakayama to $T/(3,B)T\simeq\mathbb F_3$, choosing a
complex-conjugation eigenvector as its generator, and using
$c\sigma c=\sigma^{-1}$.  Conjugation is scalar on the
resulting cyclic $\mathbb Z_3[[B]]$-module, which forces
$\sigma=1$ and hence $\lambda T=0$.

The weight-zero descent over $E_1$ is
\[
                X'/(\gamma^3-1)X'
                       =\widetilde{\operatorname{Cl}}_3(E_1).
\]
It is finite by (KC.1).  Thus $G^3-I$ has nonzero determinant,
and the snake-lemma length formula used in (KI.7) gives
\[
 \operatorname{length}_{\mathcal O}
                  \widetilde{\operatorname{Cl}}_3(E_1)
   =\min(k,3)+v_\lambda\det(G^3-I).                  \tag{KC.11}
\]
Indeed $\gamma^3-1=B^3$ on $T$, since $T$ has characteristic
three.  If $r=1$, (KC.10) gives $k\geq5$, while
$G=(g)$ with $g\in1+3\mathbb Z_3$ and $v_3(g-1)=1$.
Lifting the exponent gives $v_\lambda(g^3-1)=4$.
The right side of (KC.11) is therefore $3+4=7$.

If $r\geq2$, write $P(z)=(z-1)^r+3Q(z)$ with
$Q\in\mathbb Z_3[z]$.  Both $P(\omega)$ and
$P(\omega^2)$ belong to $3\mathcal O$, because
$\omega-1$ and $\omega^2-1$ have $\lambda$-valuation one
and $r\geq2$.  The factorization
\[
 \det(G^3-I)=\pm P(1)P(\omega)P(\omega^2)
\]
therefore has $\lambda$-valuation at least six.
The torsion contribution $\min(k,3)$ is at least one,
so (KC.11) again has length at least seven.
Both cases contradict (KC.1).  Consequently $r=0$.
Finite generation in AC-K16 now makes $X'$ finite.

Finally $W_e(E)$ is a quotient of this fixed finite module
by the twisted endomorphism $4^e\gamma-1$.
Its $\mathcal O$-length is at most that of $X'$.
Equations (KL.1)--(KL.2) prove (KC.2), and AC-K14 supplies
the equality on multiples of three.
\end{proof}

\paragraph{Scope of the fixed-field result.}
The coefficient-prime-three contribution for this one cubic
character now has a uniform linear budget in the number of
weights.  This is stronger than mere weight-by-weight
finiteness.  It does not yet estimate other coefficient
primes, the varying CM hosts in \texttt{SupplyChi}, or the
full sum required in H04.4.  The constant $C_E$ and the two
remaining classes of individual weights modulo three have
not been explicitly evaluated.

\section{An explicit exponent bound and finite weight periodicity}
\label{sec:cubic-uniform-explicit}

The finite-module conclusion of AC-K15 can be made numerical.
The only arithmetic inputs in the following argument are the
already certified base and first-layer logarithmic groups and
the generator count in AC-K16.  The remaining argument is
finite-module algebra.

\begin{theorem}[AC-K18: exponent five and nine-periodic cubic lengths]
\label{AC-K18}
For the fixed field $E=\mathbb Q(\omega,\sqrt[3]{269})$,
\[
 \lambda^5X'=0,\qquad
 \operatorname{length}_{\mathcal O}X'\leq11.           \tag{KE.1}
\]
Consequently the exponents of the individual cubic character
satisfy
\[
        3\leq m_e\leq6,\qquad m_{e+9}=m_e
                  \quad(e\geq1),                    \tag{KE.2}
\]
with $m_e=4$ for $3\mid e$.
In fact either $m_e\leq4$ for every weight, or there is one
nonzero residue class $r\pmod3$ such that
\[
 \begin{cases}
 m_e\in\{5,6\},&e\equiv r\pmod3,\\
 m_e\in\{3,4\},&e\not\equiv r\pmod3.
 \end{cases}                                         \tag{KE.3}
\]
The class $r$ and the remaining individual values are not
computed here.  For $N=t-1\geq1$, an explicit uniform budget is
\[
 3N+\lfloor N/3\rfloor
 \ \leq\sum_{e=1}^{N}m_e
 \ \leq 4N+2\lceil N/3\rceil.                       \tag{KE.4}
\]
Multiplication by $\log3$ gives the corresponding fixed
coefficient-prime-three contribution.  These statements use
no RH hypothesis.
\end{theorem}
\begin{proof}
\proofstep{1. The two conjugation components.}
Set $M=X'$, which is finite by AC-K15, and put
\[
 \delta=\omega-\omega^2,\qquad \delta^2=-3,
 \qquad B=\gamma-1,\qquad R=\mathbb Z_3[[B]].
\]
The elements $\delta$ and $\lambda$ differ by a unit.
Complex conjugation $c$ commutes with $B$ and sends
$\delta$ to $-\delta$.  The action of $B$ is nilpotent
on the finite three-group $M$: the image of the cyclotomic
pro-three group is finite, so some $\gamma^{3^h}=1$;
reduction modulo three makes $B^{3^h}=0$, and iteration
using an exponent of $M$ gives nilpotence on $M$ itself.

The quotient $M/BM$ is the degree-zero logarithmic class
group of $E$.  It has underlying group
$\mathbb Z/9\oplus\mathbb Z/3$ by AC-K13.
As a finite $\mathcal O$-module it is therefore
$\mathcal O/\delta^3$: its order is $27$, and the presence
of an element of order nine excludes every other partition
of the $\mathcal O$-length three.
Thus $M/(\delta,B)M$ is one-dimensional over $\mathbb F_3$.
Choose a lift of its generator which is a $c$-eigenvector,
by averaging with $c$ and dividing by two.  Nakayama gives
an $\mathcal O[[B]]$-generator $x$ of $M$.
Label its conjugation eigenspace by a plus sign, reversing
the two sign labels if necessary.  Then
\[
 M_+=Rx,\qquad M_-=R\delta x,\qquad M=M_+\oplus M_-.
                                                               \tag{KE.5}
\]
Indeed $\mathcal O=\mathbb Z_3\oplus\delta\mathbb Z_3$,
and averaging any coefficient with its conjugate gives
these two summands.  Multiplication by $\delta$ maps
$M_+$ onto $M_-$.
Since $(\delta^3)=(9,3\delta)$ as a $\mathbb Z_3$-lattice,
\[
             M_+/BM_+\simeq\mathbb Z/9,
       \qquad M_-/BM_-\simeq\mathbb Z/3.              \tag{KE.6}
\]
Moreover
\[
 M/\delta M\simeq M_+/3M_+,
 \qquad d:=\dim_{\mathbb F_3}(M_+/3M_+)\in\{6,7\},  \tag{KE.7}
\]
because $\delta M_+=M_-$ and
$\delta M_-=-3M_+$.  The last assertion is AC-K16.

\proofstep{2. The smaller conjugation component has order at most nine.}
Write
\[
 F(B)=(1+B)^3-1=B^3+3B^2+3B,
 \qquad C_\pm=M_\pm/F(B)M_\pm.
\]
By the first-layer descent and AC-K15,
\[
       \operatorname{length}_{\mathbb Z_3}C_+
       +\operatorname{length}_{\mathbb Z_3}C_-
                                    \leq6.           \tag{KE.8}
\]
The cyclic $R$-module $M_+/3M_+$ is
$\mathbb F_3[B]/B^d$.  As $d\geq6$ and $F(B)=B^3$
modulo three, $C_+/3C_+$ has dimension three.
Also $C_+/BC_+=M_+/BM_+\simeq\mathbb Z/9$.
Thus $C_+$ is not elementary abelian, and its length is
at least four.  Hence $C_-$ has length at most two.

Put $U=M_-$.  Its nonzero reduction modulo three is a
cyclic $\mathbb F_3[B]$-module, say
$U/3U=\mathbb F_3[B]/B^h$.
The dimension of $C_-/3C_-$ is $\min(h,3)$, so $h$ is
one or two.
If $h=2$, the length bound forces $C_-$ to be elementary
abelian of dimension two.  The surjection
$U/3U\to C_-$ is then an isomorphism, so
$F(B)U=3U$.  Since $B^2U\subset3U$,
\[
 F(B)U=B^3U+3B^2U+3BU\subset3BU.
\]
Consequently $3U=B(3U)$.  Nilpotence of $B$ forces
$3U=0$, and $U$ has order nine.

It remains to treat $h=1$.  Nakayama over $\mathbb Z_3$
makes $U$ cyclic as an abelian group: $U\simeq\mathbb Z/3^a$.
Suppose for contradiction that $a\geq3$.
The action of $B$ is multiplication by a scalar
$b\in\mathbb Z_3$.  Equation (KE.6) forces
$v_3(b)=1$.  The vector $\delta x$ generates $U$ even
as a $\mathbb Z_3$-module, and its $R$-annihilator is
$(3^a,B-b)$.
Let $H$ be the kernel of
$\delta:M_+\to U$.  It is killed by three, and as an
$\mathbb F_3[B]$-module it is generated by
\[
                    u=(B-b)x,
             \qquad z=3^a x.                          \tag{KE.9}
\]
Indeed these are the images of the generators of that
annihilator.  Also $Bz=0$:
$Bz=3^a(bx+u)=bz+3^au=0$, since $3H=0$ and $3\mid b$.
The image of $u$ in
$M_+/3M_+=\mathbb F_3[B]/B^d$ is $Bx$.
Therefore the cyclic block
$J=\mathbb F_3[B]u$ has length at least $d-1\geq5$.
Either $z\in J$, in which case it lies in the $B$-socle
of this block and hence in $B^3J$, or
$H=J\oplus\mathbb F_3z$.

Apply $F(B)$ to the exact sequence
$0\to H\to M_+\to U\to0$.
On $H$, this endomorphism is $B^3$.
Since $v_3(b)=1$, lifting the exponent gives
$F(b)=9v$ with $v\in\mathbb Z_3^\times$.
Thus $U/F(b)U$ has length two and $U[F(b)]$ is generated
by $3^{a-2}\delta x$.
The connecting map
\[
             U[F(b)]\longrightarrow H/B^3H
\]
sends this generator to the class of $vz$.
To verify the assertion, polynomial division gives
\[
 F(B)x=F(b)x+
 [B^2+Bb+b^2+3B+3b+3]u
       =F(b)x+B^2u.
\]
All terms removed in the second equality kill $u$, since
$3H=0$ and $3\mid b$.
Multiplication by $3^{a-2}$ kills the last term because
$a\geq3$, leaving $3^a vx=vz$, as claimed.

If $z\in J$, the quotient $H/B^3H$ has dimension three
and the connecting image is zero.
If $z\notin J$, that quotient has dimension four and
the connecting image has dimension one.
In either case the exact cokernel sequence gives
\[
              \operatorname{length}_{\mathbb Z_3}C_+
                         =3+2=5.
\]
But $C_-$ has length two in the present case, contradicting
(KE.8).  Hence $a\leq2$.  We have proved in both cases
\[
                          \#M_-\leq9.                \tag{KE.10}
\]

\proofstep{3. Exponent and weight consequences.}
Equation (KE.10) gives $9M_-=0$.
Since $\delta M=M_-+\delta M_-$, it follows that
$9\delta M=0$.  The identity $\delta^4=9$ now gives
$\delta^5M=0$, proving the exponent assertion in (KE.1).
Also $3M_+=-\delta M_-$ has length at most two.
By (KE.7), $M_+$ therefore has length at most $d+2$;
adding (KE.10) gives
$\operatorname{length}_{\mathcal O}M\leq d+4\leq11$.
Here the sum of the two $\mathbb Z_3$-lengths equals the
$\mathcal O$-length, since the residue field has order three.

Every $W_e(E)$ is a quotient of $M$, and AC-K12 makes it
cyclic over $\mathcal O$, of length $a_e=m_e-1$.
Thus $a_e\leq5$.  AC-K13 supplies $a_e\geq2$,
proving the bounds in (KE.2).
For $e'\equiv e\pmod9$, AC-K14 with $n=3$ identifies
$W_e/27$ and $W_{e'}/27$.  Both quotients equal the whole
group because $\lambda^5$ annihilates them and
$27$ is a unit times $\lambda^6$.  Their lengths are equal,
which proves nine-periodicity.

Finally consider the sharper alternatives.  The image of
$x$ generates $W_e$ over $\mathcal O$: after imposing
$4^e(1+B)-1=0$, the $R$-action is scalar.
Its conjugation sign is the same as that chosen in (KE.5).
Thus the minus component of the cyclic quotient
$\mathcal O/\delta^{a_e}$ has length
$\lfloor a_e/2\rfloor$.  It is also exactly
\[
                M_-/(4^e\gamma-1)M_- .              \tag{KE.11}
\]
If $M_-$ is elementary abelian, $4^e$ acts as one on it,
and (KE.6) makes (KE.11) have length one for every $e$.
Hence $a_e\in\{2,3\}$ for all $e$.
Otherwise (KE.10) forces $M_-=\mathbb Z/9$.
In this case $B$ is a scalar $b=3c\pmod9$ with
$c\not\equiv0\pmod3$, again by (KE.6).
Since $4^e\equiv1+3e\pmod9$,
\[
             4^e(1+b)-1\equiv3(e+c)\pmod9.
\]
The quotient (KE.11) has length two on the single nonzero
class $e\equiv-c\pmod3$, and length one elsewhere.
This gives (KE.3).
There are at most $\lceil N/3\rceil$ exceptional weights
among $1,\ldots,N$, proving the upper bound in (KE.4).
The lower bound uses $m_e\geq3$ and the known equality
$m_e=4$ on the $\lfloor N/3\rfloor$ multiples of three.
\end{proof}

\paragraph{What is still numerical.}
The fixed coefficient-prime-three budget is now explicit.
Determining which alternative in (KE.3) occurs requires
additional finite arithmetic data.  In either case at most
nine positive weights determine the entire sequence.
The other coefficient primes and the growing-host comparison
in H04.4 remain outside this fixed-field conclusion.

\section{The exact finite cyclotomic action}
\label{sec:cubic-finite-action}

The explicit bound in AC-K18 can be strengthened to a finite
descent statement.  This identifies the exact arithmetic data
that would determine all the remaining weights, without
performing a class-group calculation in a degree-54 field.

\begin{theorem}[AC-K20: action of order nine and a finite matrix recipe]
\label{AC-K20}
Keep $M=X'$, $\delta=\omega-\omega^2$ and
$\varepsilon_3(\gamma)=4$ as in AC-K18.  Then
\[
                    \gamma^9=1\quad\hbox{on }M,
       \qquad \gamma^3\ne1\quad\hbox{on }M.          \tag{KF.1}
\]
For
\[
 E_n=\mathbb Q(\mu_{3^{n+1}},\sqrt[3]{269}),
 \qquad n\geq0,
\]
the logarithmic descent isomorphisms consequently give
\[
 M\simeq\widetilde{\operatorname{Cl}}_3(E_n)
                      \quad(n\geq2),                 \tag{KF.2}
\]
compatibly with the cubic action and the remaining cyclotomic
action.  Thus $E_2=\mathbb Q(\mu_{27},\sqrt[3]{269})$, of
degree $54$ over $\mathbb Q$, suffices.  It is the first layer
for which the quotient map
$M\to M/(\gamma^{3^n}-1)M$ is an isomorphism.

More explicitly, suppose that a certified presentation of its
degree-zero logarithmic class group and its cyclotomic action
is given by
\[
 C=\mathbb Z^s/D\mathbb Z^s,
 \qquad D=\operatorname{diag}(3^{r_1},\ldots,3^{r_s}),
 \qquad G:C\longrightarrow C.                         \tag{KF.3}
\]
Here $G$ must represent the field automorphism fixing
$\sqrt[3]{269}$ and sending $\zeta_{27}$ to
$\zeta_{27}^4$, with its induced action on logarithmic ideal
classes.  For $j=1,\ldots,9$, put
\[
 C_j=\operatorname{coker}_{\mathbb Z}
       \begin{bmatrix}D&4^jG-I_s\end{bmatrix}.
                                                               \tag{KF.4}
\]
Then the exact exponent sequence is recovered by
\[
       m_e=1+\log_3\#C_j
       \quad\text{whenever }e\equiv j\pmod9.         \tag{KF.5}
\]
Thus nine finite Smith normal forms of one certified module
with its action determine every weight.  The invariant factors
of $C$ without $G$ do not constitute the data in this recipe.
No such degree-54 arithmetic certificate is asserted here.
\end{theorem}
\begin{proof}
\proofstep{1. A bound on the kernel of multiplication by $\delta$.}
Use the generator $x$ and the conjugation components from
AC-K18:
\[
 M_+=\mathbb Z_3[[B]]x,\qquad
 U=M_-=\mathbb Z_3[[B]]\delta x,
 \qquad B=\gamma-1.
\]
Multiplication by $\delta$ maps $M_+$ onto $U$.
Let $H$ be its kernel.  Since $\delta^2=-3$, $3H=0$.
The endomorphism $B$ is nilpotent, as proved in AC-K18.
Write $d=\dim_{\mathbb F_3}(M_+/3M_+)\in\{6,7\}$.
Because $3M_+=-\delta U$, the length of $3M_+$ is at
most the length of $U$.  The exact sequence defining $H$
therefore gives
\[
 \begin{aligned}
 \dim_{\mathbb F_3}H
 &=\operatorname{length}_{\mathbb Z_3}M_+
        -\operatorname{length}_{\mathbb Z_3}U
 \\
 &=d+\operatorname{length}_{\mathbb Z_3}(3M_+)
        -\operatorname{length}_{\mathbb Z_3}U
 \leq d\leq7.
 \end{aligned}
                                                               \tag{KF.6}
\]
In particular $B^7H=0$.

By AC-K18, the nonzero group $U$ has order at most nine
and $U/BU\simeq\mathbb Z/3$.
There are precisely the following possibilities for its
cyclic $\mathbb Z_3[[B]]$-module structure:
\[
 \begin{array}{ll}
 U=\mathbb Z/3^a,&a=1\text{ or }2,\quad
       B\text{ acts as }b\in3\mathbb Z_3;\\
 U=\mathbb F_3[B]/B^2.&
 \end{array}                                                   \tag{KF.7}
\]
In the first line, take $b=0$ if $a=1$; if $a=2$,
the quotient $U/BU$ forces $v_3(b)=1$.
In the elementary group of order nine, cyclicity and the same
quotient force the single nilpotent block in the second line.
Every case in (KF.7) is killed by
\[
                 F=(1+B)^3-1=B^3+3B^2+3B.           \tag{KF.8}
\]
Indeed this is immediate in characteristic three; for
$\mathbb Z/9$ it follows from $3\mid b$.
Consequently $FM_+\subset H$.
Since $3H=0$ and $F$ acts on $H$ as $B^3$,
\[
       (\gamma^9-1)M_+=F^3M_+=B^6FM_+.             \tag{KF.9}
\]
The first equality follows by expanding
$(1+F)^3-1=F^3+3F^2+3F$; the last two terms vanish on
$M_+$ because $FM_+\subset H$.

\proofstep{2. The two possible small components.}
Suppose first that $U=\mathbb Z/3^a$ as in the first
line of (KF.7).  The vector $\delta x$ is an abelian
generator of $U$ and has annihilator $(3^a,B-b)$ over
$\mathbb Z_3[[B]]$.  Hence $H$ is generated by
\[
                    t=(B-b)x,\qquad z=3^a x.
\]
Both lie in $H$.  Moreover
$Bz=bz+3^at=0$, because $3H=0$ and $3\mid b$.
Polynomial division, followed by $3t=0$ and $3\mid b$,
gives the exact identity
\[
             Fx=F(b)x+B^2t.                          \tag{KF.10}
\]
For $a=1$ our choice $b=0$ makes the first term zero.
For $a=2$, $F(b)$ is divisible by nine, so that the first
term is a multiple of $z$.
It follows that $B^6Fx=B^8t=0$, using (KF.6).
Since $x$ generates $M_+$ and all endomorphisms commute,
(KF.9) proves $\gamma^9=1$ on $M_+$ in this case.

In the remaining case $U=\mathbb F_3[B]/B^2$,
the annihilator of $\delta x$ is $(3,B^2)$.
Thus $H$ is generated by
\[
                         t=B^2x,\qquad z=3x.
\]
The relation $3H=0$ gives $B^2z=3t=0$.
From (KF.8) one obtains
\[
                         Fx=Bt+Bz.
\]
Therefore $B^6Fx=B^7t+B^7z=0$, again by (KF.6).
Equation (KF.9) proves the assertion on $M_+$.
On $U$, already $\gamma^3=1$, so this proves
$\gamma^9=1$ on all of $M$.

For the strict assertion in (KF.1), reduce $M_+$ modulo
three.  This is the cyclic nilpotent module
$\mathbb F_3[B]/B^d$ with $d\geq6$.
Here $\gamma^3-1=B^3$ is nonzero.  Thus the action of
$\gamma$ on $M$ has order exactly nine.

\proofstep{3. Finite descent and the matrix computation.}
The same logarithmic descent theorem used in (KL.9), applied
over $E_n$, identifies
\[
 \widetilde{\operatorname{Cl}}_3(E_n)
                 =M/(\gamma^{3^n}-1)M.              \tag{KF.11}
\]
Here the cyclotomic infinite extension and its maximal split
unramified module are unchanged; only the cyclotomic subgroup
has generator $\gamma^{3^n}$.
The exact primary-source input is the Schneider description
stated in the proof of Jaulent--Michel, Theorem~3, p.~7 of
the author's version, already checked in AC-K14:
\url{https://www.math.u-bordeaux.fr/~jjaulent/Articles/J44.pdf}.
It is an equality of compact coinvariants, not a conjectural
formula for a characteristic ideal.
Equation (KF.1) proves (KF.2) and the assertion about the
first stable quotient.  The field degree is
$3\varphi(27)=54$, since the cubic radical extension is
disjoint from the cyclotomic extension over
$\mathbb Q(\omega)$, as used in AC-K14.

Under (KF.2), the positive-weight wild kernel is
\[
              W_e(E)=C/(4^eG-1)C.                   \tag{KF.12}
\]
This is (KL.9), with the generator normalized by
$\varepsilon_3(\gamma)=4$.
If $C$ is given by (KF.3), adjoining the relations
$(4^eG-I_s)v$ to the columns of $D$ is exactly the
integer presentation (KF.4).
AC-K18 gives $\lambda^5C=0$ and therefore $27C=0$.
Since $4^9\equiv1\pmod{27}$, the quotient depends only
on $e$ modulo nine.  Finally AC-K12 and AC-G26 identify
$\log_3\#W_e=a_e=m_e-1$, which proves (KF.5).
Every step of the matrix recipe is finite exact algebra once
the stated arithmetic module and its action are certified.
\end{proof}

\paragraph{The remaining arithmetic certificate.}
A computation for this purpose must certify the degree-zero
logarithmic class module at three and the action induced by
$\zeta_{27}\mapsto\zeta_{27}^4$; an ordinary class group or
an unlabelled list of logarithmic invariant factors is
insufficient for (KF.4).  The induced action can be supplied
on a presentation by logarithmic ideal classes, with the images
of its generators verified modulo the certified relations.
The cubic action is a useful compatibility check, but its
matrix is not needed merely to take the orders in (KF.5).
Reversing the cyclotomic generator without changing the Tate
twist would permute the weight classes, so its normalization
is part of the certificate.  No uncontrolled degree-54
calculation is included in this reduction.

\paragraph{Finite-cover and even K-theory interpretation.}
The cubic and cyclotomic actions on $M$ factor through
$\operatorname{Gal}(E_2/K)\simeq C_3\times C_9$.
Thus $M$, killed by $\lambda^5$, is a finite locally constant
coefficient on $\operatorname{Spec}\mathcal O_K[1/(3\cdot269)]$
trivialized by its finite etale cover from $E_2$.
This is precisely the finite-cover/action dictionary of
SGA~1, Expos\'e~V, \S\S4 and 7, reused here as a categorical
interpretation.  Equations (KF.4) and (KF.12) describe the
finite specialization obtained by imposing the twisted
cyclotomic relation $4^e\gamma=1$.
Under the already cited even K-theory and etale-cohomology
comparison in AC-K10--AC-K12, this quotient is the wild
localization kernel in degree $2e$, rather than the entire
even K-group.  AC-K18 bounds its integral length and AC-K20
reduces all its specializations to one finite module with
action.  These statements supply no polarization or analytic
zero operator.  No passage in Grothendieck's handwritten
archive, or additional scan-level locator, is claimed.

\section{Three weight classes and the remaining degree-18 datum}
\label{sec:cubic-three-weights}

The full module first stabilizes at $E_2$, as AC-K20 proves.
Its positive-weight specializations need less information.
The same certified first-layer bound used in AC-K18
forces the following sharper result.

\begin{theorem}[AC-K21: exponents four and five with period three]
\label{AC-K21}
For the fixed cubic character over $K=\mathbb Q(\omega)$,
\[
          4\leq m_e\leq5,
       \qquad m_{e+3}=m_e\quad(e\geq1).              \tag{KJ.1}
\]
There are only three possible sequences:
\[
 \begin{array}{c|ccc}
 &e\equiv1\pmod3&e\equiv2\pmod3&e\equiv0\pmod3\\\hline
 \text{first possibility}&4&4&4\\
 \text{second possibility}&5&4&4\\
 \text{third possibility}&4&5&4.
 \end{array}                                         \tag{KJ.2}
\]
Equivalently, either every $m_e=4$, or there is one
$r\in\{1,2\}$ such that $m_e=4+\mathbf1_{e\equiv r\ (3)}$.
For $N=t-1\geq1$, the exact aggregate in these alternatives is
\[
 \sum_{e=1}^Nm_e=
 \begin{cases}
 4N,&\text{in the first possibility},\\
 4N+\left\lfloor (N+3-r)/3\right\rfloor,
                  &\text{in the possibility indexed by }r.
 \end{cases}                                         \tag{KJ.3}
\]
In particular the fixed coefficient-prime-three budget is
at most $(4N+\lceil N/3\rceil)\log3$.
All the wild kernels $W_e(E)$ are killed by nine.

The choice among (KJ.2) is determined by a specified
conjugation component of
\[
 C_1=\widetilde{\operatorname{Cl}}_3
        (\mathbb Q(\mu_9,\sqrt[3]{269})),             \tag{KJ.4}
\]
with its normalized cyclotomic action.  The criterion below
is finite and uses the degree-18 field; this theorem does
not compute which alternative occurs.
\end{theorem}
\begin{proof}
\proofstep{1. Specializing the cyclic plus component.}
Keep $M=X'$, $R=\mathbb Z_3[[B]]$, $B=\gamma-1$,
and $M_+=Rx$, $U=M_-=R\delta x$ from AC-K18.
Here the plus sign is the conjugation sign of $x$;
it need not be the numerical sign $+1$.
The already proved facts are
\[
 \begin{gathered}
 \delta M_+=U,\quad \delta U=-3M_+,\quad
 M_+/BM_+\simeq\mathbb Z/9,\quad U/BU\simeq\mathbb Z/3,\\
 M_+/3M_+\simeq\mathbb F_3[B]/B^d,
       \qquad d\in\{6,7\},\qquad \#U\leq9,\qquad27M=0.
 \end{gathered}                                      \tag{KJ.5}
\]
Write $M_+=R/I$ using the generator $x$.
Reduction modulo three and evaluation at $B=0$ give
\[
       I\bmod3=(B^d),\qquad
       \{f(0):f\in I\}=9\mathbb Z_3,
       \qquad27\in I.                               \tag{KJ.6}
\]
Consequently, for $f(B)=\sum_{k\geq0}f_kB^k\in I$,
$9\mid f_0$ and $3\mid f_k$ for $k<d$.
If $b\in3\mathbb Z_3$, the convergent evaluation therefore
satisfies
\[
                         9\mid f(b).                \tag{KJ.7}
\]
The constant and linear terms are divisible by nine;
the terms of degree at least two are divisible by nine
already from the powers of $b$.
Formal division by $B-b$, valid for $b\in3\mathbb Z_3$,
identifies the specialized plus component with
\[
 M_+/(B-b)M_+
       \simeq\mathbb Z_3/\operatorname{ev}_b(I).
                                                               \tag{KJ.8}
\]
The image here is an ideal, since evaluation is onto
$\mathbb Z_3$; it contains $27\mathbb Z_3$ and lies
inside $9\mathbb Z_3$ by (KJ.6)--(KJ.7).
Thus its quotient has length at least two.

For a positive weight $e$, put $b_e=4^{-e}-1$.
The relation $4^e\gamma-1=0$ is equivalent to $B=b_e$,
and $b_e\in3\mathbb Z_3$.
The image of $x$ generates $W_e=M/(4^e\gamma-1)M$
over $\mathcal O$, because $B$ acts as a scalar there.
By AC-K12 it is a cyclic module of length $a_e=m_e-1$.
Using the original conjugation on $M$, the lengths of
its two components are
\[
       \left\lceil a_e/2\right\rceil
       \quad\hbox{and}\quad
       \left\lfloor a_e/2\right\rfloor.             \tag{KJ.9}
\]
This follows directly from
$\mathcal O=\mathbb Z_3\oplus\delta\mathbb Z_3$
and $\delta^2=-3$.
These are the labels inherited from $M$, irrespective
of the extra sign $(-1)^e$ in the action of actual
conjugation on the Tate twist.
The lower bound two in (KJ.8) implies $a_e\geq3$.

\proofstep{2. Every specialized plus component is killed by nine.}
If $U$ is elementary abelian, then $3U=0$ and
$3M_+=-\delta U$ give $9M_+=0$.
Together with (KJ.8), this makes every specialized
plus component have length exactly two.

The remaining possibility is $U\simeq\mathbb Z/9$.
By (KJ.5), $B$ acts on $U$ as multiplication by
$b_U=3c$ modulo nine, where $c\in\mathbb F_3^\times$.
Set $F(B)=(1+B)^3-1$.
This endomorphism kills $U$.
Logarithmic descent and the certified first-layer bound
AC-K15 give
\[
 C_1=M/F(B)M=C_+\oplus U,
 \qquad \operatorname{length}_{\mathbb Z_3}C_+\leq4,
 \qquad C_+=M_+/F(B)M_+.                             \tag{KJ.10}
\]
Modulo three, $F(B)=B^3$; hence $C_+/3C_+$ has
dimension three, since $d\geq6$.
Also $C_+/BC_+=M_+/BM_+\simeq\mathbb Z/9$.
Thus $C_+$ is not elementary abelian.
Its length bound and its three generators modulo three
force
\[
                   C_+\simeq\mathbb Z/9\oplus
                                     (\mathbb Z/3)^2.
                                                               \tag{KJ.11}
\]
In particular there exists $h(B)\in R$ such that
\[
                       9x=F(B)h(B)x.                \tag{KJ.12}
\]
Reduce this identity modulo three in
$M_+/3M_+=\mathbb F_3[B]/B^d$.
It gives $B^3h(B)=0$ in that quotient, and therefore
\[
                       h(B)\in(3,B^{d-3}).          \tag{KJ.13}
\]
For every $b\in3\mathbb Z_3$, this implies
$3\mid h(b)$; also
$F(b)=3b+3b^2+b^3$ is divisible by nine.
After specialization in (KJ.12), its right side is
therefore a multiple of $27x$, which is zero by
(KJ.5).  Thus $9x=0$ in every specialized plus
component.  Equation (KJ.8) again makes its length
exactly two.

\proofstep{3. The remaining component and the three possibilities.}
If $U$ is elementary abelian, then $4^e=1$ on $U$,
and its quotient by $4^e\gamma-1$ is $U/BU$ of
length one.  The plus length is two, so $a_e=3$
and $m_e=4$ at every weight.
If $U=\mathbb Z/9$ with $B=b_U=3c$ as above, then
\[
                  4^e(1+3c)-1\equiv3(e+c)\pmod9.
                                                               \tag{KJ.14}
\]
Its quotient has length two for $e\equiv-c\pmod3$
and length one elsewhere.
The plus length is always two, so these alternatives
give respectively $a_e=4$ and $a_e=3$, proving
(KJ.1)--(KJ.2).  In particular $a_e\leq4$;
since $(\lambda^4)=(9)$, nine kills every $W_e$.
Counting the integers in the exceptional residue class
gives (KJ.3).
\end{proof}

\paragraph{The exact degree-18 criterion.}
The proof of AC-K20 shows that $F(B)U=0$ also in the
elementary cases.  Consequently the component of $C_1$
with sign opposite to that of $x$ is exactly $U$, with
its cyclotomic action.  The sign of $x$ can be identified
intrinsically: it is the unique conjugation sign whose
component of
$C_1/BC_1=M/BM\simeq\mathbb Z/9\oplus\mathbb Z/3$
has order nine.  If this opposite component is elementary,
all $m_e=4$.  If it is cyclic of order nine, compute
$\gamma=1+3c$ on it using the field automorphism
$\zeta_9\mapsto\zeta_9^4$, $\sqrt[3]{269}$ fixed.
Then $r\equiv-c\pmod3$ is the unique class with $m_e=5$.
Thus a certificate for just this component and its action
would choose among (KJ.2); a complete certified
presentation of $C_1$ with conjugation and cyclotomic
actions is sufficient.  The first-layer upper certificate
AC-K15 does not by itself supply these exact actions or
select an alternative.

The degree-54 reconstruction of the whole module in AC-K20
remains valid, as does its nine-quotient recipe.  For the
exponents alone the present theorem reduces that recipe
to three specializations and, more economically, to the
degree-18 component just specified.  The finite-cover and
even K-theory interpretation following AC-K20 applies
with these sharper integral lengths.  No new external
arithmetic theorem, numerical certification, or archive
locator enters this finite-module deduction.

\InputIfFileExists{modules/principal_regular_twist.tex}{}{}
\section{The actual theta-index normalization and the mixed-character gate}
\label{sec:theta-cohomology-interface}

The regular construction of AC-G27 can be compared with the
literal arithmetic normalization in the supplied source.  This
comparison gives a bounded lattice loss, but the cohomology
bound remains a separate input.  Throughout this section the
source version is the immutable commit
\texttt{7cbf7a0c32fca1badf05931879a37c3f8a2b03b4}.
The newly recovered files and their hashes are recorded in the
companion notes and acquisition manifest.

\begin{theorem}[AC-G28: cubic projectors in the source index units]
\label{AC-G28}
Let $R$ be the ring of integers of a finite extension of
$\mathbb Q_\ell$ containing $\omega$, let $V$ be a finite-dimensional
vector space over its fraction field, and let $M\subset V$ be a
full $R$-lattice.  Suppose $U(M)=M$ and $U^3=1$.
For $J\subset\{0,1,2\}$ put
\[
 P_J=\sum_{j\in J}\frac13\sum_{a=0}^2\omega^{-ja}U^a,
 \qquad V_J=P_JV,
 \qquad M_J^+=P_JM,\quad M_J^-=M\cap V_J.
                                                        \tag{G28.1}
\]
Then
\[
 3M_J^+\subset M_J^-\subset M_J^+,
 \qquad
 M_J^+=V_J\cap(M+V_{J^c}).                              \tag{G28.2}
\]
Writing $r_J=\dim V_J$, the normalized length used in
\texttt{DataN.lenUnit} satisfies
\[
 0\le\frac{\operatorname{length}_R(M_J^+/M_J^-)}
      {\operatorname{length}_R(R/\ell R)}
 \le\begin{cases}
 r_J,&\ell=3,\\
 0,&\ell\ne3.
 \end{cases}                                            \tag{G28.3}
\]
The same inclusion bound holds after intersecting with a
rational subspace or passing to a rational quotient.  If the
integral $C_3$-action is retained through tensor, Weyl, dual,
or other functorial integral constructions and the character
projection is performed only at the end, (G28.2)--(G28.3)
still hold with the constant $3$, independently of their degree.
In contrast, applying a homogeneous polynomial functor of degree
$d$ to the two lattices in (G28.2), after projection, gives the
general inclusion $3^d F(M_J^+)\subset F(M_J^-)$.

If $M=R^a\oplus R[C_3]^b$, with trivial action on $R^a$, then
for every nonempty proper $J$ one has the sharper exact formula
\[
 M_J^+/M_J^-\simeq(R/3R)^b,
 \qquad \operatorname{lenUnit}_3(M_J^+/M_J^-)=b
                      \quad(\ell=3).                  \tag{G28.4}
\]
The theta permutation lattice of AC-G27 therefore has this
normalized length $b=(r^3-r)/3$, where $r=d^g$, for every
nontrivial proper character selection.  This projected-to-intersected
length is a different invariant from the full Fourier determinant
index of AC-G27.
\end{theorem}
\begin{proof}
The polynomial $X^3-1$ is separable over the fraction field.
The three displayed character operators are mutually orthogonal
idempotents and sum to the identity.  Thus $V=V_J\oplus V_{J^c}$.
Every coefficient of $3P_J$ belongs to $R$, so $3P_JM\subset M$;
it also lies in $V_J$.  Conversely $P_J$ is the identity on
$M\cap V_J$.  This proves the inclusions.

If $x=P_Jm$, then $x=m-P_{J^c}m\in V_J\cap(M+V_{J^c})$.
If $x=m+y$ belongs to that intersection, with $m\in M$ and
$y\in V_{J^c}$, applying $P_J$ gives $x=P_Jm$.
This proves the second equality in (G28.2), exactly the
projected/intersected shape in \texttt{NData.MtUp} and
\texttt{NData.MtLow}.

Both $M_J^+$ and $M_J^-$ are free of rank $r_J$.
Their quotient is a quotient of $M_J^+/3M_J^+$, so its
length is at most $r_Jv_R(3)$.  The denominator in
\texttt{lenUnit} is $v_R(\ell)$.  At $\ell=3$ these
valuations agree, and at other primes $3$ is a unit.
This proves (G28.3), including finite length; no manipulation
of the source's $\mathrm{toNat}$ at infinite length is involved.

For any rational subspace $W$, multiplying an element of
$M_J^+\cap W$ by $3$ places it in $M_J^-\cap W$.
Applying a linear quotient map gives the corresponding statement
for image lattices.  For any newly constructed integral module
with the induced order-three action, the same projector proof
starts afresh and has the same denominator $3$.  No estimate of
its dimension or tensor degree enters this argument.
If instead $F$ is a homogeneous polynomial functor of degree
$d$ applied to an already projected inclusion, the maps
$3\operatorname{id}:M_J^+\to M_J^-$ and
$\operatorname{id}:M_J^-\to M_J^+$ induce maps whose composite
is $F(3\operatorname{id})=3^d\operatorname{id}$.
Interpreting their images in the common rational representation
gives the asserted degree-dependent inclusion.  The latter
claim is restricted to such polynomial functors; the
degree-independent projector assertion only requires an
integral order-three action.

It remains to compute one regular orbit.  Write its basis as
$e_0,e_1,e_2$, with $Ue_a=e_{a+1}$, the suffixes taken modulo
$3$.  For a singleton $J=\{j\}$, an eigenvector is
$v_j=\sum_a\omega^{-ja}e_a$.  Its coefficients
are units, so $M\cap V_J=Rv_j$, whereas $P_JM=(1/3)Rv_j$.
The quotient is $R/3R$.
For a two-element $J$, let $k$ be its complement.  The
linear functional defining $V_J$ has unit coefficients and
takes the value $3$ on a generator of the complementary
eigenline.  It follows that
\[
 M\big/\big((M\cap V_J)+(M\cap V_{\{k\}})\big)
              \simeq R/3R.
\]
Projection identifies this quotient with $M_J^+/M_J^-$.
The trivial summands contribute nothing, and direct sums give
(G28.4).  The integral orbit decomposition of AC-G27 now gives
the theta-lattice assertion.
\end{proof}

\begin{corollary}[The exact height and global index normalization]
\label{AC-G28-normalization}
Let $F$ be a totally imaginary number field of degree $D$,
with the same Hermitian metrics on two full $\mathcal O_F$-lattices
$L\subset M$ of rank $n>0$.  Suppose
$3^aM\subset L\subset M$, for an integer $a\ge0$.
In the convention of
\texttt{Formalism.HarderNarasimhan.StatementPl.size} and
\texttt{muDef},
\[
 0\le\mu(M)-\mu(L)
   =\frac{2}{Dn}\log[M:L]\le2a\log3.                  \tag{G28.5}
\]
The same bounds hold for the differences of the maximum
slopes and of the minimum slopes.  Thus a final cubic
character projection costs at most $2\log3$ in each of
these slope conventions.

There is also a direct conditional comparison in the precise
\texttt{LinkN\_at} units.  Suppose an actual source row has
\texttt{etaJ\_tie} identifying its local indices with comparisons
as in (G28.1), and suppose for every coefficient-field place
$w\mid3$ the sum of their selected ranks is at most
$d_0=\texttt{dZ0}\,j$.  The additional index contribution is
then at most
\[
 \sum_{w\mid3}\frac{\log Nw}{[F:\mathbb Q]}
                 e(w/3)\,\eta_J(j,w)
              \le d_0\log3.                            \tag{G28.6}
\]
It is zero at all other coefficient primes.  Hence the
additional term inside the rank-normalized bracket of
\texttt{LinkN\_at} is at most $\log3$, once that row dictionary
and rank comparison have been proved.
\end{corollary}
\begin{proof}
The two covolumes differ by the additive-group index
$[M:L]$.  All degree, discriminant, and rank factors in the
definition of \texttt{ht} are the same.  Since $F$ is totally
imaginary, $2\,\texttt{nPl}=D$; the source definitions give
\[
 \operatorname{size}(M)=\operatorname{ht}(M)^{1/(Dn)},
 \qquad \mu(M)=-2\log\operatorname{size}(M).
\]
Thus the middle equality in (G28.5) follows.  The quotient
$M/L$ is a quotient of $M/3^aM$, whose cardinality is
$3^{aDn}$, proving the inequalities.

Intersecting with any rational subspace preserves the
inclusions and the constant $3^a$.  Passing to the image in
any rational quotient also preserves them.  Apply (G28.5)
to each such intersection or quotient.  The maximum slope
is the supremum of the intersection slopes: a sublattice
with a fixed rational span has no larger slope than the
intersection with that span.  Similarly the minimum slope
is the infimum of the saturated quotient slopes.  Taking
these suprema and infima proves the corresponding bounds.

For (G28.6), (G28.3) and the assumed rank comparison imply
$\eta_J(j,w)\le d_0$.  All summands are nonnegative.
Writing $Nw=3^{f(w/3)}$ and using
$\sum_{w\mid3}e(w/3)f(w/3)=D$ yields the stated bound.
The source \texttt{etaSum} uses the weight
$\texttt{lwP}(w)\,\texttt{ramIdx}(w)$, with no factor $2$;
\texttt{DataArith.sizeOfVal} uses twice that weight.
This accounts for $\log3$ in (G28.6) and $2\log3$ in
(G28.5).  Neither is a normalization of a spectral trace.
\end{proof}

\begin{remark}[What the theta ties do and do not identify]
The exact source locators for the preceding comparison are
\texttt{DataN.lean}: \texttt{lenUnit}, \texttt{NData.MtUp},
\texttt{MtLow}, \texttt{etaLoc};
\texttt{NStageTrivialCharacterCoeffEq.lean}: the field
\texttt{NStage.etaJ\_tie}; and
\texttt{DataParams\_2}: \texttt{SizeData\_2.dZ0},
\texttt{lwP}, \texttt{sz}.  The last two displayed filenames
with an ASCII suffix refer to the supplied files whose actual
suffix is the Unicode subscript $2$.
The source \texttt{DataTheta.ThetaX.TieRowIsom} instead says
that a Petersson product scaled by \texttt{mY} equals the
given row inner product; \texttt{TieMY} expresses
\texttt{mY} as a product of explicit parabolic indices.
The source \texttt{DataCoeff.normFock} divides a Fock pairing
by the pairing of the chosen highest-weight vector with itself.
Neither of these declarations divides a determinant index
by the number of theta sections.
The inequality (G28.6) is therefore a conditional comparison
with the actual row index: the identification of the new
regular construction with those rows, including the selected
rank bound, has not been supplied by AC-G27 or by the
archimedean orthogonality of its projectors.
\end{remark}

\begin{theorem}[AC-G29: scalar cancellation and the actual cohomology gate]
\label{AC-G29}
Let $G$ be a group, $R$ a commutative ring, and
$\kappa:G\to R^\times$ a character.  For $R$-representations
$V,W$, tensoring both with the free rank-one character
$R(\kappa)$ gives an integral, $G$-equivariant identification
\[
 \operatorname{Hom}_R(V\otimes R(\kappa),
                     W\otimes R(\kappa))
       \simeq\operatorname{Hom}_R(V,W).                 \tag{G29.1}
\]
In particular, the literal family in \texttt{NData.chi\_def},
\[
          \chi_c=\chi_1\varepsilon^{-(c-1)},
                         \qquad 1\le c\le t,
\]
has pairwise coefficients
\[
 \operatorname{Hom}(R(\chi_a),R(\chi_b))
                =R(\varepsilon^{a-b}).                 \tag{G29.2}
\]
A common cubic twist of this family leaves these coefficients
unchanged, including at the coefficient prime $3$.

For comparison, put $P=R[C_3]$ with its regular action.
Under conjugation by the regular permutation,
\[
              \operatorname{End}_R(P)\simeq P^{\oplus3}
                                                        \tag{G29.3}
\]
integrally.  After extending scalars to a characteristic-zero
splitting field, the mixed family
$R(\chi_1\kappa^i\varepsilon^{-(a-1)})$, $i=0,1,2$,
has cross coefficients
\[
       R(\kappa^{j-i}\varepsilon^{a-b}).                \tag{G29.4}
\]
Thus retaining regular cubic blocks produces the required
mixed cubic coefficients, whereas a common scalar twist
cannot produce them through (G29.2).

In the supplied source the existing arithmetic tie is
\texttt{NData.sha\_eq}: at each finite length and
$0<e<t$, it identifies \texttt{Rows.shaLen} with the
strict cohomology length of
$\mathbb Q_\ell/\mathbb Z_\ell(-e)$ over the actual field
\texttt{Host.E}.  It is not a tie to (G29.4).
The source \texttt{LinkN\_at} has the direction
\[
 \eta\le d_0\left(C\,\mathcal E+
           \sum_\ell\log\ell\sum_{0<e<t}
                    \operatorname{shaLen}_\ell(e)\right),
                                                        \tag{G29.5}
\]
with an explicitly defined geometric quantity $\mathcal E$.
Consequently an upper bound on the additional index from
AC-G28 supplies no upper bound on the cohomology sum in
(G29.5).  That sum is bounded separately by
\texttt{LinkV\_at} or by the varying-field \texttt{VVle}
input to \texttt{VAt}/\texttt{VWinAt}.
\end{theorem}
\begin{proof}
Choose the standard generator of the free rank-one module
$R(\kappa)$.  It identifies the underlying Hom modules in
(G29.1).  For a map $f$, the Galois action on the twisted Hom
is
\[
 g\cdot f=(\kappa(g)\rho_W(g))\,f\,
                 (\kappa(g)\rho_V(g))^{-1}
            =\rho_W(g)f\rho_V(g)^{-1}.
\]
This proves equivariance without dividing by the order of
the character.  Inverting the tensor functor by
$R(\kappa^{-1})$ also shows that it is exact and preserves
extension classes and all restriction maps.  For rank-one
characters the Hom character is $\chi_b\chi_a^{-1}$,
which gives (G29.2).  When $b>a$, its negative Tate weight
is precisely $e=b-a$ in the source \texttt{qzTw} convention.

For (G29.3), index a basis of $P$ by $i\in\mathbb Z/3$.
The matrix units $E_{ij}$ form a basis of its endomorphism
module.  Conjugation takes $E_{ij}$ to $E_{i+1,j+1}$.
There are exactly three free orbits, distinguished by
$j-i$.  Each is an integral regular module, proving (G29.3)
over every $R$, including residue characteristic $3$.
The rank-one ratio calculation gives (G29.4).
For finite projective $V,W$, one can equivalently identify
the regular-block Hom with
\[
 \operatorname{Hom}_R(V\otimes P,W\otimes P)
       \simeq\operatorname{Hom}_R(V,W)
                    \otimes\operatorname{End}_R(P).
\]
This expresses the new mixed coefficient before any
nonintegral character splitting.

The asserted source identification is read directly from
\path{GaloisCohomology.Statement.LData}: its finite set
$S$ is required to equal the primes above $\ell$, and its
character must pull back to the cyclotomic character.
The definitions are
$\texttt{qzTw}(e)=\mathbb Q_\ell/\mathbb Z_\ell(-e)$,
$\texttt{Sha}=\texttt{Sha1}$ for the local family at $S$,
and $\texttt{shaLen}=\operatorname{length}_{\mathbb Z_\ell}
\texttt{Sha}$.
In \texttt{NData}, \texttt{ldat} is such a datum for
$\texttt{H.E}$; \texttt{sha\_eq} gives the equality to the
row field whenever the actual length equals a natural number.
This ties the row variable to arithmetic when finite.  The
abstract \texttt{Rows.Data} countermodel of AC-G25 therefore
applies to that raw structure, and is not a countermodel to
a fully tied finite \texttt{NData}.

Finally \texttt{GlobalTheoremInputs.LinkN\_at} is exactly
(G29.5), under summability of its \texttt{vvTerm}.
Even numerically, an inequality
$0\le\eta\le d_0(E+V)$ with $d_0>0$ and $E,V\ge0$
cannot imply an upper bound on $V$: for example fix
$\eta=0$, $d_0=1$, $E=0$, and let $V$ be arbitrarily large.
This is a counterexample to reversing that inequality,
not an assertion that actual cohomology lengths in a fixed
number field can be assigned arbitrarily.
The independently declared \texttt{LinkV\_at} bounds the
cohomology sum from above.  The definitions
\texttt{ExistsPageN.VVle}, \texttt{VAt}, and
\texttt{VWinAt} quantify over the source host fields and
their windows.  Neither (G28.3) nor (G29.1) changes those
quantifiers or proves the required bound.
\end{proof}

\begin{proposition}[Exact transport of the local Hecke fragment]
\label{AC-G29-Hecke}
Let $R$ be the integers of a finite extension of $\mathbb Q_3$
containing $\omega$, and write $e_R=v_R(3)$.
Let $T\subset R^X$ be an $R$-subalgebra, where $X$ is finite
and nonempty, such that $T$ is local and separates the points
of $X$.  Let $M$ be a finite free $R$-module with a faithful
$T$-action.  The regular cubic extension
\[
 T'=T[C_3]\subset R^{X\times\{0,1,2\}},
             \qquad M'=M\otimes_R R[C_3]               \tag{G29.6}
\]
again gives a local algebra separating its rational
eigencharacters, acting faithfully on a finite free module.
These are the corresponding local algebra clauses of
\texttt{NData}; no global Hecke or automorphic realization is
asserted here.

Let $P$ be the rational projector associated with any subset
of $X$, and put
\[
 A=P(M),\qquad B=M\cap\operatorname{im}P,
                 \qquad r=\dim\operatorname{im}P.
\]
Select one cubic character in (G29.6).  For the joint
projector $P\otimes P_j$ the new projected and intersected
lattices are, respectively,
\[
           (1/3)A\,v_j,\qquad B\,v_j,
\]
and consequently
\[
 \operatorname{length}_R(M'_{P,j,+}/M'_{P,j,-})
       =\operatorname{length}_R(A/B)+r e_R,
 \qquad \eta'_{P,j}=\eta_P+r.                          \tag{G29.7}
\]
At a coefficient prime different from $3$, the same joint
character comparison has exactly the old length; the
algebra $T[C_3]$ is then the product of three copies of $T$
over a splitting coefficient ring, and each character factor
is treated separately.

The old length in (G29.7) can be arbitrarily large even in
the displayed local algebra fragment.  Indeed, for a
uniformizer $\pi$ and any integer $n\ge1$, take
\[
 T_n=M_n=\{(a,b)\in R^2:a\equiv b\pmod{\pi^n}\}.
\]
For the first-coordinate Hecke projector one has
\[
 A=R(1,0),\quad B=\pi^nR(1,0),\quad
       \eta_P=n/e_R,\quad \eta'_{P,j}=n/e_R+1.          \tag{G29.8}
\]
Thus the cubic splitting contributes a controlled additive
term and does not bound the original Hecke congruence length.
\end{proposition}
\begin{proof}
Since $R$ is Noetherian, $T\subset R^X$ is a finite
torsion-free $R$-module.  It is therefore finite free.
Point separation implies that $T\otimes_R\operatorname{Frac}R$
contains every coordinate idempotent.  For a fixed $x$, choose
for each $y\ne x$ an element $f_{xy}\in T$ with
$f_{xy}(x)\ne f_{xy}(y)$; the product of
\[
       \frac{f_{xy}-f_{xy}(y)}{f_{xy}(x)-f_{xy}(y)}
\]
over $y\ne x$ is that idempotent.  Thus the rational
eigenspace projectors used in the statement exist.

Evaluation of $\sum_{a=0}^2t_a\sigma^a$ at $(x,j)$
means $\sum_a t_a(x)\omega^{ja}$.  A polynomial of degree
at most two that vanishes at all three distinct roots is
zero over the fraction field.  Evaluation is consequently
injective.  Elements of $T$ separate unequal $x$'s, and
$\sigma$ separates unequal $j$'s, so $T'$ separates all
the stated characters.
For locality, every maximal ideal of the finite integral
extension $T'$ contracts to the unique maximal ideal
$\mathfrak m_T$ of $T$.
The residue field $k_T$ has characteristic $3$, and
\[
        T'/\mathfrak m_TT'=k_T[C_3]
                    \simeq k_T[z]/(z-1)^3
\]
has just one maximal ideal.  Hence $T'$ is local.
For faithfulness, apply $\sum_a t_a\sigma^a$ to
$m\otimes1$ for every $m\in M$.
The result is $\sum_a t_am\otimes\sigma^a$; its
vanishing implies that every $t_a$ acts as zero and hence
is zero.  This proves all asserted local algebra properties.

Identify $M'$ with three copies of $M$.
On one regular cubic orbit, $P_j$ has image
$(1/3)Rv_j$, with all three coordinates of $v_j$ units.
Thus applying the joint projector gives $(1/3)A v_j$.
A vector $a v_j$ in the joint rational eigenspace lies in
$M^3$ exactly when $a\in M$; since $a$ already belongs
to $\operatorname{im}P$, its intersection is $Bv_j$.
The exact sequence
\[
 0\longrightarrow A/B\longrightarrow (1/3)A/B
       \longrightarrow (1/3)A/A\longrightarrow0
\]
has last term isomorphic to $(R/3R)^r$.
Taking lengths and dividing by $e_R$ proves (G29.7).
If $3$ is a unit, the last term is zero and the usual
Chinese remainder theorem gives the three local character
factors as stated.

Finally $T_n$ is free with basis $(1,1)$ and $(\pi^n,0)$.
A pair in $T_n$ is a unit exactly when its common residue
modulo $\pi$ is nonzero: the coordinate inverses still
agree modulo $\pi^n$.  Hence $T_n$ is local.  The element
$(\pi^n,0)$ separates its two rational characters, and
its regular action on itself is faithful.  Projection of
$(1,1)$ gives $(1,0)$, while a pair $(a,0)$ lies in
$T_n$ exactly when $a\in\pi^nR$.
This proves (G29.8), and (G29.7) gives its joint cubic
version.  The example concerns precisely this local
algebra/lattice fragment; it is neither a countermodel to
all the fields of \texttt{NData} nor a refutation of a
supplied global theorem.
\end{proof}

\begin{proposition}[The full-packet regular tensor fails the existing good-row test]
\label{AC-G29-good-row}
Let $V$ be a nonzero representation of dimension $n$ over a
field, and let $P$ be the three-dimensional regular permutation
representation through a cubic character $\kappa$.
Then $V\otimes P$ is reducible: it has the canonical proper
invariant subspace
\[
             V\otimes\langle e_0+e_1+e_2\rangle.
                                                        \tag{G29.9}
\]
This remains true in characteristic $3$.
If $\kappa$ is trivial on a subgroup $I$, then for every
$\tau\in I$,
\[
 \det\bigl(X-\tau\mid V\otimes P\bigr)
                  =\det\bigl(X-\tau\mid V\bigr)^3.     \tag{G29.10}
\]
In particular this polynomial is not squarefree.

The source \texttt{CyclotomicNorms.StatementT.IsGood}
requires each prescribed auxiliary polynomial to be
squarefree.  A literal extension of an entire packet by
the regular cubic representation therefore cannot preserve
this good-row condition at an auxiliary inertia place
where $\kappa$ is unramified, if the prescribed polynomial
is transported through the existing inertia tie.
For $\chi_{269}$, every auxiliary rational prime different
from $269$ has this unramified property.
This rules out that particular full-packet substitution;
it does not rule out a new construction using separately
twisted packets or a regular action on a multiplicity space.
\end{proposition}
\begin{proof}
The vector $e_0+e_1+e_2$ is fixed by every permutation in
the regular action and is nonzero over every field.
Tensoring its one-dimensional span with $V$ gives an
invariant subspace of dimension $n$ in a space of
dimension $3n$, proving (G29.9).
If the action on $P$ is trivial on $I$, the restricted
representation is the direct sum of three copies of
$V|_I$.  Its characteristic polynomial is the cube in
(G29.10).  A nonconstant polynomial has a root over an
algebraic closure, and that root has multiplicity at least
three in its cube.  Thus the cube is not squarefree,
in every characteristic.

The source clause is explicit: following the monicity,
degree, and nonzero-constant-term conditions, \texttt{IsGood}
requires $\forall v,\ \texttt{Squarefree}(f(v))$.
The tie is also explicit in \texttt{DataN.NTypes}:
\texttt{inertiaCharpoly\_type} expresses the inertia
polynomial at the chosen tame generator as the product
of the linear factors obtained from the roots of
\texttt{items\_2}; \texttt{teichC\_inj} is injective on
the nonzero roots.  Consequently a squarefree monic
polynomial with nonzero constant term would give distinct
inertia eigenvalues under this tie.
The clause \texttt{NAut.wdOf\_ine} identifies those
inertia factors with the characteristic polynomial of an
actual Galois representation away from its coefficient
prime.  This is incompatible with (G29.10) at any such
auxiliary place.  Equivalently, a direct polynomial
transport would give $f(v)^3$, which fails the displayed
squarefree clause immediately.

The character $\chi_{269}$ is unramified away from $269$
by AC-K04 and AC-G22; restriction to a larger base field
preserves this unramified property.  The application to
the auxiliary inertia places follows.  The statement is
restricted to tensoring the entire Galois packet.  Applying
scalar twists to three separate packets leaves each
individual auxiliary inertia representation unchanged
and does not encounter this particular cubic repetition,
but requires the new cross-block arithmetic tie already
identified in AC-G29.
\end{proof}

\begin{remark}[The narrowed remaining theorem]
To use the regular principal construction in the supplied
arithmetic argument, one needs an actual mixed-block version
of its automorphic/cohomology dictionary: the new integral
row and theta lattices must be identified with the source
objects, the selected ranks must be controlled as in
(G28.6), and the strict local conditions and arithmetic tie
must be proved for the cross coefficients (G29.4).
The normalized extra character-projection loss can then be
bounded by AC-G28.  The uniform upper bound on those
cohomology lengths over the growing host family remains
a separate arithmetic assertion.  The construction does
not replace the field \texttt{Host.E} by an etale CM algebra
in an existing theorem, and its common scalar-twist variant
does not alter the pairwise pure-Tate coefficients at all.
\end{remark}

\section{A constant quadratic enlargement of the signed hosts}
\label{sec:mixed-host-interface}

The fixed cubic coefficient is scalar over a field containing
$K=\mathbb Q(\omega)$.  The literal signed Host in AM-H17 does
not contain this field.  We first make the required enlargement
and its integral cost precise.  Here ``Host'' has exactly the
scope of \texttt{DataHost.lean}; it does not include the later
theta identities or the proposition \texttt{VWin}.

\begin{proposition}[AC-G30: adjoining the cubic coefficient field]
\label{AC-G30}
Write the AM-H17 field as
\[
 C_h=\mathbb Q(\zeta_{2^{h+2}},\zeta_7),\quad
 H=C_hK_qK_d,\quad K_q=\mathbb Q(\sqrt{-q}),\quad
 K_d=\mathbb Q(\sqrt{-d}),
\]
and put $n=3\cdot2^{h+2}$ and $g=[H:\mathbb Q]/2=2n$.
Assume $h\geq294$, the splitting bound is at least $269$,
and the AM-H17 primes are chosen large enough for the ranks
specified below.  Let $F=HK$.  Then
\[
 [F:H]=2,\qquad
 \operatorname{rd}(F)=\sqrt3\operatorname{rd}(H),
 \qquad \mu(F)=\mu_{3\cdot7\cdot2^{h+2}}.             \tag{G30.1}
\]
The AM-H17 construction gives a complete bare Host over $F$,
with the same $q,d$ and
\[
 Q'_S=3Q_S,\quad n'=2n,\quad n'_0=2^h,\quad o'=24.
                                                               \tag{G30.2}
\]
More explicitly, if the original ranks are $(k,m,t)$ and
$t\geq2k+m+4$, take $(k',m',t')=(2k,m,t)$ and require
$q,d>4t(2k+m-1)+1$.  Thus this enlargement preserves the
weight cutoff $t$; these are explicit rank and prime-size
conditions, not a claim that arbitrary old ranks still work.

Take the AM-H17 choice $\mathfrak c_K=\mathcal O_{K_d}$ and
write its different-splitting ideal and lattice as
$\mathfrak b_H$ and $AA_H=\mathfrak b_H^{-1}$.
There is a three-supported ideal $\mathfrak b_3$ of $F$ with
\[
 \mathfrak b_3\overline{\mathfrak b}_3
       =\mathfrak D_{F/H},\quad
 N_{F/\mathbb Q}\mathfrak b_3=3^g,
\]
such that the new source lattice is
\[
 AA_F=(AA_H\mathcal O_F)\mathfrak b_3^{-1},\qquad
 [AA_F:AA_H\mathcal O_F]=3^g.                         \tag{G30.3}
\]
It is unchanged relative to the extended old lattice away
from three.  With the source coefficient
$\xi_0=-1/\sqrt{-d}$ it satisfies the exact principal identity
\[
 (\xi_0)AA_F\overline{AA_F}\mathfrak D_{F/\mathbb Q}=1.
                                                               \tag{G30.4}
\]
Its Euclidean covolume is
\[
 \operatorname{covol}(AA_F)
       =\left(\frac{\sqrt d}{2}\right)^{2g},           \tag{G30.5}
\]
and the determinant of the principal trace metric makes its
metric covolume one.  The logarithm of the integral enlargement
index in (G30.3) is $g\log3$; that logarithmic cost must not be
substituted for an unlogged theta norm.

Nevertheless neither $H$ nor $F$ is a literal
\texttt{IsChiField} field for any one Dirichlet character
with these same $h,q,d$.  This last assertion concerns the
single-character definition; it does not assert that several
instances of a separately supplied theorem \texttt{VWin}
could never imply a bound for these fields.
\end{proposition}
\begin{proof}
\proofstep{1. The field and frame.}
The four factors used in AM-H17 have ramification support
in $\{2\}$, $\{7\}$, $\{q\}$ and $\{d\}$.
The primes $q,d$ exceed $269$.  Thus $H$ is unramified at
three and $H\cap K=\mathbb Q$, since $K/\mathbb Q$
ramifies at three.  The discriminants are coprime, so
\[
 \mathcal O_F=\mathcal O_H\otimes_{\mathbb Z}\mathcal O_K,
 \qquad
 |\operatorname{disc}F|
       =|\operatorname{disc}H|^2\,3^{[H:\mathbb Q]}.
\]
This is the integral-basis/discriminant formula used in
AM-H17, with the additional quadratic factor of discriminant
$-3$; see the primary locator there, Milne,
\emph{Algebraic Number Theory}, Lemma~6.5 and Remark~6.6(c),
printed pp.~100--101.  Taking the $[F:\mathbb Q]$-th root
gives (G30.1).  The roots-of-unity argument in AM-H17 applies
at every old prime, while the exact inertia index at three
is two and excludes a primitive ninth root.  It gives the
displayed root group.

Use the frame group
\[
 A'=\operatorname{Gal}(C_hK/\mathbb Q)
        \simeq C_2\times C_{2^h}\times C_6\times C_2.
\]
Project to the two-power cyclotomic factor modulo conjugation,
and use the element $5$ in that factor as the distinguished
generator, acting trivially on the other factors.  Its image
has order $2^h$, with the same injectivity on the distinguished
cyclic subgroup as in AM-H17.  Thus $|A'|=2n$ and
$o'=|A'|/2^h=24$.  The integer inequality
$12^{64}<2^{230}$ used in AM-H17 gives
$24^{64}<2^{294}\leq2^h$.

All rank clauses in \texttt{Params.Ranks} hold for
$(2k,m,t)$ under the stated extra slack.  In particular
$2n'=4n\leq2k$, the parity conditions are unchanged,
and $k'+m-1+5\leq t$ is exactly the added condition.
The remaining inequalities either stay the same or become
weaker when $k$ increases.  The strengthened lower bound
on $q,d$ is exactly the new \texttt{qd\_big} clause.

\proofstep{2. All bare Host field clauses survive.}
Use the rebasing map of AM-H17 with the conjugation of
$C_hK$ in place of that of $C_h$.  It again sends full
complex conjugation to $((-1,-1),1)$ and pure $d$-inertia
to $((1,-1),1)$.  Restriction from the cyclotomic field of
conductor $3Q_Sqd$ supplies the embedding, Artin map,
quadratic-coordinate identities, and Frobenius, inertia
and decomposition descriptions in \texttt{HostRaw} and
\texttt{HostDer}.  The new conductor is exact at three
because its quadratic factor acts nontrivially; exactness
at the old primes is preserved by fixing the added factor.
The inequalities $Q'_S\leq(n')^2$ and the required
root-discriminant bounds follow from
$Q'_S=3Q_S$, $n'=2n$, and (G30.1).

There is only one new local check.  Three splits in $K_d$
by the stipulated small-prime splitting condition.  Hence
its decomposition group fixes $K_d$ and does not contain
complex conjugation.  The old ramified-prime checks remain
valid: if conjugation belonged to a new decomposition
group, its restriction would belong to the old one.
Thus the CM extension $F/F^+$ is unramified at every finite
place, and the different at the new prime can be separated
into conjugate halves.  The old proofs of the ordinary-type,
quadratic splitting and sign-character clauses use the same
$K_d$ and the same two quadratic coordinates, so they apply
without alteration to the added factor.

The old four auxiliary primes in the AM-H17 construction
were chosen congruent to one modulo $15Q_Sqd$.
They are therefore already congruent to one modulo three
and split in the added quadratic factor.  They retain all
bare Host level and coprimality requirements.  One can keep
$\mathfrak c_K=\mathcal O_{K_d}$, and the choices of empty
$TT,T_2$ allowed for the bare structure.  As in AM-H17, this
last observation does not claim the later quantitative
lower bounds on these sets.

The exact declarations checked here are
\texttt{HostRaw}, \texttt{HostDer}, \texttt{HostAx} in
\path{Formalism/SharpRestriction/RamificationFrobenius/TotientGaloisClasses/ConjEMemHostFIff.lean},
and \texttt{Host} in \texttt{DataHost.lean}.
The separately defined property
\texttt{HostRaw.CoprimeFifteen} is not a field of this Host;
the new conductor is divisible by three and does not satisfy
that separate property.

\proofstep{3. The exact principal lattice at three.}
Since the discriminants are coprime,
\[
 \mathfrak D_{F/H}=(\sqrt{-3})\mathcal O_F.
\]
Choose one of the two primes of $K_d$ above three, and let
$\mathfrak b_3$ be the part of this relative different above
that chosen prime.  Conjugation exchanges the two supports,
and every automorphism fixing $K_d$ preserves each support.
Consequently
\[
 (\mathfrak b_3,\overline{\mathfrak b}_3)=1,
 \qquad
 \mathfrak b_F=\mathfrak b_H\mathcal O_F\,\mathfrak b_3
\]
is the required stable half of $\mathfrak D_{F/K_d}$.
Its inverse is (G30.3).  The norm of the relative different
is $3^{2g}$; its two conjugate halves have equal norm, so
$N\mathfrak b_3=3^g$.  For invertible fractional ideals,
$[I\mathfrak b_3^{-1}:I]=N\mathfrak b_3$, which proves
the index.  Transitivity of differents and the old principal
identity give (G30.4) directly.

The extended old lattice has Euclidean covolume
\[
 2^{-2g}\sqrt{|\operatorname{disc}F|}\,N(AA_H)^2
       =3^g\left(\frac{\sqrt d}{2}\right)^{2g}.
\]
Division by its enlargement index gives (G30.5).
The old positive trace metric has diagonal coefficient
$2/\sqrt d$ in each of the $2g$ complex coordinates, so
its volume density cancels this covolume.  For the signed
metric obtained by pulling back the real unit $\kappa$,
the product of $|\tau(\kappa)|$ remains one, and the
same determinant calculation applies.  These statements
compare the actual trace lattices and determinant metrics;
they do not identify normalized algebraic theta sections.

Finally the roots of unity in $F$ are those in $C_hK$.
Their norm to $K_d$ equals their rational norm from that
CM cyclotomic factor, squared, and is one.  This verifies
the remaining \texttt{norm\_mu} clause of the bare Host.

\proofstep{4. The single-character obstruction.}
Put
\[
 B_h=\mathbb Q(\zeta_{2^{h+2}})^+K_qK_d.
\]
The literal definition \texttt{IsChiField} in
\path{Development/AutomorphicArithmetic/CharacterBoxes/ExistsPageN.lean}
requires its field to be
\[
                           B_hL_\psi,                 \tag{G30.6}
\]
where $L_\psi$ is the cyclic field cut out by the one finite
Dirichlet character $\psi$.  Indeed its displayed stabilizer
is the intersection of the kernel of $\psi$ with the three
stabilizers defining the real two-power field and the two
quadratic fields.  Galois correspondence identifies this
intersection with the compositum in (G30.6), regardless
of which sufficiently divisible ambient conductor is used.
The finite character image is cyclic, so
$\operatorname{Gal}(B_hL_\psi/B_h)$ is cyclic.

By disjoint ramification of the factors,
\[
 \operatorname{Gal}(H/B_h)=C_2\times C_6,
 \qquad
 \operatorname{Gal}(F/B_h)=C_2\times C_6\times C_2.
                                                               \tag{G30.7}
\]
Neither group is cyclic.  Thus neither field can satisfy
the same single-character definition.  This proves the
precise obstruction claimed, while leaving open a comparison
using several cyclic subfields and several quantitative
instances of \texttt{VWin}.
\end{proof}

\begin{remark}[The scalar coefficient and the one-Speh-block gate]
Over $H$, the splitting field of $X^3-269$ still has group
$S_3$: its only possible nontrivial abelian intersection
with the abelian field $H$ is $K$, and $K\not\subset H$.
The rational rank-two Artin representation therefore remains
irreducible over $H$, even after extending its coefficient
field.  Over $F=HK$ it splits into the two cubic characters.
The CM-linear Hom construction of AC-G23 is scalar over
this latter base.

On $F$, scalar twists of separate packets preserve their
inertia types at the two old auxiliary primes, since the
cubic character is unramified there.  This avoids the
repeated-inertia-polynomial defect of the full regular
tensor in AC-G29.  It does not repair the source's
arithmetic interface.  The exact declarations
\texttt{NData.chi\_def}, \texttt{muPart},
\texttt{HasParamX}, \texttt{ZN} and \texttt{IsRhoBar}
in \texttt{LevelLetterCocycles/DataN.lean} use the one
distinguished Speh block
$\langle1,\mu_C,t\rangle$ and its characters
$\eta_c=\eta_1\varepsilon_\ell^{-(c-1)}$.
All pairwise ratios inside that block are pure Tate.

For example, replacing these characters by
$\eta_c=\eta_0\chi^{i_c}\varepsilon_\ell^{-(c-1)}$
can satisfy that same \texttt{chi\_def} only if all
$i_c$ agree modulo three.  Evaluate the ratio at the
three-primary tame inertia at $269$: $\chi$ is surjective,
whereas $\varepsilon_\ell$ is trivial there.  For
$\ell\ne269$ this is the usual unramified cyclotomic
action; for $\ell=269$ it follows because
$\mathbb Z_{269}^\times$ has no nontrivial pro-three
subgroup.  Thus $\omega^{i_c-i_1}=1$ for every $c$.
Separate scalar packets consequently require a new
comparison for extensions between distinct distinguished
blocks, together with its integral theta and Hecke ties.
The bare Host enlargement above supplies none of those
additional identities.
\end{remark}

\begin{proposition}[AC-G30-Fourier: fixed finite-character Fourier defects]
\label{AC-G30-Fourier}
For the extension $F/B_h$ in (G30.7), let
$G=\operatorname{Gal}(F/B_h)\simeq C_2^3\times C_3$.
Its characters are the restrictions of the $24$ Dirichlet
characters of modulus $84$.  Let
\[
 \mathcal F:\mathcal O[G]\longrightarrow
       \bigoplus_{\eta\in\widehat G}\mathcal O(\eta)
\]
be the integral Fourier map.
At three, take $\mathcal O=\mathbb Z_3[\omega]$.
The cokernel is killed by three and has
$\mathcal O$-length $24$.
At two, take the unramified quadratic ring
$\mathcal O=\mathbb Z_2[\omega]$.
The elementary divisors are
\[
 1\ (3\text{ times}),\qquad
 2\ (9\text{ times}),\qquad
 4\ (9\text{ times}),\qquad
 8\ (3\text{ times}).                                \tag{G30.8}
\]
Thus this cokernel is killed by eight and has
$\mathcal O$-length $36$, or underlying
$\mathbb Z_2$-length $72$.
At all other coefficient primes the Fourier map is an
integral isomorphism after an unramified coefficient
extension containing the sixth roots of unity.

These are fixed finite defects, independent of $h,q,d$
and of a common Tate twist.  Their residual characters
need not be trivial.  The common fixed extension
$\mathbb Q(\zeta_{168})$ trivializes their finite actions
and the required finite Tate twists after base change
from $B_h$.
\end{proposition}
\begin{proof}
The full two-power cyclotomic field is obtained from its
maximal real subfield by adjoining $i$.  Consequently
\[
 F=B_h\mathbb Q(i,\zeta_7,\omega)
                  =B_h\mathbb Q(\zeta_{84}).
\]
Its relative degree is $24$ by (G30.7), equal to
$\varphi(84)$, so restriction identifies the two Galois
groups.  This proves the character assertion.

Fourier evaluation on a product of groups is the tensor
product of their evaluation matrices.  Over any ring in
which two is not a unit, the $C_2$ matrix
$\left(\begin{smallmatrix}1&1\\1&-1\end{smallmatrix}\right)$
has Smith factors $1,2$, by integral row and column
operations.  The threefold tensor product has factors
$2^j$ with multiplicities $\binom3j$.  At two the $C_3$
Fourier matrix is invertible, since three is a unit, and
adds three copies of each factor.  This gives (G30.8).
The residue field of $\mathbb Z_2[\omega]$ has degree two,
which gives the stated underlying length.

At three all three $C_2$ matrices are invertible.
AM-H06 computes the $C_3$ matrix factors
$1,\lambda,\lambda^2$, of length three and killed by
$\lambda^2$, a unit times three.  Eight copies give the
length $24$.  If the coefficient prime divides neither
two nor three, every group order is a unit, and the usual
inverse Fourier matrix is integral.

Twisting both lattices by a common character changes
none of their underlying Smith factors.  Their untwisted
Galois actions factor through $\mathbb Q(\zeta_{84})$.
The three-primary quotient is killed by three, and the
two-primary quotient by eight; adjoining $\mu_3$ and
$\mu_8$ therefore kills every residual Tate action as
well.  The single field $\mathbb Q(\zeta_{168})$ contains
these fields.  This last argument retains the nontrivial
prime-to-coefficient characters, rather than declaring
every composition factor trivial on the original base.
\end{proof}

\section{What descent actually preserves in the source windows}
\label{sec:finite-window-descent}

\begin{proposition}[AC-G31: exact two-power descent and its window cost]
\label{AC-G31}
For a number field $L$ and $e\geq1$, let
$S_\ell(L,e)$ be the strict group with coefficient
$\mathbb Q_\ell/\mathbb Z_\ell(-e)$: it is unramified
outside $\ell$ and zero at each completion above $\ell$.
Suppose $L'/L$ is cyclic of degree $2^s$ and all the
groups over $L'$ in the range $1\leq e<t$ are finite.
Restriction maps strict groups to strict groups, is
injective for odd $\ell$, and satisfies
\[
 \operatorname{length}_{\mathbb Z_2}
       \ker\bigl(S_2(L,e)\longrightarrow S_2(L',e)\bigr)
                  \leq s.                            \tag{G31.1}
\]
In particular a bound \texttt{VVle}$(L',t,A)$ implies
the mathematical bound
\[
 \texttt{VVle}\bigl(L,t,A+s(t-1)\log2\bigr).         \tag{G31.2}
\]
Here the groups are the actual strict Tate groups identified
by the source's \texttt{LData}; no Artin coefficient has
been inserted into that definition.

Fix a finite set of Dirichlet characters $\psi_i$, a single
$0<\epsilon<1/52$, and the monotone page function
\[
 a(x)=\left\lceil(2+8\epsilon)\log_2 x\right\rceil
                         \quad(x\geq1).
\]
Assume the corresponding source propositions
\texttt{VWinAt}$(\psi_i,401)$ have been supplied, with
window constants $J_i$.  Put $J_* =\max(2,\max_i J_i)$.
For all sufficiently large $Z$, and any one pair $q,d$
satisfying \texttt{WinPair}$(\epsilon,Z,q,d)$, let
$L_i(Z)$ be its cyclic \texttt{IsChiField} at page $a(Z)$.
Then each such field satisfies, with
$t_Z=\lfloor Z^{401}\rfloor$,
\[
 \texttt{VVle}\left(L_i(Z),t_Z,
   \exp\left(\tfrac12 Z^{\epsilon J_*/2}\right)
      +(t_Z-1)\left((2+8\epsilon)(J_*/2-1)\log Z
                                      +\log2\right)\right).
                                                               \tag{G31.3}
\]
This is a common lower-page comparison.  It does not give
the required source budget
$\exp(\tfrac12 Z^\epsilon)$ without further control of
the window constants or a simultaneous-window theorem.
\end{proposition}
\begin{proof}
Restriction preserves the local zero conditions and the
unramified conditions: each new inertia group maps into
the old inertia group, even at places ramified in $L'/L$.
Inflation--restriction embeds the kernel of global
restriction into
\[
 H^1\left(\operatorname{Gal}(L'/L),
  (\mathbb Q_\ell/\mathbb Z_\ell(-e))^{G_{L'}}\right).
                                                               \tag{G31.4}
\]
For odd $\ell$, the group order is a unit on the
$\ell$-primary coefficient module, so averaging makes
this group zero.  For $\ell=2$, cyclic cohomology expresses
(G31.4) as a subquotient of the indicated invariant group
and annihilates it by $2^s$.  The invariant group is a
subgroup of the rank-one group $\mathbb Q_2/\mathbb Z_2$.
Every bounded-exponent subquotient of such a group is
cyclic; if killed by $2^s$ it has order at most $2^s$.
This proves (G31.1).  Finiteness descends, and summing the
length bounds for $1\leq e<t$, with the prime weight
$\log\ell$, gives (G31.2) for every finite prime set.
The universal choice of \texttt{LData} in that definition
causes no extra condition: its set of places is fixed,
its character pulls back to the specified cyclotomic
character, and changing a chosen place in the algebraic
closure conjugates the local restriction map.  These
changes preserve the strict kernels and their lengths.

For each $i$ apply its \texttt{VWinAt} at
$X_i=Z^{J_i}$.  Its \texttt{WinPair} input is exactly
the common pair at $X_i^{1/J_i}=Z$.
It supplies a number $Y_i$ with
\[
 Z\leq Y_i\leq Z^{J_i/2}\leq Z^{J_*/2}
\]
and the \texttt{VAt} bound at $x=Y_i$.
If a stated $J_i<2$ occurs, these inequalities exclude
such a pair for $Z>1$; there is no need to assume its
existence in that vacuous case.
The supplied bound is
$\exp(\tfrac12Y_i^\epsilon)$ for weights up to
$\lfloor Y_i^{401}\rfloor$, which includes $t_Z$.

By the compositum description (G30.6), increasing the page
from $a(Z)$ to $a(Y_i)$ extends only the real two-power
cyclotomic factor.  Thus
$L_i(Y_i)/L_i(Z)$ is cyclic of degree $2^{s_i}$, with
\[
 s_i\leq a(Y_i)-a(Z)
 \leq(2+8\epsilon)(J_*/2-1)\log_2Z+1.
\]
Apply (G31.2) and use
$Y_i^\epsilon\leq Z^{\epsilon J_*/2}$ to obtain
(G31.3).  No intersection of the different supplied
windows has been assumed.
\end{proof}

\begin{remark}[The remaining quantitative condition]
The exact source definitions are \texttt{VVle},
\texttt{IsChiField}, \texttt{WinPair}, \texttt{PageN},
\texttt{VAt} and \texttt{VWinAt} in
\texttt{CharacterBoxes/ExistsPageN.lean}.
They use $t=\lfloor x^{401}\rfloor$ and budget
$\sqrt{\exp(x^\epsilon)}$.
In (G31.3), the additional restriction-kernel term is
polynomial in $Z$ and $\log Z$ and would be negligible
relative to that budget.  The loss in the exponent of the
first term is the actual unresolved issue.

Replacing $\epsilon$ by a smaller value does not by itself
resolve it: the supplied $J_i$ depends on that value,
and the definition gives no upper control on
$\epsilon J_i(\epsilon)$.  There is also a second
constraint on naive reparameterization.  Changing
$Z$ to $Z^{\epsilon/\epsilon'}$ preserves the scale
$Z^{\epsilon/5}$ that bounds $\log q$, but it multiplies
the central scale for $\log d$ by
$(\epsilon/\epsilon')^2$, because \texttt{WinPair}
also contains $(\log Z)^2$.
If that multiplier exceeds two, the two prescribed
factor-two intervals for $\log d$ are disjoint.
Thus neither independent windows nor a change of epsilon
has been promoted to a simultaneous supply theorem.
Finite character decompositions and their integral errors
could still be used after such a theorem is proved.
\end{remark}

\section{A smaller host with one literal cyclic character}
\label{sec:cyclic-cm-host}

The noncyclicity in AC-G30 belongs to the chosen signed
Jacobian host.  It is not an obstruction to every bare Host
containing the cubic coefficient field.  The following smaller
construction removes that particular field-definition gate.
It retains an explicit integral CM-type gate.

\begin{proposition}[AC-G32: a cubic-compatible cyclic Host]
\label{AC-G32}
Let $h\geq64$ and put
\[
 R_h=\mathbb Q(\zeta_{2^{h+2}})^+,
 \qquad K=\mathbb Q(\omega),\qquad C_h=R_hK,
 \qquad n=2^{h+1},\qquad Q_S=3\cdot2^{h+2}.
                                                               \tag{G32.1}
\]
Use the frame
$A=\operatorname{Gal}(C_h/\mathbb Q)\simeq C_{2^h}\times C_2$,
projecting to its first factor and with distinguished generator
induced by $5$.  Then $n_0=2^h$ and $o=2$.
Take any source \texttt{Ranks} for this frame.  For example
$k=2n$, $m=2$, $t=k+7$ are permissible.
Let $B\geq269$, and choose distinct primes
$q,d>\max(B,4t(k+m-1)+1,11)$ satisfying
\[
 \begin{gathered}
 q\equiv3\pmod8,\quad d\equiv7\pmod {16},\quad
 d\equiv-1\pmod q,\\
 d\equiv-1\pmod\ell
       \quad\text{for every odd prime }\ell\leq B.
 \end{gathered}
                                                               \tag{G32.2}
\]
Then
\[
 E_h=C_h\mathbb Q(\sqrt{-q},\sqrt{-d})
                                                               \tag{G32.3}
\]
admits every field of the bare source \texttt{Host}, with
$\mathfrak c_K=\mathcal O_{\mathbb Q(\sqrt{-d})}$.
Its exact numerical data include
\[
 \begin{split}
 [E_h:\mathbb Q]&=2^{h+3}=4n,\\
 \operatorname{rd}(E_h)
   &=2^{h+1-2^{-h}}\sqrt{3qd},\qquad
 \mu(E_h)=\mu_6.
 \end{split}                                                   \tag{G32.4}
\]
The source half-different gives the same principal identity
\[
 AA=\mathfrak b^{-1},\quad \xi_0=-1/\sqrt{-d},\qquad
        (\xi_0)AA\overline{AA}\mathfrak D_{E_h/\mathbb Q}=1.
                                                               \tag{G32.5}
\]

Let $\psi_{-3}$ be the nontrivial Dirichlet character of
conductor three.  This field satisfies the literal definition
\[
             \texttt{IsChiField}(\psi_{-3},h,q,d,E_h).
                                                               \tag{G32.6}
\]
Consequently any separately supplied instance
$\texttt{VAt}(\psi_{-3},\epsilon,401,q,d,h,x)$ applies
directly to the actual strict \emph{untwisted} groups over
$E_h$, with the exact cutoff $\lfloor x^{401}\rfloor$
and budget $\sqrt{\exp(x^\epsilon)}$.
No simultaneous-character window or quadratic descent is
needed for this particular application.  The proposition
neither proves \texttt{VWin} from A nor substitutes a
cubic Artin coefficient into its pure-Tate definition.
\end{proposition}
\begin{proof}
\proofstep{1. The real cyclotomic factor and the frame.}
Write $N=2^{h+1}$ and $z=\zeta_{2^{h+2}}$.
The integral basis $1,z,\ldots,z^{N-1}$ has minimal
polynomial $X^N+1$.  Comparing its coefficients under
conjugation shows
\[
 \mathcal O_{R_h}=\mathbb Z[z+z^{-1}].
\]
Indeed the conjugation-fixed integral basis is $1$ and
$z^j+z^{-j}$ for $1\leq j<N/2$; these are monic integral
polynomials of successive degrees in $z+z^{-1}$.
Also $\mathcal O_{\mathbb Q(z)}$ is free on $1,z$ over
this real ring, by the relation
$z^2-(z+z^{-1})z+1=0$.
Its relative discriminant is generated by
$(z-z^{-1})^2$.  Its absolute norm from $R_h$ is four:
$z^2$ is a primitive $2^{h+1}$-st root and
$|N_{\mathbb Q(z)/\mathbb Q}(1-z^2)|=\Phi_{2^{h+1}}(1)^2=4$.
The derivative of $X^N+1$ gives the full cyclotomic
absolute discriminant $2^{(h+1)2^{h+1}}$.
Transitivity therefore gives
\[
 |\operatorname{disc}R_h|=2^{(h+1)2^h-1},
 \qquad \operatorname{rd}(R_h)=2^{h+1-2^{-h}}.
                                                               \tag{G32.7}
\]
The integral cyclotomic basis used here is the standard
one in Milne, \emph{Algebraic Number Theory},
Theorem~6.4 and Remark~6.6, printed pp.~99--101.
The relative calculation above proves the precise
maximal-real exponent needed here.

The groups of $R_h$ and $K$ are disjoint by ramification.
The image of $5$ generates the cyclic real two-power
quotient, with order $2^h$, exactly as in AM-H17.
Its chosen lift acts trivially on $K$, so the kernel
condition on its cyclic subgroup holds.  This proves
the frame data, and $o^{64}=2^{64}\leq2^h$.
The displayed example of ranks satisfies all inequalities
in \texttt{Params.Ranks}.

\proofstep{2. Field coordinates and the local clauses.}
The four factors $R_h,K,K_q,K_d$ have disjoint conductor
supports $2,3,q,d$.  They are linearly disjoint, and their
integral bases tensor because the discriminants are
coprime.  Thus the degree and root discriminant in
(G32.4) follow from (G32.7) and the quadratic discriminants
$-3,-q,-d$.  These quadratic discriminants have the stated
form because $q,d\equiv3\pmod4$.

Let $c_C$ be conjugation on $C_h$; it is trivial on $R_h$
and nontrivial on $K$.  In physical coordinates use
\[
 (a,\epsilon_q,\epsilon_d)\longmapsto
 \bigl((\epsilon_q,\epsilon_d),
        a c_C^{(1-\epsilon_q)/2}\bigr).              \tag{G32.8}
\]
It takes full conjugation to $((-1,-1),1)$ and pure
$d$-inertia to $((1,-1),1)$.  Since $c_C$ projects to
one in the real two-power factor, it preserves the exact
\texttt{art\_h} condition.  Cyclotomic restriction from
conductor $Q_Sqd$ supplies the Artin map, its surjectivity,
the two quadratic coordinate characters, and the same
inertia, decomposition and Frobenius dictionary as AM-H17.
This dictionary includes the ramified primes: one removes
the prime-power part of the conductor before imposing the
congruence describing inertia, decomposition or Frobenius.
The exact declarations are \texttt{HostRaw} and
\texttt{HostDer} in
\texttt{ConjEMemHostFIff.lean}.

The conductor is exact at three by its nontrivial quadratic
factor.  At two the residue $1+2^{h+1}$ modulo
$2^{h+2}$ is neither $1$ nor $-1$ and acts nontrivially
on $R_h$; extending it by one at the other factors proves
the literal \texttt{QS\_exact} condition.  The clauses
$2^{h+2}\mid Q_S$, $(Q_S,qd)=1$ and $Q_S\leq n^2$
hold, and the only odd prime in $Q_S$ is three, with
$4\cdot3\leq B$.

At $2,3,q$, the decomposition group fixes $K_d$.
For two this follows from $d\equiv7\pmod8$; for three
and $q$ it follows from (G32.2).  Hence these groups
preserve the source CM type and exclude conjugation.
Quadratic reciprocity gives $(-q/d)=-1$:
$(d/q)=(-1/q)=-1$, hence $(q/d)=1$, and $(-1/d)=-1$.
If conjugation lay in decomposition at $d$, its restriction
to $K_q$ would force the Frobenius power to be odd.
On the real subfield of $\mathbb Q(\zeta_{16})$ this
odd power acts by the class of $7$ modulo $\{\pm1\}$,
which is nontrivial; conjugation acts trivially there.
This contradiction excludes it at $d$ as well.
Thus $E_h/E_h^+$ is unramified at every finite place,
and all the \texttt{decomp\_typeI}, \texttt{decompConj},
\texttt{unramEF}, \texttt{eta}, \texttt{ordinary} and
\texttt{splitK} clauses follow by the identical group
interpretations used in AM-H17.  The type in these
clauses is still the source type induced from $K_d$.

The displayed discriminant proves all its support clauses.
The required lower root-discriminant bound follows by
dropping $\sqrt{3q}$ from (G32.4); its upper bound follows
from
$2^{-2^{-h}}\sqrt3\leq3$.
Finally $\operatorname{rd}(E_h)^2\geq4n$, since already
$2h+2-2^{1-h}\geq h+3$ for $h\geq2$.
These check the three inequalities in \texttt{HostDer}.

\proofstep{3. The integral data and the roots of unity.}
The relative different of $E_h/K_d$ is supported at
$2,3,q$, all split in $K_d$.  At each such prime choose
one of its two primes in $K_d$ and put in $\mathfrak b$
all factors of the relative different above that choice,
with their full exponents.  Conjugation exchanges the
supports.  Thus
\[
 \mathfrak b\overline{\mathfrak b}=\mathfrak D_{E_h/K_d},
 \quad(\mathfrak b,\overline{\mathfrak b})=1,
 \quad \sigma\mathfrak b=\mathfrak b
                  \quad(\sigma\in\operatorname{Gal}(E_h/K_d)).
\]
The trace-dual ideal is exactly $\mathfrak b^{-1}$,
which verifies the source unimodularity clause for every
ideal with this factorization.  Transitivity of differents
and $\mathfrak D_{K_d/\mathbb Q}=(\sqrt{-d})$ give
(G32.5).

The unique quadratic subfield of $R_h$ is $\mathbb Q(\sqrt2)$.
Hence all quadratic subfields of $E_h$ have squareclasses
in the span of $2,-3,-q,-d$.  The squareclass $-1$ is
not in that span: valuations at $3,q,d$ first force their
three exponents to be zero, and neither $1$ nor $2$ is
$-1$ modulo rational squares.  In particular $i\notin E_h$.
A higher two-power root of unity is therefore impossible.
At three the exact ramification index is two, excluding
$\mu_9$; at $q$ and $d$ it is two, excluding their
primitive roots as $q,d>3$.  Every other odd prime is
unramified.  This proves $\mu(E_h)=\mu_6$.
The norm of each such root from $K$ to $\mathbb Q$ is
one, and its norm from $E_h$ to $K_d$ is a power of that
norm.  This verifies \texttt{norm\_mu}.

Take $\mathfrak c_K=\mathcal O_{K_d}$ and four distinct
level primes congruent to one modulo $15Q_Sqd$.
The same choices as in AM-H17 verify primality,
coprimality, membership and the two special-prime
congruences.  The source embedding is chosen positive
on $\sqrt{-d}/i$.  The bare choices
$TT=T_2=\varnothing$, $u_0=0$, $L_w=2$, $C_{T_2}=0$
verify the remaining stored fields and their vacuous
range conditions.  No later size lower bound on these
sets is asserted.  Thus all fields of \texttt{Host}
have been supplied.  As in AC-G30, the separately defined
\texttt{CoprimeFifteen} property fails because $3\mid Q_S$.

For completeness, primes satisfying (G32.2) exist without
any window assertion: choose $q>B$ in its residue class,
then choose $d$ in the coprime CRT class given there,
using Dirichlet's theorem.  The choice in the source
\texttt{WinPair} has these same residue conditions for
all odd primes below its stipulated splitting bound;
its much stronger size and window assertions are separate.

\proofstep{4. The literal cyclic field.}
Embed $E_h$ in the cyclotomic field of conductor $Q_Sqd$.
An automorphism given by a unit residue $a$ fixes it
exactly when
\[
 a\equiv\pm1\pmod {2^{h+2}},\quad
 \psi_{-3}(a)=1,\quad (a/q)=(a/d)=1.
\]
These are precisely the four conditions in
\texttt{ExistsPageN.IsChiField}, including the ambient
conductor divisibilities.  This proves (G32.6).
The assertion about \texttt{VAt} is its direct definition
applied to this field, not a new theorem supplying its
hypothesis.  Its \texttt{LData} remains the strict Tate
coefficient of \texttt{GaloisCohomology/Statement.lean}.
\end{proof}

\begin{proposition}[AC-G32-signature: the exact fixed-lattice CM-type gate]
\label{AC-G32-signature}
The maximal real subfield of (G32.3) is
\[
 E_h^+=R_h\mathbb Q(\sqrt{3q},\sqrt{3d}).             \tag{G32.9}
\]
There is a purely imaginary trace coefficient $\beta$
with $(\beta)=(\xi_0)$ making the unchanged lattice
$AA$ in (G32.5) positive for the type induced from $K$
if and only if there is a unit
$\kappa\in\mathcal O_{E_h^+}^{\times}$ satisfying
\[
 \operatorname{sign}\tau(\kappa)
       =\operatorname{sign}\tau(\sqrt{3d})
             \quad\text{for every real embedding }\tau.
                                                               \tag{G32.10}
\]
Here the imaginary square roots and the positive root of
$3d$ are chosen compatibly.  Such a unit cannot lie in
the quadratic subfield $\mathbb Q(\sqrt{3d})$.
This last obstruction alone makes no assertion about
existence in the larger field (G32.9).
\end{proposition}
\begin{proof}
Complex conjugation changes the signs of all three
imaginary quadratic roots and fixes $R_h$.  Its fixed
field is generated by their pairwise products, giving
(G32.9).

If $\beta$ has the stated ideal and is purely imaginary,
then $\kappa=\xi_0/\beta$ is a real integral unit;
conversely every real unit gives such a $\beta$.
Put $\beta_3=-1/\sqrt{-3}$, positive imaginary for the
chosen type on $K$.  Then
\[
 \frac{\beta}{\beta_3}
       =\frac{\sqrt{-3}}{\sqrt{-d}\,\kappa}
       =\frac{\sqrt{3d}}{d\,\kappa}.                 \tag{G32.11}
\]
The right-hand side is real.  It is positive at every
embedding exactly when (G32.10) holds.  This is precisely
the positivity condition for the induced type.
The equality of ideals leaves (G32.5) unchanged, so
when the condition holds the corresponding alternating
trace form is integral and principal.  This verifies
both directions of the stated fixed-lattice equivalence.

Now $d\equiv7\pmod {16}$ and $d\equiv-1\pmod3$.
The squarefree number $3d$ is $1$ modulo four, so a
unit of its real quadratic field is
$(a+b\sqrt{3d})/2$ with integers $a,b$ of matching parity.
Opposite signs at its two embeddings would force norm
$-1$, hence
\[
                         a^2-3db^2=-4.
\]
Reducing modulo three gives $a^2=2$, impossible in
$\mathbb F_3$.  This proves the quadratic obstruction.
The extension from that quadratic field to (G32.9) has
even degree; a norm from it squares this signature.
Thus the same argument cannot be used to exclude a
unit in the larger field.
\end{proof}

\paragraph{What is and is not repaired.}
The smaller construction supplies a cubic-compatible bare Host
and one exact cyclic-character field.  It removes the specific
multi-character field-definition obstruction of (G30.7) for
this choice of host.  It does not retain the inherited
$\mathbb Q(\zeta_7)$ Jacobian and its specified sign unit:
neither $\zeta_7$ nor $i$ lies in the new field.
The AM-H16 units live in explicitly larger real subfields;
weak approximation on signs does not replace them by units
in (G32.9).  AC-G33 below constructs the required unit in (G32.10)
with an explicit logarithmic height bound.  The equivalence
here isolates exactly the condition that construction meets.
The cross-block coefficient tie of AC-G29--G30 is
still needed for the cubic Artin character.  The source
\texttt{VVle} hypothesis obtained from a supplied
\texttt{VWin} continues to bound pure-Tate strict groups
over this varying field, not those mixed coefficients.

\paragraph{Algebraic-geometric dictionary.}
The finite cyclic character in (G32.6) corresponds to a finite
etale covering after the conductor primes are inverted, in the
Galois-category dictionary of SGA~1, Expos\'e~V, \S\S4 and 7,
already used in AC-K19.  This records the finite coefficient
and its field; the principal CM lattice and the later theta
comparison are additional structures.  No new archive cote
is asserted here.

\section{An explicit sign unit inside the smaller cyclic host}
\label{sec:cyclic-cm-unit}

The larger-field unit question isolated in AC-G32 has a
constructive answer.  Two conjugates of one cyclotomic norm
suffice.  The norm is taken before the unit is used, so the
field of the smaller Host is unchanged.

\begin{theorem}[AC-G33: the missing signature unit and elliptic principal model]
\label{AC-G33}
Let $d$ be as in AC-G32 and put $D=3d$ and
$R=\mathbb Q(\zeta_{16})^+$.  There is an explicitly
specified integral unit
\[
 \kappa_d\in\mathcal O_{R(\sqrt D)}^\times
              \subset\mathcal O_{E_h^+}^\times
\]
with
\[
 \operatorname{sign}\tau(\kappa_d)
     =\operatorname{sign}\tau(\sqrt D)
                 \quad\text{at every real embedding},        \tag{G33.1}
\]
and absolute logarithmic Weil height
\[
             h(\kappa_d)\leq2(d-1)\log2+\log105.               \tag{G33.2}
\]
It requires only a cyclotomic norm, one fixed conjugate,
and at most one copy of each of three fixed circular units,
with an optional minus sign.

Consequently $\beta=\xi_0/\kappa_d$ makes the unchanged
source lattice $AA$ of AC-G32 a principally polarized
CM lattice of type induced from $K=\mathbb Q(\omega)$.
With $g=[E_h:K]=2^{h+2}$, its unpolarized abelian variety is
\[
                       C_0^g,\qquad C_0:y^2=x^3+1.             \tag{G33.3}
\]
The induced integral $\mathcal O_{E_h}$-action and the
principal polarization can be realized on this power over
$K$.  The polarization need not be its product polarization.

For any fixed $A,a\geq0$, the source window bound
$\log d\leq Y_0^{\epsilon/5}(\log Y_0)^2$, with
$x\geq Y_0\geq e$ and fixed $\epsilon>0$, gives
\[
        x^A(1+h(\kappa_d))^a
                    =o\bigl(\exp(x^\epsilon/2)\bigr).
                                                               \tag{G33.4}
\]
This controls the displayed logarithmic height and fixed
polynomial expressions in it.  It does not identify that
quantity with an unlogged theta norm or prove the mixed-block
arithmetic tie.
\end{theorem}
\begin{proof}
\proofstep{1. A real unit from the quadratic cyclotomic norm.}
Put $t=\zeta_{16}$ and $z=\zeta_D$.
The primitive quadratic character
\[
        \chi_D(a)=(a/3)(a/d)
                 \quad(a\in(\mathbb Z/D\mathbb Z)^\times)
\]
is even and cuts out $\mathbb Q(\sqrt D)$.
Indeed it is the product of the characters of quadratic
discriminants $-3$ and $-d$, whose compositum has this real
quadratic subfield.  The fields of conductors $16$ and $D$
are linearly disjoint.  Define
\[
 \begin{split}
 P&=\prod_{\chi_D(a)=1}(1-tz^a)
    =N_{\mathbb Q(t,z)/\mathbb Q(t,\sqrt D)}(1-tz),\\
 v&=t^{-(d-1)/2}P.
 \end{split}\tag{G33.5}
\]
The product has $d-1$ factors.  Each $tz^a$ is a primitive
$16D$-th root, and $1-tz^a$ is an integral unit because
$\Phi_{16D}(1)=1$.
For this last equality one may use
$\Phi_{16D}(X)=\Phi_{2D}(X^8)$ and
$\Phi_{2D}(1)=\Phi_D(-1)=\Phi_d((-1)^3)/\Phi_d(-1)=1$.
Thus both $P$ and $v$ are units.

The sum of the residues $a$ occurring in (G33.5) is zero
modulo $D$.  Modulo $d$ they run once through every nonzero
residue: its quadratic sign determines the unique nonzero
residue modulo three.  Modulo three there are $(d-1)/2$
of each of its two nonzero residues, since the two quadratic
signs modulo $d$ occur equally often.  Both sums are zero.
As $d-1$ is even, complex conjugation consequently gives
\[
 \overline P=(-t^{-1})^{d-1}P=t^{-(d-1)}P,
                       \qquad\overline v=v.          \tag{G33.6}
\]
The fixed real field of $\mathbb Q(t,\sqrt D)$ is
$R(\sqrt D)$, so $v$ lies in precisely this field.
The auxiliary root $z$ has disappeared through the norm;
no extension of the Host has been made.

\proofstep{2. Two conjugates give opposite signs in every quadratic pair.}
Let $\sigma_D$ fix $t$ and negate $\sqrt D$.
It exchanges the two character cosets in (G33.5), hence
\[
 n:=v\sigma_D(v)=t^{-(d-1)}\Phi_D(t).                 \tag{G33.7}
\]
The polynomial identity
\[
 \Phi_{3d}(X)
       =\frac{(X^{3d}-1)(X-1)}{(X^3-1)(X^d-1)}
\]
and $d\equiv7\pmod {16}$ show that, at the real embedding
of $R$ indexed by $b\in\{1,3,5,7\}$,
\[
 n_b=-\frac{\sin(5b\pi/16)\sin(b\pi/16)}
                  {\sin(3b\pi/16)\sin(7b\pi/16)}.
                                                               \tag{G33.8}
\]
For example, substituting $t^d=t^7$ and $t^{3d}=t^5$
in (G33.7) leaves the phase $t^{-8}=-1$;
this fixes its sign without a choice of square root.
None of the displayed sines vanishes, and their signs give
\[
 \begin{array}{c|rrrr}
 b&1&3&5&7\\\hline
 \operatorname{sign}(n_b)&-&+&+&-
 \end{array}.                                                   \tag{G33.9}
\]
Let $\tau_3$ send $t$ to $t^3$ and fix $\sqrt D$.
On the four real embeddings it permutes the indices by
multiplication by three modulo $\{\pm1\}$.
Inspection of (G33.9) shows
$n_b n_{3b}<0$ at all four embeddings.
Therefore the unit
\[
                             w=v\tau_3(v)                      \tag{G33.10}
\]
has totally negative norm from $R(\sqrt D)$ to $R$.
Its two real values above every embedding of $R$ have
opposite signs.

\proofstep{3. Aligning those four pairs by fixed circular units.}
For $a=3,5,7$, put
\[
 c_a=t^{(1-a)/2}\frac{1-t^a}{1-t}\in\mathcal O_R^\times.
                                                               \tag{G33.11}
\]
The quotient is an integral geometric sum.  Its inverse
is integral as well: choose $b$ with $ab\equiv1\pmod {16}$
and use $(1-(t^a)^b)/(1-t^a)$, adjusting by a root of unity.
Conjugation fixes the expression, so it is a real unit.
Its values at the four real embeddings are the sine ratios
$\sin(ab\pi/16)/\sin(b\pi/16)$, with signs
\[
 \begin{array}{c|rrrr}
 &1&3&5&7\\\hline
 -1&-&-&-&-\\
 c_3&+&+&+&-\\
 c_5&+&+&-&+\\
 c_7&+&-&+&-
 \end{array}.                                                   \tag{G33.12}
\]
These four sign vectors form a basis over $\mathbb F_2$.
This is the same explicit four-dimensional calculation as
AM-H16, (H16.7)--(H16.9); no general unit-signature theorem
is needed.

Above each $b$, take the extension sending $\sqrt D$ to
its positive real value, and let $a_b\in\{0,1\}$ record
whether $w$ is negative there.  Reduce the following
exponents modulo two and set
\[
 w_0=(-1)^{a_1}c_7^{a_1+a_3}c_5^{a_1+a_5}c_3^{a_3+a_7},
                  \qquad \kappa_d=w/w_0.             \tag{G33.13}
\]
The table makes $w_0$ and $w$ have the same sign at
these four selected extensions.  On the other member of
each pair $w_0$ is unchanged and $w$ reverses sign by
(G33.10).  This proves (G33.1).
All sign tests involve nonzero real algebraic numbers,
so exact real-root isolation specifies this construction
without a numerical choice.

There are $d-1$ factors of height at most $\log2$ in
(G33.5); roots of unity have height zero.  Thus
$h(v)\leq(d-1)\log2$ and
$h(w)\leq2(d-1)\log2$.
Each $c_a$ is a sum of $a$ roots of unity, giving
$h(c_a)\leq\log a$.  Formula (G33.13) uses at most one
of each, so $h(w_0)\leq\log105$.
Subadditivity of height under multiplication and
$h(w_0^{-1})=h(w_0)$ prove (G33.2).

\proofstep{4. The unchanged principal lattice and the elliptic power.}
The field $R(\sqrt D)$ is contained in (G32.9).
Apply AC-G32-signature to the constructed unit.
For every embedding $\phi$ in the induced $K$-type,
write $\phi(\beta)=ib_\phi$ with $b_\phi>0$.
The alternating form
\[
                     (x,y)\longmapsto
                  \operatorname{Tr}_{E_h/\mathbb Q}(\beta x\bar y)
\]
is integral and unimodular by (G32.5), and its value on
$(x,ix)$ is $2\sum_\phi b_\phi|x_\phi|^2>0$.
These are the Riemann relations, giving the principal
polarization on $\mathbb C^{\Phi}/\Phi(AA)$.
Multiplication by $\mathcal O_{E_h}$ preserves $AA$;
the trace identity shows that its Rosati involution is
complex conjugation.  This is exactly the CM lattice
argument of AM-H17, step~6, with its primary locators
Milne, \emph{Complex Multiplication}, Example~2.9,
pp.~25--26, and Proposition~3.17, p.~31.

The Eisenstein ring $\mathcal O_K=\mathbb Z[\omega]$
is norm-Euclidean: rounding first the imaginary coordinate
and then the real coordinate leaves squared absolute value
at most $3/16+1/4<1$.  Consequently the torsion-free
finite module $AA$ is free over $\mathcal O_K$, of rank
$g=[E_h:K]$.
An $\mathcal O_K$-basis identifies the complex lattice of
its induced type with $\mathcal O_K^g$, up to the common
complex scaling that identifies $\mathbb C/\mathcal O_K$
with $C_0$.  This proves (G33.3).
After transporting the polarization and the action to
$C_0^g$, their homomorphisms have entries in
$\operatorname{End}_{\mathbb C}(C_0)=\mathcal O_K$.
Every such endomorphism is defined over $K$, generated by
integer multiplications and $(x,y)\mapsto(\omega x,y)$.
Thus the transported action and polarization morphism are
defined over $K$.  Positivity and principality can be checked
after embedding in $\mathbb C$, so this is the asserted
principal abelian model.  This reasoning does not identify
its polarization with the product one.

\proofstep{5. The quantitative statement has logarithmic normalization.}
For $x\geq Y_0\geq e$, the given upper bound on $\log d$
is at most $x^{\epsilon/5}(\log x)^2$.
For fixed $A,a$, (G33.2) therefore gives
\[
 \log\bigl(x^A(1+h(\kappa_d))^a\bigr)
     \leq A\log x+a x^{\epsilon/5}(\log x)^2+O_a(1)
                         =o(x^\epsilon).
\]
Subtracting $x^\epsilon/2$ proves (G33.4).
The assertion is about this explicitly normalized height,
not its exponential or a separately defined theta quantity.
\end{proof}

\paragraph{Source scope after this construction.}
The signature-unit existence gate for AC-G32 is now closed.
The proof gives a principal elliptic CM model on the exact
source lattice, with a controlled logarithmic sign-unit height,
inside a host that is one literal \texttt{IsChiField}.
It still changes the CM type from the source type induced
from $K_d$ to the one induced from $K$.
The source ordinary clause remains a statement about the
former type.  For instance, at eleven the latter elliptic
curve has trace zero: cubing permutes $\mathbb F_{11}$,
so its quadratic-character point-count sum is zero.
The ordinary and theta identities for the new type therefore
require their actual local hypotheses and comparison proofs.
The cross-block character tie, the main cubic-twisted
cohomology supply, and its exact theta normalization are
not consequences of the sign unit or of (G33.4).

\paragraph{Algebraic-geometric dictionary.}
The principal lattice statement uses the audited CM lattice
and Riemann-form dictionary in Milne,
\emph{Complex Multiplication}, Example~2.9, pp.~25--26,
and Proposition~3.17, p.~31, also used in AM-H17.
The finite character field remains the finite etale-cover
object of SGA~1, Expos\'e~V, \S\S4 and 7, after inversion
of its conductor primes.  The explicit unit supplies the
additional integral polarization data; it does not identify
a finite-cover coefficient with a theta or cohomology
comparison.  These are the already audited locators;
no new Grothendieck archive cote is asserted.

\begin{corollary}[AC-G33-twist: two principal scalar cubic packets]
\label{AC-G33-twist}
The principal model $A_0=C_0^g$ of AC-G33 has a scalar
$\mu_3$-action preserving its polarization and commuting
with its integral $\mathcal O_{E_h}$-action.
Put $p=269$, $C_p:y^2=x^3+p^2$ and $A_p=C_p^g$.
There is a principal polarization on $A_p$, defined over
$K$, and the same integral $\mathcal O_{E_h}$-action,
with the following properties.  If $\alpha^3=p$ and
$\chi(\sigma)=\sigma(\alpha)/\alpha$, then the product
isomorphism
\[
 \phi_g:A_0\longrightarrow A_p,
       \qquad (x,y)\longmapsto(\alpha^2x,py)
                  \quad\text{on each factor}                 \tag{G33.14}
\]
is $\mathcal O_{E_h}$-linear and polarized over $K(\alpha)$.
Its Galois conjugates satisfy
\[
                  \sigma(\phi_g)=\phi_g\,[\chi(\sigma)^2].
                                                               \tag{G33.15}
\]
Consequently every covariant Tate module is twisted by
$\chi^2$ integrally, including at coefficient prime three.
Its reverse $\mathcal O_{E_h}$-linear Hom is a rank-one
integral coefficient with character $\chi$.
The two principal varieties have identical geometric CM
data and stable Faltings heights.  This construction invokes
no Fourier projector and no full regular tensor.
\end{corollary}
\begin{proof}
Identify $A_0$ with $C_0^g$ as in AC-G33, and use the
product principal polarization only to identify its dual.
The constructed principal polarization is represented by
an integral positive Hermitian matrix
$H\in M_g(\mathcal O_K)$.
Every scalar $u\in\mu_3$ satisfies $\bar u u=1$, so
$\bar u H u=H$.  This proves the asserted polarization
invariance; the scalar action commutes with all matrices
representing $\mathcal O_{E_h}$.

On $C_p$, the same Eisenstein endomorphisms are defined over
$K$, and (G33.14) commutes with them.  An elliptic group
isomorphism preserves the principal polarization defined by
the origin.  Therefore the same matrix $H$ on $C_p^g$
pulls back under $\phi_g$ to the polarization on $C_0^g$.
It is defined over $K$; positivity and degree one are
checked after this isomorphism.  The same integral action
matrices also commute with $\phi_g$, proving the statement
about $\mathcal O_{E_h}$.  Taking Galois conjugates of the
coordinate formula gives (G33.15) directly.

For a covariant Tate vector $v$, formula (G33.15) reads
$\sigma(\phi_g(v))=\phi_g(\chi(\sigma)^2\sigma v)$.
This proves the integral twist, without division by three.
The original Tate lattice is rank one projective over
$\mathcal O_{E_h}\otimes\mathbb Z_\ell$: the characteristic-zero
integral comparison identifies it with the invertible
$\mathcal O_{E_h}$-ideal $AA$ tensored with $\mathbb Z_\ell$.
The latter is locally free of rank one over this product
of local rings.  Thus its endomorphism coefficient is
$\mathcal O_{E_h}\otimes\mathbb Z_\ell$; reversing the Hom
replaces $\chi^2$ by $\chi$.  This also fixes the variance
at $\ell=3$.
Finally the varieties are geometrically isomorphic, which
preserves their CM data and stable Faltings heights.
This is the concrete scalar instance of AC-G24 on the
principal model just constructed.
\end{proof}

The corollary supplies separate principal scalar packets.
A common twist within the one distinguished source Speh
block still cancels from its pairwise Hom coefficients,
as proved in AC-G29.  To use the nontrivial Hom between
these two packets in the source bound requires the actual
cross-block theta and cohomology comparison.  The polarized
isomorphism above does not assert that comparison or a
comparison of ramified integral theta lattices.

\section{The finite integral comparison error on growing hosts}
\label{sec:finite-gluing-budget}

The integral Fourier comparison in ALP\_MACHINE, AM-H06,
leaves finite-coefficient cohomology terms.  These terms admit
a uniform elementary bound.  Their control does not supply
the main infinite-coefficient strict Selmer lengths.

Write $n_F=[F:\mathbb Q]$, $D_F=|\operatorname{disc}F|$,
and $s_F=|S_{\mathrm{fin}}|$.  Throughout this section $F$
is totally imaginary and $S$ contains the archimedean places
and all places over the coefficient prime $\ell$.
All logarithms without a subscript are natural.

\begin{theorem}[AC-K19: finite gluing errors fit the growing-host budget]
\label{AC-K19}
Suppose first that $\mu_\ell\subset F$.  Define
\[
 \begin{split}
 d_i(F,S)&=\dim_{\mathbb F_\ell}H^i(G_{F,S},\mathbb F_\ell)
                  \quad(i=1,2),\\
 d_{\mathrm{loc}}(F,S)&=
 \sum_{v\mid\ell}\dim_{\mathbb F_\ell}H^1(F_v,\mathbb F_\ell)
 +\sum_{v\in S_{\mathrm{fin}},\ v\nmid\ell}
                \dim_{\mathbb F_\ell}H^1(I_v,\mathbb F_\ell).
 \end{split}
\]
Then these dimensions are finite and
\[
 d_1+d_2+d_{\mathrm{loc}}
 \leq B_\ell(F,S):=
 2\log_\ell D_F+(3+2\log_\ell8)n_F+5s_F.       \tag{FG.1}
\]

More generally, let $M$ be a finite $G_{F,S}$-module,
killed by $\ell^a$ and generated by at most $r$ elements
as an abelian group.  Let $L/F$ be a finite Galois extension
unramified outside $S$, trivializing $M$ and containing
$\mu_{\ell^a}$.  Put $D=[L:F]$, and let $S_L$ be the
places above $S$.  The quantity
\[
 \begin{split}
 \mathcal C_{F,S}(M)={}&\log|H^1(G_{F,S},M)|
      +\log|H^2(G_{F,S},M)|\\
 &+\sum_{v\mid\ell}\log|H^1(F_v,M)|
      +\sum_{v\in S_{\mathrm{fin}},\ v\nmid\ell}
                         \log|H^1(I_v,M)|
 \end{split}
\]
is finite and satisfies
\[
 \begin{split}
 \mathcal C_{F,S}(M)\leq{}&2r\log D_L+2rn_L\log8\\
 &+ar\log\ell\,[2(D+n_L+s_L)+n_F+3s_F].
 \end{split}
 \tag{FG.2}
\]
In particular, a fixed finite Artin coefficient, with an
arbitrary integral Tate twist, has
$\mathcal C=O(\log D_F+n_F+s_F)$ when trivialized after
compositum with one fixed number field; the constant is
independent of the twist.

For the growing signed hosts of AM-H17, adjoin the cubic
coefficient field and write
\[
 F=\mathbb Q(\zeta_{2^{h+2}},\zeta_7,
               \sqrt{-q},\sqrt{-d},\omega),
 \qquad F^{\mathrm c}=F(\sqrt[3]{269}).          \tag{FG.3}
\]
Here $q,d>7$ are the distinct primes in the source's
$\operatorname{WinPair}(\varepsilon,Y_0)$, $x\geq Y_0\geq\exp(1)$,
and $h$ satisfies its $\operatorname{PageN}$ condition
\[
 x^{2+8\varepsilon}\leq2^h\leq2x^{2+8\varepsilon},
 \qquad 0<\varepsilon<1/52.
\]
Take $S_{\mathrm{fin}}$ to be the primes of $F$ over
$3\cdot269$ and $t=\lfloor x^{401}\rfloor$.
For a positive weight $j$, let $a_{F^{\mathrm c}}(j)$
and $b_F(j)$ be the two sides of AM-H06: the untwisted
strict $\mathbb Z_3$-length over $F^{\mathrm c}$, and the
sum of the three character lengths over $F$, respectively.
Then
\[
 \begin{split}
 \sum_{1\leq j<t}|2a_{F^{\mathrm c}}(j)-b_F(j)|\log3
 &\leq3(t-1)B_3(F,S)\log3\\
 &=O_\varepsilon\!\left(
 x^{403+41\varepsilon/5}(\log x)^2\right)\\
 &=o\!\left(e^{x^\varepsilon/2}\right).
 \end{split}
 \tag{FG.4}
\]
Thus this finite integral comparison error consumes less
than any prescribed positive fraction of the source budget
for sufficiently large $x$.  Absorbing it into an exact
budget requires that positive relative slack in the
main arithmetic estimate.  No such main estimate is
asserted by this theorem.
\end{theorem}

\begin{proof}
\proofstep{1. A coarse class-number bound.}
Let $H_F$ be the ordinary ideal class number.
Minkowski's bound supplies an integral ideal in every
class of norm at most
\[
 \frac{n_F!}{n_F^{n_F}}
       \left(\frac4\pi\right)^{r_2(F)}\sqrt{D_F}
       \leq2^{n_F}\sqrt{D_F}=:U_F.
\]
For the source statement, see Milne,
\emph{Algebraic Number Theory}, v3.08 (2020),
Theorem~4.3, p.~70.
The Euler factors give $\zeta_F(2)\leq\zeta(2)^{n_F}<2^{n_F}$:
at each rational prime there are at most $n_F$ factors,
and every residue degree is at least one.
Counting all ideals of norm at most $U_F$ therefore gives
\[
 H_F\leq U_F^2\sum_{\mathfrak a\ne0}(N\mathfrak a)^{-2}
             \leq8^{n_F}D_F.                         \tag{FG.5}
\]
This estimate uses no hypothesis about zeros.

\proofstep{2. Constant finite coefficients.}
Choose an identification $\mathbb F_\ell\simeq\mu_\ell$
over $F$.  Kummer theory gives
\[
 0\longrightarrow\mathcal O_{F,S}^{\times}/
                    (\mathcal O_{F,S}^{\times})^\ell
 \longrightarrow H^1(G_{F,S},\mu_\ell)
 \longrightarrow\operatorname{Cl}(\mathcal O_{F,S})[\ell]
 \longrightarrow0.                                  \tag{FG.6}
\]
Indeed, outside $S$ the residue characteristic differs
from $\ell$, and a Kummer class is unramified exactly
when its valuation is a multiple of $\ell$.
Dividing its divisor outside $S$ by $\ell$ gives the
right-hand map.  A trivial divisor class leaves an
$S$-unit, and any $\ell$-torsion ideal class lifts by
choosing a generator of its $\ell$th power.  This proves
the displayed exact form from field Kummer theory.

The $S$-unit rank is $r_1(F)+r_2(F)-1+s_F$; its torsion
subgroup is cyclic.  Hence its quotient by $\ell$th
powers has dimension at most $n_F+s_F$.
Here the known unit theorem is Milne, loc. cit.,
Theorem~5.11, p.~90.
The $S$-class group is a quotient of the ordinary class
group, so
\[
             d_1\leq n_F+s_F+\log_\ell H_F.           \tag{FG.7}
\]

The last part of the finite Poitou--Tate sequence is
\[
 H^1(G_{F,S},\mu_\ell)^\vee
 \longrightarrow H^2(G_{F,S},\mathbb F_\ell)
 \longrightarrow\bigoplus_{v\in S}H^2(F_v,\mathbb F_\ell).
\]
Finite local duality bounds the dimension of each finite-place
term on the right by one; complex places contribute zero.
Consequently $d_2\leq d_1+s_F$.
The exact sources are Sharifi, \emph{Iwasawa Theory},
Appendix~A, Theorems~A.0.3 and A.0.8; these are the
finite-module statements, with $S$ containing all places
over $\ell$ and infinity.

Local duality and the local Euler characteristic give,
for $v\mid\ell$,
\[
      \dim H^1(F_v,\mathbb F_\ell)=[F_v:\mathbb Q_\ell]+2.
\]
The Euler-characteristic source is Milne,
\emph{Arithmetic Duality Theorems}, second edition (2006),
I, Theorem~2.8, p.~31.
If $v\nmid\ell$, wild inertia is pro-$p$ with $p\ne\ell$
and tame inertia has maximal pro-$\ell$ quotient
$\mathbb Z_\ell(1)$.  Thus
$\dim H^1(I_v,\mathbb F_\ell)=1$.
Summing the local degrees yields
\[
 d_{\mathrm{loc}}\leq n_F+2s_F.
\]
Combining this with (FG.7), $d_2\leq d_1+s_F$, and
(FG.5) proves (FG.1).

\proofstep{3. Arbitrary bounded finite coefficients.}
Set $M^D=\operatorname{Hom}(M,\mu_{\ell^a})$.
It has the same elementary divisors as $M$, is trivial
over $L$, and is the finite Tate dual used by duality.
Inflation--restriction and the elementary bound on
one-cocycles of the finite group $\operatorname{Gal}(L/F)$
give
\[
 |H^1(G_{F,S},M)|
 \leq |M|^D\,|H^1(G_{L,S_L},M)|.
\]
Over $L$, decompose $M$ into at most $r$ cyclic groups
of orders $\ell^{b_i}$, where $b_i\leq a$.
Since $\mu_{\ell^a}\subset L$, the same Kummer exact
sequence as (FG.6), now with exponent $\ell^{b_i}$,
bounds each factor by
$\ell^{b_i(n_L+s_L)}H_L$.  Therefore, for both $M$ and
$M^D$,
\[
 \log|H^1(G_{F,S},M)|
 \leq ar\log\ell\,(D+n_L+s_L)+r\log H_L.           \tag{FG.8}
\]
Finite Poitou--Tate and local duality give
\[
 \log|H^2(G_{F,S},M)|
 \leq\log|H^1(G_{F,S},M^D)|+s_F\log|M|.           \tag{FG.9}
\]
For $v\mid\ell$, the same local Euler characteristic
and duality imply
$|H^1(F_v,M)|\leq|M|^{[F_v:\mathbb Q_\ell]+2}$.
For $v\nmid\ell$, first taking wild-inertia invariants
and then the invariants of the prime-to-$\ell$ part of
tame inertia is exact on finite $\ell$-groups.
The remaining first cohomology is a quotient
$N/(\tau-1)N$ for a submodule $N$ of $M$ and a
generator of $\mathbb Z_\ell$.  It has order at most
$|M|$.  The two local sums in $\mathcal C$ are thus at
most $(n_F+2s_F)\log|M|$.
Substitute (FG.8)--(FG.9), $|M|\leq\ell^{ar}$ and
(FG.5) for $L$.  This is (FG.2).

For the fixed-field corollary, take a fixed finite
Galois field $L_0$ trivializing the Artin action and
containing $\mu_{\ell^a}$.  It trivializes all Tate
twists as well.  Put $L=FL_0$, and include the fixed
ramified primes in $S$.  If $C=[L_0:\mathbb Q]$, then
$D\leq C$, $n_L\leq Cn_F$ and $s_L\leq Cs_F$.
Moreover
\[
             D_L\leq D_F^C D_{L_0}^{n_F}.            \tag{FG.10}
\]
For a direct proof valid even when the fields intersect,
the finite etale algebra $F\otimes_{\mathbb Q}L_0$ is
a product of number fields, one of them $L$.
The tensor order $\mathcal O_F\otimes\mathcal O_{L_0}$
has discriminant $D_F^C D_{L_0}^{n_F}$ in absolute value,
by the determinant of the tensor trace matrix.
The product maximal order has discriminant the product
of the positive absolute field discriminants; this
divides the tensor-order discriminant by the square
of the index.  Each factor is at least one, proving
(FG.10).  Equation (FG.2) now has the stated uniform order.

\proofstep{4. The actual cubic growing hosts.}
The five fields in (FG.3) have pairwise coprime
conductors $2^{h+2},7,q,d,3$.
They are linearly disjoint: their Galois groups occur
in distinct factors of the cyclotomic Chinese-remainder
product, and the quadratic fields use the $q$ and $d$
factors.  Their discriminants are also pairwise coprime.
Milne, \emph{Algebraic Number Theory}, Lemma~6.5,
p.~100, and Remark~6.6(c), p.~101, identify the product
integral basis and its discriminant.  The cyclotomic
discriminant formula in that remark gives
\[
 \begin{split}
 n_F&=96\cdot2^h,\\
 \log D_F/n_F&=(h+1)\log2+\tfrac56\log7
                 +\tfrac12\log q+\tfrac12\log d
                 +\tfrac12\log3.
 \end{split}
 \tag{FG.11}
\]
The quadratic discriminants are $-q$ and $-d$, because
the source congruences give $q,d\equiv3\pmod4$.
There are at most $2n_F$ places in $S_{\mathrm{fin}}$.
The field $F$ is abelian over $\mathbb Q$, so it cannot
contain the nonabelian splitting field
$\mathbb Q(\omega,\sqrt[3]{269})$.
Thus $F^{\mathrm c}/F$ is indeed cyclic of degree three.

The proof of AM-H06 applies over any base containing
$\mu_3$; its local conditions explicitly retain inertia
at $269$.  With $c_3=3$, its equation (H06.3) gives
\[
 -3(d_2+d_{\mathrm{loc}})
 \leq2a_{F^{\mathrm c}}(j)-b_F(j)\leq3d_1.
\]
For positive weights the untwisted group is finite by
AC-K10, and AM-H06 then makes all the displayed lengths
finite.  Equation (FG.1) bounds the absolute difference
by $3B_3(F,S)$, proving the first line of (FG.4).

The exact source window inequalities are
\[
 \log q\leq Y_0^{\varepsilon/5},\qquad
 \log d\leq Y_0^{\varepsilon/5}(\log Y_0)^2.
\]
For $x\geq Y_0\geq e$ these imply the same upper bounds
with $x$ in place of $Y_0$; also $h=O_\varepsilon(\log x)$.
Using $n_F\leq192x^{2+8\varepsilon}$ in (FG.11) yields
\[
 B_3(F,S)=O_\varepsilon\!\left(
 x^{2+8\varepsilon+\varepsilon/5}(\log x)^2\right).
\]
There are fewer than $x^{401}$ weights.
This proves the polynomial estimate in (FG.4).
Finally, for every fixed $A$ and fixed $\varepsilon>0$,
$A\log x+2\log\log x-x^\varepsilon/2\to-\infty$.
The last assertion of (FG.4) and the statement about
arbitrary positive relative slack follow.
\end{proof}

\paragraph{Application to the full finite-character decomposition.}
The two defects of AC-G30-Fourier satisfy (FG.2) with
$(\ell,a,r)=(3,1,24)$ and $(2,3,42)$, respectively.
For the second count, the $21$ nonunit Smith factors
over the unramified quadratic coefficient ring each
have two generators over $\mathbb Z_2$.
The common trivializing field is
$L_0=\mathbb Q(\zeta_{168})$, of degree $48$.
On the base $B_h$ of AC-G30 include the fixed primes
$2,3,7$ in $S$, and also $269$ for the cubic comparison.
The degree and log discriminant of $B_h$ are bounded
by those of $F$, by the tower discriminant formula.
Thus (FG.2), (FG.10) and (FG.11) give the same
polynomial order in (FG.4) for the sum of these two
finite defects over all the weights.  The long exact
global and local coefficient sequences bound the
kernel and cokernel of strict localization exactly as
in AM-H06.  The normalized lengths times $\log\ell$
are at most the logarithms of the finite group orders,
so the bound also applies to those comparison errors.
This controls the full finite Fourier cost without
assuming trivial residual characters on $B_h$; it
does not synchronize the main character windows in
AC-G31.

\paragraph{Source interfaces and remaining gate.}
The exact inherited comparison is
\path{RETURN_ALP_MACHINE_checkpoint_04_20261121T2304Z/parts/hecke/section.tex},
AM-H06, equations (H06.2)--(H06.4); its proof was read
in full for this application.  The window and host
quantifiers are those of
\path{Development/AutomorphicArithmetic/CharacterBoxes/ExistsPageN.lean},
declarations \texttt{WinPair}, \texttt{PageN},
\texttt{VVle}, \texttt{VAt} and \texttt{VWinAt}.
The source's previously unbounded $d_i,d_{\mathrm{loc}}$
are now quantitatively controlled on these hosts.
Neither (FG.1) nor (FG.2) bounds the length of an
infinite-coefficient strict group from its residual
cohomology.  That main depth estimate, its all-prime
sum, and the arithmetic theta tie remain separate.

\paragraph{Algebraic-geometric dictionary and source access.}
The finite module $M$ is a locally constant finite
etale coefficient on $\operatorname{Spec}\mathcal O_F[S^{-1}]$;
the finite covering that trivializes it belongs to the
Galois-category dictionary of SGA~1, Expos\'e~V,
\S\S4 and 7.  The Fourier cokernel is a finite integral
coefficient supported at the bad coefficient primes.
Equations (FG.1)--(FG.4) bound the cost of that coefficient
in the global and local cohomology of the arithmetic
curve.  They do not produce a surface square or a
polarization of the infinite-coefficient object.
No new Grothendieck archive scan is used here; the
archive access qualifications in Section~\ref{sec:groth-source}
remain unchanged.  The primary source URLs for the
number-theoretic inputs are
\url{https://www.jmilne.org/math/CourseNotes/ANT.pdf},
\url{https://www.jmilne.org/math/Books/ADTnot.pdf}, and
\url{https://math.ucla.edu/~sharifi/notes/iwasawa-app-a.html}.


\part{Unproved comparisons and source-access limits}
\section{The central arithmetic gate has not been removed}
The strongest cubic input carried from ALP\_MACHINE is the
generic pole/window implication with the actual full local
coefficients. The new ray-class integral model resolves
the algebra and its ramified Frobenius lifts. AC-G22 now
proves that the fixed positive-weight $3$-adic twisted
group is finite cyclic, while AC-K10--AC-K11 identify its
length with an even K-theory size up to explicit local
terms. AC-G26 sharpens the character comparison to the
exact formula $m_e=\kappa_e-2(1+v_3(e))+1$. The exponent
is now uniformly bounded for this fixed field and coefficient
prime three: AC-K13--AC-K21 combine certified logarithmic
class data, genus theory and conjugation to prove finiteness
of the split cyclotomic Iwasawa module and the explicit
bound $4\leq m_e\leq5$, with $m_{e+3}=m_e$ and
$m_e=4$ when $3\mid e$. This closes that
specific growth question. It does not give the all-prime
aggregate or the uniform bound over the actual growing
CM hosts in the source's window theorem. AC-G28--AC-G29
also retain the exact mixed-character arithmetic tie still
needed for the regular CM construction. AC-G30 supplies
the complete bare Host enlargement containing the cubic
coefficient field, with its exact integral lattice cost.
AC-K19 bounds the finite Fourier comparison errors over
these growing hosts by a polynomial, negligible relative
to the arithmetic window budget when positive relative
slack is available. This removes the finite coefficient
error from the remaining growth obligations; it does not
bound the infinite-coefficient Selmer depths.
AC-G32 then supplies a smaller host satisfying one literal
\texttt{IsChiField}, and AC-G33 constructs its required
signature unit and a principal elliptic CM model on the
unchanged integral lattice. This closes those field and unit
existence gates. The changed CM type still needs its actual
local and theta comparisons; the mixed-block cohomology tie
and the growing-host Selmer budget remain separate.
The bounded test-function window of AC-K09 and the
unbounded one-sided arithmetic windows of H04.4 have
different variables and quantifiers. Conflating them
would be a gap. The final ledger and implications register
will state the exact gate reached at the last review.

\section{Spectral and geometric comparisons still requiring proof}
AC-T39 now proves exact vague convergence of the actual positive
prolate defect measures to the zero measure, including multiplicities,
for every nonprincipal primitive Dirichlet sector under A.
The original full trace-test passage still requires a uniform
weighted high-frequency tail estimate; compact frequency convergence
does not supply it. The actual adelic cutoff dictionary retains
the scope stated in AC-T11 and AC-T30.
No proof here identifies a specified prolate eigenvalue
with an individual zeta ordinate merely from a counting
asymptotic. The original prolate counting comparison is
recorded in \path{modules/spectral_prolate_unchecked.md},
including the versions whose logarithmic terms require
reconciliation. It is not used in any proved finding.
No proof here constructs the full arithmetic square
from the numerical Hodge inequality. An isometry of a
marked numerical kernel is not a scheme, a motive, or
a polarized cohomology theory. The unconditional
Riemann--Roch statements cited in the body retain their
actual spaces, dimension functions and exceptional
terms; A does not silently change those definitions.

\section{Grothendieck source dictionary and access levels}
\label{sec:groth-source}
The geometric interpretations in the body identify the
category in which each construction lives. A framework
reference is not an attribution of a new analytic theorem
to Grothendieck. The following exact source entries were
read in the indicated form in this run.
\begin{itemize}
\item SGA~1, expos\'e~V, \S4, ``Conditions axiomatiques
d'une th\'eorie de Galois'', and \S7, ``Cas des
pr\'esch\'emas'': finite continuous Galois sets and
finite \'etale covers. The online French transcription
was read. This is the framework for AC-K01--AC-K06,
AC-B03--AC-B06 and the finite abelian monodromy results.
\item SGA~1, expos\'e~XI, \S6, ``Application aux
rev\^etements principaux: th\'eories de Kummer et
d'Artin--Schreier'', Proposition~1, Definition~2,
Corollary~3, Proposition~4 and Corollary~5: the group
$\mu_n$, its \'etaleness away from $n$, and the
unit/Picard description of Kummer torsors. The online
French transcription was read, including the exact
sequence. This directly describes the cubic covers
after inverting $3p$ and their coefficient-field
character projectors.
\item SGA~4, expos\'e~II, Theorem~4.1: existence and
left exactness of associated sheaves. The online
French transcription was read. This is the applicable
framework for AC-G09, not a statement that the
sectionwise split-zero presheaf was already a sheaf.
\item SGA~6, expos\'e~V, 2.2(b) and Lemma~2.2.1:
exterior-power operations and their representation-ring
identities. The online French transcription was read.
These locate the $\lambda$-ring and tensor language
for the Artin--Tate constructions and AC-K08. They
do not imply global zero bounds for nonabelian
Artin functions.
\end{itemize}

The Montpellier catalogue was consulted for cotes 2
(Riemann--Roch), 10 (tensor structures), 11
(correspondences), 14 (Hodge--Tate--Lefschetz), 16
(standard conjectures), and 19 (topoi). These are
\emph{catalogue-level thematic references}. The scans
were not opened successfully here, so no new
scan-level verification is claimed. Earlier supplied
documentation contains readings of cote54, scans003,
005 and007, and cote4, scans003,005 and006. Those
records are carried as historical source records,
not as scans reread by this run. Any requested
archival implication requiring such a reading
remains a source-access gate. The explicit source
manifest and source-check evidence preserve the
distinction.

The identification $\Delta=\chi_{12}\circ u$ in the
updated intake note has not been independently
verified: the underlying DELTA23 intake and definition
of $u$ were not located in the supplied material.
AC-B06 uses only its fully proved character table.

\appendix
\section{Evidence, reproducibility, and licensing}
Every input file is hashed in
\path{evidence/source_manifest.json} or
\path{evidence/followup_1_source_manifest.json}.
The extraction was performed once for each attachment.
The Python scripts in \path{scripts/} are finite
certificates of the expressly stated algebraic and
interval computations. They are not formal proofs
of A, a Lean build of the inherited archives, or
certificates of cubic zero location. Independent
whole-proof readings are recorded in \path{reviews/}
and the check register.

Checkpoint timing and the scheduled movement reviews
are in \path{LEDGER.md}; actual clock times are used.
No earlier checkpoint ZIP is included inside a later
one. The source ZIPs remain separate inputs.
Text newly written by the AI collaboration is
dedicated to the public domain under CC0~1.0, as
requested. This does not change the rights in the
attached literature or third-party source material.

\end{document}
